\documentclass[11pt, b5paper]{book}
\usepackage[b5paper]{geometry}
\usepackage[slovak]{babel}
\usepackage[utf8]{inputenc}
\usepackage[T1]{fontenc}
\usepackage{tabularx}
\usepackage{amsfonts}
\usepackage{amssymb}
\usepackage{amsmath}
\usepackage{epsfig}
\usepackage{epigraph}
\setlength{\epigraphwidth}{.6\textwidth}
\usepackage{color}
\usepackage{mathrsfs}
\usepackage{verbatim}
\usepackage{pifont}
\usepackage[unicode]{hyperref}
\usepackage[autostyle]{csquotes}
\usepackage[inline]{enumitem}
\usepackage{probsoln}
\usepackage{multicol}
\usepackage{theorem}
\usepackage{interval}
\usepackage{tasks}
\usepackage{xcolor}
\usepackage[framemethod=TikZ]{mdframed}

\usepackage{pgfplots}
\usepgfplotslibrary{polar}
\usepgflibrary{shapes.geometric}
\usetikzlibrary{calc}
\pgfplotsset{my style/.append style={axis x line=middle, axis y line=middle, xlabel={$x$}, ylabel={$y$}, axis equal }}

% Justified citácie na začiatku kapitol
\renewcommand{\textflush}{flushepinormal}

% Menší znak ekvivalencie
\DeclareRobustCommand\iff{\;\Leftrightarrow\;}

% Štýl úvodzoviek
%\DeclareQuoteAlias{dutch}{slovak}

% probsoln package Solution label redefinition
\renewcommand{\solutionname}{Riešenie}

% šedé pozadie riešení
\renewenvironment{solution}{
  \par
  \begin{mdframed}[backgroundcolor=lightgray]
  \noindent
  \textbf{\solutionname}:
  \ignorespaces
}{
  \ignorespacesafterend\end{mdframed}
}

% sign function
\newcommand{\sgn}{\operatorname{sgn}}
\DeclareMathOperator{\arcsec}{arcsec}
\DeclareMathOperator{\arccot}{arccot}
\DeclareMathOperator{\arccsc}{arccsc}

\settasks{
  %debug=true
}


% Author: Michal Foríšek, misof ~at~ ksp.sk
\catcode`\@=11

% definicie znamych tried jazykov
%\def\R{\mathscr{R}}   % R_0 je \R_0
\def\Llin{\mathscr{L}_{\it lin}}
\def\Ldpda{\mathscr{L}_{\it DPDA}}
\def\Ledpda{\mathscr{L}_{\it eDPDA}}
\def\Lcf{\mathscr{L}_{\it CF}}
\def\Lcs{\mathscr{L}_{\it CS}}
\def\Lecs{\mathscr{L}_{\it ECS}}
\def\Lrec{\mathscr{L}_{\it rec}}
\def\Lre{\mathscr{L}_{\it RE}}

\def\L{\mathscr{L}}   % vseobecna trieda jazykov
%\def\pisD{\mathscr{D}} %pisane D - vseobecna trieda automatov
%\def\L{{\cal L}}   % vseobecna trieda jazykov
%\def\Ldcf{{\cal L}_{\it dCF}}
%\def\Lcf{{\cal L}_{\it CF}}
%\def\Lcs{{\cal L}_{\it CS}}
%\def\Lecs{{\cal L}_{\it ECS}}
%\def\Lre{{\cal L}_{\it RE}}
%\def\Lrec{{\cal L}_{\it rec}}
\def\R{{\cal R}}
\def\cent{{c\kern-0.3em|\kern0.1em}}
\def\N{{N}} % FIXME FIXME

%\def\cent{%
% \mathchoice%
%  {\rm c\kern-0.21em\vrule width0.022em depth0.05em height0.50em\kern-0.022em\kern0.21em}%
%  {\rm c\kern-0.21em\vrule width0.022em depth0.05em height0.50em\kern-0.022em\kern0.21em}%
%  {\rm c\kern-0.15em\vrule width0.0205em depth0.035em height0.335em\kern-0.0205em\kern0.15em}%
%  {\rm c\kern-0.12em\vrule width0.018em depth0.030em height0.25em\kern-0.018em\kern0.12em}%
%}

\let\eps=\varepsilon

\def\relupdown#1#2#3{\mathrel{\mathop{#1}\limits^{#2}_{#3}}}

\let\then=\Rightarrow
\let\neht=\Leftarrow

\def\krok#1{\relupdown{\Longrightarrow}{}{#1}}
\def\thenrm{\relupdown{\Longrightarrow}{}{rm}}

\def\bicik{\upharpoonright}
\def\B{{\mathbf B}}
\def\kodTS#1{{\tt <}#1{\tt >}}
\def\bigdot{\textrm{\Huge .}}

\newtheorem{definicia}{Definícia}[section]
\newtheorem{HLPpoznamka}{Poznámka}[section]
\newtheorem{HLPpriklad}{Príklad}[section]
\newtheorem{HLPcvicenie}[HLPpriklad]{Cvičenie}
\newtheorem{zadanie}{Úloha}[section]
\newenvironment{poznamka}{\begin{HLPpoznamka}\rm}{\end{HLPpoznamka}}
\newenvironment{priklad}{\begin{HLPpriklad}\rm}{\end{HLPpriklad}}
\newenvironment{cvicenie}{\begin{HLPcvicenie}\rm}{\end{HLPcvicenie}}
\newtheorem{veta}{Veta}
\newtheorem{lema}[veta]{Lema}
\newtheorem{dosledok}[veta]{Dôsledok}
\newtheorem{teza}[veta]{Téza}
% \newtheorem{dokaz}{Dôkaz}[section]

\long\def\odsadene#1{
\leftskip=\parindent
\parindent=0pt
\vskip-5pt

\parskip=5pt
#1
\parskip=0pt

\parindent=\leftskip
\leftskip=0pt

} % end \odsadene




%%%%%%%%%%% PROSTREDIE PRE PISANIE KOMENTAROV

%\newenvironment{komentar}{%
%\vskip\baselineskip
%\tabularx{0.95\textwidth}{|X|}
%\sl
%}
%{\endtabularx
%\vskip\baselineskip
%}

\newenvironment{komentar}{%
\vskip\baselineskip\noindent
\tabularx{\textwidth}{>{\hsize=.2\hsize}X>{\hsize=1.8\hsize}X}
\sl ~ & \sl
}
{\endtabularx
\vskip\baselineskip
}

%\newenvironment{komentar}{%
%\vskip\baselineskip
%\trivlist\vspace{-4pt}\raggedleft\item\relax\tabularx{0.9\textwidth}{X}\sl}
%{\endtabularx\vspace{-4pt}\endtrivlist
%\vskip\baselineskip
%}

\newenvironment{dokaz}{\trivlist
  \item[\hskip \labelsep{\bfseries Dôkaz.}]}{\endtrivlist}

%\newenvironment{dokaz}{%
%\vskip\baselineskip\noindent
%\tabularx{\textwidth}{||X||}
%\sl
%}
%{\endtabularx
%\vskip\baselineskip
%}

%%%%%%%%%%% PROSTREDIE PRE PISANIE KODU

\newcounter{program}[chapter]
\renewcommand \theprogram
   {\ifnum \c@chapter>\z@ \thechapter.\fi \@arabic\c@program}
\def\fps@program{tbp}
\def\ftype@program{1}
\def\ext@program{lop}
\def\fnum@program{Program~\theprogram}

\newenvironment{program}{\@float{program}}{\end@float}

\newcounter{schema}[chapter]
\renewcommand \theschema
   {\ifnum \c@chapter>\z@ \thechapter.\fi \@arabic\c@schema}
\def\fps@schema{tbs}
\def\ftype@schema{1}
\def\ext@schema{los}
\def\fnum@schema{Schéma~\theschema}

\newenvironment{schema}{\@float{schema}}{\end@float}

\newlength{\myboxwidth}
\myboxwidth=\textwidth \advance\myboxwidth -13pt

\long\def\mybox#1{\noindent%
\fcolorbox{black}{lgray}%
%  {\parbox{\myboxwidth}{ \vbox { % \vskip -0.74cm
  {\parbox{\myboxwidth}{ \begin{center}\begin{minipage}{0.95\myboxwidth} #1 % \vskip -0.74cm
\end{minipage}\end{center}}}}

\def\VlozKod#1#2{

\begin{program}[!bhtp]
\noindent%
\fcolorbox{black}{lgray}%
  {\parbox{\myboxwidth}{ \vbox { \InputC{code/#1} \vskip -0.74cm }} }
\caption{#2}
\end{program}

}

%%%%%%%%%%% PROSTREDIE PRE MOJE ITEMIZE

\newenvironment{myitemize}{%
\begin{itemize}
\itemsep-3pt
}
{\end{itemize}
}

%%%%%%%%%%% MATICKE MAKRA

\font\tenrm=csr10

\def\eps{\varepsilon}
% \def\R{{\mathbb R}}
\def\lvec#1{\overrightarrow{#1}}
\def\uhol{{\measuredangle}}
\def\then{\Rightarrow}
% \def\lg{{\rm lg}}
\def\lg{\log_2}
%\def\div{\mathbin{\rm div}}
\def\div{{\rm div}}

%%%%%%%%%%% PDF

\newif\ifpdf
\ifx\pdfoutput\undefined
  \pdffalse
\else
  \pdfoutput=1 \pdftrue
\fi

%%%%%%%%%%% OBRAZKY

\newcommand{\myincludegraphics}[2][]{\includegraphics[#1]{images/#2}}

%%%%%%%%%%% SLOVNICEK

\openout2=\jobname.slo

\newcommand{\definuj}[3][]{%
\def\tmpvoid{}\def\tmpfirst{#1}%
\ifx\tmpvoid\tmpfirst%
  {\sl #2}\label{definicia:#2}\write2{#2 & #3 & \pageref{definicia:#2} \cr}%
\else%
  {\sl #2}\label{definicia:#2}\write2{#1 & #3 & \pageref{definicia:#2} \cr}%
\fi}

\newcommand{\definujsilent}[2]{%
\label{definicia:#1}\write2{#1 & #2 & \pageref{definicia:#1} \cr}%
}

\newcommand\myglossary{
  \immediate\closeout2
  %\if@twocolumn\@restonecoltrue\onecolumn\else\@restonecolfalse\fi
  \chapter{Slovníček pojmov}
  \begin{tabular}{|l|l|r|}
  \hline
  {\bfseries slovenský pojem} & {\bfseries anglický preklad} & {\bfseries str.} \\
  \hline
  \InputIfFileExists{\jobname.srs}{}{~ & ~ & ~ \\}
  \hline
  \end{tabular}
  %\if@restonecol\twocolumn\fi
}

%%%%%%%%%%% UVODZOVKY

\catcode`\"=13
\def "{\begingroup\clqq\def "{\endgroup\crqq}}
\def\dospecials{\do\ \do\\\do\{\do\}\do\$\do\&%
  \do\#\do\^\do\^^K\do\_\do\^^A\do\%\do\~\do\"}

%%%%%%%%%%% DANGER BENDS

\font\manual=manfnt % font used for the METAFONT logo, etc.
\def\dbend{{\manual\char127}} % dangerous bend sign

\newlength{\bendwidth}   \settowidth{\bendwidth}{\dbend}    \newlength{\hangwidth}

\def\hangone{%
  \hangwidth=\bendwidth%
  \advance\hangwidth 5pt%
  \hangindent\hangwidth%
}
\def\hangtwo{%
  \hangwidth=\bendwidth%
  \multiply\hangwidth 2%
  \advance\hangwidth 6pt%
  \hangindent\hangwidth%
}

\def\medbreak{\par\ifdim\lastskip<\medskipamount \removelastskip\penalty-100\medskip\fi}
\let\endgraf=\par

\def\d@nger{\medbreak\begingroup\clubpenalty=10000
%\def\d@nger{\begingroup\clubpenalty=10000
%  \def\par{\endgraf\endgroup\medbreak} \noindent\hangone\hangafter=-2
  \def\par{\endgraf\endgroup} \noindent\hangone\hangafter=-2
  \hbox to0pt{\hskip-\hangindent\dbend\hfill}}
\outer\def\danger{\d@nger}

\def\dd@nger{\medbreak\begingroup\clubpenalty=10000
%  \def\par{\endgraf\endgroup\medbreak} \noindent\hangtwo\hangafter=-2
  \def\par{\endgraf\endgroup} \noindent\hangtwo\hangafter=-2
  \hbox to0pt{\hskip-\hangindent\dbend\kern1pt\dbend\hfill}}
\outer\def\ddanger{\dd@nger}

\def\enddanger{\endgraf\endgroup} % omits the \medbreak
\def\enddangerhop{\endgraf\endgroup\medbreak}




\def\@nakedcite#1#2{{#1\if@tempswa , #2\fi}}
\DeclareRobustCommand\nakedcite{%
  \@ifnextchar [{\@tempswatrue\@nakedcitex}{\@tempswafalse\@nakedcitex[]}}
\def\@nakedcitex[#1]#2{%
  \let\@citea\@empty
  \@nakedcite{\@for\@citeb:=#2\do
    {\@citea\def\@citea{,\penalty\@m\ }%
     \edef\@citeb{\expandafter\@firstofone\@citeb\@empty}%
     \if@filesw\immediate\write\@auxout{\string\citation{\@citeb}}\fi
     \@ifundefined{b@\@citeb}{\mbox{\reset@font\bfseries ?}%
       \G@refundefinedtrue
       \@latex@warning
         {Citation `\@citeb' on page \thepage \space undefined}}%
       {\hbox{\csname b@\@citeb\endcsname}}}}{#1}}

\long\def\FIXME#1{
  \begin{center}
  \begin{minipage}{0.8\textwidth}
  {\bf FIXME:~}\sl #1
  \end{minipage}
  \end{center}
}


\catcode`\@=12


%\includeonly{
%  ./tex/05-diferencialny-pocet/main
%}

\begin{document}

\intervalconfig{
soft  open  fences ,
}

\hyphenation{in-ver-zná}
\lccode`\'=39\hyphenation{l'Ho-spi-ta-lo-vho}

\thispagestyle{empty}
\begin{minipage}{0.25\textwidth}
  \includegraphics[width=0.9\textwidth]{img/komlogo-new}
\end{minipage}
\begin{minipage}{0.69\textwidth}
  \begin{center}
    \sc Katedra Informatiky \\
    Fakulta Matematiky, Fyziky a Informatiky \\
    Univerzita Komenského, Bratislava
  \end{center}
\end{minipage}
\vfill
\begin{center}
  \begin{minipage}{0.8\textwidth}
    \hrule
    \bigskip\bigskip
    \centerline{\LARGE\sc Cvičenia z matematickej analýzy I}
    \bigskip\bigskip
    \centerline{\large\sc Zbyněk Kubáček}
    \bigskip
    \centerline{\large\sc Ján Valášek}
    \bigskip\bigskip
    \hrule
  \end{minipage}
\end{center}
\vfill
{~}
\hfill {\bf\today}

\hfill \textbf{Editori}

\hfill Ivona Hrivová, Jakub Šimo
\eject % EOP i

\chapter*{Predslov}

\setlength\epigraphrule{0pt}
\epigraph{
  \enquote{Človek, ktorý nedokáže hovoriť tak,
  aby mu ostatní rozumeli, je idiot.
  Chápeš ma?} \\
  \enquote{Nie, otec.}
}{}

Predslov je pre zúfalých autorov často poslednou príležitosťou napraviť,
čo sa ešte dá. Chápeme sa preto možnosti vysvetliť zachmúrenému čitateľovi,
ako sme to vlastne mysleli.

Skriptá \emph{Cvičenia z matematickej analýzy I} sú určené pre študentov prvého
ročníka odboru Matematika a tých študentov prvého ročníka odboru Učiteľstvo
všeobecnovzdelávacích predmetov, ktorých aprobačným predmetom je matematika;
obsahujú cvičenia k látke preberanej v prvom semestri.

Dvojitou čiarou sú v týchto skriptách označené definície a tie tvrdenia, ktoré
sa čitateľ môže - bez toho, že by ich dokazoval - odvolávať pri riešení za nimi
nasledujúcich úloh. Teda:
\begin{enumerate*}[label=\alph*)]
  \item pri riešení príkladu sa nemožno odvolávať na tvrdenia uvedené až za ním
  \item riešenie musí obsahovať dôkazy všetkých ostatných použitých tvrdení.
\end{enumerate*}

Jednoduchou čiarou sú označené riešenia príkladov uvedené priamo v texte.

V každej kapitole možno rozlíšiť dve časti; druhá z nich (doplňujúca) má vždy
názov \textit{Ďaľšie príklady}. Do prvej (hlavnej) časti každej kapitoly sme
zaradili úlohy, ktoré pokladáme za základné, odsek \textit{Ďaľšie príklady}
obsahuje aj trocha náročnejšie cvičenia. (Skoro v každej študíjnej skupine sa
totiž nájdu študenti, ktorí príklady preberané na cvičeniach už ovládajú alebo
si to aspoň myslia. Nech si teda ticho prepočítavajú úlohy z odseku
\textit{Ďaľšie príklady} a nedeprimujú svojím znudeným výrazom cvičiaceho. Tým
samozrejme nechceme povedať, že by tieto úlohy boli neprístupné pre ostatných
študentov.)

V poslednej časti skrípt sú uvedené výsledky všetkých výpočtových príkladov a
návody k mnohým dôkazovým. Najmä druhá z týchto skutočností by mohla pôsobiť
demoralizujúco, preto sa pri nej zastavme. Cieľom zbierky úloh nie je podľa
nášho názoru študenta odradiť, ale niečo ho naučiť. Či sa niečo naučí neborák,
štyri hodiny bezvýsledne riešiaci nejaké cvičenie, je otázne.

Samozrejme, že najnevhodnejším možným postupom je prečítať si zadanie a hneď si
vzadu nalistovať návod; človek sa tak iste veľa nenaučí. Morálne právo uchýliť
sa k návodu dáva až predchádzajúca neúspašná snaha o nájdenie riešenia. (Celý
tento odstavec sa napokon rovnako vzťahuje na riešené príklady.) Vzadu uvedené
návody nie sú samozrejme jediné možné, iné postupy môžu byť rovnako správne
(nehovoriac už o chybách a nedopatreniach, ktoré sa vo výsledkoch a návodoch
zákonite vyskytnú).

Na úlohu, ktorá je súčasťou príkladu pozostávajúceho z viacerých zadaní, sa
odvolávame vždy ( aj v rámci toho istého príkladu) dvojčíslom: napr. na druhú
úlohu v príklade 4 sa odvolávame ako na pr. 4.2. Pokiaľ sme u niektorého
príkladu nepovažovali za vhodné zaradiť do poslednej časti skrípt návod na jeho
riešenie, má jeho číslo index $_0$. Pokiaľ je niektorý príklad formulovaný ako
tvrdenie, očakávame, že ho čitateľ dokáže, preto nie všetky dôkazové príklady
sa začínajú "Dokážte, že ...".

Napokon ešte jedna poznámka: goniometrické, exponenciálne a logaritmické
funkcie sú zväčša na cvičeniach potrebné skôr, než sa ich podarí na prednáške,
resp. v učebnici korektne zaviesť. Pri zostavovaní príkladov sme preto od
začiatku považovali tieto funkcie a ich vlastnosti (grafy, goniometrické vzorce
a pod.) za známe.

Ďakujeme na tomto mieste všetkým, ktorí príspeli k výslednej podobe popisu
týchto skrípt. Sú to predovšetkým recenzenti doc. RNDr. Jozef Vencko, CSc. a
RNDr. Igor Book, CSc., ďalej RNDr. Zuzana Ondrejková, RNDr. Ivan Kupka, Ing.
Jiří Kubáček, CSc., RNDr. Iveta Kundraciková a RNDr. Viera Čerňanová. Za
mimoriadnu ochotu, s ktorou nám vyšiel v ústrety, sme zaviazaní vedúcemu
Vydavateľského oddelenia RUK Ing. Antonínovi Skovajsíkovi.


\tableofcontents

\chapter*{Úvod (výroky, dôkazy)}
% Pridáva TOC entry pre nečíslovanú kapitolu
\addcontentsline{toc}{chapter}{Úvod}
% Opravuje to, že nečíslovaná kapitola berie header z predchádzajúcej kapitoly
\markboth{ÚVOD}{}

\loadallproblems[uvod]{./tex/00-uvod/ulohy.tex}
\epigraph{
  \enquote{Prosím také, aby čtenář nepovažoval tyto vědomě neúplné poznámky za
nějaký (sebe elementárnejší) úvod do matematickej logiky. Šlo mně jenom o to,
abychom si ujasnili způsob vjadřovaní, kterého ostatně nadaný matematik užívá
intuitivně s neomylnou správností.}
}{\textit{Vojtěch Jarník, Diferenciální počet (II)}}

V matematike sa stretávame s tvrdeniami, o ktorých má
zmysel povedať, že sú pravdivé alebo nepravdivé. Také tvrdenia sa nazývajú
\emph{výroky}.

Z výrokov A, B môžeme pomocou operácií negácie ($\neg$), konjukcie
($\land$), disjunkcie ($\lor$), implikácie ($\Rightarrow$) a ekvivalencie
($\iff$) vytvoriť nové, tzv. \emph{zložené výroky}:

\begin{center}
  \begin{tabular}{|m{0.15\linewidth} | p{0.6\linewidth}|}
    \hline
    $\neg A$ & negácia výroku $A$ \\
    \hline
    \bigskip
    $A \land B$
      & \enquote{$A$ a súčasne $B$},
        \enquote{$A$ a $B$};

        namiesto $A \land B$ a často píše $A, B$ \\
    \hline
    \bigskip
    $A \lor B$
      & \enquote{$A$ alebo $B$},
        \enquote{alebo $A$ alebo $B$} \\
    \hline
    \bigskip
    $A \Rightarrow B$
      & \enquote{ak $A$, tak $B$},
        \enquote{nech $A$, potom $B$},

        \enquote{$A$ je postačujúca podmienka pre $B$},

        \enquote{$B$ je nutná podmienka pre $A$} \\
    \hline
    \bigskip
    $A \iff B$
      & \enquote{$A$ vtedy a len vtedy, keď $B$},

        \enquote{$A$ práve vtedy, keď $B$},

        \enquote{$A$ je nutná a postačujúca podmienka pre $B$} \\
    \hline
  \end{tabular}
\end{center}

Pravdivosť týchto zložených výrokov závisí len na pravdivosti výrokov $A, B$,
a to následovne (1 označuje, že daný výrok je pravdivý; 0, že je nepravdivý):
\begin{center}
  \begin{tabular}{|c|c|}
    \hline
    $A$ & $\neg A$ \\
    \hline
    1 & 0 \\
    0 & 1 \\
    \hline
  \end{tabular}
  \quad
  \begin{tabular}{|c|c|c|c|c|c|}
    \hline
    $A$ & $B$ & $A \land B$ & $A \lor B$ & $A \Rightarrow B$ & $A \iff B$ \\
    \hline
    1 & 1 & 1 & 1 & 1 & 1 \\
    1 & 0 & 0 & 1 & 0 & 0 \\
    0 & 1 & 0 & 1 & 1 & 0 \\
    0 & 0 & 0 & 0 & 1 & 1 \\
    \hline
  \end{tabular}
\end{center}

Pre negáciu zložených výrokov platia tieto pravidlá:
\begin{center}
  \begin{tabular}{c | c}
    výrok & negácia \\
    \hline
    $\neg A$ & $A$ \\
    $A \land B$ & $\neg A \lor \neg B$ \\
    $A \lor B$ & $\neg A \land \neg B$ \\
    $A \Rightarrow B$ & $A \land \neg B$ \\
    $A \iff B$ & $(A \land \neg B) \lor (B \land \neg A)$ \\
  \end{tabular}
\end{center}

Často budeme využívať, že pre ľubovoľné výroky $A, B$ je:
\begin{enumerate}
  \item
    implikácia rovnocenná s obmenou (kontrapozíciou)
    \[
      A \Rightarrow B \iff \neg B \Rightarrow \neg A
    \]
  \item
    ekvivalencia rovnocenná s výrokom
    \[
      (A \iff B) \iff (A \Rightarrow B) \land (B \Rightarrow A)
    \]
\end{enumerate}

Každá definícia v matematike má podobu ekvivalencie. Je však zvykom nevyjadrovať
v definícií symbol \enquote{$\iff$} slovami \enquote{vtedy a len vtedy}, ale
slovom \enquote{ak}. Teda napr. definíciu v slovnej podobe \enquote{funkcia sa
nazýva rastúca, ak platí ...} treba chápať ako ekvivalenciu \enquote{funkcia je
rastúca $\iff$ ...}.

\subsubsection*{Výrokové formy. Všeobecné a existenčné výroky a ich negovanie}

Výroková forma je výraz, ktorý síce nie je výrokom, ale obsahuje premenné,
ktorých vhodným nahradením vznikne výrok (napr. výroková forma $x^2 + y \geq 0
\land y < 0$ obsahuje dve také premenné). Tvorenie zložených výrokových foriem
je analogické tvoreniu zložených výrokov.

Nech $\varphi (x)$ je výroková forma, v ktorej je možno dosadzovať za jedinú
premennú $x$, a $A$ je ľubovoľná neprázdna množina obsahujúca len také prvky,
ktorých dosadením do $\varphi (x)$ vznikne výrok.

Výrok \enquote{pre všetky $x \in A$ platí $\varphi (x)$} -- symbolicky ho
zapisujeme
\[
  \forall x \in A_1: \varphi (x)
\]
sa nazýva všobecný symbol $\forall$ zastupujúci slová \enquote{pre všetky} je
tzv. všeobecný (alebo veľký) kvantifikátor. (Častá slovná podoba všeobecného
výroku $\forall x \in A : \varphi (x)$ je \enquote{nech $x \in A$, potom
(platí) $\varphi (x)$} alebo \enquote{ak $x \in A$, tak (platí) $\varphi
(x)$}.)

Výrok \enquote{existuje $x \in A$ také, že platí $\varphi (x)$} -- symbolicky
\[
  \exists x \in A : \varphi (x)
\]
sa nazýva existenčný výrok; symbol $\exists$ vyjadrujúci slovo
\enquote{existuje} (\enquote{pre niektoré}) je tzv. existenčný (malý)
kvantifikátor.

Pre negáciu všeobecných a existenčných kvantifikátorov platia tieto pravidlá:
\begin{center}
  \begin{tabular}{c | c}
    výrok & negácia \\
    \hline
    $\forall x \in A : \varphi (x)$ & $\exists x \in A : \neg \varphi (x)$ \\
    $\exists x \in A : \varphi (x)$ & $\forall x \in A : \neg \varphi (x)$ \\
  \end{tabular}
\end{center}

Ak uplatníme všeobecný (alebo existenčný) kvantifikátor na premennú $x_1$ vo
výrokovej forme $\varphi(x_1, x_2, \ldots, x_n)$ -- v ktorej možno dosadzovať za
$n$ premenných $x_1, \ldots, x_n$ kde $(n \geq 2)$ -- dostaneme výrokovú formu
$\forall x_1 \in A: \varphi (x_1, \ldots, x_n) (\exists x_1 \in A: \varphi
(x_1, \ldots, x_n))$, v ktorej možno dosadzovať už len za $n - 1$ premenných
$x_2, \ldots, x_n$ za premennú viazanú kvantifikátorom už dosadzovať nemožno,
predpokladáme samozrejme, že $A \neq B$ a že pre každé $a \in A$ je výraz
$\varphi (a, x_2, \ldots, x_n)$ opäť výrokovou formou). Na premenné $x_2$ opäť
možeme uplatniť kvantifikátor atď. Ak v poradí posledný použitý -- t.j. pri
zápise prvý -- kvantifikátor je veľký (malý), dostaneme takýmto postupom
všeobecný (existenčný) výrok.

Dôkaz ekvivalencie $A \iff B$ spravidla prevádzame na dôkaz dvoch implikácií
$A \Rightarrow B$ a $B \Rightarrow A$

\subsubsection*{Podrobnejšie o dôkazoch všeobecných výrokov:}
Ako príklad priameho dôkazu všobecného výroku nám poslúži dôkaz tvrdenia
\enquote{ak celé číslo $n$ nie je deliteľné 3, tak $n^2 - 1$ je deliteľné 3}
(symbolicky $\forall n \in \mathbb{Z}: (3 \nmid n \Rightarrow 3 \mid
(n^2 - 1))$), ktorý prebieha nasledovne: ak $n$ nie je deliteľné 3, tak
$n - 1$ alebo $n + 1$ je deliteľné 3, teda $n^2 - 1$ je deliteľné 3. Symbolicky
možno tento dôkaz zapísať takto:
\[
3 \nmid n
\Rightarrow (3 \mid (n - 1) \lor 3 \mid (n + 1))
\Rightarrow 3 \mid (n^2 - 1)
\]

Keby sme v uvedených úvahách namiesto $n$ písali napr. číslo 5, dostali by sme
dôkaz implikácie $3 \nmid 5 \Rightarrow 3 \mid (5^2 - 1)$. To je pre tento typ
dôkazu charakteristické; priamy dôkaz výroku $\forall x \in A: \varphi (x)$ je
vlastne schémou, z ktorej možno dostať dôkaz každého jednotlivého výroku
$\varphi (a)$, ak v nej za $x$ dosadíme prvok $a \in A$.

Zvláštnym typom priameho dôkazu všeobecného výroku je dôkaz matematickou
indukciou. Výroky v tvare $\forall n \in N: \varphi (n)$ sa spravidla nedokazujú
bezprostredne, namiesto toho stačí dokázať výroky
\begin{enumerate}[label=(\Alph*)]
  \item platí $\varphi (1)$
  \item $\forall n \in \mathbb{N}: \varphi (n) \Rightarrow \varphi (n + 1)$
\end{enumerate}
(výrok $\varphi (n)$ v implikácií $\varphi (n) \Rightarrow \varphi (n + 1)$ sa
nazýva indukčný predpoklad).Z (A) a (B) možno totiž reťazcom úsudkov odvodiť
ľubovoľný jednotlivý výrok $\varphi (n)$ takto:

platí $\varphi (1)$

\underline{platí $\varphi (1) \Rightarrow \varphi (2)$}

\hspace{1em} teda platí $\varphi (2)$;

\hspace{3.2em} \underline{platí $\varphi(2) \Rightarrow \varphi(3)$}

\hspace{4.4em} teda platí $\varphi(3)$;

\hspace{5cm} $\ldots$

\hspace{6cm} platí $\varphi (n - 1)$

\hspace{6cm} \underline{platí $\varphi (n - 1) \Rightarrow \varphi (n)$}

\hspace{7.2cm} teda platí $\varphi (n)$

Takýmto spôsobom možno dokazovať aj tvrdenia platné pre všetky prirodzené čísla
väčšie než dané číslo $p$, pre všetky celé čísla, pre všetky párne čísla a pod.

Pri nepriamom dôkaze všeobecného výroku odvodzujeme z jeho negácie nepravdivé
tvrdenie.

\medskip

\subsubsection*{Dôkazy existenčných výrokov}
Výrok $\exists ~ x \in A: \varphi (x)$
môžeme dokázať tak, že priamo určíme (skonštruujeme) prvok $a \in A$, pre ktorý
je výrok $\varphi (a)$ pravdivý. Taký dôkaz sa nazýva konštruktívny. Často však
odvodzujeme existenčný výrok z iných existenčných viet bez toho, že by sme
skutočne určili prvok s požadovanou vlastnosťou.

Existenčné výroky možno dokazovať aj nepriamo, teda odvodením nepravdivého
tvrdenia z ich negácie.

\begin{comment}
\begin{enumerate}
  \item Určte negácie následujúcich výrokov:

  \item Namiesto bodiek doplňte slová \enquote{je nutné}, \enquote{stačí},
        \enquote{je nutné a stačí} tak, aby vznikol pravdivý výrok:
  \begin{enumerate}[label=\arabic*.]
    \item aby súčet dvoch celých čísel bol deliteľný 2, $\ldots$ , aby každý
          sčítanec bol deliteľný 2;
    \item aby celé číslo bolo deliteľné 100, $\ldots$ , aby bolo deliteľné 10;
    \item aby celé číslo bolo deliteľné 100, $\ldots$ , aby bolo deliteľné
          1000;
    \item aby platila nerovnosť $\frac{1}{x} < 1, \ldots$ , aby bolo $x > 1$;
    \item aby platilo $\frac{1}{x} < 1, \ldots$ , aby bolo $x > 1$ alebo
          $x < 0$;
  \end{enumerate}
\end{enumerate}
\end{comment}

\begin{enumerate}
  \item \useproblem[uvod]{uvod-1}
  \item \useproblem[uvod]{uvod-2}
\end{enumerate}



\chapter{Množiny}

\loadallproblems[realne-cisla]{./tex/01-mnoziny/ulohy-realne-cisla.tex}
\loadallproblems[supremum-infimum]{./tex/01-mnoziny/ulohy-supremum-infimum.tex}
\loadallproblems[mohutnost]{./tex/01-mnoziny/ulohy-mohutnost.tex}
\loadallproblems[kompaktne]{./tex/01-mnoziny/ulohy-otvorene-uzavrete-kompaktne.tex}

Ďalej budeme používať tieto označenia:
\begin{itemize}[label=]
  \item $\mathbb{N}$ množina všetkých prirodzených čísel ($= \{ 1, 2, 3, \ldots \}$)
  \item $\mathbb{Z}$ množina všetkých celých čísel ($= \{ 0, -1, 1, 2, -2, \ldots \}$)
  \item $\mathbb{Q}$ množina všetkých racionálnych čísel
    ($= \{ \frac{p}{q}; p \in \mathbb{Z} \land q \in \mathbb{N} \}$)
  \item $\mathbb{R}$ množina všetkých reálny čísel
  \item $\mathbb{R}^+$ množina všetkých kladných reálny čísel ($= ( 0, \infty )$)
  \item $\mathbb{R}^+_0$ množina všetkých nezáporných reálny čísel
    ($= \interval[open right]{0}{\infty}$)
\end{itemize}
Ak pre niektorý prvok $a$ neprázdnej množiny $A \subset \mathbb{R}$ platí
\[
  (\forall x \in A) : x \leq a ((\forall x \in A): q \geq a)
\]
nazývame tento prvok \emph{maximum (minimum) množiny $A$} a označujeme ho
$\max(A)$ ($\min(A)$)

\section{Reálne čísla}
\begin{enumerate}[resume]
  \foreachproblem[realne-cisla]{\item\thisproblem}
\end{enumerate}


\section{Ohraničené množiny, supremum a infimum}
Neprázdna množina $A \subset \mathbb{R}$ sa nazýva \textit{zhora (zdola)
ohraničená}, ak platí
\[
  (\exists K \in  \mathbb{R})
    (\forall x \in A):
      x \leq K
\]
\[
  ((\exists K \in \mathbb{R})
    (\forall x \in A):
      x \geq K)
\]

Číslo $K$ s uvedenou vlastnosťou sa nazýva \textit{horné (dolné)
ohraničenie množiny $A$.} $\emptyset$ považujeme za ohraničenú zhora aj zdola.

Množina, ktorá je zhora aj zdola ohraničená, sa nazýva \textit{ohraničená}.
Množina, ktorá nie je ohraničená sa nazýva \textit{neohraničená}.

\begin{enumerate}[resume]
\item \useproblem[supremum-infimum]{supremum-infimum-14}
\item \useproblem[supremum-infimum]{supremum-infimum-15}
\item \useproblem[supremum-infimum]{supremum-infimum-16}
\end{enumerate}


Číslo $\alpha \in \mathbb{R}$ sa nazýva \textit{supremum množiny $A$}
$\subset \mathbb{R}, A \neq \emptyset$, ak
\[
  (\forall x \in A): x \leq \alpha \label{eq:supremum-i}
\]
\[
  (\forall \varepsilon > 0) (\exists \: x_\varepsilon \in A):
  x_\varepsilon > \alpha - \varepsilon
\]

Podľa \ding{34} je $\alpha$ horné ohraničenie množiny $A$; \ding{37} je negácia
výroku ``pre niektoré $\varepsilon > 0$ je číslo $\alpha - \varepsilon$ horným
ohraničením množiny $A$'', hovorí teda, že neexistuje horné ohraničenie množiny
$A$, ktoré by bolo menšie než $\alpha$. Teda $\alpha$ je najmenšie horné
ohraničenie množiny $A$. Supremum množiny $A$ označujeme $\sup \{ A \}$.

Číslo $\beta \in \mathbb{R}$ sa nazýva \textit{infimum množiny} $A \subset
\mathbb{R}, A \neq \emptyset$, ak:

\begin{align*}
\forall x \in A&: x \geq \beta \tag{\ding{34}} \\
(\forall \varepsilon > 0) (\exists \: x_\varepsilon \in A)&:
x_\varepsilon < \beta + \varepsilon \tag{\ding{37}}
\end{align*}

To znamená, že $\beta$ je najväčšie z dolných ohraniční množiny $A$. Infimum
množiny $A$ označujeme $\inf\{A\}$

Ak usporiadané pole $\mathbb{R}$ reálnych čísel konštruujeme z poľa
$\mathbb{Q}$ racionálnych čísel pomocou Dedekindových rezov, môžeme dokázať
nasledujúce dve ekvivalentné tvrdenia:

\begin{veta}
Každá neprázdna zhora ohraničená množina reálnych čísel má supremum.
\end{veta}

\begin{veta}
Každá neprázdna zdola ohraničená množina reálnych čísel má infimum.
\end{veta}

Ak usporiadané pole $\mathbb{R}$ zavádzame axiomaticky, považujeme prvé z
uvedených tvrdení za axiómu, z nej možno odvodiť vetu o existencií infima.

\showanswers
\begin{enumerate}[resume]
\item \useproblem[supremum-infimum]{supremum-infimum-17}
\hideanswers
\item \useproblem[supremum-infimum]{supremum-infimum-18}
\item \useproblem[supremum-infimum]{supremum-infimum-19}
\item \useproblem[supremum-infimum]{supremum-infimum-20}
\item \useproblem[supremum-infimum]{supremum-infimum-21}
\end{enumerate}


\section{Mohutnosť množín}
Neprázdna množina $A$ sa nazýva konečná, ak pre niektoré prirodzené číslo $n$
existuje bijekcia $\varphi:\{1,...,n\}\rightarrow A$ (t.j. $\varphi$ je prosté
zobrazenie
%\footnote{Ak každému prvku a neprázdnej množiny $A$ priradíme práve
%jeden prvok neprázdnej množiny $B$, ktorý označíme $\varphi (a)$ (množiny $A,B$
%nemusia byť číselné), tak hovoríme, že $\varphi$ je zobrazenie množiny $A$ do
%množiny $B$. Špeciálne, ak $A=\mathbb{N}$, hovoríme, že $\varphi$ je postupnosť
%prvkov z $B$. Zrejme pojmy funkcie a postupnosti uvedené v odseku $1.3$ sú
%špeciálnymi prípadmi tu uvedených pojmov zobrazenie a postupnosť.}
a $\varphi(\{1,...,n\})=A$). Prázdnu množinu pokladáme za konečnú. Množina,
ktorá nie je konečná, sa nazýva nekonečná.

Množina $A$ sa nazýva nekonečne spočítateľná, ak existuje bijekcia
$\varphi:\mathbb{N}\rightarrow A$ (t.j. ak prvky množiny $A$ možno zoradiť do
prostej postupnosti). Konečné a nekonečne spočítateľné množiny sa označujú
spoločným názvom spočítateľné. Množina, ktorá nie je spočítateľná, sa nazýva
nespočítateľná.

\begin{veta}
Každý nedegenerovaný interval je nespočítateľná množina (degenerovanými
intervalmi sa nazývajú jednoprvkové množiny).
\end{veta}

\textbf{Poznámka:} Namiesto dvojice pojmov nekonečne spočítateľná - spočítateľná
sa často v tom istom význame používa dvojica spočítateľná - najviac
spočítateľná.

\begin{enumerate}[resume]
  \item \useproblem[mohutnost]{mohutnost-90}
  \item \useproblem[mohutnost]{mohutnost-91}
  \item \useproblem[mohutnost]{mohutnost-92}
  \item \useproblem[mohutnost]{mohutnost-93}
  \item \useproblem[mohutnost]{mohutnost-94}
  \item \useproblem[mohutnost]{mohutnost-95}
  \item \useproblem[mohutnost]{mohutnost-96}
  \item \useproblem[mohutnost]{mohutnost-97}
  \item \useproblem[mohutnost]{mohutnost-98}
  \item \useproblem[mohutnost]{mohutnost-99}
\end{enumerate}


\section{Otvorené, uzavreté a kompaktné množiny}
Bod $a \in \mathbb{R}$ sa nazýva \textit{vnútorný bod množiny $a \subset
\mathbb{R}$}, ak existuje také jeho okolie $O(a)$, že platí $O(a) \subset A$.

Neprázdna množina $a \subset \mathbb{R}$ sa nazýva \textit{otvorená}, ak každý
jej prvok je jej vnútorným bodom. Ďalej sa dohodneme, že prázdnu množinu budeme
pokladať za otvorenú. Množina $A \subset \mathbb{R}$ sa nazýva
\textit{uzavretá}, ak je $\mathbb{R} \setminus A$ otvorená.

\begin{veta}
Pre ľubovoľnú množinu $A \subset \mathbb{R}$ sú nasledujúce tvrdenia
ekvivalentné:
\begin{enumerate}
\item
    $A$ je uzavretá množina
\item
    ak $a \in \mathbb{R}$ je hromadný bod množiny $A$, tak $a \in A$
\item
    ak všetky členy konvergentnej postupnosti ${\{a_n\}}_{n=1}^\infty$ sú
    prvkami množiny $A$, tak $\lim\limits_{n \rightarrow \infty} a_n \in A$
\end{enumerate}
\end{veta}

\textit{Poznámka:}
Často sa pojem uzavretej množiny definuje pomocou vlastnosti $2.$ alebo $3.$.
Otvorená množina sa niekedy definuje pomocou uzavretej, a to tak, že sa najprv
pomocou niektorej z vlastností $2.$,$3.$ zavedie pojem uzavretej množiny a za
otvorené sa potom vy.hlásia tie množiny $A$, ktorých doplnky $\mathbb{R}
\setminus A$ uzavreté.
\begin{enumerate}[resume]
    \item \useproblem[kompaktne]{kompaktne-208}
    \item \useproblem[kompaktne]{kompaktne-209}
    \item \useproblem[kompaktne]{kompaktne-210}
    \item \useproblem[kompaktne]{kompaktne-211}
    \item \useproblem[kompaktne]{kompaktne-212}
    \item \useproblem[kompaktne]{kompaktne-213}
\end{enumerate}

Systém $\{A_t: t \in I\}$ množín reálnych čísel (neprázdna množina I - nie nutne
číselná - sa nazýva indexová) sa nazýva \textit{pokrytie množiny $A \subset
\mathbb{R}$}, ak
\[
    A \subset \bigcup_{t \in I} A_t (:= \{x \in \mathbb{R};
    (\exists t \in I): x \in A_t\})
\]
Ak je naviac každá z množín $A_t$, $t \in I$,
otvorená, nazýva sa tento systém \textit{otvorené pokrytie množiny A}; termín
\textit{konečné pokrytie} používame, ak je množina I konečná.

Množina $A \subset \mathbb{R}$ sa nazýva \textit{kompaktná množina (kompakt)},
ak z každého jej otvoreného pokrytia $\{a_t: t \in I\}$ možno vybrať konečné
pokrytie (t.j. existuje konečná neprázdna množina $I' \subset I$ tak, že $A
\subset \bigcup_{t \in I} A_t$).

\begin{veta}
Pre ľubovoľnú množinu $A \subset \mathbb{R}$ sú nasledujúce tvrdenia
ekvivalentné:
\begin{enumerate}
\item $A$ je kompaktná množina
\item $A$ je uzavretá a ohraničená množina
\item
    ak všetky členy postupnosti ${\{a_n\}}_{n=1}^\infty$ sú prvkami množiny
    $A$, tak z ${\{a_n\}}_{n=1}^\infty$ možno vybrať konvergentnú podpostupnosť,
    ktorej milita je prvkom množiny $A$

\textit{Poznámka:}
Často sa kompaktná množina definuje pomocou vlastnosti $3.$.
\end{enumerate}
\end{veta}

\begin{enumerate}[resume]
    \item \useproblem[kompaktne]{kompaktne-214}
    \item \useproblem[kompaktne]{kompaktne-215}
    \item \useproblem[kompaktne]{kompaktne-216}
    \item \useproblem[kompaktne]{kompaktne-217}
    \item \useproblem[kompaktne]{kompaktne-218}
    \item \useproblem[kompaktne]{kompaktne-219}
    \item \useproblem[kompaktne]{kompaktne-220}
\end{enumerate}



\chapter{Funkcie}

\loadallproblems[funkcie]{./tex/02-funkcie/ulohy.tex}

\section{Definícia funkcie}
Nech $A\subset \mathbb{R}$ je neprázdna množina. Ak je každému číslu $x\in A$
priradené práve jedno číslo $y\in\mathbb{R}$, ktoré označíme $f(x)$, hovoríme,
že $f$ je funkcia (funkcia definovaná na množine $A$). Číslo $f(x)$ sa nazýva
funkčná hodnota (v bode $x$), množina $A$ definičný obor funkcie $f$ (túto
množinu budeme označovať $D(f)$). Na označenie funkcií budeme používať písmená
latinskej a gréckej abecedy. Ak chceme zdôraznič, že definičným oborom funkcie
$f$ je množina $A$, použijeme zápis $f:A \rightarrow \mathbb{R}$ (prípadne $f:A
\rightarrow B$, ak pre každé $x\in A$ platí $f(x)\in B$) alebo $f(x)$, $x\in A$.
Okrem označení \enquote{funkcia $f$}, \enquote{funkcia $f:A \rightarrow
\mathbb{R}^n$} sa možno stretnúť aj so spojeniami \enquote{funkcia $y=f(x)$}
alebo \enquote{funkcia $f(x)$} (istá nepresnosť posledných dvoch spojení spočíva
v tom, že symbol $f(x)$ sa zvykne označovať funkčná hodnota v danom bode $x$;
mnohí autori preto rozlišujú označenie $f(x)$ pre funkčnú hodnotu a $f(.)$ pre
funkciu).

Hovoríme, že funkcie $f$ a $g$ sa rovnajú, ak $D(f)=D(g)$ a pre každé $x\in
D(f)$ platí $f(x)=g(x)$ (teda funkcia je jednoznačne určená predpisom
priradenia a definičným oborom).

Funkciu $a$, ktorej definičným oborom je množina $N$, nazývame postupnosť a
označujeme ju spravidla $\{a_n\}_{n=1}^\infty$; funkčná hodnota v bode $n$ sa
nazýva $n$-tý člen postupnosti $\{a_n\}_{n=1}^\infty$ a označuje sa $a_n$.

Ak $f:A \rightarrow \mathbb{R}$ je funkcia a $B \subset A$ neprázdna množina,
tak množina $f(B):=\{f(x);x\in B \}$ sa nazýva obraz množiny $B$ (pri zobrazení
$f$). Špeciálne množina $f(A)$ sa nazýva obor hodnôt funkcie $f$.


\section{Zložené a elementárne funkcie}
Nech sú dané funkcie $f$, $g$.

\begin{itemize}
  \item
    Ak je množina $D_1:=D(f)\cap D(g)$ neprázdna, nazývajú sa funkcie $p, q, r$
    definované na množine $D_1$ predpismi
    \begin{align*}
      p(x) &= f(x)+g(x) \\
      q(x) &= f(x)-g(x) \\
      r(x) &= f(x)\cdot g(x)
    \end{align*}
    súčet, rozdiel a súčin funkcií $f, g$ a označujú sa $f+g$, $f-g$, $f\cdot g$.
  \item
    Ak je množina $D_2:=D(f)\cap \{x\in D(g);g(x)\neq 0\}$ neprázdna, nazýva sa
    funkcia $s:D_2 \rightarrow \mathbb{R}$ definovaná predpisom
    \[
      a(x)=\frac{f(x)}{g(x)}
    \]
    podiel funkcií $f$, $g$ a označuje sa $\frac{f}{g}$.
  \item
    Ak je množina $D_3:=\{x\in D(f);f(x)\in D(g)\}$ neprázdna, nazýva sa funkcia
    $t:D_3\rightarrow\mathbb{R}$ daná predpisom
    $$t(x)=g(f(x))$$
    zložená funkcia z funkcií $f$, $g$ (superpozícia funkcií $f$ a $g$) a označuje
    sa $g \circ f$. Funkcia $f$ sa nazýva vnútorná zložka, funkcia $g$ vonkajšia
    zložka funkcie $g \circ f$.
\end{itemize}

Základnými elementárnymi funkciami nazývame nasledujúce funkcie:

\noindent
% FIXME odkazy na footnotes v tabulke
\begin{tabular}{|p{0.21\textwidth}|p{0.22\textwidth}|p{0.45\textwidth}|}
  \hline
  \textbf{Názov} & \textbf{Predpis} & \textbf{Definičný obor} \\
  \hline
  konštantné &
  $\underset{a \in \mathbb{R}}{f(x)}\equiv a$  & $\mathbb{R}$ \\
  \hline
  mocninové &
  $\underset{a \in \mathbb{R} \setminus {0}}{f(x)}\equiv a$ &
  \begin{enumerate}[leftmargin=0.5cm]
    \item ak $a>0$:
      \begin{enumerate}[leftmargin=0.5cm]
        \item ak $a=\frac{p}{q},p,q\in\mathbb{N},$ $p$ je párne alebo $p$ aj $q$
              sú nepárne: $\mathbb{R}$
        \item vo všetkých ostatných prípadoch: $\interval[open right]{0}{\infty}$
      \end{enumerate}
    \item ak $a<0$:
      \begin{enumerate}[leftmargin=0.5cm]
        \item ak $a=-\frac{p}{q},p,q\in\mathbb{N}$, $p$ je párne lebo $p$ aj $q$
              sú nepárne: $\mathbb{R} \setminus \{0\}$
        \item vo všetkých ostatných prípadoch: $\interval[open]{0}{\infty}$
      \end{enumerate}
  \end{enumerate}
  \\
  \hline
  exponenciálne
    & $f(x)=a^x, $\newline$ a>0,a\neq 1$
    & $\mathbb{R}$ \\
  \hline
  logaritmické
    & $f(x)=\log_a x,$\newline$ a>0,a\neq 1$ & $\mathbb{R}^{+}$ \\
  \hline
  goniometrické & $f(x)=\sin x$ & $\mathbb{R}$ \\
    & $f(x)=\cos x$ & $\mathbb{R}$ \\
    & $f(x)=\tan x$ & $\mathbb{R}\ \{\frac{\pi}{2}+k\pi;k\in\mathbb{Z}\}$ \\
    & $f(x)=\cot x$ & $\mathbb{R}\ \{k\pi;k\in\mathbb{Z}\}$ \\
  \hline
  cyklometrické
    & $f(x)=\arcsin x$ & $\interval{-1}{1}$ \\
    & $f(x)=\arccos x$ & $\interval{-1}{1}$ \\
    & $f(x)=\arctan x$ & $\mathbb{R}$ \\
    & $f(x)=\arccot x$ & $\mathbb{R}$ \\
  \hline
\end{tabular}

Funkcie, ktoré vzniknú zo základných elementárnych funkcií len použitím
operácií súčtu, rozdielu, súčinu, podielu a superpozície funkcií, sa nazývajú
elementárne funkcie.

Všimnime si, že definičný obor funkcie, ktorá je súčtom, rozdielom, súčinom,
podielom alebo superpozíciou daných funkcií $f$ a $g$, je jednoznačne určený
množinami $D(f)$ a $D(g)$. Teda ak funkcia $h$ vznikne z funkcií
$f_1,\ldots,f_n$ len použitím operácií súčtu, rozdielu, súčinu, podielu a
superpozície funkcií, je množina $D(h)$ jednoznačne určená množinami
$D(f_1),\ldots,D(f_n)$. Preto ak napíšeme predpis takejto funkcie $h$ bez toho,
aby sme výslovne určili jej definičný obor, považujeme funkciu $h$ za
definovanú práve na tej množine, ktorá je určená množinami
$D(f_1),\ldots,D(f_n)$ na základe definícií súčtu, rozdielu, súčinu, podielu a
superpozície funkcií. (Teda trocha nepresne povedané, za definičný obor takejto
funkcie $h$ považujeme množinu všetkých tých $x\in\mathbb{R}$, pre ktoré má
predpis funkcie $h$ \enquote{zmysel}.)

\begin{enumerate}[resume]
  \item \useproblem[funkcie]{funkcie-22}
  \item \useproblem[funkcie]{funkcie-23}
  \item \useproblem[funkcie]{funkcie-24}
  \item \useproblem[funkcie]{funkcie-25}
\end{enumerate}

Ak $f: A \rightarrow \mathbb{R}, g: B \rightarrow \mathbb{R}$ sú funkcie, $A
\subset B$ a pre všetky $x \in A$ platí $f(x) = g(x)$, hovoríme, že funkcia $f$
je zúženie funkcie $g$ na množinu $A$ ($f$ je funkcia $g$ zúžená na množinu
$A$) a označujeme $f = \frac{g}{A}$.

% má tam byť číslovanie v príklade
\begin{enumerate}[resume]
  \item \useproblem[funkcie]{funkcie-26}
\end{enumerate}


\section{Graf funkcie}
Nech je v rovine daná pravouhlá súradnicová sústava, pričom jednotky dĺžky na
súradnicových osiach $Ox$ a $Oy$ sú rovnaké. Množina
\[
  \{(x,f(x)); x \in A\}
\]
bodov roviny, kde $f:A \rightarrow \mathbb{R}$ je daná funkcia, sa nazýva
\emph{graf funkcie $f$}, $((x,f(x))$ je zápis bodu roviny pomocou jeho súradníc
v danej súradnicovej sústave).

V nasledujúcej tabuľke sú opísané elementárne transformácie grafov funkcií:

\begin{center}
  \begin{tabular}{|p{0.2\textwidth}|p{0.6\textwidth}|}
    \hline
    {\bf funkcia $y=g(x)$} & {\bf transformácia grafu funkcie $y=f(x)$}  \\
    \hline
    $y = f(x) + c$     & posunutie o $c$ v smere osi $Oy$ \\
    $y = f(x - c)$     & posunutie o $c$ v smere osi $Ox$ \\
    $y = f(-x)$        & symetria podľa osi $Oy$ \\
    $y = -f(x)$        & symetria podľa osi $Ox$ \\
    $y = a \cdot f(x)$ & vynásobenie každej $y$-ovej súradnice číslom $a$ \\
    $y = f(ax)$        & vydelenie každej $x$-ovej súradnice číslom $a$, $(a \neq 0)$ \\
    \hline
  \end{tabular}
\end{center}

\begin{enumerate}[resume]
  \showanswers
  \item \useproblem[funkcie]{funkcie-27}
  \hideanswers
  \item \useproblem[funkcie]{funkcie-28}
  \item \useproblem[funkcie]{funkcie-29}
  \showanswers
  \item \useproblem[funkcie]{funkcie-30}
  \item \useproblem[funkcie]{funkcie-31}
  \hideanswers
  \item \useproblem[funkcie]{funkcie-32}
  \item \useproblem[funkcie]{funkcie-33}
  \item \useproblem[funkcie]{funkcie-34}
\end{enumerate}


\section{Niektoré vlastnosti funkcií}
Nech $A,B\subset\mathbb{R}$ sú neprázdne množiny, $B\subset A$. Funkcia
$f:A\rightarrow\mathbb{R}$ sa nazýva ohraničená na množine $B$, ak je ohraničená
množina $f(B)$. Funkcia ohraničená na svojom definičnom obore sa nazýva
ohraničená. Analogicky sa zavádza pojem funkcie ohraničenej zhora, resp. zdola a
funkcie ohraničenej zhora, resp. zdola na množine $B$.

Ak je množina $f(B)$ zhora (zdola) ohraničená, tak jej supremum (infimum) sa
nazýva supremum (infimum) funkcie $f$ na množine $B$ a označuje sa
$\sup\limits_{x\in B} f(x)$ ($\inf\limits_{x\in B} f(x)$). Ak existuje maximum
(minimum) množiny $f(B)$, nazýva sa toto číslo maximum (minimum) funkcie $f$ na
množine $B$ (často sa používa názov globálne maximum (globálne minimum) funkcie
$f$ na množine $B$) a označuje sa $\max\limits_{x\in B}f(x)$ ($\min\limits_{x\in
B}f(x)$).

%Ak je množina $f(B)$ zdola ohraničená, tak jej infimum sa nazýva infimum funkcie
%$f$ na množine $B$ a označuje sa $\inf_{x\in B} f(x)$. Ak existuje minimum
%množiny $f(B)$, nazýva sa toto číslo minimum funkcie $f$ na množine $B$ (často
%sa používa názov globálne minimum funkcie $f$ na množine $B$) a označuje sa
%$\min_{x\in B}f(x)$.

%Ak je množina $f(B)$ zhora ohraničená, tak jej supremum sa nazýva supremum
%funkcie $f$ na množine $B$ a označuje sa $\sup_{x\in B} f(x)$. Ak existuje
%maximum množiny $f(B)$, nazýva sa toto číslo maximum funkcie $f$ na množine $B$
%(často sa používa názov globálne maximum funkcie $f$ na množine $B$) a označuje
%sa $\max_{x\in B}f(x)$.

\begin{enumerate}[resume]
  \item \useproblem[funkcie]{funkcie-35}
  \item \useproblem[funkcie]{funkcie-36}
  \item \useproblem[funkcie]{funkcie-37}
  \item \useproblem[funkcie]{funkcie-38}
  \item \useproblem[funkcie]{funkcie-39}
\end{enumerate}

Na $A,B\subset\mathbb{R}$ sú neprázdne množiny, $B\subset A$. Funkcia
$f:A\rightarrow\mathbb{R}$ sa nazýva rastúca na množine $B$, resp. neklesajúca
na množine $B$, ak platí
\[
  (\forall x,y \in B):
    x < y \Rightarrow f(x) < f(y)
\]
resp.
\[
  (\forall x,y \in B):
    x < y \Rightarrow f(x) \leq f(y)
\]

Ak má funkcia $f$ niektorú z uvedených štyroch vlastností, nazýva sa monotónna
na množine $B$; funkcia, ktorá je rastúca na množine $B$ alebo klesajúca na
množine $B$, sa nazýva rýdzomonotónna na množine $B$.

Funkcia rastúca (klesajúca, nerastúca, neklesajúca, monotónna, rýdzomonotónna)
na svojom definičnom obore sa nazýva rastúca (klesajúca, nerastúca, neklesajúca,
monotónna, rýdzomonotónna).

\begin{enumerate}[resume]
  \item \useproblem[funkcie]{funkcie-40}
  \item \useproblem[funkcie]{funkcie-41}
  \item \useproblem[funkcie]{funkcie-42}
  \showanswers
  \item \useproblem[funkcie]{funkcie-43}
  \hideanswers
  \item \useproblem[funkcie]{funkcie-44}
  \item \useproblem[funkcie]{funkcie-45}
  \item \useproblem[funkcie]{funkcie-46}
  \item \useproblem[funkcie]{funkcie-47}
\end{enumerate}

Funkcia $f$ sa nazýva periodická, ak existuje číslo $T>0$ tak, že platí
\begin{itemize}
\item
  $
    (\forall a\in D(f)):
      d(f)\cap \interval{a+T}{a+2T}=\{x+T;x\in D(f)\cap \interval{a}{a + T}\}
  $
\item  $(\forall a \in D(f)):f(a+T)=f(a)$
\end{itemize}

Každé číslo $T$ s uvedenými vlastnosťami sa nazýva perióda funkcie $f$; ak
existuje najmenšie také $T>0$, nazýva sa najmenšia perióda funkcie $f$ (možno sa
stretnúť aj s terminológiou, v ktorej sa pojem perióda používa výlučne v zmysle
tu zavedeného pojmu najmenšej periódy).

\begin{enumerate}[resume]
  \item \useproblem[funkcie]{funkcie-48}
  \item \useproblem[funkcie]{funkcie-49}
  \item \useproblem[funkcie]{funkcie-50}
  \item \useproblem[funkcie]{funkcie-51}
  \item \useproblem[funkcie]{funkcie-52}
\end{enumerate}

Funkcia $f:A \rightarrow\mathbb{R}$ sa nazýva prostá (injektívna, jednoznačná),
ak platí:
\[
  (\forall x,y\in A):
    x \neq y \Rightarrow f(x) \neq f(y)
\]

Ak je funkcia $f:A \rightarrow \mathbb{R}$ prostá, tak funkcia s definičným
oborom $f(A)$, ktorý každému číslu $x\in f(A)$ priradí to číslo $y\in A$, pre
ktoré $f(y)=x$, sa nazýva inverzná funkcia k funkcii $f$ a označuje sa $f^{-1}$.

\begin{enumerate}[resume]
  \showanswers
    \item \useproblem[funkcie]{funkcie-53}
  \hideanswers
  \item \useproblem[funkcie]{funkcie-54}
  \item \useproblem[funkcie]{funkcie-55}
  \item \useproblem[funkcie]{funkcie-56}
\end{enumerate}


\section{Cyklometrické a hyperbolické funkcie}
Inverzné funkcie k funkciám
$\sin / \interval{-\frac{\pi}{2}}{\frac{\pi}{2}}$,
$\cos / \interval{0}{\pi}$,
$\tan / \interval[open]{-\frac{\pi}{2}}{\frac{\pi}{2})}$,
$\cot / \interval[open]{0}{\pi}$
sa nazývajú arkussínus, arkuskosínus, arkustangens a arkuskotangens a označujú
sa $\arcsin, \arccos, \arctan, \arccot$. Tieto funkcie majú spoločný názov
cyklometrické.

\begin{enumerate}[resume]
  \item \useproblem[funkcie]{funkcie-57}
  \item \useproblem[funkcie]{funkcie-58}
  \item \useproblem[funkcie]{funkcie-59}
  \showanswers
  \item \useproblem[funkcie]{funkcie-60}
  \item \useproblem[funkcie]{funkcie-61}
  \hideanswers
  \item \useproblem[funkcie]{funkcie-62}
\end{enumerate}

Funkcie definované predpismi
$\sinh x=\frac{e^x-e^{-x}}{2}$ (sínus hyberbolický),
$\cosh x=\frac{e^x+e^{-x}}{2}$ (kosínus hyperbolický),
$\tanh x=\frac{\sinh x}{\cosh x}$ (tangens hyperbolický) a
$\coth x=\frac{\cosh x}{\sinh x}$ (kotangens hyperbolický) sa
nazývajú hyperbolické funkcie.

\begin{enumerate}[resume]
  \item \useproblem[funkcie]{funkcie-63}
  \item \useproblem[funkcie]{funkcie-64}
\end{enumerate}


\section{Ďalšie príklady}
\begin{enumerate}[resume]
  \item \useproblem[funkcie]{funkcie-65}
  \item \useproblem[funkcie]{funkcie-66}
  \item \useproblem[funkcie]{funkcie-67}
  \item \useproblem[funkcie]{funkcie-68}
  \item \useproblem[funkcie]{funkcie-69}
  \item \useproblem[funkcie]{funkcie-70}
  \item \useproblem[funkcie]{funkcie-71}
  \item \useproblem[funkcie]{funkcie-72}
  \item \useproblem[funkcie]{funkcie-73}
  \item \useproblem[funkcie]{funkcie-74}
  \item \useproblem[funkcie]{funkcie-75}
  \item \useproblem[funkcie]{funkcie-76}
  \item \useproblem[funkcie]{funkcie-77}
  \item \useproblem[funkcie]{funkcie-78}
  \item \useproblem[funkcie]{funkcie-79}
  \item \useproblem[funkcie]{funkcie-80}
  \item \useproblem[funkcie]{funkcie-81}
  \item \useproblem[funkcie]{funkcie-82}
  \item \useproblem[funkcie]{funkcie-83}
  \item \useproblem[funkcie]{funkcie-84}
  \item \useproblem[funkcie]{funkcie-85}
  \item \useproblem[funkcie]{funkcie-86}
  \item \useproblem[funkcie]{funkcie-87}
  \item \useproblem[funkcie]{funkcie-88}
  \item \useproblem[funkcie]{funkcie-89}
\end{enumerate}



\chapter{Limita funkcie}%\label{chapter:gramatiky}

\loadallproblems[limita-funkcie]{./tex/03-limita-funkcie/ulohy.tex}

\section{Okolia a hromadné body}
Nech $a \in \mathbb{R};$ každý interval
$\interval[open]{a - \varepsilon}{a + \varepsilon}$,
kde $\varepsilon > 0$, sa nazýva \textit{okolie bodu $a$}. Číslo $\varepsilon$
sa nazýva \textit{polomer okolia}
$\interval[open]{a - \varepsilon}{a + \varepsilon};$

ak chceme zdôrazniť, že dané okolie bodu $a$ má polomer $\varepsilon$, hovoríme
o $\varepsilon-$okolí bodu $a$. Okolím bodu $\infty$ sa nazýva každý interval
$\interval[open]{K}{\infty}$, kde $K \in \mathbb{R};$ okolím bodu $- \infty$ je
každý interval $(- \infty,K)$, kde $K \in \mathbb{R}.$ Okolie bodu $b \in
\mathbb{R^*}$ ($\mathbb{R^*}$ sa nazýva rozšírená množina reálnych čísel a
pozostáva zo všetkých reálnych čísel a symbolov
$\interval[open]{+\infty}{-\infty}$) budeme označovať $O(b)$, symbol
$O_\varepsilon (b)$ budeme používať pre $\varepsilon-$okolie bodu $b \in
\mathbb{R}$.

\textit{Poznámka:} Namiesto symbolu $+\infty$ sa často používa symbol $\infty$,
niektorí autori však zavádzajú symbol $\infty$ s iným významom, k tomu pozri
poznámku na konci odseku $2.5$

Bod $a \in \mathbb{R^*}$ sa nazýva hromadný bod množiny $M \subset \mathbb{R}$,
ak každé jeho okolie $O(a)$ obsahuje aspoň jeden prvok množiny $M$, rôzny od
bodu $a$, t.j. ak platí
\[
  (\forall O(a)): (O(a) \setminus \{ a \}) \cap M \neq \emptyset
\]
(tento výrok možno zapísať aj v tvare $(\forall O(a))(\exists x\in M); x \neq a
\wedge x \in O(a);$ je ekvivalentný s výrokom: každé okolie bodu $a$ obsahuje
nekonečne veľa prvkov množiny $M$).

Množinu všetkých hromadných bodov množiny $M$ budeme označovať $M'$.
\begin{enumerate}[resume]
  \item \useproblem[limita-funkcie]{limita-100}
  \item \useproblem[limita-funkcie]{limita-101}
  \item \useproblem[limita-funkcie]{limita-102}
  \item \useproblem[limita-funkcie]{limita-103}
  \item \useproblem[limita-funkcie]{limita-104}
\end{enumerate}


\section{Definícia limity}
Nech $a \in \mathbb{R^*}$ je hromadný bod definičného oboru funkcie $f$. Bod $b
\in \mathbb{R^*}$ sa nazýva \textit{limita funkcie $f$ v bode $a$} ak platí
\[
  (\forall O(b)) (\exists O(a)) (\forall x \in (O(a)\setminus \{ a \}) \cap D(f) ):
  f(x) \in O(b)
\]
(Ak $b \in \mathbb{R}$, hovoríme o \textit{vlastnej (alebo konečnej) limite}; ak
$b=\infty$ alebo $b=-\infty$, o \textit{nevlastnej limite}; ak $a=\infty$ alebo
$a=-\infty$, používane názov \textit{limita v nevlastnom bode}.) Zapisujeme
$\lim_{x \rightarrow a} f(x)=b$ alebo $f(x) \rightarrow b$ pre $x \rightarrow
a$.

Všimnime si teraz jednotlivé prípady, ktoré zahŕňa uvedená definícia; začneme
postupnosťami:
\begin{itemize}
\item
  Nech $b \in \mathbb{R}$; hovoríme, že postupnosť ${\{a_n\}}_{n=1}^\infty$
  konverguje k (číslu) $b$, ak $\lim\limits_{n \rightarrow \infty} a_n=b$, t.j.
  ak platí
  \[
    (\forall \varepsilon > 0)
      (\exists n_0 \in \mathbb{N})
        (\forall n \in \mathbb{N}) (n > n_0): |a_n-b| < \varepsilon
  \]

  \textit{Poznámka:} Výroky
  \[
    (\forall \varepsilon > 0)
      (\exists n_0 \in \mathbb{N})
        (\forall n \in \mathbb{N})
          (n > n_0): |a_N - b| < \varepsilon
  \]
  a
  \[
    (\forall \varepsilon > 0)
      (\exists n_0 \in \mathbb{R})
        (\forall n \in \mathbb{N})
          (n > n_0): |a_N - b| < \varepsilon
  \]
  sú ekvivalentné. Pretože v definícii konečnej limity postupností je zvykom
  žiadať $n_0 \in \mathbb{N}$ (a nie $n_0 \in \mathbb{R}$), budeme ju v takej
  podobe používať aj my (hoci - ako uvidíme v príklade $105$ - sa tým niekedy
  komplikuje vyjadrenie závislosti čísla $n_0$ na čísle $\varepsilon$).
  Analogická poznámka sa vzťahuje aj na definície nevlastných limít postupností.
\item
  Hovoríme, že postupnosť ${\{a_n\}}_{n=1}^\infty$ diverguje k $+\infty$, ak
  $\lim\limits_{n \rightarrow \infty} a_N = +\infty$, t.j. ak platí
  \[
    (\forall K \in	\mathbb{R})
      (\exists n_0 \in \mathbb{N})
        (\forall n \in \mathbb{N})
        (n > n_0): a_n>K
  \]
\item
  Hovoríme, že postupnosť ${\{a_n\}}_{n=1}^\infty$ diverguje k $-\infty$, ak
  $\lim\limits_{n \rightarrow \infty} a_N = -\infty$, t.j. ak platí
  \[
    (\forall K \in \mathbb{R})
      (\exists n_0 \in \mathbb{N})
        (\forall n \in \mathbb{N})
          (n > n_0): a_n < K
  \]
  Postupnosti, ktoré majú vlastnú limitu, sa nazývajú \textit{konvergentné}; ak
  postupnosť nemá limitu alebo má nevlastnú limitu, nazýva sa
  \textit{divergentná}.
\end{itemize}

\begin{enumerate}[resume]
  \item \useproblem[limita-funkcie]{limita-105}
  \item \useproblem[limita-funkcie]{limita-106}
  \item \useproblem[limita-funkcie]{limita-107}
  \item \useproblem[limita-funkcie]{limita-108}
  \item \useproblem[limita-funkcie]{limita-109}
  \item \useproblem[limita-funkcie]{limita-110}
\end{enumerate}

Ak $x \in \mathbb{R^*}$, nastáva práve jedna z troch možností:
$x \in \mathbb{R}$, $x = +\infty$, $x = -\infty$.

Ak v definícii limity funkcie rozlíšime pre body $a, b \in \mathbb{R^*}$ tieto
možnosti, dostaneme nasledujúcich deväť špeciálnych prípadov (v nich už $a, b$
označujú len reálne čísla):

\begin{multicols}{2}
\begin{enumerate}
    \item $\lim\limits_{x \rightarrow a} f(x) = b$
    \item $\lim\limits_{x \rightarrow a} f(x) = +\infty$
    \item $\lim\limits_{x \rightarrow a} f(x) = -\infty$
    \item $\lim\limits_{x \rightarrow \infty} f(x) = b$
    \item $\lim\limits_{x \rightarrow \infty} f(x) = +\infty$
    \item $\lim\limits_{x \rightarrow \infty} f(x) = -\infty$
    \item $\lim\limits_{x \rightarrow -\infty} f(x) = b$
    \item $\lim\limits_{x \rightarrow -\infty} f(x) = +\infty$
    \item $\lim\limits_{x \rightarrow -\infty} f(x) = -\infty$
\end{enumerate}
\end{multicols}

\begin{enumerate}[resume]
  \item \useproblem[limita-funkcie]{limita-111}
\end{enumerate}

\begin{veta}[Cauchyho-Bolzanovo kritérium konvergencie]
Nech $a \in \mathbb{R}^*$ je hromadný bod definičného oboru funkcie $f$. Funkcia
$f$ má v bode $a$ konečnú limitu práve vtedy, keď platí:
\[
  (\forall \varepsilon > 0)
    (\exists O(a))
      (\forall x,y \in (O(a) \setminus \{ a \}) \cap D(f))
        : |f(x)-f(y)| < \varepsilon
\]
\end{veta}

\begin{enumerate}[resume]
  \item \useproblem[limita-funkcie]{limita-112}
  \item \useproblem[limita-funkcie]{limita-113}
  \item \useproblem[limita-funkcie]{limita-114}
  \item \useproblem[limita-funkcie]{limita-115}
\end{enumerate}


\section{Vety o limitách I}
\begin{veta}
\textbf{o limite skalárneho násobku, súčtu, rozdielu, súčinu a podielu}

Nech sú dané funkcie $f,g$, nech $a \in \mathbb{R^*}$ je hromadný bod množiny
$D(f) \cap D(g)$. Ak existujú konečné $\lim\limits_{x \rightarrow a} f(x)=A$,
$\lim\limits_{x \rightarrow a} g(x)=B$, tak existuje aj
$\lim\limits_{x \rightarrow a} cf(x)$ ($c \in \mathbb{R}$ je konštanta),
$\lim\limits_{x \rightarrow a} (f(x)+g(x)),
\lim\limits_{x \rightarrow a}(f(x)-g(x)),
\lim\limits_{x \rightarrow a} (f(x) \cdot g(x))$
a platí
\begin{multicols}{2}
\begin{enumerate}[label=]
    \item $\lim\limits_{x \rightarrow a} c \cdot f(x) = cA$
    \item $\lim\limits_{x \rightarrow a} (f(x)+g(x)) = A+B$
    \item $\lim\limits_{x \rightarrow a} (f(x)-g(x)) = A-B$
    \item $\lim\limits_{x \rightarrow a} (f(x) \cdot g(x)) = A \cdot B$
    \item $( = c \cdot \lim\limits_{x \rightarrow a} f(x))$
    \item $( = \lim\limits_{x \rightarrow a} f(x)+\lim\limits_{x \rightarrow a} g(x))$
    \item $( = \lim\limits_{x \rightarrow a} f(x)-\lim\limits_{x \rightarrow a} g(x))$
    \item $( = \lim\limits_{x \rightarrow a} f(x) \cdot \lim\limits_{x \rightarrow a} g(x))$
\end{enumerate}
\end{multicols}
Ak naviac $B \neq 0$, tak existuje aj $\lim\limits_{x \rightarrow a} \frac{f(x)}{g(x)}$
a platí
\begin{multicols}{2}
\begin{enumerate}[label=]
\item $\lim\limits_{x \rightarrow a} \frac{f(x)}{g(x)}=\frac{A}{B}$
\item $(=\frac{\lim\limits_{x \rightarrow a} f(x)}{\lim\limits_{x \rightarrow a} g(x)})$.
\end{enumerate}
\end{multicols}
\end{veta}

\begin{veta} \textbf{o limite zloženej funkcie}
Nech sú dané funkcie $f, g$ nech $a \in \mathbb{R^*}$ je hromadný bod množiny
$D(f \circ g)$. Ak $\lim\limits_{x \rightarrow a} g(x) = A$ $(A \in
\mathbb{R^*})$, pričom je splnená podmienka
\[
  (\exists O(a))
    (\forall x \in O(a)):
      x \neq a \Rightarrow g(x) \neq A
\]
a
\[
  \lim\limits_{x \rightarrow A} f(x) = B \quad (B \in \mathbb{R^*})
\]
tak
\[
  \lim\limits_{x \rightarrow a} f(g(x)) = B
\]
\end{veta}

\textit{Poznámka:}
Ak $A$ je hromadným bodom $D(f)$, ale $A \notin D(f)$, tak uvedená veta platí aj
vtedy, keď nie je splnená podmienka (*).

Niektoré limity možno nájsť len na základe definície, v ostatných prípadoch je
však oveľa efektívnejšie použiť vety o limitách. Pritom je potrebné osvojiť si
zdôvodňovanie jednotlivých krokov výpočtu, inak sa nenaučíme odlišovať správne
postupy od nesprávnych. Na ilustráciu podrobne popíšme nasledujúci výpočet

\[
  \begin{split}
    \lim_{x \rightarrow \infty} \frac{2x^2 + 3x + 5}{3x^2 - 7}
      &= \lim_{x \rightarrow \infty}
          \frac{
            2
            + \frac{3}{x}
            + \frac{5}{x^2}}{3 - \frac{7}{x^2}}
      = \frac{\lim\limits_{x \rightarrow \infty}
        (
          2
          + \frac{3}{x}
          + \frac{5}{x^2}
        )}{
          \lim\limits_{x \rightarrow \infty} (3 - \frac{7}{x^2})
        } \\
      &= \frac{
          \lim\limits_{x \rightarrow \infty}
          2
          + \lim\limits_{x \rightarrow \infty} \frac{3}{x}
          + \lim\limits_{x \rightarrow \infty} \frac{5}{x^2}
        }{
          \lim\limits_{x \rightarrow \infty} 3
          -
          \lim\limits_{x \rightarrow \infty}
            \frac{7}{x^2}
        } \\
      &= \frac{2+0+0}{3-0}
      = \frac{2}{3}.
  \end{split}
\]

\textbf{1. krok:}
Na intervale $\interval[open]{0}{\infty}$ iste platí
\[
  \frac{2x^2 + 3x + 5}{3x^2 - 7}
  = \frac{2 + \frac{3}{x} + \frac{5}{x^2}}{3 - \frac{7}{x^2}}
\]
(zlomok vľavo stačí rozšíriť výrazom $\frac{1}{x^2}$), preto: ak existuje limita
na pravej strane prvej rovnosti, tak existuje aj limita na jej ľavej strane a
tieto limity sa rovnajú.

\textit{Poznámka}: Táto elementárna, ale veľmi častá úvaha sa vo všeobecnosti
formuluje takto: Nech $a \in \mathbb{R^*}$ je hromadný bod $D(f)$, nech existuje
$O(a)$ tak, že:
\[
  D(f) \cap (O(a) \setminus \{ a\})=D(g) \cap (O(a) \setminus \{ a\})
\]
a $(\forall x \in D(f) \cap (O(a) \setminus \{ a\})$ platí $f(x)=g(x)$
Ak existuje $\lim\limits_{x \rightarrow a} g(x)=b$ $(\in \mathbb{R^*})$, tak
platí $\lim\limits_{x \rightarrow a} f(x)=b$.

Ďalej sa teda snažíme zistiť, či
existuje
\[
  \lim\limits_{x \rightarrow \infty}
    \frac{2+\frac{3}{x} + \frac{5}{x^2}}{3-\frac{7}{x^2}}
\]

\textbf{2. krok:}
tu sme použili vetu o limite podielu (zatiaľ ovšem len \enquote{na čestné
slovo}); presnejšie povedané: ak ukážeme, že limita v čitateli aj v menovateli
existujú a sú konečné, pričom limita v menovateli je nenulová, tak podľa vety o
limite podielu bude platiť druhá rovnosť;

\textbf{3. krok:}
v čitateli sme použili vetu o limite súčtu (tú možno indukciou rozšítiť na
ľubovoľný konečný počet sčítancov), v menovateli vetu o limite rozdielu (obidve
zatiaľ tiež len \enquote{na čestné slovo});

\textbf{4. krok:}
teraz už ľahko overíme, že druhá a tretia rovnosť skutočne platia, pretože
$\lim\limits_{x \rightarrow \infty} 2$,
$\lim\limits_{x \rightarrow \infty} \frac{3}{x}$,
$\lim\limits_{x \rightarrow \infty} \frac{5}{x^2}$,
$\lim\limits_{x \rightarrow \infty} 3$
a
$\lim\limits_{x \rightarrow \infty} \frac{7}{x^2}$
existujú a sú konečné (to ľahko dokážeme priamo z definície), bolo použitie viet
o limite súčtu a rozdielu v tretej rovnosti oprávnené (a preto $\lim\limits_{x
\rightarrow \infty} (2 + \frac{3}{x} + \frac{5}{x^2}) = 2$, $\lim\limits_{x
\rightarrow \infty} (3 - \frac{7}{x^2}) = 3$); rovnako oprávnené bolo použitie
vety o limite podielu v druhom kroku (limita v čitateli aj v menovateli - ako
sme sa práve presvedčili - skutočne existujú, sú konečné a limita v menovateli
je naviac nenulová).

Z platnosti prvej, druhej a tretej rovnosti vyplýva $\lim\limits_{x \rightarrow
\infty} \frac{2x^2 + 3x + 5}{3x^2 - 7} = \frac{2}{3}$.

Takéto zdôvodnenie (vykonané ovšem len v duchu alebo ústne) by malo byť súčasťou
výpočtu každej limity; po získaní istej prace budú zápisy aj argumentácia
podstatne stručnejšie

\begin{veta}
Ak $f$ je elementárna funkcia a bod $a \in D(f)$ je hromadný bod množiny $D(f)$
je hromadný bod množiny $D(f)$, tak $\lim\limits_{x \rightarrow a} f(x) = f(a)$.
(Toto tvrdenie veľmi úzko súvisí s pojmom spojitosti (pozri kap. 3).)
\end{veta}

\begin{enumerate}[resume]
  \item \useproblem[limita-funkcie]{limita-116}
  \item \useproblem[limita-funkcie]{limita-117}
  \item \useproblem[limita-funkcie]{limita-118}
  \item \useproblem[limita-funkcie]{limita-119}
  \item \useproblem[limita-funkcie]{limita-120}
  \item \useproblem[limita-funkcie]{limita-121}
  \item \useproblem[limita-funkcie]{limita-122}
  \item \useproblem[limita-funkcie]{limita-123}
  \item \useproblem[limita-funkcie]{limita-124}
  \item \useproblem[limita-funkcie]{limita-125}
\end{enumerate}

\begin{veta}
$\lim\limits_{x \rightarrow 0} \frac{\sin{x}}{x}=1$.
\end{veta}

\begin{enumerate}[resume]
  \item \useproblem[limita-funkcie]{limita-126}
  \item \useproblem[limita-funkcie]{limita-127}
  \item \useproblem[limita-funkcie]{limita-128}
  \item \useproblem[limita-funkcie]{limita-129}
  \item \useproblem[limita-funkcie]{limita-130}
\end{enumerate}


\section[Vety o limitách II]{Vety o limitách II (porovnávacie vety, vety o nevlastných limitách, jednostranné limity)}
\begin{definicia}
Prstencovým okolím bodu $a$  sa nazýva množina $O(a) \setminus \{ a\}$, kde
$O(a)$ je okolie bodu $a$. (Teda každé okolie bodov $+\infty,-\infty$ je súčasne
aj ich prstencovým okolím.). Označujeme znakom $\mathcal{P}(a)$.
\end{definicia}

\begin{veta}
Nech sú dané sunkcie $f, g, h$ nech $a \in \mathbb{R}^*$ je hromadný bod množiny
$D(g)$ a nech pre niektoré jeho prstencové okolie $O^*(a)$ platí $O^*(a) \cap
D(f)=O^*(a)\cap D(g)=O^*(a) \cap D(h)=D$. Ak $\lim\limits_{x \rightarrow a}
f(x)=\lim\limits_{x \rightarrow a} h(x)=b$ $( \in \mathbb{R})$ a pre všetky $x
\in D$ je $f(x) \leq g(x) \leq h(x)$, tak existuje aj $\lim\limits_{x
\rightarrow a} g(x)$ a rovná sa $b$.
\end{veta}

\begin{veta}
Nech sú dané funkcie $f, g$, nech $a \in \mathbb{R}^*$ je hromadný bod množiny
$D(f) \cap D(g)$. Ak $\lim\limits_{x \rightarrow a} f(x) = 0$ a funkcia $g$ je
ohraničená v niektorom prstencovom okolí bodu $a$ (t.j. na niektorej z množín
$(O(a) \setminus \{ a\})\cap D(g)$, tak $\lim\limits_{x \rightarrow a} f(x)
\cdot g(x)=0$.
\end{veta}

\begin{enumerate}[resume]
  \item \useproblem[limita-funkcie]{limita-131}
  \item \useproblem[limita-funkcie]{limita-132}
  \item \useproblem[limita-funkcie]{limita-133}
  \item \useproblem[limita-funkcie]{limita-134}
  \item \useproblem[limita-funkcie]{limita-135}
  \item \useproblem[limita-funkcie]{limita-136}
\end{enumerate}

\begin{veta}
Nech sú dané funkcie $f,g$, nech $a \in \mathbb{R^*}$ je hromadný bod množiny
$D(f)$ a nech pre niektoré jeho rýche okolie $O^*(a)$ a platí $O^*(a)\cap
D(g)=:D$. Ak $\lim\limits_{x \rightarrow a} g(x)=\infty$  $(-\infty)$ a pre
všetky $x \in D$ platí $f(x)\geq g(x)$  $(f(x)\leq g(x))$ tak existuje aj
$\lim\limits_{x \rightarrow a} f(x)$ a rovná sa $\infty$  $(-\infty)$.
\end{veta}

\begin{veta}
Nech sú dané funkcie $f,g$, nech $a \in \mathbb{R^*}$ je hromadný bod množiny
$D(f)\cap D(g)$, nech existuje $\lim\limits_{x \rightarrow a} f(x)=:A$,
$\lim\limits_{x \rightarrow a} g(x)=:B$. Potom
\begin{enumerate}
\item ak $A=+\infty$, $B \in \mathbb{R}$ alebo $B=+\infty$, tak $\lim\limits_{x \rightarrow a}(f(x)+g(x))=+\infty$;
\item ak $A=+\infty$, $B >0$ alebo $B=+\infty$, tak $\lim\limits_{x \rightarrow a}(f(x)\cdot g(x))=+\infty$.
\end{enumerate}
\end{veta}

\begin{veta}
Nech $a \in \mathbb{R^*}$ je hromadný bod definičného oboru funkcie $f$. Ak
$\lim\limits_{x \rightarrow a} f(x)=+\infty$   $(-\infty)$, tak
$\lim\limits_{\frac{1}{f(x)}=0}$.
\end{veta}

\textit{Poznámka:}
Predchádzajúce tvrdenie (spolu s ďalšími, analogickými) si ľahko zapamätáme
pomocou nasledujúcich rovností (ktoré ovšem považujeme len za mnemotechnickú
pomôcku):
\begin{multicols}{2}
\begin{enumerate}
    \item $A +/- \infty = +/- \infty$;
    \item $+/- \infty +/- \infty = +/- \infty$;
    \item $\frac{1}{+/- \infty} = 0$;
    \item $A \cdot (+/- \infty) = \left\{ \begin{array}{r@{\quad}c}
    +/- \infty, & $ak$ A>0 \\
    -/+ \infty, & $ak$ A<0 \\ \end{array} \right.
    $;
    \item $\infty \cdot(+/- \infty)=+/- \infty$.
\end{enumerate}
\end{multicols}
($A$ označuje reálne číslo.) Všimnime si, že žiadna z uvedených viet sa
nevzťahuje na limity funkcií typu $+\inf\ -\infty, 0 \cdot ( +/-
\infty),\frac{+/- \infty}{ +/- \infty}, \frac{0}{0}$. Také funkcie budeme
nazývať neurčitými výrazmi práve im je venovaná väčšina príkladov na výpočet
limít.

\begin{veta}
Nech je daná funkcia $f$, nech $a \in \mathbb{R^*}$ je hromadný bod množiny
$D:=\{ x \in D(f): f(x)\neq 0 \}$, nech $\lim\limits_{x \rightarrow a} f(x)=0$.
Ak existuje také prstencové okolie $O^*(a)$ bodu $a$, že pre všetky $x \in D
\cap O^*(a)$ platí $f(x)>0$  $(f(x)<0)$, tak existuje $\lim\limits_{x
\rightarrow a}\frac{1}{f(x)}$ a rovná sa $+\infty)(-\infty)$
\end{veta}

\begin{enumerate}[resume]
  \item \useproblem[limita-funkcie]{limita-137}
  \item \useproblem[limita-funkcie]{limita-138}
  \item \useproblem[limita-funkcie]{limita-139}
  \item \useproblem[limita-funkcie]{limita-140}
  \item \useproblem[limita-funkcie]{limita-141}
\end{enumerate}

Nech je daná funkcia $f$, nech $a \in \mathbb{R}$ je hromadný bod množiny
$D^+:=D(f)\cap (a,\infty)$ (množiny $D^-:=D(f)\cap (-\infty,a)$); označme $f'$
zúženie funkcie na množinu $D^+$ (na množinu $D^-$). Ak existuje limita funkcie
$f'$ v bode $a$, nazývame ju \textit{limitou funkcie $f$ v bode $a$ sprava
(zľava)} a označujeme $\lim\limits_{x \rightarrow a+} f(x)$ $(\lim\limits_{x
\rightarrow a-} f(x))$. Pre limity sprava a zľava sa používa súhrnný názov
\textit{jednostranné limity}.

\begin{veta}
Nech je daná funkcia $f$, nech $a$ je hromadný bod množín $D(f) \cap
(-\infty,a)$ a $D(f) \cap (a,\infty)$. Potom $\lim\limits_{x \rightarrow a}
f(x)=\lim\limits_{x \rightarrow -a} f(x)$; pritom $\lim\limits_{x \rightarrow a}
f(x)$ sa rovná spoločnej hodnote týchto jednostranných limít.
\end{veta}

\begin{enumerate}[resume]
  \item \useproblem[limita-funkcie]{limita-142}
  \item \useproblem[limita-funkcie]{limita-143}
  \item \useproblem[limita-funkcie]{limita-144}
  \item \useproblem[limita-funkcie]{limita-145}
  \item \useproblem[limita-funkcie]{limita-146}
\end{enumerate}


\section{Limity mocninovo-exponenciálnych funkcií}
\begin{veta}
$\lim\limits_{x \rightarrow +\infty} (1+\frac{1}{x})^x=e,\lim\limits_{x \rightarrow 0}
(1+x)^{\frac{1}{x}}=e$.
\end{veta}

\begin{enumerate}[resume]
  \item \useproblem[limita-funkcie]{limita-147}
  \item \useproblem[limita-funkcie]{limita-148}
  \item \useproblem[limita-funkcie]{limita-149}
  \item \useproblem[limita-funkcie]{limita-150}
  \item \useproblem[limita-funkcie]{limita-151}
  \item \useproblem[limita-funkcie]{limita-152}
  \item \useproblem[limita-funkcie]{limita-153}
  \item \useproblem[limita-funkcie]{limita-154}
\end{enumerate}


\section{Limity monotónnych postupností}
\begin{veta}
Každá zhora ohraničená neklesajúca postupnosť ${\{a_n\}}_{n=1}^\infty$ je
konvergentná a platí $\lim\limits_{n \rightarrow \infty} a_n=\sup {a_n;n \in
\mathbb{N}}$. Ak ${\{a_n\}}_{n=1}^\infty$ je zhora neohraničená neklesajúca
postupnosť, tak $\lim\limits_{n \rightarrow \infty} a_n=+\infty$. Analogicky tvrdenie
pre nerastúcu postupnosť ${\{a_n\}}_{n=1}^\infty$  dostaneme, ak predchádzajúcu
vetu aplikujeme na postupnosť ${\{-a_n\}}_{n=1}^\infty$.
\end{veta}

\begin{enumerate}[resume]
  \item \useproblem[limita-funkcie]{limita-155}
  \item \useproblem[limita-funkcie]{limita-156}
  \item \useproblem[limita-funkcie]{limita-157}
\end{enumerate}


\section{Heiného definícia limity}
\begin{veta}
Nech $a \in \mathbb{R^*}$ je hromadný bod definičného oboru $D(f)$ funkcie $f$a
$b \in \mathbb{R^*}$. Potom sú nasledujúce výroky ekvivalentné:
\begin{itemize}
\item $\lim\limits_{x \rightarrow a} f(x)=b$
\item pre každú postupnosť ${\{a_n\}}_{n=1}^\infty$ prvkov z $D(f)\setminus \{
a\}$, ktorej limitou je $a$, platí $\lim\limits_{n \rightarrow \infty} f(a_n)=b$
\end{itemize}
\end{veta}

Limitu funkcie možno definovať aj bez použitia pojmu okolia; v takomto prípade
sa zavedie len pojem vlastnej a nevlastnej limity postupnosti a na definíciu
limity funkcie $f$ v bode $a$ sa použije vlastnosť $2$ z uvedenej vety.
Definícia limity funkcie v takejto podobe sa nazýva \textit{Heineho definíciou
limity}. Veta $15$ teda hovorí, že definícia limity pomocou okolia a Hejného
definícia limity sú ekvivalentné.

\begin{enumerate}[resume]
  \item \useproblem[limita-funkcie]{limita-158}
  \item \useproblem[limita-funkcie]{limita-159}
\end{enumerate}

\section[Hromadné hodnoty]{Hromadné hodnoty, limes inferior a limes superior postupností}
Postupnosť ${\{a_n\}}_{n=1}^\infty$ sa nazýva \textit{podpostupnosť (vybranou
postupnosťou z) postupnosti ${\{a_n\}}_{n=1}^\infty$}, ak existuje rastúce
zobrazenie $k: \mathbb{N} \rightarrow \mathbb{N}$ také, že $b_n=a_{k(n)}$.

Bod $a \in \mathbb{R^*}$ sa nazýva \textit{hromadnou hodnotou postupnosti
${\{a_n\}}_{n=1}^\infty$}, ak je limitou niektorej jej podpostupnosti (to je
ekvivalentné s podmienkou: pre každé okolie $O(a)$ bodu $a$ je množina $\{n \in
\mathbb{N}: a_n \in O(a) \}$). Supremum (infimum) množiny $H$ v množine
$\mathbb{R^*}$ sa nazýva \textit{limes superior (limer inferior) postupnosti
${\{a_n\}}_{n=1}^\infty$} a označuje sa $\lim\limits_{n \rightarrow \infty}\sup
a_n$ $(\lim\limits_{n \rightarrow \infty} \inf a_n)$.

Suprémum (infimum) množiny $H$ v $\mathbb{R^*}$ sa definuje ako najmenšie
(najväčšie) z jej horných (dolných) ohraničení, pritom ako horné a dolné
ohraničenia prichádzajú do úvahy aj body $+\infty,-\infty$. Teda napr.
supremom zhora neohraničenej množiny $H \subset \mathbb{R}$ je v množine
$\mathbb{R^*}$ bod $+\infty$; $\sup\{-\infty\}$=$-\infty$ a pod. Pre zhora
(zdola) ohraničenú množinu $H \subset \mathbb{R}$ je jej supremum (infimum) v
$\mathbb{R^*}$ zhodné s jej supremom (infimom).

\begin{veta}
Limes superior a limes ingerior postupnosti ${\{a_n\}}_{n=1}^\infty$ sú jej
hromadnými hodnotami.
\end{veta}

\begin{veta}
Nech je daná postupnosť ${\{a_n\}}_{n=1}^\infty$, nech $b \in \mathbb{R}$.
Rovnosť
\begin{align*}
  \lim\limits_{n \rightarrow \infty} \sup a_n =b
  && (\lim\limits_{n \rightarrow \infty} \inf a_n =b)
\end{align*}
platí práve vtedy, keď:
\begin{itemize}
\item $b$ je hromadná hodnota postupnosti ${\{a_n\}}_{n=1}^\infty$
\item
  pre každé $\varepsilon > 0$ je množina
  \begin{align*}
    \{ n \in \mathbb{N}: a_n \geq b +\varepsilon \}
    && (\{  n \in \mathbb{N}: a_n \leq b +\varepsilon \} )
  \end{align*}
  konečná.
\end{itemize}
\end{veta}

\begin{veta}
Postupnosť ${\{a_n\}}_{n=1}^\infty$ má limitu práve vtedy, keď:
\[
  \lim\limits_{n \rightarrow \infty} \sup a_n = \lim\limits_{n \rightarrow \infty} \inf a_n
\]
Limitou je pritom spoločná hodnota $\lim\limits_{n \rightarrow \infty} \sup a_n$
a $\lim\limits_{n \rightarrow \infty} \inf a_n$.
\end{veta}

\begin{enumerate}[resume]
  \item \useproblem[limita-funkcie]{limita-160}
  \item \useproblem[limita-funkcie]{limita-161}
  \item \useproblem[limita-funkcie]{limita-162}
  \item \useproblem[limita-funkcie]{limita-163}
  \item \useproblem[limita-funkcie]{limita-164}
  \item \useproblem[limita-funkcie]{limita-165}
  \par
  \needspace{2\baselineskip}
  \item \useproblem[limita-funkcie]{limita-166}
\end{enumerate}


\section{Ďalšie príklady}
\begin{enumerate}[resume]
  \item \useproblem[limita-funkcie]{limita-167}
  \item \useproblem[limita-funkcie]{limita-168}
  \item \useproblem[limita-funkcie]{limita-169}
  \item \useproblem[limita-funkcie]{limita-170}
  \item \useproblem[limita-funkcie]{limita-171}
  \item \useproblem[limita-funkcie]{limita-172}
  \item \useproblem[limita-funkcie]{limita-173}
  \item \useproblem[limita-funkcie]{limita-174}
  \item \useproblem[limita-funkcie]{limita-175}
  \item \useproblem[limita-funkcie]{limita-176}
  \item \useproblem[limita-funkcie]{limita-177}
  \item \useproblem[limita-funkcie]{limita-178}
  \item \useproblem[limita-funkcie]{limita-179}
  \item \useproblem[limita-funkcie]{limita-180}
  \item \useproblem[limita-funkcie]{limita-181}
  \item \useproblem[limita-funkcie]{limita-182}
  \item \useproblem[limita-funkcie]{limita-183}
  \item \useproblem[limita-funkcie]{limita-184}
  \item \useproblem[limita-funkcie]{limita-185}
  \item \useproblem[limita-funkcie]{limita-186}
  \item \useproblem[limita-funkcie]{limita-187}
  \item \useproblem[limita-funkcie]{limita-188}
  \item \useproblem[limita-funkcie]{limita-189}
  \item \useproblem[limita-funkcie]{limita-190}
  \item \useproblem[limita-funkcie]{limita-191}
  \item \useproblem[limita-funkcie]{limita-192}
  \item \useproblem[limita-funkcie]{limita-193}
  \item \useproblem[limita-funkcie]{limita-194}
  \item \useproblem[limita-funkcie]{limita-195}
  \item \useproblem[limita-funkcie]{limita-196}
  \item \useproblem[limita-funkcie]{limita-197}
  \item \useproblem[limita-funkcie]{limita-198}
  \item \useproblem[limita-funkcie]{limita-199}
  \item \useproblem[limita-funkcie]{limita-200}
  \item \useproblem[limita-funkcie]{limita-201}
  \item \useproblem[limita-funkcie]{limita-202}
  \item \useproblem[limita-funkcie]{limita-203}
  \item \useproblem[limita-funkcie]{limita-204}
  \item \useproblem[limita-funkcie]{limita-205}
  \item \useproblem[limita-funkcie]{limita-206}
  \item \useproblem[limita-funkcie]{limita-207}
\end{enumerate}



\chapter{Spojitosť funkcie}%\label{chapter:gramatiky}

\loadallproblems[spojitost-funkcie]{./tex/04-spojitost-funkcie/ulohy.tex}

\section{Definícia spojitosti v bode}
Hovoríme, že funkcia $f$ je spojitá v bode $a \in D(f)$, ak platí
\[
  (\forall \varepsilon > 0)
    (\exists \delta > 0)
      (\forall x \in D(f)):
        |x-a| < \delta \Rightarrow |f(x)-f(a)| < \varepsilon
\]
\textit{Poznámka:} Všimnime si, že bod $a$ nemusí byť hromadným bodom množiny
$D(f)$. Môžu teda nastať dva prípady:
\begin{itemize}
\item
  ak $a \in D(f)$ je hromadným bodom množiny $D(f)$, tak podmienka vyššie je
  ekvivalentná s podmienkou $\lim\limits_{x \rightarrow a} f(x)=f(a)$;
\item
  ak $a \in D(f)$ nie je hromadným bodom množiny $D(f)$ (t.j. existuje $\eta > 0
  $ tak, že $D(f) \cup \interval[open]{a-\eta}{a+\eta}=\{ a\}$; taký bod sa
  nazýva izolovaným bodom množiny $D(f)$), je podmienka vyššie splnená (stačí
  položiť $\delta=\eta$ pre ľubovoľné $\varepsilon > 0$). Teda funkcia $f$ je
  splnená v každom izolovanom bode množiny $D(f)$.
\end{itemize}

\begin{veta}
Ak funkcie $f,g,f-g,f \cdot g$ sú spojité v bode $a$. Ak naviac $g(a) \neq 0$,
tak aj funkcia $\frac{f}{g}$ je spojitá v bode $a$.
\end{veta}

\begin{veta}
Ak funkcia $f$ je spojitá v bode $a$, funkcia $g$ je spojitá v bode $f(a)$, tak
funkcia $g \circ f$ je spojitá v bode $a$. Hovoríme, že funkcia $f$ je spojitá,
ak je spojitá v každom bode svojho definičného oboru. Hovoríme, že funkcia $f$
je spojitá na množine $A$ ($\emptyset\neq a\subset D(f)$), ak je spojitá funkcia
$f \setminus A$.
\end{veta}

\begin{veta}
Každá základná elementárna funkcia je spojitá.
\end{veta}
(Z viet $1,2$ a $3$ vyplýva veta $4$ z kapitoly $2.$)
\begin{enumerate}[resume]
  \item \useproblem[spojitost-funkcie]{s-funkcie-221}
  \item \useproblem[spojitost-funkcie]{s-funkcie-222}
  \item \useproblem[spojitost-funkcie]{s-funkcie-223}
  \item \useproblem[spojitost-funkcie]{s-funkcie-224}
  \item \useproblem[spojitost-funkcie]{s-funkcie-225}
  \item \useproblem[spojitost-funkcie]{s-funkcie-226}
  \item \useproblem[spojitost-funkcie]{s-funkcie-227}
  \item \useproblem[spojitost-funkcie]{s-funkcie-228}
  \item \useproblem[spojitost-funkcie]{s-funkcie-229}
  \item \useproblem[spojitost-funkcie]{s-funkcie-230}
  \item \useproblem[spojitost-funkcie]{s-funkcie-231}
\end{enumerate}

Číslo $a \in \mathbb{R}$, ktoré je hromadným bodom definičného oboru $D(f)$
funkcie $f$, sa nazýva bodom nespojitosti funkcie $f$, ak je splnená niektorá z
nasledujúcich podmienok:
\begin{itemize}
\item $a \notin D(f)$;
\item neexistuje $\lim\limits_{x \rightarrow a}f(x)$;
\item $a \in D(f)$ a existuje $\lim\limits_{x \rightarrow a}f(x)$, ale neplatí
      $\lim\limits_{x \rightarrow a}f(x)=f(a)$.
\end{itemize}

Ak $a \in \mathbb{R}$ je bod nespojitosti funkcie $f$ a existuje konečná
$\lim\limits_{x \rightarrow a}f(x)$, nazýva sa bodom odstrániteľnej
nespojitosti. Ak neexistuje vlastná $\lim\limits_{x \rightarrow a}f(x)$ nazýva
sa $a$ bodom neodstrániteľnej nespojitosti.

Nech $a \in \mathbb{R}$ je bod nespojitosti funkcie $f$, pričom $a$ je hromadným
bodom množín $D(f) \cap (-\infty,a)$ aj $D(f)\cap (a,+\infty)$. Ak v bode $a$
existujú vlastné jednostranné limity a $\lim\limits_{x \rightarrow a-}f(x) \neq
\lim\limits_{x \rightarrow a+}f(x)$ nazýva sa $a$ bodom nespojitosti $1.$ druhu
(rozdiel $\lim\limits_{x \rightarrow] a+}f(x)-\lim\limits_{x \rightarrow
a-}f(x)$ sa potom nazýva skokom funkcie $f$ v bode $a$). Ak aspoň jedna z
jednostranných limít v bode $a$ neexistuje alebo je nevlastná, nayzva sa $a$
bodom nespojitosti $2.$ druhu.

\textit{Poznámka:}
Klasifikácia bodov nespojitosti nie je v matematickej literatúre jednotná; napr.
niekedy sa hromadné body množiny $D(f)$, ktoré nie sú prvkami $D(f)$, nepovažujú
za body nespojitosti funkcie $f$; v definícii bodu nespojitosti $1.$ druhu s
aniekedy nepovažuje splnenie podmienky $\lim\limits_{x \rightarrow a+}f(x) \neq
\lim\limits_{x \rightarrow a-}f(x)$ atď.

\begin{enumerate}[resume]
  \item \useproblem[spojitost-funkcie]{s-funkcie-232}
  \item \useproblem[spojitost-funkcie]{s-funkcie-233}
\end{enumerate}


\section{Vlastnosti spojitých funkcií}
Funkcia $f$ definovaná na intervale $I$ sa nazýva darbouxovská na $I$, ak pre
Každé $a,b \in I,a<b$ platí: na intervale $\interval{a}{b}$ nabudúce funkcia $f$
všetky hodnoty medzi $f(a)$ a $f(b)$; t.j. ak platí
\[
  \forall a, b \in I, a < b:
    f(\left[ a, b \right] ) \supset X
\]
\[
  X = \{y \in \mathbb{R}; \min \{f(a),f(b)\} \leq y \leq \max \{f(a),f(b)\}\}
\]

\begin{veta}
Ak je funkcia $f$ spojitá na intervale $I$, tak je darbouxovská na $I$.
\end{veta}

\begin{enumerate}[resume]
  \item \useproblem[spojitost-funkcie]{s-funkcie-234}
  \item \useproblem[spojitost-funkcie]{s-funkcie-235}
  \item \useproblem[spojitost-funkcie]{s-funkcie-236}
  \item \useproblem[spojitost-funkcie]{s-funkcie-237}
  \item \useproblem[spojitost-funkcie]{s-funkcie-238}
  \item \useproblem[spojitost-funkcie]{s-funkcie-239}
\end{enumerate}

\begin{veta}
Ak je funkcia $f$ spojitá na kompaktnej množine $A$ (alebo špeciálne na
uzavretom ohraničenom intervale), tak je na $A$ ohraničená a existujú
$\max\limits_{x \in A}f(x)$ a $\min\limits_{x \in A}f(x)$.
\end{veta}

\begin{enumerate}[resume]
  \item \useproblem[spojitost-funkcie]{s-funkcie-240}
  \item \useproblem[spojitost-funkcie]{s-funkcie-241}
  \item \useproblem[spojitost-funkcie]{s-funkcie-242}
  \item \useproblem[spojitost-funkcie]{s-funkcie-243}
  \item \useproblem[spojitost-funkcie]{s-funkcie-244}
  \item \useproblem[spojitost-funkcie]{s-funkcie-245}
  \item \useproblem[spojitost-funkcie]{s-funkcie-246}
\end{enumerate}


\section{Rovnomerná spojitosť}
Funkcia $f$ sa nazýva rovnomerne spojitá na množine $A \subset D(f)$ $(A \neq
\emptyset)$, ak platí
\[
  (\forall \varepsilon > 0)
    (\exists \delta > 0)
      (\forall x,y\in A):
        |x-y| < \delta \Rightarrow |f(x)-f(y)| < \varepsilon
\]
Funkcia $f$ rovnomerne spojitá na svojom definičnom obore sa nazýva rovnomerne
spojitá.

\begin{enumerate}[resume]
  \item \useproblem[spojitost-funkcie]{s-funkcie-247}
  \item \useproblem[spojitost-funkcie]{s-funkcie-248}
  \item \useproblem[spojitost-funkcie]{s-funkcie-249}
  \item \useproblem[spojitost-funkcie]{s-funkcie-250}
\end{enumerate}

\begin{veta}
Funkcia spojitá na kompaktnej množine $K$ je rovnomerne spojitá na $K$.
\end{veta}

\begin{enumerate}[resume]
  \item \useproblem[spojitost-funkcie]{s-funkcie-251}
  \item \useproblem[spojitost-funkcie]{s-funkcie-252}
  \item \useproblem[spojitost-funkcie]{s-funkcie-253}
  \item \useproblem[spojitost-funkcie]{s-funkcie-254}
  \item \useproblem[spojitost-funkcie]{s-funkcie-255}
  \item \useproblem[spojitost-funkcie]{s-funkcie-256}
\end{enumerate}


\section{Ďalšie príklady}
\begin{enumerate}[resume]
  \item \useproblem[spojitost-funkcie]{s-funkcie-257}
  \item \useproblem[spojitost-funkcie]{s-funkcie-258}
  \item \useproblem[spojitost-funkcie]{s-funkcie-259}
  \item \useproblem[spojitost-funkcie]{s-funkcie-260}
  \item \useproblem[spojitost-funkcie]{s-funkcie-261}
  \item \useproblem[spojitost-funkcie]{s-funkcie-262}
  \item \useproblem[spojitost-funkcie]{s-funkcie-263}
  \item \useproblem[spojitost-funkcie]{s-funkcie-264}
  \item \useproblem[spojitost-funkcie]{s-funkcie-265}
  \item \useproblem[spojitost-funkcie]{s-funkcie-266}
  \item \useproblem[spojitost-funkcie]{s-funkcie-267}
  \item \useproblem[spojitost-funkcie]{s-funkcie-268}
  \item \useproblem[spojitost-funkcie]{s-funkcie-269}
  \item \useproblem[spojitost-funkcie]{s-funkcie-270}
  \item \useproblem[spojitost-funkcie]{s-funkcie-271}
  \item \useproblem[spojitost-funkcie]{s-funkcie-272}
  \item \useproblem[spojitost-funkcie]{s-funkcie-273}
  \item \useproblem[spojitost-funkcie]{s-funkcie-274}
  \item \useproblem[spojitost-funkcie]{s-funkcie-275}
  \item \useproblem[spojitost-funkcie]{s-funkcie-276}
  \item \useproblem[spojitost-funkcie]{s-funkcie-277}
  \item \useproblem[spojitost-funkcie]{s-funkcie-278}
\end{enumerate}



\chapter[Diferenciálny počet]{Diferenciálny počet funkcií jednej premennej}
\loadallproblems[diferencialny-pocet]{./tex/05-diferencialny-pocet/ulohy.tex}

\section{Definícia derivácie}
Hovoríme, že funkcia $f$ (definovaná v okolí bodu $a$ \footnote{t.j. pre
niektoré $\varepsilon>0$ platí $\interval[open]{a-\varepsilon}{a+\varepsilon}
\subset D(f)$; pojmy derivácie a nevlastnej derivácie a nevlastnej derivácie v
bode $a$ by bolo možné zaviesť aj za slabšieho predpokladu \enquote{$a \in D(f)$
je hromadný bod množiny $D(f)$} (odstránili by sa tým aj niektoré ťažkosti so
zavedením pojmu derivácie ako funkcie), definícia predpokladajúca, že $a$ je
vnútorný bod množiny $D(f)$, je však v literatúre najčastejšia}) má deriváciu v
bode $a,(a \in \mathbb{R})$, ak existuje konečná $\lim\limits_{x \rightarrow
a}\frac{f(x)-f(a)}{x-a}$. Ak existuje nevlastná $\lim\limits_{x \rightarrow
a}\frac{f(x)-f(a)}{x-a}$, hovoríme, že funkcia $f$ má nevlastnú (alebo
nekonečnú) deriváciu v bode $a$. Hodnotu $\lim\limits_{x \rightarrow
a}\frac{f(x)-f(a)}{x-a}$ v obidvoch týchto prípadoch označujeme $f'(a)$. Ak nás
nezaujíma, či je $\lim\limits_{x \rightarrow a}\frac{f(x)-f(a)}{x-a}$ konečná
alebo nekonečná. používame spoločný názov vlastná alebo nevlastná derivácia v
bode $a$. Pojmy derivácie sprava a nevlastnej derivácie sprava, resp. derivácie
zľava a nevlastnej derivácie zľava dostaneme, ak v predchádzajúcich definíciách
$\lim\limits_{x \rightarrow a}\frac{f(x)-f(a)}{x-a}$ nahradíme limitou
$\lim\limits_{x \rightarrow a+}\frac{f(x)-f(a)}{x-a}$, resp. $\lim\limits_{x
\rightarrow a-}\frac{f(x)-f(a)}{x-a}$ (v takom prípade stačí predpokladať, že
definičný obor funkcie $f$ obsahuje interval $\interval[open
right]{a}{a+\varepsilon}, resp. (a-\varepsilon,a\rangle$ pre niektoré
$\varepsilon>0$). Hodnoty $\lim\limits_{x \rightarrow a+}\frac{f(x)-f(a)}{x-a}$,
resp. $\lim\limits_{x \rightarrow a-}\frac{f(x)-f(a)}{x-a}$ označujeme
$f'_+(a)$,resp. $f'_-(a)$.

\begin{veta}
Funkcia $f$ (definovaná v okolí bodu $a \in \mathbb{R}$) má vlastnú alebo
nevlastnú deriváciu $f'(a)$ práve vtedy, keď existujú $f'_+(a),f'_-(a)$ a platí
$f'_+(a)=f'_-(a)$. Hodnota $f'(a)$ sa pritom rovná spoločnej hodnote $f'_+(a)$ a
$f'_-(a)$.
\end{veta}

Pojem derivácie ako funkcie sa vo všeobecnosti definuje nasledovne: Nech $M$ je
množina všetkých bodov definičného oboru $D(f)$, v ktorých má funkcia $f$
deriváciu. Funkcia $f':M \rightarrow \mathbb{R}$, ktorá každému bodu $a$  $M$
priradí hodnotu $f'(a)$ derivácie funkcie $f$ v bode $a$, sa nazýva derivácia
funkcie $f$. V špeciálnych prípadoch, ktoré ilustruje nasledujúca poznámka, sa
niekedy pojem derivácie ako funkcie chápe širšie: ak definičným oborom funkcie
$f$ je niektorý z intervalov $\interval{a}{b}, \interval[open right]{a}{b},
\interval[open left]{a}{b}$, pričom v jeho koncovom bode existuje jednostranná
derivácia (v prvom prípade prichádzajú do úvahy body $a,b$, v druhom bod $a$, v
trečom bod $b$), tak funkciu $f'$ považujeme za definovanú aj v tomto bode, jej
funkčnou hodnotou je hodnota príslušnej jednostrannej derivácie.

\begin{enumerate}[resume]
  \item \useproblem[diferencialny-pocet]{diferencialny-pocet-279}
  \item \useproblem[diferencialny-pocet]{diferencialny-pocet-280}
  \item \useproblem[diferencialny-pocet]{diferencialny-pocet-281}
  \item \useproblem[diferencialny-pocet]{diferencialny-pocet-282}
  \item \useproblem[diferencialny-pocet]{diferencialny-pocet-283}
  \item \useproblem[diferencialny-pocet]{diferencialny-pocet-284}
  \item \useproblem[diferencialny-pocet]{diferencialny-pocet-285}
\end{enumerate}

\begin{veta}
Ak funkcie $f,g$ majú derivácie v bode $a$, tak aj funkcie $c \cdot f$ ($c$ je
reálna konštanta), $f+g,f-g,f\cdot g$ majú v bode $a$ derivácie a platí:
\begin{enumerate}
\item $(c\cdot f)'(a)=c\cdot f'(a)$,
\item $(f+g)'(a)=f'(a)+g'(a)$,
\item $(f-g)'(a)=f'(a)-g'(a)$,
\item $(f\cdot g)'(a)=f'(a)\cdot g'(a)$.
\end{enumerate}
Ak naviac $g(a)\neq 0$ majú v bode $a$ deriváciu aj funkcie
$\frac{1}{g},\frac{f}{g}$ a platí
\begin{enumerate}
\item $(\frac{1}{g})'(a)=-\frac{g'(a)}{g^2(a)}$,
\item $(\frac{f}{g})'(a)=\frac{f'(a)\cdot g(a)-f(a)\cdot g'(a)}{g^2 (a)}$.
\end{enumerate}
\end{veta}

\begin{veta}
Ak funkcia $f$ má deriváciu v bode $a$, funkcia $g$ v bode $f(a)$ a zložená
funkcia $h=g \circ f$ je definovaná v okolí bodu $a$, tak $h$ má v bode $a$
deriváciu a platí $$h'(a)=f'(a)\cdot g'(f(a)).$$ (Analogické vety možno dokázať
aj pre jednostranné derivácie.)
\end{veta}

V nasledujúcej tabuľke sú derivácie základných elementárnych funkcií ($c$ je
reálna konštanta):
\begin{itemize}
\item $c'=0$,
\item $(x^n)'=n\cdot x^{n-1},(n \in \mathbb{R} \setminus \{0\})$,
\item $(\sin x)'=\cos x$,
\item $(\cos x)'=-\sin x$,
\item $(\tan x)'=\frac{1}{\cos^2 x}$,
\item $(a^x)'=a^x \cdot \ln a$, špeciálne $(e^x)'=e^x$,
\item $(\log_{a} x)'=\frac{1}{x\cdot \ln a},x>0$, špeciálne $(\ln x)'=\frac{1}{x},x>0$,
\item $(\arcsin x)'=\frac{1}{\sqrt{1-x^2}}$,
\item $(\arccos x)'=-\frac{1}{\sqrt{1-x^2}}$,
\item $(\arctan x)'=\frac{1}{1+x^2}$,
\item $(arcctan x)'=-\frac{1}{1+x^2}$.
\end{itemize}

Pre deriváciu hyperbolických funkcií platia nasledujúce vzorce
\begin{multicols}{2}
\begin{itemize}
  \item $(\sinh x)'=\cosh x$,
  \item $(\cosh x)'=\sinh x$,
  \item $(\tanh x)'=\frac{1}{ch^2 x}$,
  \item $(\coth x)'=-\frac{1}{sh^2 x}$.
\end{itemize}
\end{multicols}

Ak pri hľadaní derivácie funkcie využívame len znalosť derivácií základných
elementárnych funkcií a vety o derivácii súčtu, rozdielu, súčinu, podielu a
zloženej funkcie, nazýva sa taký postup tabuľkovým derivovaním.

\begin{enumerate}[resume]
  \item \useproblem[diferencialny-pocet]{diferencialny-pocet-286}
  \item \useproblem[diferencialny-pocet]{diferencialny-pocet-287}
  \item \useproblem[diferencialny-pocet]{diferencialny-pocet-288}
  \item \useproblem[diferencialny-pocet]{diferencialny-pocet-289}
  \item \useproblem[diferencialny-pocet]{diferencialny-pocet-290}
  \item \useproblem[diferencialny-pocet]{diferencialny-pocet-291}
  \item \useproblem[diferencialny-pocet]{diferencialny-pocet-292}
  \item \useproblem[diferencialny-pocet]{diferencialny-pocet-293}
  \item \useproblem[diferencialny-pocet]{diferencialny-pocet-294}
  \item \useproblem[diferencialny-pocet]{diferencialny-pocet-295}
  \item \useproblem[diferencialny-pocet]{diferencialny-pocet-296}
  \item \useproblem[diferencialny-pocet]{diferencialny-pocet-297}
  \item \useproblem[diferencialny-pocet]{diferencialny-pocet-298}
  \item \useproblem[diferencialny-pocet]{diferencialny-pocet-299}
  \item \useproblem[diferencialny-pocet]{diferencialny-pocet-300}
  \item \useproblem[diferencialny-pocet]{diferencialny-pocet-301}
  \item \useproblem[diferencialny-pocet]{diferencialny-pocet-302}
  \item \useproblem[diferencialny-pocet]{diferencialny-pocet-303}
  \item \useproblem[diferencialny-pocet]{diferencialny-pocet-304}
  \item \useproblem[diferencialny-pocet]{diferencialny-pocet-305}
  \item \useproblem[diferencialny-pocet]{diferencialny-pocet-306}
\end{enumerate}

\begin{veta}
Ak funkcia $f$ má v bode $a$ deriváciu, tak $f$ je v tomto bode spojitá.
\end{veta}

\begin{enumerate}[resume]
  \item \useproblem[diferencialny-pocet]{diferencialny-pocet-307}
  \item \useproblem[diferencialny-pocet]{diferencialny-pocet-308}
  \item \useproblem[diferencialny-pocet]{diferencialny-pocet-309}
\end{enumerate}


\section{Derivácia inverznej funkcie}
\begin{veta}
Nech funkcia $f:\mathbb{I}\rightarrow\mathbb{R}$, rýdzomonotónna na intervale
$I$, má deriváciu v bode $a$. Potom inverzná funkcia $f^{-1}$ má vlastnú alebo
nevlastnú deriváciu v bode $f(a)$, pričom platí
\begin{enumerate}
\item ak $f'(a)\neq 0$, tak $(f^{-1})'(f(a))=\frac{1}{f'(a)}$
\item ak $f'(a)=0$ a $f$ je rastúca, tak $(f^{-1})'(f(a))=+\infty$
\item ak $f'(a)=0$ a $f$ je klesajúca, tak $(f^{-1})'(f(a))=-\infty$
\end{enumerate}
\end{veta}

\begin{enumerate}[resume]
  \item \useproblem[diferencialny-pocet]{diferencialny-pocet-310}
  \item \useproblem[diferencialny-pocet]{diferencialny-pocet-311}
  \item \useproblem[diferencialny-pocet]{diferencialny-pocet-312}
  \item \useproblem[diferencialny-pocet]{diferencialny-pocet-313}
\end{enumerate}

\section{Diferenciál}
Hovoríme, že funkcia $f$ (definovaná v okolí bodu $a$ má v bode $a$
diferenciál (je diferencovateľná v bode $a$), ak existuje reálna konštanta $A$
taká, že pre funkciu $\omega$, definovanú vzťahom
\[
  f(x)=f(a)+A(x-a)+\omega(x)
\]
platí $\lim\limits_{x \rightarrow a}\frac{\omega(x)}{x-a}=0$. Funkcia definovaná
predpisom $y=A(x-a)$ sa v takom prípade označuje $df(a)$ a nazýva sa diferenciál
funkcie $f$ v bode $a$. Funkcia $df(a)$ sa zvyčajnezapisuje v tvare $df(a)=A$
$dx(a)$ \footnote{písmeno $a$ sa v zápisoch často vynecháva,preto sa možno
stretnúť aj so zápisom $df=A$ $dx$}, kde symbol $dx(a)$ (označujúci diferenciál
funkcie $g(x)=x$ v bode $a$,t.j. funkciu danú predpisom $y=x-a$) sa nazýva
diferenciál nezávislej premennej.
\begin{veta}
Funkcia $f$ je diferencovateľná v bode $a$ práve vtedy, keď $f$ má v bode $a$
deriváciu; pritom platí $A=f'(a)$, kde $A$ je konštanta z definície diferenciálu
funkcie $f$ v bode $a$.
\end{veta}

Graf funkcie $y=f(a)+f'(a)(x-a)$ je dotyčnicou v bodu $(a,f(a))$ ku grafu
funkcir $f$. Ak je funkcia $f$ spojitá v bode $a$, pričom $f'(a)$ je nevlastná,
je dotyčnicou v bode $(a,f(a))$ ku grafu funkcie $f$ priamka $x=a$.

\begin{enumerate}[resume]
  \item \useproblem[diferencialny-pocet]{diferencialny-pocet-314}
  \item \useproblem[diferencialny-pocet]{diferencialny-pocet-315}
  \item \useproblem[diferencialny-pocet]{diferencialny-pocet-316}
  \item \useproblem[diferencialny-pocet]{diferencialny-pocet-317}
  \item \useproblem[diferencialny-pocet]{diferencialny-pocet-318}
  \item \useproblem[diferencialny-pocet]{diferencialny-pocet-319}
  \item \useproblem[diferencialny-pocet]{diferencialny-pocet-320}
\end{enumerate}


\section{Derivácie vyšších rádov}
Derivácie vyšších rádov definujeme rekurentne: Nech funkcia $f$ je definovaná v okolí bodu $a\in\mathbb{R}$; označme $f^{(0)}:=f$. Hovoríme,že funkcia $f$ má n-tú deriváciu v bode $a$, reps. nevlastnú n-tú deriváciu v bode $a$, ak existuje derivácia, resp. nevlastná derivácia funkcie $f^{(n-1)}$ v bode $a$. n-tú deriváciu v bode $a$ aj nevlastnú n-tú deriváciu v bode $a$   označujeme $f^{(n)}(a)$. (Analogicky sa definujú vlastné a nevlastné jednostranné n-té derivácie v bode $a$.) Derivácia funkcie $f^{(n-1)}$ sa nazýva n-tá derivácia funkcie $f$ a označuje sa $f^{(n)}$. Okrem označení $f^{(1)},f^{(2)},f^{(3)},f^{(4)},f^{(5)},...$ sa používajú aj označenia $f',f'',f''',f''''$ (alebo $f^{VI}$),$f^{V},...$ .
Platia nasledujúce vzťahy:

\begin{itemize}
    \item $(x^m)^{(n)}=
        \begin{cases}
            \frac{m!x^{m-n}}{(m-n)!},& $ak $ n\leq m \\
            0, &  $ak $ n>m
        \end{cases}
        $,$(m\in\mathbb{N})$
    \item $(\sin x)^{(n)}=\sin (x+\frac{n\pi}{2})$
    \item $(\cos x)^{(n)}=\cos (x+\frac{n\pi}{2})$
    \item $(a^x)^{(n)}=a^x\ln^n a$
    \item $(e^x)^{(n)}=e^x$
    \item $(\log_a x)^{(n)}=\frac{(-1)^{n-1}(n-1)!}{x^n\ln^n a},x>0$
    \item $(\ln x)^{(n)}=\frac{(-1)^{n-1}(n-1)!}{x^n},x>0$
\end{itemize}

\begin{veta}
\textit{Leibnitzov vzorec}
Ak funkcie $f,g$ majú n-tú deriváciu v bode $a$, tak existuje $(f\cdot
g)^{(n)}(a)$ a platí
\[
    (f\cdot g)^{(n)}(a)=\sum_{k=0}^n {n \choose k} f^{k}(a)\cdot g^{(n-k)}(a)
\]
\end{veta}

\begin{enumerate}[resume]
	\item \useproblem[diferencialny-pocet]{diferencialny-pocet-321}
	\item \useproblem[diferencialny-pocet]{diferencialny-pocet-322}
	\item \useproblem[diferencialny-pocet]{diferencialny-pocet-323}
	\item \useproblem[diferencialny-pocet]{diferencialny-pocet-324}
	\item \useproblem[diferencialny-pocet]{diferencialny-pocet-325}
	\item \useproblem[diferencialny-pocet]{diferencialny-pocet-326}
	\item \useproblem[diferencialny-pocet]{diferencialny-pocet-327}
	\item \useproblem[diferencialny-pocet]{diferencialny-pocet-328}
\end{enumerate}


\section{Základné vety diferenciálneho počtu}
\begin{veta}
\textit{(Rolle).}
Nech funkcia $f:\interval{a}{b} \rightarrow \mathbb{R}$ vyhovuje nasledujúcim
podmienkam:
\begin{itemize}
\item $f$ je spojitá na intervale $\interval{a}{b}$
\item
  v každom bode $x\in \interval[open]{a}{b}$ existuje vlastná alebo nevlastná
  $f'(x)$
\item $f(a)=f(b)$
\end{itemize}
Potom existuje bod $c\in \interval[open]{a}{b}$, v ktorom $f'(c)=0$.
\end{veta}

\begin{enumerate}[resume]
  \item \useproblem[diferencialny-pocet]{diferencialny-pocet-329}
  \item \useproblem[diferencialny-pocet]{diferencialny-pocet-330}
  \item \useproblem[diferencialny-pocet]{diferencialny-pocet-331}
  \item \useproblem[diferencialny-pocet]{diferencialny-pocet-332}
  \item \useproblem[diferencialny-pocet]{diferencialny-pocet-333}
  \item \useproblem[diferencialny-pocet]{diferencialny-pocet-334}
  \item \useproblem[diferencialny-pocet]{diferencialny-pocet-335}
\end{enumerate}

\begin{veta}
\textit{(Lagrangeova veta o strednej hodnote.)}
Nech funkcia $f:\interval{a}{b} \rightarrow\mathbb{R}$ vyhovuje nasledujúcim
podmienkam:
\begin{itemize}
\item $f$ je spojitá;
\item
  v každom bode $x\in \interval[open]{a}{b}$ existuje vlastná alebo nevlastná
  $f'(x)$. Potom existuje bod $c\in \interval[open]{a}{b}$, v ktorom
  $f'(c)=\frac{f(b)-f(a)}{b-a}.$
\end{itemize}
\end{veta}

\begin{enumerate}[resume]
  \item \useproblem[diferencialny-pocet]{diferencialny-pocet-336}
  \item \useproblem[diferencialny-pocet]{diferencialny-pocet-337}
  \item \useproblem[diferencialny-pocet]{diferencialny-pocet-338}
  \item \useproblem[diferencialny-pocet]{diferencialny-pocet-339}
  \item \useproblem[diferencialny-pocet]{diferencialny-pocet-340}
  \item \useproblem[diferencialny-pocet]{diferencialny-pocet-341}
  \item \useproblem[diferencialny-pocet]{diferencialny-pocet-342}
  \item \useproblem[diferencialny-pocet]{diferencialny-pocet-343}
  \item \useproblem[diferencialny-pocet]{diferencialny-pocet-344}
  \item \useproblem[diferencialny-pocet]{diferencialny-pocet-345}
\end{enumerate}

\begin{veta}
\textit{(Cauchyho veta o strednej hodnote).}
Nech funkcie $f: \interval{a}{b} \rightarrow \mathbb{R}$ a $g: \interval{a}{b}
\rightarrow \mathbb{R}$ spĺňajú nasledujúce podmienky:
\begin{itemize}
\item
  $f$ a $g$ sú spojité na intervale $\interval{a}{b}$
\item
  v každom bode $x\in \interval[open]{a}{b}$ existujú vlastná alebo nevlastná
  $f'(x)$ a vlastná $g'(x)$
\end{itemize}
Potom existuje bod $c\in \interval[open]{a}{b}$, v ktorom platí
\[
  (f(b)-f(a))g'(c)=(g(b)-g(a))f'(c)
\]
Ak sú naviac splnené predpoklaady
\begin{itemize}
\item $(f'(x))^2+(g'(x))^2>0$ Pre každé $x\in \interval[open]{a}{b}$
\item $g(b)\neq g(a)$
\end{itemize}
možno uvedenú rovnosť písať v tvare
\[
  \frac{f(b)-f(a)}{g(b)-g(a)}=\frac{f'(c)}{g'(c)}
\]
\end{veta}

\begin{enumerate}[resume]
  \item \useproblem[diferencialny-pocet]{diferencialny-pocet-346}
  \item \useproblem[diferencialny-pocet]{diferencialny-pocet-347}
  \item \useproblem[diferencialny-pocet]{diferencialny-pocet-348}
\end{enumerate}


\section[Vyšetrovanie vlastností funkcií]{Vyšetrovanie vlastností funkcií pomocou diferenciálneho počtu}
\subsection{Monotónnosť}
\begin{veta}
Nech funkcia $f: I \rightarrow\mathbb{R}$ je spojitá na intervale $I$ a má
deriváciu v každom jeho vnútornom bode. Potom
\begin{itemize}
\item
  ak $f'>0,(f'\geq 0)$ vnútri intervalu $I$ (t.j. ak pre každý vnútorný bod
  $x$ intervalu $I$ platí $f'(x)>0,(f'(x)\geq 0)$), tak $f$ je rastúca
  (neklesajúca) na $I$
\item
  ak $f'<0,(f'\leq 0)$ vnútri intervalu $I$, tak $f$ je klesajúca (nerastúca) na
  $I$
\end{itemize}
\end{veta}

\begin{enumerate}[resume]
  \item \useproblem[diferencialny-pocet]{diferencialny-pocet-349}
  \item \useproblem[diferencialny-pocet]{diferencialny-pocet-350}
  \item \useproblem[diferencialny-pocet]{diferencialny-pocet-351}
\end{enumerate}

\begin{veta}
Nech funkcie $f,g$ sú $n$-krát diferencovateľné na intervale $I$, nech v bode
$a\in I$ platí $f(a)=g(a),f'(a)=g'(a),...,f^{(n-1)}(a)=g^{(n-1)}(a)$ (teda ak
$n=1$, predpokladáme len $f(a)=g(a)$). Potom
\begin{enumerate}
\item
  ak $f^{(n)}(x)>g^{(n)}(x)$, pre všetky $x\in I \cap \interval[open]{a}{\infty}$, tak
  $f(x)>g(x)$ Pre všetky $x\in I \cap \interval[open]{a}{\infty}$ (pritom samozrejme
  predpokladáme, že $I \cap \interval[open]{a}{\infty} \neq \emptyset$ )
\item
  \begin{itemize}
    \item
      ak $f^{(n)}(x)>g^{(n)}(x)$ pre všetky $x\in I \cap
      \interval[open]{-\infty}{a}$ a $n$ je párne, tak $f(x)>g(x)$ pre všetky
      $x\in I \cap \interval[open]{-\infty}{a}$;
    \item
      ak $f^{(n)}(x)>g^{(n)}(x)$ pre všetky $x\in I \cap
      \interval[open]{-\infty}{a}$ a $n$ je nepárne, tak $f(x)<g(x)$ pre všetky
      $x\in I \cap \interval[open]{-\infty}{a}$ (pritom predpokladáme $I \cap
      \interval[open]{-\infty}{a}\neq \emptyset$)
  \end{itemize}
\end{enumerate}
\end{veta}

\begin{enumerate}[resume]
  \item \useproblem[diferencialny-pocet]{diferencialny-pocet-352}
  \item \useproblem[diferencialny-pocet]{diferencialny-pocet-353}
  \item \useproblem[diferencialny-pocet]{diferencialny-pocet-354}
  \item \useproblem[diferencialny-pocet]{diferencialny-pocet-355}
  \item \useproblem[diferencialny-pocet]{diferencialny-pocet-356}
\end{enumerate}

\subsection{Konvexnosť a konkávnosť. Inflexné body}
Funkcia $f$ sa nazýva rýdzo konvexná (konvexná) na intervale $I \subset D(f)$,
ak platí $(*)$
\[
  \forall x,y\in I,x\neq y \forall p,q>0,p+q=1;f(px+qy)<pf(x)+qf(y)
\]
\[
  (\forall x,y\in I,x\neq y \forall p,q>0,p+q=1;f(px+qy)\leq pf(x)+qf(y))
\]
Funkcia $f$ sa nazýva rýdzo konkávna (konkávna) na intervale $I \subset D(f)$,
ak platí
\[
  \forall x,y\in I,x\neq y \forall p,q>0,p+q=1;f(px+qy)>pf(x)+qf(y)
\]
\[
  (\forall x,y\in I,x\neq y \forall p,q>0,p+q=1;f(px+qy)\geq pf(x)+qf(y))
\]
(Geometricky možno výrok $(*)$ interpretovať takto: pre ľubovoľné $2$ čísla
$x,y\in I,x<y$, leží úsečka spájajúca body $(x,f(x))$ nad grafom funkcie
$f/(x,y)$.)

\begin{veta}
Nech funkcia $f$ je spojitá na intervale $I$ a dvakrát diferencovateľná v každom
jeho vnútornom bode. Potom
\begin{itemize}
\item
  $f''>0,(f''\geq 0)$ vnútri intervalu $I$, tak $f$ je rýdzo konvexná (konvexná)
  na $I$
\item
  $f''<0,(f''\leq 0)$ vnútri intervalu $I$, tak $f$ je rýdzo konkávna (konkávna)
  na $I$
\end{itemize}
Vnútorný bod $a$ množiny $D(f)$ sa nazýva inflexný bod funkcie $f$, ak $f$ má v
bode $a$ deriváciu a existuje $\varepsilon >0$ tak, že funkcia $f$ je rýdzo
konvexná na jednej z množín $\interval[open left]{a-\varepsilon}{a}$,
$\interval[open right]{a}{a + \varepsilon}$ a rýdzo konkávna na druhej z nich.
\end{veta}

\begin{veta}
Nech funkcia $f$ je trikrát diferencovateľná v bode $a$ a dvakrát
diferencovateľná v niektorom jeho okolí. Ak $f''(a)=0,f''(a)\neq 0$, tak $a$ je
inflexný bod funkcie $f$.
\end{veta}

\textit{Poznámka:}
Existujú aj iné definície rýdzej konvexnosti, rýdzej konkávnosti a inflexného
bodu, ktoré nie sú ekvivalentné tu uvedeným. Všetky v matematickej literatúre
používané definície týchto pojmov sú však volené tak, že vety $13$ a $14$
zostanú v platnosti.

\begin{enumerate}[resume]
  \item \useproblem[diferencialny-pocet]{diferencialny-pocet-357}
  \item \useproblem[diferencialny-pocet]{diferencialny-pocet-358}
  \item \useproblem[diferencialny-pocet]{diferencialny-pocet-359}
  \item \useproblem[diferencialny-pocet]{diferencialny-pocet-360}
  \item \useproblem[diferencialny-pocet]{diferencialny-pocet-361}
  \item \useproblem[diferencialny-pocet]{diferencialny-pocet-362}
  \item \useproblem[diferencialny-pocet]{diferencialny-pocet-363}
\end{enumerate}

\textit{Poznámka:}
Na vlastnostiach odvodených v príklade $362$ sa zakladajú nasledujúce definície,
odlišné od definícií používaných v tomto odstavci:
\begin{itemize}
\item
  Funkcia $f$ diferencovateľná na intervale $I$ sa nazýva rýdzo konvexná (rýdzo
  konkávna) na $I$, ak pre každé $a\in I$ platí: na množine $I \setminus \{a\}$
  leží graf funkcie $f$ nad (pod) svojou dotyčnicou v bode $(a,f(a))$.
\item
  Bod $a$ sa nazýva inflexný bod funkcie $f$, ak existuje konečná $f'(a)$ a graf
  funkcie $f$ prechádza v bode $a$ z jednej strany svojej dotyčnice na druhú.
\end{itemize}
Prvá z týchto definícií predstavuje užšie chápanie pojmov rýdzo konvexná a rýdzo
konkávna funkcia (to sme tu nedokazovali, ale môžete to skúsiť sami), druhá
naopak širšie chápanie pojmu inflexný bod (to ukazujú príklady $362$ a $363$).

\subsection{Extrémy}
Nech definičným oborom funkcie $f$ je interval $I$. Hovoríme, že funkcia $f$ má
v bode $a$ $I$ lokálne maximum (lokálne minimum), ak existuje okolie $O(a)$ bodu
$a$ tak, že platí:
\[
  \forall x\in (O(a)\setminus \{a\})\cap I:f(x)\leq f(a)
\]
\[
  (\forall x\in (O(a)\setminus \{a\})\cap I:f(x)\geq f(a))
\]
Definíciu ostrého lokálneho maxima (ostrého lokálneho minima) dostaneme, ak v
predchádzajúcej definícii zameníme znak $\leq$ ($\geq$) znakom $<(>)$. Lokálne
maximá a lokálne minimá sa súhrnne nazývajú lokálnymi extrémami.

\begin{veta}
Ak funkcia $f$ má lokálny extrém vo vnútornom bode $a$ svojho definičného oboru,
tak buď neexistuje vlastná ani nevlastná $f'(a)$, alebo $f'(a)=0$. Bod $a$ sa
nazýva stacionálny bod funkcie $f$, ak $f'(a)=0$. Pri hľadaní lokálnych extrémov
funkcie $f:I\rightarrow\mathbb{R}$, treba teda vyšetriť:
\begin{enumerate}
\item všetky jej stacionárne body
\item všetky body $a\in I$, v ktorých neexistuhe $f'(a)$
\item všetky body $a\in I$, ktoré nie sú vnútornými bodmi intervalu $I$
\end{enumerate}
\end{veta}

\begin{veta}
Ak funkcia $f$ je dvakrát diferencovateľná vo vnútornom bode $a$ množiny $D(f)$
a platí $f'(a)=0,f''(a)>0$  $(f'(a)=0,f''(a)<0)$, tak $f$ má v bode $a$ ostré
lokálne minimum (ostré lokálne maximum).
\end{veta}

\begin{veta}
Nech funkcia $f$ je $n$-krát $(n\geq 2)$ diferencovateľná vo vnútornom bode $a$
množiny $D(f)$, nech $f'(a)=...=f^{(n-1)}(a)=0,f^{(n)}\neq 0$.

Ak $n$ je párne a $f^{(n)}(a)>0$  $(f^{(n)}(a)<0)$, tak funkcia $f$ má v bode
$a$ ostré lokálne minimum (ostré lokálne maximum).

Ak $n$ je nepárne, nemá funkcia $f$ v bode $a$ lokálny extrém.
\end{veta}

\begin{enumerate}[resume]
  \item \useproblem[diferencialny-pocet]{diferencialny-pocet-364}
  \item \useproblem[diferencialny-pocet]{diferencialny-pocet-365}
  \item \useproblem[diferencialny-pocet]{diferencialny-pocet-366}
  \item \useproblem[diferencialny-pocet]{diferencialny-pocet-367}
  \item \useproblem[diferencialny-pocet]{diferencialny-pocet-368}
  \item \useproblem[diferencialny-pocet]{diferencialny-pocet-369}
  \item \useproblem[diferencialny-pocet]{diferencialny-pocet-370}
  \item \useproblem[diferencialny-pocet]{diferencialny-pocet-371}
  \item \useproblem[diferencialny-pocet]{diferencialny-pocet-372}
  \item \useproblem[diferencialny-pocet]{diferencialny-pocet-373}
  \item \useproblem[diferencialny-pocet]{diferencialny-pocet-374}
  \item \useproblem[diferencialny-pocet]{diferencialny-pocet-375}
  \item \useproblem[diferencialny-pocet]{diferencialny-pocet-376}
\end{enumerate}


\section{L' Hospitalovo pravidlo}
\begin{veta}
Nech
\begin{enumerate}
\item
  funkcia $f,g$ sú diferencovateľné v niektorom prstencovom okolí $O^*(a)$ bodu
  $a\in\mathbb{R^*}$
\item
  $\forall x\in O^*(a):g'(x)\neq 0$
\item
  $\lim\limits_{x\rightarrow a}f(x)=\lim\limits_{x\rightarrow a}g(x)=0$ alebo
  $\lim\limits_{x\rightarrow a}|g(x)|=+\infty$
\item
  existuje vlastná alebo nevlastná $\lim\limits_{x\rightarrow a}\frac{f'(x)}{g'(x)}$.
  Potom existuje aj $\lim\limits_{x\rightarrow a}\frac{f(x)}{g(x)}$ a platí
  $\lim\limits_{x\rightarrow a}\frac{f(x)}{g(x)}=\lim\limits_{x\rightarrow
  a}\frac{f'(x)}{g'(x)}$

(Analogické tvrdenia možno sformulovať aj pre jednostranné limity.)
\end{enumerate}
\end{veta}

\begin{enumerate}[resume]
  \item \useproblem[diferencialny-pocet]{diferencialny-pocet-377}
  \item \useproblem[diferencialny-pocet]{diferencialny-pocet-378}
  \item \useproblem[diferencialny-pocet]{diferencialny-pocet-379}
  \item \useproblem[diferencialny-pocet]{diferencialny-pocet-380}
  \item \useproblem[diferencialny-pocet]{diferencialny-pocet-381}
  \item \useproblem[diferencialny-pocet]{diferencialny-pocet-382}
  \item \useproblem[diferencialny-pocet]{diferencialny-pocet-383}
  \item \useproblem[diferencialny-pocet]{diferencialny-pocet-384}
\end{enumerate}


\section{Taylorov polynóm}
Nech funkcia $f$ je $n$-krát diferencovateľná v bode $a\in\mathbb{R}$. Potom
Taylorovým polynómom stupňa $n$ funkcie $f$ v bode $a$ sa nazýva polynóm (v
premennej $x$)
\[
  f(a)+\frac{f'(a)}{1!}(x-a)+...+\frac{f^{(n)}(a)}{n!}(x-a)^n
\]
Ak špeciálne $a=0$, používa sa namiesto názvu Taylorov polynóm označenie
Maclaurinov polynóm.

\begin{veta}
Ak funkcia $f$ je $n$-krát diferencovateľná v bode $a\in\mathbb{R}$ a $T_n$ je
jej Taylorov polynóm stupňa $n$ v bode $a$, tak
\[
  \lim_{x \rightarrow a}\frac{f(x)-T_n(x)}{(x-a)^n}=0
\]
(Rozdiel $f-T_n$ sa nazýva zvyšok Taylorovho polynómu stupňa $n$ funkcie $f$ v
bode $a$.)
\end{veta}

Nech funkcia $g$ je definovaná v niektorom rýdzom okolí bodu $a\in\mathbb{R^*}$
a nenadobúda tam nulové hodnoty. Znakom $c(g)$ budeme označovať triedu všetkých
funkcií $f$ takých, že
\begin{itemize}
\item ich definičný obor obsahuje niektoré rýdze okolie bodu $a$
\item platí $\lim\limits_{x\rightarrow a}\frac{f(x)}{g(x)}=0$
\end{itemize}

Namiesto zápisu $f\in c(g)$ sa používa zápis $f=c(g)$. (Pokiaľ by z kontextu
nebolo jasné, ktorého $a\in\mathbb{R^*}$ sa vzťah $f=c(g)$ týka, používa sa
zápis $f=c(g)$ $(x\rightarrow a).$) Zápis $f=h+c(g)$ treba chápať nasledovne:
funkcia $f$ je súčtom funkcie $h$ a niektorej funkcie z triedy $c(g)$.

Ak funkcia $f$ je $n$-krát diferencovateľná v bode $a\in\mathbb{R}$, patrí podľa
vety $19$ zvyšok jej Taylorovho polynómu stupňa $n$ v bode $a$ do triedy
$c((x-a)^n)$, možno teda písať
\[
  f(x)=f(a)+\frac{f'(a)}{1!}(x-a)+...+\frac{f^{(n)}(a)}{n!}c((x-a)^n)
\]
tento zápis sa nazýva Taylorovým vzorcom so zvyškom v Peanovom tvare. V
nasledujúcich príkladoch budeme používať tieto tvrdenia (v zápisoch všade
vynechávame $x\rightarrow 0;m,n$ sú prirodzené čísla):
\begin{itemize}
\item ak $f(x)=c(x^n)$ a $g(x)=c(x^n)$, tak $f(x)+g(x)=c(x^n)$
\item ak $m>n$ a $f(x)=c(x^m)$, tak $f(x)=c(x^n)$
\item ak $f(x)=c(x^n)$, tak $f^m(x)=c(x^{m\cdot n})$
\item ak $f(x)=c(x^n)$, tak $x^m\cdot f(x)=c(x^{m+n})$
\item ak $f(x)=c(x^n)$ a $g(x)=c(x^m)$, tak $f(x)\cdot g(x)=c(x^{m+n})$
\end{itemize}

Uvedené implikácie budeme zapisovať nasledovne:
\begin{itemize}
\item $c(x^n)+c(x^n)=c(x^n)$
\item $c(x^m)=c(x^n)$ pre $m>n$
\item $c^m(x^n)=c(x^{m\cdot n})$
\item $x^m\cdot c(x^n)=c(x^{m+n})$
\item $c(x^n)\cdot c(x^m)=c(x^{m+})$
\end{itemize}

Používanie týchto zápisov vyžaduje istú opatrnosť, pretože ide o symbolické
vyjadrenie implikácií, možno (na rozdiel od skutočných rovností) tieto
\enquote{rovnosti} čítať sprava doľava.

\begin{enumerate}[resume]
  \item \useproblem[diferencialny-pocet]{diferencialny-pocet-385}
  \item \useproblem[diferencialny-pocet]{diferencialny-pocet-386}
\end{enumerate}

V nasledujúcich príkladoch využijeme tvrdenie z príkladu $386$ a znalosť
Maclaurinových polynómov týchto funkcií (možno ich zostrojiť výpočtom
príslušných derivácií):
\begin{itemize}
\item
  $e^x=1+x+\frac{x^2}{2!}+...+\frac{x^n}{n!}+c(x^n)$
\item
  $\sin
  x=x-\frac{x^3}{3!}+...+(-1)^{n-1}\frac{x^{2n-1}}{(2n-1)!}+c(x^{2n})$\footnote{pretože
  $\sin^{(2n)}(0)=0$, majú Maclaurinov polynóm stupňa $2n-1$ a Maclaurinov
  polynóm stupňa $2n$ rovnaký tvar, preto aj zvyšok Maclaurinovho polynómu
  stupňa $2n-1$ funkcie $\sin$ patrí do triedy $c(x^{2n})$; podobná poznáma
  platí o funkcii $\cos$}
\item
  $\cos x=1-\frac{x^2}{2!}+...+(-1)^{n}\frac{x^{2n}}{(2n)!}+c(x^{2n+1})$
\item
  $\ln (1+x)=x-\frac{x^2}{2}+...+(-1)^{n-1}\frac{x^n}{n}+c(x^n)$
\item
  $\frac{1}{1-x}=1+x+x^2+...+\frac{x^n}{n!}+c(x^n)$
\item
  $(1+x)^{\alpha}=1+x+\frac{\alpha(\alpha-1)}{2!}x^2+...+\frac{(\alpha-1) ...
  (\alpha-n+1)}{n!}x^n+c(x^n),(\alpha\in\mathbb{R})$
\end{itemize}

\begin{enumerate}[resume]
  \item \useproblem[diferencialny-pocet]{diferencialny-pocet-387}
  \item \useproblem[diferencialny-pocet]{diferencialny-pocet-388}
  \item \useproblem[diferencialny-pocet]{diferencialny-pocet-389}
  \item \useproblem[diferencialny-pocet]{diferencialny-pocet-390}
  \item \useproblem[diferencialny-pocet]{diferencialny-pocet-391}
\end{enumerate}

\begin{veta}
Nech je daná funkcia $f$, nech funkcia $f^{(n)}$ je definovaná a spojitá v
niektorom okolí $O(a)$ bodu $a\in\mathbb{R}$, nech pre každé $x\in O(a)\setminus
\{a\}$ existuje vlastná $f^{(n+1)}(x)$. Nech $T_n$ je Taylorov polynóm stupňa
$n$ funkcie $f$ v bode $a$. Potom
\begin{itemize}
  \item
    pre každé $x\in O(a),x>a$ $(x\in O(a),x<a)$ existuje číslo $\vartheta (x)\in
    (a,x)$  $(\vartheta (x)\in (x,a))$ také, že:
    \[
      f(x)-T_n(x)=\frac{f^{(n+1)}(\vartheta (x))}{(n+1)!}(x-a)^{n+1}
    \]
    (tzv. Lagrangeov tvar zvyšku)
  \item
    pre každé $x\in O(a),x>a$  $(x\in O(a),x<a)$ existuje číslo $\vartheta
    (x)\in (a,x)$  $\vartheta (x)\in (a,x)$  $(\vartheta (x)\in (x,a))$ také,
    že:
    \[
      f(x)-T_n(x)=\frac{f^{(n+1)}(\vartheta (x))}{n!}(x-a)(x-\vartheta (x))^n
    \]
    (tzv. Cauchyho tvar zvyšku)
\end{itemize}
\end{veta}

\begin{enumerate}[resume]
  \item \useproblem[diferencialny-pocet]{diferencialny-pocet-392}
  \item \useproblem[diferencialny-pocet]{diferencialny-pocet-393}
\end{enumerate}


\section[Zostrojovanie grafov]{Použitie diferenciálneho počtu pri zostrojovaní grafov funkcií}
Pri zostrojovaní grafu funkcie $f$ postupujeme spravidla nasledovne:
\begin{itemize}
\item
  určíme $D(f)$
\item
  nájdeme všetky hodnoty $x$, pre ktoré $f(x)=0$ (t.j. priesečníky grafu funkcie
  $f$ s osou $Ox$)
\item
  vyšetríme spojitosť funkcie $f$ a jej správanie sa v bodoch nespokojnosti
\item
  zistíme, na ktorých intervaloch je $f$ monotónna, a nájdeme body, v ktorých
  nadobúda lokálne extrémy
\item
  vyšetríme konvexnosť a konkávnosť funkcie $f$, nájdeme inflexné body
\item
  nájdeme asymptoty grafu funkcie (definíciu asymptoty pozri ďalej)
\end{itemize}
Zostrojenie grafu funkcie $f$ môžu uľahčiť niektoré jej špeciálne vlastnosti:
pri párnej alebo nepárnej funkcii stačí zostrojiť graf funkcie $f/D(f)\cap
\interval[open right]{0}{+\infty}$, v prípade periodickej funkcie $f$ graf
funkcie $f/\interval{a}{a + T} \cap D(f)$ a $T$ je niektorá perióda funkcie $f$.

(Nech je daná funkcia $f$, nech $a\in\mathbb{R}$ je hromadný bod množiny
$D(f)\cap (a,+\infty)$) (množiny $D(f)\cap (-\infty,a)$). Ak existuje
$\lim\limits_{x\rightarrow a+}f(x)$  $(\lim\limits_{x\rightarrow a-}f(x))$ a je nevlastná,
nazýva sa priamka $x=ay$ asymptotou bez smernice grafu funkcie $f$.

Nech bod $+\infty$ je hromadný bod definičného oboru funkcie $f$. Priamka
$y=kx+q$ sa nazýva asymptota so smernicou grafu funkcie $f$ v bode $+\infty$, ak
$\lim\limits_{x\rightarrow \infty}(f(x)-kx-q)=0$.

Analogicky sa definuje asymptota so smernicou grafu funkcie $f$ v bode
$-\infty$.

\begin{veta}
Nech bod $+\infty$ je hromadný bod definičného oboru funkcie $f$. Priamka
$y=kx+q$ je asymptotou so smernicou grafu funkcie $f$ v bode $+\infty$ práve
vtedy, keď:
\begin{itemize}
\item
  existuje konečná $\lim\limits_{x\rightarrow \infty} \frac{f(x)}{x}$ a platí
  $\lim\limits_{x\rightarrow \infty} \frac{f(x)}{x}=k$
\item
  existuje konečná $\lim\limits_{x\rightarrow \infty} (f(x)-kx)$ a rovná sa
  číslu $q$

(Analogická veta platí pre asymptotu so smernicou grafu funkcie $f$ v bode $-\infty$.)
\end{itemize}
\end{veta}

Zostrojte grafy nasledujúcich funkcií:
\begin{enumerate}[resume]
  \begin{multicols}{2}
    \item \useproblem[diferencialny-pocet]{diferencialny-pocet-394}
    \item \useproblem[diferencialny-pocet]{diferencialny-pocet-395}
    \item \useproblem[diferencialny-pocet]{diferencialny-pocet-396}
    \item \useproblem[diferencialny-pocet]{diferencialny-pocet-397}
    \item \useproblem[diferencialny-pocet]{diferencialny-pocet-398}
    \item \useproblem[diferencialny-pocet]{diferencialny-pocet-399}
    \item \useproblem[diferencialny-pocet]{diferencialny-pocet-400}
    \item \useproblem[diferencialny-pocet]{diferencialny-pocet-401}
    \item \useproblem[diferencialny-pocet]{diferencialny-pocet-402}
    \item \useproblem[diferencialny-pocet]{diferencialny-pocet-403}
    \item \useproblem[diferencialny-pocet]{diferencialny-pocet-404}
    \item \useproblem[diferencialny-pocet]{diferencialny-pocet-405}
    \item \useproblem[diferencialny-pocet]{diferencialny-pocet-406}
    \item \useproblem[diferencialny-pocet]{diferencialny-pocet-407}
    \item \useproblem[diferencialny-pocet]{diferencialny-pocet-408}
    \item \useproblem[diferencialny-pocet]{diferencialny-pocet-409}
    \item \useproblem[diferencialny-pocet]{diferencialny-pocet-410}
    \item \useproblem[diferencialny-pocet]{diferencialny-pocet-411}
  \end{multicols}
\end{enumerate}


\section{Ďalšie príklady}
\begin{enumerate}[resume]
	\item \useproblem[diferencialny-pocet]{diferencialny-pocet-412}
	\item \useproblem[diferencialny-pocet]{diferencialny-pocet-413}
	\item \useproblem[diferencialny-pocet]{diferencialny-pocet-414}
	\item \useproblem[diferencialny-pocet]{diferencialny-pocet-415}
	\item \useproblem[diferencialny-pocet]{diferencialny-pocet-416}
	\item \useproblem[diferencialny-pocet]{diferencialny-pocet-417}
	\item \useproblem[diferencialny-pocet]{diferencialny-pocet-418}
	\item \useproblem[diferencialny-pocet]{diferencialny-pocet-419}
	\item \useproblem[diferencialny-pocet]{diferencialny-pocet-420}
	\item \useproblem[diferencialny-pocet]{diferencialny-pocet-421}
	\item \useproblem[diferencialny-pocet]{diferencialny-pocet-422}
	\item \useproblem[diferencialny-pocet]{diferencialny-pocet-423}
	\item \useproblem[diferencialny-pocet]{diferencialny-pocet-424}
	\item \useproblem[diferencialny-pocet]{diferencialny-pocet-425}
	\item \useproblem[diferencialny-pocet]{diferencialny-pocet-426}
	\item \useproblem[diferencialny-pocet]{diferencialny-pocet-427}
	\item \useproblem[diferencialny-pocet]{diferencialny-pocet-428}
	\item \useproblem[diferencialny-pocet]{diferencialny-pocet-429}
	\item \useproblem[diferencialny-pocet]{diferencialny-pocet-430}
	\item \useproblem[diferencialny-pocet]{diferencialny-pocet-431}
	\item \useproblem[diferencialny-pocet]{diferencialny-pocet-432}
	\item \useproblem[diferencialny-pocet]{diferencialny-pocet-433}
	\item \useproblem[diferencialny-pocet]{diferencialny-pocet-434}
	\item \useproblem[diferencialny-pocet]{diferencialny-pocet-435}
	\item \useproblem[diferencialny-pocet]{diferencialny-pocet-436}
	\item \useproblem[diferencialny-pocet]{diferencialny-pocet-437}
	\item \useproblem[diferencialny-pocet]{diferencialny-pocet-438}
	\item \useproblem[diferencialny-pocet]{diferencialny-pocet-439}
	\item \useproblem[diferencialny-pocet]{diferencialny-pocet-440}
	\item \useproblem[diferencialny-pocet]{diferencialny-pocet-441}
	\item \useproblem[diferencialny-pocet]{diferencialny-pocet-442}
	\item \useproblem[diferencialny-pocet]{diferencialny-pocet-443}
	\item \useproblem[diferencialny-pocet]{diferencialny-pocet-444}
	\item \useproblem[diferencialny-pocet]{diferencialny-pocet-445}
	\item \useproblem[diferencialny-pocet]{diferencialny-pocet-446}
	\item \useproblem[diferencialny-pocet]{diferencialny-pocet-447}
	\item \useproblem[diferencialny-pocet]{diferencialny-pocet-448}
	\item \useproblem[diferencialny-pocet]{diferencialny-pocet-449}
	\item \useproblem[diferencialny-pocet]{diferencialny-pocet-450}
	\item \useproblem[diferencialny-pocet]{diferencialny-pocet-451}
	\item \useproblem[diferencialny-pocet]{diferencialny-pocet-452}
	\item \useproblem[diferencialny-pocet]{diferencialny-pocet-453}
	\item \useproblem[diferencialny-pocet]{diferencialny-pocet-454}
	\item \useproblem[diferencialny-pocet]{diferencialny-pocet-455}
	\item \useproblem[diferencialny-pocet]{diferencialny-pocet-456}
	\item \useproblem[diferencialny-pocet]{diferencialny-pocet-457}
	\item \useproblem[diferencialny-pocet]{diferencialny-pocet-458}
	\item \useproblem[diferencialny-pocet]{diferencialny-pocet-459}
	\item \useproblem[diferencialny-pocet]{diferencialny-pocet-460}
	\item \useproblem[diferencialny-pocet]{diferencialny-pocet-461}
	\item \useproblem[diferencialny-pocet]{diferencialny-pocet-462}
	\item \useproblem[diferencialny-pocet]{diferencialny-pocet-463}
	\item \useproblem[diferencialny-pocet]{diferencialny-pocet-464}
	\item \useproblem[diferencialny-pocet]{diferencialny-pocet-465}
	\item \useproblem[diferencialny-pocet]{diferencialny-pocet-466}
	\item \useproblem[diferencialny-pocet]{diferencialny-pocet-467}
	\item \useproblem[diferencialny-pocet]{diferencialny-pocet-468}
	\item \useproblem[diferencialny-pocet]{diferencialny-pocet-469}
	\item \useproblem[diferencialny-pocet]{diferencialny-pocet-470}
	\item \useproblem[diferencialny-pocet]{diferencialny-pocet-471}
	\item \useproblem[diferencialny-pocet]{diferencialny-pocet-472}
	\item \useproblem[diferencialny-pocet]{diferencialny-pocet-473}
	\item \useproblem[diferencialny-pocet]{diferencialny-pocet-474}
	\item \useproblem[diferencialny-pocet]{diferencialny-pocet-475}
	\item \useproblem[diferencialny-pocet]{diferencialny-pocet-476}
	\item \useproblem[diferencialny-pocet]{diferencialny-pocet-477}
	\item \useproblem[diferencialny-pocet]{diferencialny-pocet-478}
\end{enumerate}



\chapter{Riešenia}%\label{chapter:gramatiky}

\epigraph{
  \enquote{Jíž pohled na ty pítvorné vzorce,\\
  hemžící se odmocninami a trigonometrií,\\
  je s to vyhnat veškerou nezpednost z těla.}
}{\textit{K. Poláček: Hostinec U kamenného stolu}}

\section{Úvod}
$\boxed{1.}$
$\boldsymbol{1.}\ \exists x\in \mathbb{R}:x^2<x;\\
\boldsymbol{2.}\ \forall x\in \mathbb{R}\ \exists y\in \mathbb{N}:y=x;\\
\boldsymbol{3.}\ \exists x,y\in \mathbb{R}:x>y \wedge	x^2 \leq y^2;\\
\boldsymbol{4.}\ \exists x,y\in \mathbb{R}\ \forall y\in \mathbb{N}: x\geq z \vee y\geq z; \\
\boldsymbol{5.}\ $symbolický zápis uvedeného výroku je 

$\forall a\in \mathbb{R}\setminus \{0\}\ \forall b\in \mathbb{R}\setminus \{0\}:((\ln |a|>\ln |b|) \Rightarrow (a>b)),$ 

negácia je $\exists a,b\in \mathbb{R}\setminus \{0\}: (\ln|a|>\ln|b|\vee a\leq b);\\
\boldsymbol{6.}\ \exists a\in \mathbb{R^+}:(|\ln(a)|<|\ln(b)|\wedge a\geq b);\\
\boldsymbol{7.}\ \exists A\subset \mathbb{R},A\neq \emptyset \exists B\subset \mathbb{R},B\neq \emptyset:$

$((\forall x\in A\ \forall y\in B\ \exists c>0:|x-y|>c)\wedge A\cap B\neq \emptyset),$ 

teda slovne: existujú neprázdne množiny $A,B\subset \mathbb{R}$ také, že platí 

$A\cap B\neq \emptyset$ a súčasne $\forall x\in A\ \forall y\in B\ \exists c>0:|x-y|>c;\\
\boldsymbol{8.}\ \exists A\subset \mathbb{R},A\neq \emptyset\ \exists B\subset \mathbb{R},B\neq \emptyset\ \exists c\subset \mathbb{{R_{0}}^+},C\neq \emptyset:$

$((\forall x\in A\ \exists y\in B\ \forall x\in C:x+c<y)\wedge(\exists x\in A\ \exists y\in B: x\geq y));\\
\boldsymbol{9.}\ \exists A\subset \mathbb{N},A\neq \emptyset:$

$((\forall n\in A:(2\nmid n\vee 3\nmid n))\wedge (\exists n\in A: 2|n)\wedge (\exists n\in A:3|n));\\
\boxed{2.}
\boldsymbol{1.}\ $platí tvrdenie $\forall\ r,s\in\mathbb{Z}:((2|r\wedge 2|s)\Rightarrow 2|r+s),$ ale neplatí tvrdenie s opačnou implikáciou, t.j. $\forall r,s\in\mathbb{Z}:((2|r+s)\Rightarrow (2|r \wedge 2|s))$, treba preto doplniť ``stačí'' (ale nemožno doplniť ``je nutné'' ani ``je nutné a stačí'');\\
$\boldsymbol{2.}\ $``je nutné'';
$\boldsymbol{3.}\ $``stačí'';
$\boldsymbol{4.}\ $``stačí'';\\
$\boldsymbol{5.}\ $platí výrok $\forall x\in\mathbb{R}\setminus \{0\}:(\frac{1}{x}<1\iff (x>1 \vee x<0))$ (a teda platia aj výroky $\forall x\in\mathbb{R}\setminus \{0\}:(\frac{1}{x}<1\Rightarrow (x>1 \vee x<0))$ a 

$\forall x\in \mathbb{R}\setminus \{0\}:((x>` \vee x<0)\Rightarrow \frac{1}{x}<1)$), preto doplnením každého z uvedených výrazov vznikne pravdivý výrok;


\section{Množiny}
$\boxed{3.}$
nepriamo z negácie (t.j. z výroku $\exists a\in\mathbb{Q},b\in\mathbb{R}\setminus \mathbb{Q}:a+b\in\mathbb{Q}$) vyplýva (pretože rozdiel racionálnych čísel $a+b$, $a$ je racionálne číslo) $\exists b\in \mathbb{R}:b\in \mathbb{R}\setminus \mathbb{Q}\wedge b\in \mathbb{Q}$, čo je nepravdivý výrok;\\
$\boxed{4.}$
nepriamo; využite rovnosť $\sqrt{ab}=\frac{1}{2}((\sqrt{a})+\sqrt{b})^2-a-b$;\\
$\boxed{5.}$
$\boldsymbol{1.},\ \boldsymbol{2.}$ možno postupovať ako v známom dôkaze iracionálnosti čísla $\sqrt{2}$; $\boldsymbol{3.},\ \boldsymbol{4.}$ využite príklad č.$4$; $\boldsymbol{5.}$ nepriamo ($\sqrt{3}=\frac{6a+3}{4},$ kde $a=\frac{4\sqrt{3}-3}{6}$);\\
$\boxed{6.}$
pozor, odvodenie pravdivého výroku z danej nerovnosti nie je ešte jej dôkazom;
$\boldsymbol{4.}$ do nerovnosti $x^2+y^2\geq 2xy$ dosaďte za $(x,y)$ postupne $(a,b),(a,c),(b,c)$ a získané nerovnosti sčítajte;
$\boldsymbol{6.}$ vyplýva z nerovnosti $(\sqrt{x^2+1}-1)^2 \geq 0$;
$\boxed{7.}$ 
ak $m\ (M)$ je najmenší (najväčší) zlomok, tak $mb_1\leq a_1\leq Mb_1,...,mb_n\leq a_n\leq Mb_n$; súčet týchto nerovností vydeľte $(b_1+...+b_n)$;
$\boxed{8.}$
$\boldsymbol{3.}$ matematickou indukciou, v jej druhom kroku treba využiť nerovnosť $\frac{\sqrt{n}+1}{\sqrt{n+1}}$;
$\boldsymbol{4.}$ v druhom kroku indukcie stačí pripočítaním a odpočítaním čísla $\frac{1}{n+1}$ upraviť ľavú stranu nerovnosti tak, aby sme mohli využiť indukčný predpoklad $\frac{1}{n+2}+...+\frac{1}{2n+2}=(\frac{1}{n+1}+...+\frac{1}{2n})+(\frac{1}{2n+1}+\frac{1}{2n+2}-\frac{1}{n+1})$, pritom druhá zo zátvoriek je kladné číslo;\\
$\boxed{9.}$
$\boldsymbol{2.}$ pre $n=1$ má dokazovaná nerovnosť tvar $|x_1|\leq |x_1|$, čo zrejme platí; dôkaz pre $n=2$: podľa pr. $9.1$ je $x_1\leq |x_1|,-x_1\leq |x_1|,x_2\leq |x_2|,-x_2\leq |x_2|$; odtiaľ vyplýva $x_1+x_2 \leq |x_1|+|x_2|,-x_1-x_2\leq |x_1|+|x_2|$; preto $\max \{x_1+x_2,-x_1-x_2\}\leq |x_1|+|x_2|$; t.j. (podľa pr. $9.1$) $|x_1+x_2|\leq |x_1|+|x_2|$; pre $n>2$ možno postupovať analogicky alebo použiť matematickú indukciu: $|x_1+...+x_{n+1}|=|(x_1+...+x_n)+x_{n+1}|\leq |x_1+...+x_n|+|x_{n+1}|$ (iná možnosť dôkazu pre $n=2:|x_1+x_2|^2=|x_1|^2+2x_1x_2+|x_2|^2\leq (|x_1|+|x_2|)^2$);$3.$ do nerovnosti $|x+y|\leq |x|+|y|$ dosaďte za $(x,y)$ postupne $(a-b,b)$, $(b-a,a)$ a použite pr. $9.1$;\\
$\boxed{11.}$
$\boldsymbol{4.}$ podľa návodu stačí dokázať nerovnosť $2\leq (\frac{(n+2)}{n+1})^{n+1}=(1+\frac{1}{n+1})^{n+1}$, tá vyplýva z binomickej vety (alebo z pr. $10.1$); $\boldsymbol{5.}$ nerovnosť $((2n+2)!=)$ $(n+2)!$ $[(n+3)(n+4)...(2n+2)]\geq (n+2)!$ $[(n+2)^n]$ možno dokázať priamo: $k$-ty člen v hranatej zátvorke vľavo je väčší ako $k$-ty člen v hranatej zátvorke vpravo, $k=1,...,n$;
$\boldsymbol{6.}$ nerovnosť možno dokázať priamo použitím binomickej vety, iná možnosť: podľa návodu stačí dokázať nerovnosť $(\frac{(n+1)^2}{n*(n+2)})^{n+1}\geq 1$, na jej dôkaz možno použiť pr. $10.1$;\\
$\boxed{12.}$ 
$\boldsymbol{1.}$ $x_1(1-x_2)>1-x_2$; $\boldsymbol{2.}$ matematickou indukciou; ak neplatí $x_1=x_2=...=x_{n+1}=1$ (vtedy totiž dokazovaná nerovnosť zrejme platí), je aspoň jedno z čísel $x_1,x_2,...,x_{n+1}$ väčšie a aspoň jedno menšie než $1$; preznačme ich tak, aby $x_n>1,x_{n+1}<1$; ak na $n$ čísel $x+1,...,x_{n-1},x_n*x_{n+1}$ použijeme indukčný predpoklad (ktorý má podobu $\forall a_1,...,a_n\in \mathbb{R^+}:a_1*a_2*...*a_n=1\Rightarrow a_1+...+a_n\geq n$), dostaneme $x_1+...+x_{n-1}+x_n*x_{n+1}\geq n$; podľa pr. $12.1$ je potom $x_1+...+x_{n-1}+x_m+x_{n+1}>x_1+...+x_{n-1}+x_n*x_{n+1}+1\geq n+1$;
$3.$ čísla $\frac{x_1}{x_2},...,\frac{x_n}{x_1}$, resp. $\frac{x_1}{\sqrt[n]{x_1*...*x_n}},...,\frac{x_n}{\sqrt[n]{x_1*...*x_n}}$ vyhovujú predpokladom tvrdenia z pr. $12.2$; 
$\boxed{13.}$
$\boldsymbol{1.}$ do nerovnosti $\frac{1}{k^2}<\frac{1}{k(k-1)}=\frac{1}{k-1}-\frac{1}{k}$ $(k\geq 2)$ dosaďte postupne $k=2,...,k=n$ a sčítajte; $\boldsymbol{2.}$ analogicky využite nerovnosť $\frac{1}{n!}<\frac{1}{2^{n-1}}$ $(n\geq 3)$ (možno ju dokázať matematickou indukciou);
$\boxed{14.}$
$\boldsymbol{1.}$ zdola ohraničená $(\forall m\in A:\sqrt{2}\leq m)$; platí $\forall K\in\mathbb{R}\ \exists m\in A:m>K$ (pre dané $K$ stačí položiť $a=1,b=n^4$, kde $n\in\mathbb{N}\setminus \{1\}$ je číslo väčšie než $K$; existenciu takého $n$ zaručuje Archimedov princíp - pozri $[18,$ str. $31]$), teda $A$ nie je ohraničená zhora; $\boldsymbol{2.}$ ohraničená $(\forall b\in B: 0<b\leq \frac{1}{2})$; $\boldsymbol{3.}$ ohraničená (pre každé $x\in\mathbb{R}$, teda aj pre $x=n!$, platí $|\sin x|\leq 1$); $\boldsymbol{4.}$ ohraničená (pre všetky $x\in (2,3)$, teda aj pre všetky $x\in\mathbb{Q}\cap (2,3)$ paltí $\frac{\sqrt{2}}{\sqrt[3]{3}+\sqrt[4]{3}}\leq \frac{\sqrt{x}}{\sqrt[3]{x}+\sqrt[4]{x}}\leq \frac{\sqrt{3}}{\sqrt[3]{2}+\sqrt[4]{2}}$); $\boldsymbol{5.}$ platí $\mathbb{Z}\subset E$ $(x^2-z^2$, kde $z\in\mathbb{Z}$, je polynóm druhého stupňa s racionálnymi koeficientami, ktorého koreňom je číslo $z$), preto $E$ nie je zhora ani zdola ohraničená;
$\boxed{15.}$
nepriamo; z výroku $ \exists a\in A \quad \exists \beta > 0  \quad \forall b \in A : \quad \mid a - b \mid < \beta \quad $ vyplýva na základe nerovnosti $ \mid b \mid \leq \mid a \mid + \mid b - a \mid  $ tvrdenie $ \exists a\in A \quad \exists \beta > 0  \quad \forall b \in A : \quad \mid b \mid < \mid a \mid + \beta$ , čo znamená, že $ A $ je ohraničená množina; 
$\boxed{16.}$ 
sporom; ak platí $ ( \star ) $ a $ A $ je ohraničená, tak z toho a nerovnosti $ \mid x \mid \leq  \mid y \mid + \mid x - y \mid $ vyplýva, že aj $ B $ je ohraničená;
$\boxed{17.}$
$\boldsymbol{2.}$
platí (pri dôkaze vlastnosti (ii) vyhovuje príslušnej nerovnosti pre $\varepsilon > 1$ každé reálne číslo $ x_{\varepsilon}  $, pre $ 0 < \varepsilon \leq  1 $ každé číslo  $ x_{\varepsilon}  $  s vlastnosťou $ \mid x_{\varepsilon} \mid > \sqrt{ \frac{1 - \varepsilon }{2 \varepsilon}} $;
$\boldsymbol{3.}$ $ 2x^{2}+8x+1=2.(x+2)^{2} -7 $, odtiaľ vidno, že $ -2 $ nie je dolným ohraničením danej množiny  ( teda nevyhovuje podmienke (i), ale nevyhovuje podmienke (ii) z definície  suprema (existuje totiž menšie horné ohraničenie ako číslo 12), uvedená rovnosť teda neplatí;
$\boxed{18.}$
$\boldsymbol{1.}$ $ sup \, A = 3 $, $ inf \, A = 2 $;
$\boldsymbol{2.}$ $ inf \, B = 0 $, $ sup \, B = 1,\overline{1} =\frac{10}{9} $;
$\boldsymbol{3.}$ $ C = \left\lbrace  1, -1 \right\rbrace  $, teda $ inf \, C = -1 $, $ sup \, C = 1 $;
$\boxed{19.}$
označme $ \beta := sup \, B $, potom pre všetky $ x \in B $ a teda aj pre všetky $ x \in A $, platí $ x \leq \beta $; teda $ \beta $ je horné ohraničenie množiny $ A $; pretože $ sup \, A $ je najmenšie z horných ohraničení množiny $ A $, platí $ sup \, A \leq \beta $, tvrdenie pre infíma znie: Nech $ A \subset B  \subset \R  $, pričom $ A $ je neprázdna a $ B $ zdola ohraničená množina, potom $ A $ je zdola ohraničená a platí $ inf \, A \geq inf  \, B $. 
$\boxed{21.}$ 
označme $ \alpha := sup \, A $,  $ \beta := sup \, B $; pretože $ x \leq \alpha $ pre $ x \in A $, $ y \leq \beta $ pre $ y \in B $, je $ x+y \leq \alpha + \beta $ pre každé $ x \in A $ a každé $ y \in B $ ( tým je dokázaná vlastnosť (i)), ak $ x_{\varepsilon } > \alpha - \frac{\varepsilon}{2} $,  $ y_{\varepsilon } > \beta - \frac{\varepsilon}{2} $   ( $ \alpha, \, \beta $ majú vlastnosť (ii), preto také $ x_{\varepsilon } \in A $, $ y_{ \varepsilon }$  $ \in B $ pre dané ${\varepsilon } > 0 $ existujú), tak $ x_{\varepsilon } + y_{\varepsilon } > \alpha + \beta - \varepsilon$, pričom $ x_{\varepsilon } + y_{\varepsilon } \in A + B $ (tým je dokázaná vlastnosť (ii)), pre infíma platí  $ inf \, A + inf \,B = inf \, ( A + B ) $; 
$\boxed{22.}$ 
$\boldsymbol{1.}$
$ y = \sin x , u=y^{3} $;
$\boldsymbol{2.}$ 
$ y = x^{3}, u= \sin y $
$\boldsymbol{3.}$ 
$ y = \tan x, u= y^{2}, v = 5^{u} $;
$\boldsymbol{4.}$ 
$ y = b^{x}, u= \sqrt{y}, v= \sin u, w = v^{2}, z = \log_{3} w $;
$\boldsymbol{5.}$ 
$ y = \sin x , u= 2^{y}, v= \cos u, w = \sqrt{v}$;
$\boxed{23.}$ 
$\boldsymbol{1.}$
$ (- \infty , - \sqrt{3}\rangle \cup \langle 0 , \sqrt{3}\rangle $;
$\boldsymbol{2.}$
$ (- \infty , -2 ) \cup ( 2 , \infty) $;
$\boldsymbol{3.}$
$\bigcup\limits_{n=0}^\infty \langle 4 n^{2} \pi ^{2}, (2n+1)^{2} \pi ^{2}\rangle $;
$\boldsymbol{4.}$
$\bigcup\limits_{k \in Z} (e^{-x/2}.e^{2kx},e^{x/2}.e^{2kx}) $;
$\boldsymbol{5.}$
$ \lbrace (2k + 1) \pi  ; k \in Z  \rbrace \cup \bigcup\limits_{n  \in Z} (\langle 2n  \pi , \frac{\pi}{3} + 2n \pi\rangle \cup \langle \frac{4 \pi}{3}+2n \pi , \frac{3 \pi}{2} + 2n \pi \rangle) $;
$\boldsymbol{6.}$
$ \langle - \sqrt{ \frac{ \pi}{2}},\sqrt{ \frac{ \pi}{2} \rangle} \cup \bigcup\limits_{n=1}^\infty  
\langle  \sqrt{- \frac{ \pi}{2}+ 2n \pi} ,  \sqrt{\frac{ \pi}{2}+ 2n \pi}\rangle \cup \bigcup\limits_{n=1}^\infty  
 \langle -\sqrt{ \frac{ \pi}{2}+ 2n \pi} , -\sqrt{-\frac{ \pi}{2}+ 2n \pi}\rangle $;
 $\boldsymbol{7.}$
 $ (- \frac{1}{3}, \infty) $;
 $\boldsymbol{8.}$
 pre $ f(x) =3. \sin^{2}x - 4, \, g(x)= \ln x $ je množina $ D_{3} $ z definície zloženej funkcie prázdna, preto týmto predpisom nie je definovaná žiadna funkcia; 
 $\boldsymbol{9.}$
 $ \langle -3, -1) \cup (2, 3\rangle $;
 $\boldsymbol{10.}$
 $ \langle 1, 4 \rangle $;
 $\boldsymbol{11.}$ 
 $ (0, \frac{1}{4}) $;
 $\boldsymbol{12.}$
 $ \mathbb{Z} $;
 $\boldsymbol{13.}$
 $ \bigcup\limits_{k \in Z} ((- \frac{\pi}{2}+ 2k \pi, 2k \pi) \cup ( 2k \pi, \frac{\pi}{2}+ 2k \pi)) $;
 $\boldsymbol{14.}$
 $ (- \infty, 1) \cup \langle 2, \infty) $;
 $\boldsymbol{15.}$
 $ (2,3) $;
 $\boldsymbol{16.}$
 $ \bigcup\limits_{n \in Z} (( \frac{\pi}{4}+ n \pi,  \frac{3}{4} \pi + n \pi) $ ( pre $n$ musí platiť alebo $ \sin x + \cos x = 0 $ alebo $ 0 < (\sin x + \cos x)(\sin x - \cos x)= - \cos 2x $; využili sme, že podiel $ \frac{a}{b} $ je kladný práve vtedy, keď $ a.b > 0 $ ); 
 $\boldsymbol{17.}$
 $ (-6,\frac{-5 \pi}{3}\rangle \cup \langle \frac{- \pi}{3}, \frac{1}{6 }) $; 
 $\boxed{24.}$ 
 $\boldsymbol{1.}$
$ f(f(x)) = x^{4} $, $ f(g(x)) = 4^{x} $, $ g(f(x)) = 2^{x^{2}} $, $ g(g(x)) = 2^{2^{x}} $;\\
\begin{align*} 
\boldsymbol{2.}\ f(g(x))= g(f(x))= \left\{ \begin{array}{cc} 
                1 & \hspace{5mm} ak\ x > 0 \\
                -1 & \hspace{5mm} ak\ x < 0 , \\
                \end{array} \right.
\end{align*}
$ g(g(x))= x \, x \ne 0 $; \\
$\boldsymbol{3.}$
$ f(f(x)) = f(x) $, $ f(g(x)) = 0 $, $ g(f(x)) = g(x) $, $ g(g(x)) = 0 $;
\begin{align*} 
\boldsymbol{4.}\ f(f(x))= \left\{ \begin{array}{cc} 
                x^{4} & \hspace{5mm} ak\ x \in \langle0,1\rangle \\
                9x & \hspace{5mm} ak\ x \notin \langle0,1\rangle , \\
                \end{array} \right.
\end{align*}, 
\begin{align*} 
 f(f(x))= \left\{ \begin{array}{cc} 
                4x^{2} & \hspace{5mm} ak\ x \in \langle0, \frac{1}{2}\rangle \\
                6x & \hspace{5mm} ak\ x \in ( \frac{1}{2},1\rangle \\
                12x-6 & \hspace{5mm} ak\ x \notin \langle0,1\rangle , \\
                \end{array} \right.
\end{align*}
\begin{align*} 
g(f(x))= \left\{ \begin{array}{cc} 
                2x^{2} & \hspace{5mm} ak\ x \in \langle0,1\rangle \\
                12x-2 & \hspace{5mm} ak\ x \notin \langle0,1\rangle , \\
                \end{array} \right.
\end{align*}, 
\begin{align*} 
 f(f(x))= \left\{ \begin{array}{cc} 
                4x & \hspace{5mm} ak\ x \in \langle0, \frac{1}{2}\rangle \\
                8x-2 & \hspace{5mm} ak\ x \in ( \frac{1}{2},1\rangle \\
                16x-10 & \hspace{5mm} ak\ x \notin \langle0,1\rangle , \\
                \end{array} \right.
\end{align*}
\begin{align*} 
\boldsymbol{5.}\ f(f(x))= f(x), f(g(x))= \left\{ \begin{array}{cc} 
                0 & \hspace{5mm} ak\ x^{2} \notin \mathbb{Q} \\
                1 & \hspace{5mm} ak\ x^{2} \in \mathbb{Q}-\lbrace0\rbrace \\
                \end{array} \right.
\end{align*}
$ g(f(x))=1, \, x \in \mathbb{Q}, \, g(g(x))=x^{4}, \, x \ne 0 $;
\begin{align*} 
  f(f(x))= \left\{ \begin{array}{cc} 
                4x & \hspace{5mm} ak\ x \in \langle0, \frac{1}{2}\rangle \\
                8x-2 & \hspace{5mm} ak\ x \in ( \frac{1}{2},1\rangle \\
                16x-10 & \hspace{5mm} ak\ x \notin \langle0,1\rangle , \\
                \end{array} \right.
\end{align*}
\begin{align*} 
\boldsymbol{6.}\ f(f(x))= 0, \, f(g(x))= \left\{ \begin{array}{cc} 
                1 & \hspace{5mm} ak\ x \in \langle -2, -\sqrt{3} ) \cup (-1,1) \cup ( \sqrt{3, 2}\rangle \\
                0 & \hspace{5mm} ak\ x \in ( - \infty, -2) \cup \langle - \sqrt{3},2\rangle \cup \langle1, \sqrt{3}\rangle \cup (2, \infty), \\
                \end{array} \right.
\end{align*}, 
\begin{align*} 
g(f(x))= \left\{ \begin{array}{cc} 
                -2 & \hspace{5mm} ak\ \vert x \vert\leq 1  \\
                -1 & \hspace{5mm} ak\ \vert x \vert > 1 , \\
                \end{array} \right.
\end{align*}, 
\begin{align*} 
g(g(x))= \left\{ \begin{array}{cc} 
                x^{4}-4x^{2}+2 & \hspace{5mm} ak\ \vert x \vert\leq 2  \\
                -1 & \hspace{5mm} ak\ \vert x \vert > 2 , \\
                \end{array} \right.
\end{align*} 
$\boxed{25.}$ 
$\boldsymbol{1.}$
$ f_{n}(x)= \frac{x}{\sqrt{1+nx^{2}}} $;
$\boldsymbol{2.}$
$ f_{n}(x)= a. \frac{b^{n}-1}{b-1}+b^{n}x \quad pre \quad b \ne 1$ (použite vzťah $ 1+b+.....+b^{n-1}=\frac{b^{n}-1}{b-1} $ ),  $ f_{n}(x)= na+x \, \quad pre \quad b=1 $;
$\boldsymbol{3.}$
$ f_{n}(x)= \frac{x}{a. \frac{b^{n}-1}{b-1}x+b^{n}}, \quad D(f_{n}) = \mathbb{R}- \lbrace - \frac{b^{k}}{a. \frac{b^{n}-1}{b-1}}\, , k=1,.....,n\rbrace $ ( všetky uvedené predpisy sa dokážu matematickou indukciou );
$\boxed{26.}$ 
$\boldsymbol{1.}$
nerovnajú sa $ D(f) = (-\infty, 0\rangle \cup \langle 1, \infty) \, \ne \, \langle 1, \infty) = D(g), \quad A=\langle 1, \infty) $;
$\boldsymbol{2.}$
nerovnajú sa $ A = (2, \infty) $;
$\boldsymbol{3.}$
nerovnajú sa $ A = \mathbb{R}- \lbrace \frac{k \pi}{2}; \, k \in \mathbb{Z}\rbrace $;
$\boldsymbol{4.}$
rovnajú sa $ D(f)=D(g)= \mathbb{R} $, zlomok $ \vert \frac{\sqrt{x^{2}+1}-x}{\sqrt{x^{2}+1}+x}\vert  $ treba rozšíriť výrazom $ \vert\sqrt{x^{2}+1}+x\vert $);
$\boldsymbol{5.}$
rovnajú sa ($ f(x) = \vert x+2 \vert - \vert x-4 \vert $),
$\boxed{27.}$ 
$\boldsymbol{7.}$
$ y=-1+ \frac{2}{1-x} $;
$\boldsymbol{12.}$
$ y= \frac{2}{3}( \frac{1}{8})^{x}+2 $;
$\boxed{29.}$ 
$\boldsymbol{1.}$
$ g(x)=f(2x_{0}-x) $;
$\boldsymbol{2.}$
$ g(x)=2y_{0}-f(x) $;
$\boxed{31.}$ 
$\boldsymbol{2.}$
$ y=  \sqrt{2}\, \sin (\frac{\pi}{4}-x)=-\sqrt{2}\, \sin (x-\frac{\pi}{4}) $;
$\boldsymbol{2.}$
$ y=10 \, \sin (x+ \alpha) $, kde $ \alpha $ je určená podmienkami $ \alpha \in \langle 0, \frac{\pi}{2}\rangle, \,  \sin \alpha = 0,6 $;
$\boxed{32.}$ 
$ g(x)=f(x-a)-2a^{3} $;
$\boxed{35.}$ 
$\boldsymbol{1.}$
pretože $ \vert \sin \frac{1}{x}  \vert \leq 1 $ pre každé $ x \in (0,1) $, je $ f $ ohraničená na 
$ (0,1) $;
$\boldsymbol{2.}$
pre všetky $ x \geq 2$ platí $ 0  < \frac{1}{4+x^{2}  < \frac{1}{8}} $, teda $ f $ je ohraničená na $ \langle2, \infty) $;
$\boldsymbol{3.}$
pre každé $ x \in (1,2 \rangle $ je $ \frac{x+5}{x-1} >0 $, tj. $ f $ je zdola ohraničená na $ (1,2 \rangle  $; $ f $ nie je zhora ohraničená na $ (1,2 \rangle  $  ( nerovnica $ \frac{x+5}{x-1} >K $ má v intervale $ (1,2 \rangle  $ riešenie pre každé $ K \in \mathbb{R} $;
$\boxed{36.}$
$\exists x \in B: f(x) > f(x_{0})$;
$\boxed{38.}$
ak $ \alpha:= sup_{x \in \mathbb{R}} \, f(x), \quad \beta :=sup_{x \in \mathbb{R}} \, g(x) $, tak pre všetky $ x \in \mathbb{R} $ je $ f(x) \leq \alpha $, $ g(x) \leq \beta $, odtiaľ $ f(x)+g(x) \leq \alpha + \beta $; teda $ \alpha + \beta $ je horné ohraničenie množiny $ (f+g)(\mathbb{R}) $; pretože $ \gamma := sup_{x \in \mathbb{R}}  (f(x)+g(x))$ je jej najmenším horným ohraničením, musí platiť  $ \gamma\leq\alpha+\beta $ (ak si uvedomíme, že tento dôkaz je založený na rovnakých myšlienkach ako dôkazy z pr. 19 a 21, môžeme pomocou nich naše riešenie zapísať aj takto: Označme $ A := (f+g)(\mathbb{R}), \, B:= \, f(\mathbb{R})+g(\mathbb{R}) $, potom $ A \subset B $. Podľa pr. 19 $ sup \, A \leq sup \, B $ a podľa pr. 21 $ sup \, B  = sup \, f( \mathbb{R}) + g( \mathbb{R}) $ teda $ sup_{x \in \mathbb{R}} (f(x)+g(x)) = sup \,A \leq  sup \,B =  sup_{x \in \mathbb{R}}f (x) +  sup_{\pi \in \mathbb{R}} g(x).) $;
$\boxed{40.}$
$\boldsymbol{1.}$
$ f(x)= \sqrt{2}\sin (x+\frac{x}{4}) $ (pozri pr. 31), teda $ f $ je rastúca na každej z  množín$ \langle -\frac{3 \pi}{4}+2k \pi, \frac{\pi}{4}+2k \pi \rangle ,\, k \in \mathbb{Z} $ ale nie je rastúca na zjednotení týchto množín), $ f $ je klesajúca na každej z  množín $ \langle \frac{ \pi}{4}+2k \pi, \frac{5 \pi}{4}+2k \pi \rangle ,\, k \in \mathbb{Z} $  ale nie je klesajúca na zjednotení týchto množín);
$\boldsymbol{2.}$
$ f(x) =( \cos^{2} x+ \sin^{2} x) ^{2} -2 \sin^{2} \cos^{2} x = 1- \frac{1}{2} \sin^{2} 2x = 1- \frac{1}{2}( \frac{1- \cos 4x}{2}) $; $ f $ rastie na každej z množín $ \langle (2k-1) \pi, 2k \pi\rangle , \, k\in \mathbb{Z} $, klesá  na každej z množín $ \langle 2k \pi, (2k+1) \pi \rangle   \, k\in \mathbb{Z} $;
$\boxed{41.}$
dokazovanie sa zjednoduší, ak si uvedomíme, že postupvosť $ \lbrace x_{n}\rbrace_{n=1}^\infty $  je rastúca práve vtedy, keď platí $ \forall n \in \mathbb{N}: \, x_{n}< x_{n+1} $; analogické tvrdenia platia pre všetky monotónne postupnosti;
$\boxed{42.}$
$\boldsymbol{1.}$
$ \exists x,y \in M: \, x<y \, \land \, f(x)\leq f(y) $;
$\boldsymbol{2.}$
$( \exists x,y \in M: \, x<y \, \land \, f(x)<f(y)) \,\land \, ( \exists x,y \in M: \, x<y \, \land \, f(x)>f(y))  $ ( funkcia nie je monotónna práve vtedy, keď nie je nerastúca a nie je neklesajúca - dokážte!);
$\boxed{45.}$
$\boldsymbol{1.}$
rastúca;
$\boldsymbol{2.}$
klesajúca;
$\boldsymbol{3.}$
klesajúca;
$\boldsymbol{4.}$
rastúca;
$\boxed{46.}$
$\boldsymbol{1.}$
mnemotechnická pomôcka: priraďte rastúcej funkcii číslo 1, klesajúcej číslo -1, potom číslo priradené superpozícii rýdzomonotónnych funkcií je súčinom čísel priradených k jednotlivým zložkám; 
$\boldsymbol{2a/}$
$ x^{2}-6x+10=(x-3)^{2}+1  $ klesá na $ (-\infty,3 \rangle $, rastie na $ \langle3, \infty) $, $log u $ rastie na   $(0, \infty) $, teda podľa príkladu 46.1 $ f $ klesá na $ (-\infty,3 \rangle $, rastie na $ \langle3, \infty) $;
$\boldsymbol{2b/}$
$ f $ rastie na $ (-\infty,4 \rangle $, klesá na  $ \langle4, \infty) $;
$\boxed{48.}$
$\boldsymbol{2.}$
sporom, keby $ a \notin \mathbb{Q}, a>0 $ bola perióda funkcie $ \chi $, muselo by platiť $ 0=\chi(1-a)= \chi((1-a)+a)=1 $;
$\boxed{49.}$
$ g(x+ \frac{T}{a})=f(ax+b+T)=f(ax+b)=g(x), \, x \in \mathbb{R} $ tada $ \frac{T}{a} $ je perióda; keby nebola najmenšia, tak zo vzťahu $f(x)=g(\frac{x}{a}-\frac{b}{a}) $  vyplýva, že $ f $ má niektorú periódu menšiu ako $ T $, čo je spor;
$\boxed{51.}$
$\boldsymbol{1.}$
$ f(x+2T)=-f(x+T)=f(x), \, x\in \mathbb{R} $,  teda $ 2T $ je perióda;
$\boldsymbol{2.}$
$ f(x+2T)=f(x) $;
$\boldsymbol{3.}$
$ f(x+2T)=f(x) $;
$\boldsymbol{4.}$
$ f(x+3T)=f(x) $;
$\boxed{52.}$
ak pre začiatok pomôžeme svojej predstavivosti obrázkom, zistíme, že $ f $ musí mať periódu $ 2(b-a) $; dôkaz: podľa pr. 29.1 pre všetky $ x \in \mathbb{R} $  platí $ f(x)=f(2a-x) $ a $ f(x)=f(2b-x) $, ak na vyjadrenie $ f(2a-x) $ použijeme druhú z týchto rovností, dostaneme $ f(x)=f(2b-(2a-x))=f(x+2(b-a)), \, x \in \mathbb{R} $
$\boxed{53.}$
$\boldsymbol{1.}$
$ f^{-1}(x)= \frac{1-x}{1+x}; \, x \in \mathbb{R}-\lbrace-1\rbrace $;
$\boldsymbol{2.}$
$ f $ nie je prostá (rovnica $ f(y)=x $ má dve rôzne riešenia pre každé $ x \in (- \frac{1}{\sqrt{8}}, \frac{1}{\sqrt{8}}) - \lbrace 0 \rbrace $;
$\boldsymbol{4.}$
$ f^{-1}(x)=\frac{x}{1+\sqrt{1-x^{2}}}, \, x \in \langle-1,1\rangle $, 
ak $ x \in \mathbb{R} $ je dané a má platnosť $ f(y)=x $, tak $ y $ je riešenie rovnice $ ay^{2}-2y+x=0; \, pre x=0 $ je 
$ y=0 $, pre $ x \in \langle -1,1 \rangle - \lbrace 0 \rbrace $ sú jej riešeniami $ y_{1}= \frac{1+\sqrt{1-x^{2}}}{x} $ a  $ y_{2}= \frac{1-\sqrt{1-x^{2}}}{x} $, pre $ \vert x\vert > 1$ riešenie neexistuje; podmienka $ \vert y_{2} \vert \leq 1 $ je splnená pre všetky 
$ x \in \langle -1,1 \rangle - \lbrace 0 \rbrace $, $ \vert y_{1} \vert \leq 1  $ len pre 
$ x=1,-1 $ ale  pre  $ x=1,-1 $ je $ y _{1} = y_{2} $;
teda $ f^{-1}(x)=\frac{1-\sqrt{1-x^{2}}}{x} $ pre $ x \in \langle -1,1 \rangle - \lbrace 0 \rbrace $, $ f^{-1}(0)=0 $; to možno napísať v spoločnom tvare $ f^{-1}(x) = \frac{x}{1+ \sqrt{1-x^{2}}}, \, \vert x \vert \leq 1 $;
$\boldsymbol{5.}$
$ f^{-1}(x) = 1+\frac{1}{log_{3}x -1},x \ne 3 \,  $;
$\boldsymbol{6.}$
$ f^{-1}(x)=1- \sqrt{1+log_{2}x}, \, x\geq \frac{1}{2} $;
$\boldsymbol{7.}$
$ f^{-1}(x)= \frac{1}{2} log \frac{x}{2-x}, \, x \in (0,2) $;
$\boldsymbol{8.}$
$ ^{-1} (x)= \frac{1}{2}ln (x^{2}-2x-2), \, x>1+ \sqrt{3} $ (rovnica $ x=1+ \sqrt{3+e^{2y}} $ má riešenie len pre $ x>1+ \sqrt{3}) $;
$\boldsymbol{9.}$
$ f^{-1} (x)=10^{1/x}$, $x\ne 0  $ (platí $ log_{x}10 = \frac{1}{log x} $; 
$\boldsymbol{10.}$
$ f^{-1} (x)=2 + e^{2 log_{2}\frac{x}{2}}$, $x> 0  $;
$\boldsymbol{11.}$
$ f^{-1} (x)=x + \sqrt{x^{2}+1}$, $x \in \mathbb{R}  $;
$\boldsymbol{12.}$
$ f^{-1} (x)=x - \sqrt{x^{2}+1}$, $x \in \mathbb{R}  $;
$\boldsymbol{13.}$
$ f^{-1} (x)= \frac{1}{2}(x+ \frac{1}{x})$, $x \in \langle-1,0) \cup \langle1, \infty)  $; ( rovnicu $ x = y + \sqrt{y^{2}-1} $ vynásobte výrazom $ y -  \sqrt{y^{2}-1} $, potom možno vyjadriť $ y= \frac{1}{2}(x+ \frac{1}{x}) $, pri skúške správnosti zistíme, že vyhovujú len tie riešenia, pre ktoré je $ x - \frac{1}{x}\geq 0  $;
$\boldsymbol{14.}$
$ f^{-1} (x)= \frac{1}{2}( 2^{x}- 2^{-x})$, $x \in \mathbb{R}  $;
$\boldsymbol{15.}$
\begin{align*} 
f^{-1}(x)= \left\{ \begin{array}{cc} 
                x & \hspace{5mm} ak\ x < 0 \\
                 \frac{x}{2} & \hspace{5mm} ak\ x \geq 0  \\
                \end{array} \right.
\end{align*}

$\boldsymbol{16.}$
\begin{align*}
f^{-1}(x)= \left\{ \begin{array}{cc} 
                x & \hspace{5mm} ak\ x  \in  \mathbb{Q} \\
                 1-x & \hspace{5mm} ak\ x \notin  \mathbb{Q}  \\
                \end{array} \right.
\end{align*}
$\boldsymbol{17.}$
\begin{align*}
f^{-1}(x)= \left\{ \begin{array}{cc} 
                x +1& \hspace{5mm} ak\ x < 0 \\
                  \sqrt{x}& \hspace{5mm} ak\ x \in  \langle1,16\rangle  \\
                  log _{2} x& \hspace{5mm} ak\ x> 16 
                \end{array} \right.
\end{align*}
$\boxed{54.}$ 
treba ukázať, že platí rovnosť $ f^{-1}=f $, (graf funkcie $ f $ je symetrický podľa osi $ y = x $ s grafom funkcie  $ f^{-1} $) ;
$\boxed{55.}$ 
$\boldsymbol{2.}$
napr. funkcia z pr. 53.16 ;
$\boxed{56.}$
neplatí ( vhodný protipríklad už musíte nájsť sami );
$\boxed{54.}$
$\boldsymbol{1.}$
$ x/3 $;
$\boldsymbol{2.}$
$  \frac{1}{\sqrt{2}} $;
$\boldsymbol{3.}$
$  \frac{x}{3} $;
$\boldsymbol{4.}$
$  \frac{x}{6} $;
$\boldsymbol{5.}$
$  \sqrt{3} $;
$\boxed{58.}$
$\boldsymbol{1a/}$
$ \arcsin x+2 \pi $;
$\boldsymbol{1b/}$
$  \pi - \arcsin x $;
$\boxed{59.}$
$\boldsymbol{1.}$
$ f^{-1}(x)= \arcsin \sqrt[3]{x} $, $ \vert x \vert \leq1 $;
$\boldsymbol{2.}$
$ f^{-1}(x)= \pi - \arcsin \sqrt[3]{x} $, $ \vert x \vert \leq1 $ ( ak je $ x $ dané a má platiť $ f(y)=x $, tak postupne dostávame $ x = \sin^{3} y $, $  \sqrt[3]{x } = \sin y $, $ \arcsin \sqrt[3]{x} = \arcsin (\sin y) = \pi -y $);
$\boldsymbol{3.}$
$ f^{-1}(x)= 2( \pi + \arcsin \sqrt{x}) $, $  x \in \langle0,1\rangle $  (postupne $ x = \sin ^{2} \frac{y}{2} $, $ \sqrt{x}= \vert \sin \frac{y}{2}\vert = - \sin \frac{y}{2} $, $ \arcsin \sqrt{x}= \arcsin (- \sin \frac{y}{2})= - \arcsin ( \sin \frac{y}{2}) = - ( \pi - \frac{y}{2}))$;
$\boldsymbol{4.}$
$ f^{-1}(x)= \arctan x- \pi $ pre $ x\geq 0 $, $ f^{-1} (x)= \arctan x$ pre $ x<0 $;
$\boldsymbol{5.}$
$f^{-1}(x)= \cot x $, $0<\vert x \vert < \frac{ \pi}{2}$;
$\boldsymbol{6.}$
$ f^{-1}(x)= \sqrt[3]{1- \sin ^{2} \frac{\pi}{4}} $, $ x \in \langle0, 2 \pi \rangle $;
$\boldsymbol{7.}$
$ f^{-1}(x)= 1+ \cos  \frac{x-3}{4} $, $ x \in \langle 3, 3+4 \pi \rangle $;
$\boldsymbol{8.}$
$ f^{-1}(x)= \tan (log _{2} \frac{\pi}{8} $, $ x \in ( 2^{3- \frac{\pi}{2}}, 2^{3+ \frac{\pi}{2}}) $;
$\boldsymbol{9.}$
$ f^{-1} (x) = \frac{1+ \arcsin x}{3}$, $ \vert x \vert<1  $ ($ D (f^{-1}) $ sa v týchto príkladoch nájde ak $ f(D(f)) $);
$\boxed{60.}$
$\boldsymbol{1.}$
$ \sqrt{1-x^{2}} $;
$\boldsymbol{2.}$
$ \sqrt{1-x^{2}} $ (treba použiť vzorec $ \sin 2 \alpha = 2 \sin \alpha \cos \alpha $);
$\boldsymbol{3.}$
$ \frac{1}{1+x^{2}} $ (pre $ u \in (-\frac{\pi}{2},\frac{\pi}{2} ) $ platí pre $ \cos ^{2 }u = \frac{1}{1+\tan^{2}u} $);
$\boldsymbol{4.}$
$ \frac{1}{1+x^{2}} $;
$\boldsymbol{5.}$
$ \frac{x}{\sqrt{1+x^{2}}} $;
$\boldsymbol{6.}$
$ \frac{3x-x^{3}}{1-3x^{2}} $, $ x \ne \frac{1}{\sqrt{3}} $;
$\boxed{61.}$
$\boldsymbol{2.}$
$ x \in \langle-1,0\rangle $ (funkčné hodnoty obidvoch strán ležia v $ \langle - \frac{\pi}{2},\frac{\pi}{2}\rangle  $, možno využiť injektívnosť funkcie $ \sin/\langle - \frac{\pi}{2},\frac{\pi}{2}\rangle ) $;
$\boldsymbol{3.}$
$ x\geq 0 $ pre $ x<0 $ je $ \arccos \frac{1-x^{2}}{1+x^{2}}  \in (0,\pi\rangle $, $ 2 \arctan x \in (-\pi,0) $; pre $ x\geq 0 $ možno využiť injektívnosť funkcie $ \cos/\langle0,\pi) $ a vzťah $ \cos^{2} \alpha = \frac{1}{1+ \tan^{2}\alpha} $, $ \alpha \ne k \pi, \, k \in \mathbb{Z} $;
$\boldsymbol{4.}$
$ x \in \mathbb{R} $;
$\boldsymbol{5.}$
$ x > 1 $;
$\boxed{62.}$
$\boldsymbol{1.}$
čísla $ \arccos x, \, \frac{\pi}{2}-\arcsin x $ ležia v $ \langle0, \pi\rangle $, stačí využiť injektívnosť funkcie $ \cos/\langle0,\pi) $;
$\boldsymbol{1.}$
pre $ x>0 $ použite injektívnosť funkcie $ \cos/\langle0,\pi) $, pre $ x<0 $  injektívnosť funkcie $ \cos/( -\pi,0\rangle $;
$\boxed{64.}$
$\boldsymbol{1.}$
$ sh \, x\,\, ch\, y + sh \,y\,\, ch\, x $;
$\boldsymbol{2.}$
$ ch \, x\,\, ch\, y + sh \,y\,\, sh\, x $;
$\boldsymbol{3.}$
$ \sgn x. \sqrt{\frac{ch\, x -1}{2}} $ (stačí sčítať rovnosti 63.1 a 63.2, kam dosadíme $ \frac{x}{2} $ namiesto $ x $);
$\boldsymbol{4.}$
$ 2 \, ch \, \frac{x+y}{2}. ch\, \frac{x-y}{2}$ použite rovnosti $ x= \frac{x+y}{2}+ \frac{x-y}{2} $,  $ y= \frac{x+y}{2}- \frac{x-y}{2} $ a  vzorec pre $ ch \, ( \alpha\pm \beta) $);
$\boldsymbol{5.}$
$ 2 \, sh \, \frac{x+y}{2}. ch\, \frac{x-y}{2}$;
$\boxed{65.}$
platí
$\boxed{66.}$
neplatí; platí teda negácia uvedeného tvrdenia (na overenie jej platnosti stačí položiť $ A = \mathbb{R} $, $ B= \lbrace0\rbrace $);
$\boxed{68.}$
podľa (*) je každé $ x\in A $ dolné ohraničenie množiny $ B $, každé $ y \in B $ horné ohraničenie  množiny $ A $, preto $ A $ je zhora , $ B $ zdola ohraničená a platí $ \forall x \in A: \, x\leq \inf B $; teda $ \inf B $ je horné ohraničenie $ A $, odtiaľ $ \sup A\leq \inf B $; ďalej postupujeme sporom: ak $ \sup A > \inf B $, položme $ \varepsilon_{1}:= \sup A - \inf B $  zrejme , $ \varepsilon_{1}>0 $, z toho odvodíme spor s (**);
$\boxed{69.}$ 
ak $ k\geq 0 $, tak $ \sup kA = k \sup A $, $ \inf kA = k \inf A $; ak $ k\geq 0 $, tak $ \sup kA = k \inf A $, $ \inf kA = k \sup A $; (tieto tvrdenia treba ovšem dokázať, pritom dôkaz pre $ k<0 $ možno vykonať pomocou tvrdenia z pr. 20 a už dokázaných vzťahov pre $ k\geq 0 $, pretože $ kA = -((-k).A)) $;
$\boxed{71.}$
platí (treba dokázať, že číslo $ \sup _{x\in I} f(x) - \inf _{x\in I} f(x)  $ má obidve vlastnosti, ktoré charakterizujú suprémum množiny $ \lbrace \vert f(x)-f(y)\vert ;x,y \in I\rbrace $; iná možnosť: dokázať rovnosť $ \sup \lbrace \vert f(x)-f(y)\vert ;x,y \in I\rbrace =  \lbrace f(x)-f(y) ;x,y \in I\rbrace $, potom využiť rovnosť $ \sup (A-B) = \sup A - \inf B $);
$\boxed{72.}$
$\boldsymbol{2.}$
tento vzťah možno odvodiť z 72.1 rovnako, ako sa odvodzuje nerovnosť $ \vert \vert a \vert - \vert b \vert\vert \leq \vert a-b \vert $ z nerovnosti $ \vert a+b \vert \leq \vert a \vert + \vert b \vert $;
$\boxed{73.}$
napr. $ f(x) = 0 $ pre $ x \notin \mathbb{Q} $, $ f(x) = m $, ak $ x=\frac{n}{m} $, kde $ m \in \mathbb{N} $ a $ n \in \mathbb{Z} $ sú nesúdeliteľné; resp. $ g(x)=0 $ pre $ x \notin \mathbb{Q} $, $ g(x)=(-1)^{m}.m $, ak $ x= \frac{n}{m} $;
$\boxed{76.}$
ak $ T $ je perióda, tak $ 1= \sin 0 + \cos a.0 = \sin T + \cos aT = \sin (-T) + \cos a.(-T) = -\sin T + \cos aT $; odtiaľ $ \sin T = 0 $, tj. $ T $  je v tvare $ k \pi\ \,(k \in \mathbb{N})$; $ \cos aT = 1$, tj. $ aT = 2n \pi \, (n \in \mathbb{Z}) $, odtiaľ $ a=\frac{2n}{k} \in \mathbb{Q} $;
$\boxed{77.}$
neexistuje ; stačí dokázať, že funkcia, ktorej periódou je každé kladné iracionálne číslo, je konštantná (porovnaj pr. 77 a pr. 48);
$\boxed{79.}$
pozri pr. 29 a 52;
$\boxed{80.}$
$\boldsymbol{1.}$
$ (-1,1) $;
$\boldsymbol{2.}$
$ \bigcup\limits_{k \in \mathbb{Z}} (2k \pi , (2k+1)\pi) $;
$\boldsymbol{3.}$
$ (1,e) $;
$\boxed{81.}$
$\boldsymbol{1.}$
$ f(t) = t^{2} -5t+6 $ (ak je  dané $ t \in \mathbb{R} $ a chceme nájsť $ f(t) $, treba voliť $ x $ tak, aby platilo $ x+1=t $);
$\boldsymbol{2.}$
$ f(t)=t^{2}-2 $ pre $ \vert t \vert \geq 2 $, $ f(t) $ je ľubovoľne určené pre $ \vert t \vert < 2 $ (pri úpravách využite rovnosť  $ \frac{t+ \sqrt{t^{2}-4}}{2}= \frac{2}{t - \sqrt{t^{2}-4}}$);
$\boldsymbol{3.}$
$ f(t)=\frac{1}{t}+ \sqrt{1+ \frac{1}{t^{2}}}$,  pre  $t\ne 0$, $ f(0) $ je určené ľubovoľne;
$\boldsymbol{4.}$
$ f(t)=\frac{t^{2}}{(1-t)^{2}} $ pre $ t \ne 1 $, $ f(1) $ je určené ľubovoľne;
$\boxed{82.}$
nie;
$\boxed{83.}$
$\boldsymbol{1.}$
áno, napr. $ f(x) = \chi (x) $, $ g(x) = \chi (x)+1 $;
$\boldsymbol{2.}$
áno, pre $ g(\mathbb{R}) $ musí platiť: ak $ I \subset \mathbb{R} $ je interval, tak $ I\not\subset g(\mathbb{R} ) $; zoraďme spočitateľnú množinu $ \mathbb{Q} $ (definíciu spočitateľnosti pozri v Dodatku 1; k spočitateľnosti  množiny $ \mathbb{Q} $ pozri pr. 94) do prostej postupnosti $ \lbrace a _{n}\rbrace_{n=1}^\infty  $, potom funkcie definované predpismi $ g(x)=x $ pre $ x \notin \mathbb{Q} $, $ g(x)=n $ pre $ x=a_{n} $; $ f(x)=0 $ pre $ x \in \mathbb{N} \cup (\mathbb{R}-\mathbb{Q}) $, $ f(x)=1 $ pre $x \in \mathbb{Q} - \mathbb{N} $ vyhovujú zadaniu;
 $\boxed{84.}$
$\boldsymbol{1.}$
$ f(x)=3(\cos x - \frac{\sqrt{3}}{2})^{2}+ \frac{7}{4} $, funkcia $ g(u)=3(u - \frac{\sqrt{3}}{2})^{2}+ \frac{7}{4} $ má na $ \langle-1,1\rangle $ minimum v bode $u =\frac{\sqrt{3}}{2} $ maximum v bode $ u=-1 $; teda $ \frac{7}{4}= \min_{x\in \mathbb{R}} f(x) $ sa nadobúda pre $ x=\pm \frac{\pi}{6}+2k \pi \, k\in \mathbb{Z} $, $ 7+3\sqrt{3}= \max_{x\in \mathbb{R}} f(x) $ sa nadobúda pre $ x=(2k-1)\pi \, (k\in \mathbb{Z}) $; 
$\boldsymbol{2.}$
$ 1=3^{0}=\min_{x\in \mathbb{R}} f(x) $ pre $ x=0 $, $ \max_{x\in \mathbb{R}} f(x) $ neexistuje $ g(u)=3^{0} $ je rastúca a zhora neohraničená na $ \langle0,\infty) $);
$\boldsymbol{3.}$
$ \sqrt{a^{2}+b^{2}}=\max_{x\in \mathbb{R}}f(x) $ pre $ x=\alpha + \frac{\pi}{2}+2k \pi \, (k\in \mathbb{Z}) $, $ -\sqrt{a^{2}+b^{2}}=\min_{x\in \mathbb{R}}f(x) $ pre $ x=\alpha - \frac{\pi}{2}+2k \pi \, (k\in \mathbb{Z}) $, kde $ \alpha $ je určená podmienkami $ \sin \alpha = - \frac{a}{\sqrt{a^{2}+b^{2}}} $, $ \cos \alpha =  \frac{b}{\sqrt{a^{2}+b^{2}}} $ (pozri pr. 31);
$\boldsymbol{4.}$
$ f(x)=\frac{1}{\sqrt{2}} \sin (2x+\frac{\pi}{4})+\frac{1}{2} $; $ \max_{x\in \mathbb{R}}f(x)= \frac{\sqrt{2}+1}{2} $, $ \min_{x\in \mathbb{R}}f(x)=\frac{1-\sqrt{2}}{2} $;
$\boxed{85.}$
(pozri pr. 46) f rastie na $ (-\infty,-1\rangle $, klesá na $ \langle-1,0\rangle $, rastie na $ \langle0,1\rangle $, klesá na $ \langle1,\infty) $;
$\boxed{87.}$
$\boldsymbol{1.}$
pre $ x>0 $, $ y>0 $ využite injektívnosť $ cos/(0,\pi) $, pre $ x<0 $, $ y<0 $ injektívnosť $ cos/(-\pi,0) $, pre $ xy\leq 0 $ injektívnosť $ sin/(-\frac{\pi}{2},\frac{\pi}{2}) $, pritom tiež zistíte, že $ \varepsilon =0 $ pre $ xy<1 $, $ \varepsilon =1 $ pre $x>0 \, \land \, y > 0 \, \land \, xy>1 $, $ \varepsilon =-1 $ pre $x<0 \, \land \, y < 0 \, \land \, xy>1 $;
$\boxed{88.}$
$ f^{-1}(x)=\frac{1-2\pi+\arcsin(y-1)/2}{1+2\pi+\arcsin(y-1)/2} $, $ x \in \langle-1,3\rangle $;
$\boxed{89.}$
platí;



\section{Dodatok 1}
$\boxed{90.}$
$ \mathbb{Z} $ možno zoradiť do postupnosti $ 0,1,-1,2,-2,3,-3,..... $;
$\boxed{91.}$
zoraďme $ B $ do prostej postupnosti $ b_{1}, b_{2}..... $; "vyškrtaním" prnkov nepatriacich do $ A $ dostaneme prostú postupnosť, ktorá obsahuje všetky prvky z $ A $; pri konkrétnom zápise tejto "voľnej" úvahy využite princíp konštrukcie postupnosti matematickou indukciou: a/ určí sa prvý člen (nech $ k_{1}:= min \lbrace n \in \mathbb{N} ;  b_{n} \in A \rbrace $ , potom $ a_{1}=b_{k_{1}} $); b/ udá sa spôsob, akým určiť člen $ a_{n+1} $, ak sú už určené členy $ a_{1},...,a_{n} $ (ak $ a_{n}=b_{k_{n}} $, potom $ b_{n+1}:= \min \lbrace n \in \mathbb{N}; n>k_{n} \, \land b_{n} \in A \rbrace $, potom $ a_{n+1}:=b_{k_{n+1}} $);
$\boxed{92.}$
pre $ A\cap B=\varnothing $: všimnite si, ako sme "zoradili" množiny $ \lbrace1,2,3,...\rbrace $ a $\lbrace 0,-1,-2,...\rbrace $ v pr. 90; pre $ A \cap B \ne \varnothing $ použite ten istý postup na $ A $ a $ A-B $ (pozor  $ A-B $ môže byť aj nekonečná);
$\boxed{93.}$
ak $ A_{i}\cap B_{i}=\varnothing $ pre $ i \ne j $, $ i,j \in \mathbb{N} $, zoraďte $ A $ podľa schémy

\newpage


vo všeobecnom prípade využite rovnosť $ A=A_{1} \cup (A_{2}-A_{1}) \cup (A_{3}- (A_{1}\cup A_{2})... $;
$\boxed{94.}$
využite pr. 92 a 93 $ ( \mathbb{Q}-\mathbb{Q^{-} \cup \mathbb{Q^{+}_{O}}} $, kde $\mathbb{Q^{+}_{O}} =\bigcup\limits_{n \in \mathbb{R}} Q_{n} $, $  Q_{n}=\lbrace \frac{m}{n}; \, m\in \lbrace 0 \rbrace \cup \mathbb{N}\rbrace )$; 
$\boxed{95.}$
nepriamo, ak $ E-\langle-c,c\rangle $ je spočitateľná pre každé $ c $, tak $ E= \bigcup\limits_{n \in \mathbb{R}} (\langle-n,n\rangle \cap E)$  je spočitateľná;
$\boxed{96.}$
nepriamo, ak $ E_{1} $ je nespočitateľná, tak existuje $ n \in \mathbb{N} $ tak, že $ E_{1} \cap ( \frac{1}{n}, \infty) $ je nespočitateľná, teda nekonečná; potom $ S $ nie je zhora ohraničená;
$\boxed{97.}$
áno; $ A:= \lbrace x-y;  x,y \in E \rbrace $ je spočitateľná; ak $ a \in \mathbb{R}-A $   ($ \mathbb{R}-A $ je neprázdna množina, lebo $ \mathbb{R} $ je nespočitateľná), tak $ E+a \cap E = \varnothing $;
$\boxed{99.}$
z každého intervalu z $ \varphi $ vyberme jedno racionálne číslo (zrejme čísla vybrané z rôznych intervalov sú navzájom rôzne), takto vytvorená množina čísel je podmnožinou množiny $ \mathbb{Q} $, preto je spočitateľná;

\section{Limita funkcie}
$\boxed{100.}$
$\boldsymbol{1.}$
$ A=\lbrace 1/k \pi; \, k \in \mathbb{Z}- \lbrace 0 \rbrace \rbrace $, treba dokázať $ \forall O_{\varepsilon} (0) $ $ \exists x \in A; x \ne 0 \land x \in O_{\varepsilon} (0) $, čo možno zapísať v tvare $ \forall \varepsilon > 0 $ $ \exists k \in \mathbb{Z}- \lbrace0\rbrace : \vert 1/k \pi \vert < \varepsilon $, existencia riešenia poslednej nerovnosti vyplýva z Archimedovho princípu;
$\boldsymbol{3.}$
využite, že do $ A $ patria čísla $ 0,1;0,11;0,111...... $;
$\boxed{101.}$
$\boldsymbol{1.}$
$ A^{)} \langle0,1\rangle $ (nestačí len dokázať, že každý prvok z $ \langle0,1\rangle $ je hromadným bodom množiny $ A $, treba aj dokázať, že žiadny prvok z $ \mathbb{R}^{\chi} - \langle0,1\rangle $ už nepatrí do $ A^{)} $);
$\boldsymbol{2.}$
$ B^{)}= \lbrace0,\infty \rbrace $;
$\boldsymbol{3.}$
$ C ^{)}=\langle-2,+ \infty) \cup \lbrace+\infty\rbrace $;
$\boldsymbol{4.}$
$ D- \mathbb{Q} \cap (0,1) $, $ D^{)} = \langle0,1\rangle$ (pri dôkaze poslednej rovnosti využite fakt, že (nedegenerovaný) interval obsahuje aspoň jedno racionálne číslo);
$\boldsymbol{5.}$
$ E= \mathbb{Q}-\lbrace1\rbrace $, $ E^{)}=\mathbb{R}^{\chi} $;
$\boldsymbol{6.}$
$ F^{)}=\langle1,2\rangle $;
$\boxed{102.}$
$\boldsymbol{1.}$
$ \exists O(+ \infty):M \cap (+ \infty)=\varnothing $, to možno zapísať aj v tvare $ \exists x \in \mathbb{R} \quad \forall x \in M; x\leq K $;
$\boldsymbol{2.}$
$ \forall a \in \mathbb{R}^{\chi} \exists O(a); \quad (O(a)-\lbrace a \rbrace ) \cap M= \varnothing \ $ to možno zapísať aj v tvare $ \forall a \in \mathbb{R}^{\chi} \quad  \exists O(a) \quad \forall x \in M: \quad x=a \lor x \notin O(a)  $;
$\boxed{103.}$
$\boldsymbol{1.}$
ľubovoľná konečná neprázdna množina má túto vlastnosť;
$\boldsymbol{2.}$
napr. $ A= \mathbb{N} \cup \lbrace \frac{n+1}{n}; n\in \mathbb{N} \rbrace$;
$\boldsymbol{3.}$
napr. $ A= \mathbb{Z} $;
$\boldsymbol{4.}$
na $ A $ stačí zvoliť ľubovoľný interval;
$\boldsymbol{5.}$
napr. $ A= \lbrace m+ \frac{1}{n}; \, n \in \mathbb{N} \rbrace $, potom $ A^{)}= \mathbb{N} \cup \lbrace \infty \rbrace $;
$\boxed{104.}$
pri dôkaze využite druhú vlastnosť suprema;
$\boxed{105.}$
$\boldsymbol{1.}$
nerovnosti $ \vert \frac{3n^{2}+1}{5n^{2}-1} - \frac{3}{5} \vert < \varepsilon $ (pre dané $ \varepsilon> 0 $) vyhovujú všetky tie $ n \in \mathbb{N} $, pre ktoré platí $ n > \sqrt{\frac{5\varepsilon +8}{25 \varepsilon}} $; stačí teda položiť napr. $ n_{0}= [  \sqrt{\frac{5\varepsilon +8}{25 \varepsilon}}]   $ pre $ \varepsilon \in (0, \frac{2}{5}\rangle $ a $ n_{0}= 1 $ pre $ \varepsilon > \frac{2}{5} $ (pre $ \varepsilon > \frac{2}{5} $ je totiž $ [  \sqrt{\frac{5\varepsilon +8}{25 \varepsilon}}] =0 $ keby  sme v definícii limity postupnosti namiesto $n \in \mathbb{N}  $ požadovali $ n \in \mathbb{R} $ stačilo by pre dané $ \varepsilon > 0 $ položiť $ n_{0} \geq   \frac{5\varepsilon +8}{25 \varepsilon }$); a/ pre všetky $  n \in \mathbb{N} $; b/ pre všetky $ n\geq 8 $; c/ pre všetky $ n\geq 80 $;
$\boldsymbol{3.}$
nerovnosti $ \frac{n^{2}}{n+8}>K $ vyhovujú pre $ K > 0 $ všetky čísla $ n \in \mathbb{N} $ pre $ K \geq 0 $ všetky tie $ n \in \mathbb{N} $, pre ktoré platí $ n > \frac{1}{2} (K+ \sqrt{K^{2}+32K)} $ stačí teda položiť napr. $ n_{0}=1 $ pre $ K < 0 $, $ n_{0}= [ \frac{1}{2} (K+ \sqrt{K^{2}+32K)}] +1 $ pre $ K\geq 0 $ pripočítaním čísla 1 zaručíme, že aj v prípade $ [ \frac{1}{2} (K+ \sqrt{K^{2}+32K)}]=0 $ bude platiť $ n_{0} \in \mathbb{N}; \frac{n^{2}}{n+8}> 10^{3} $ platí pre všetky $ n > 1007 $;
$\boldsymbol{4.}$
nerovnosti $ \frac{5}{n}-n < K $ vyhovujú všetky tie $ n \in \mathbb{N} $, pre ktoré $ n > \frac{1}{2} ( \sqrt{K^{2}+20} -K)  $, stačí teda pre dané $ K \in \mathbb{R} $ položiť $ n_{0} = [ \frac{1}{2} ( \sqrt{K^{2}+20}-K )]+1 $; iné riešenie : pretože  $ \frac{5}{n}-n \leq 5-n $
pre $ n \in \mathbb{N} $ je každé riešenie nerovnice $ \frac{5}{n}-n < K $, teda nerovnici $ \frac{5}{n}-n < K $ iste vyhovujú všetky tie $ n $, pre ktoré $ a> 5-K $ (podobnými únahami možno zjednodušiť aj riešenia pr. 105.1-3);
$\boldsymbol{5.}$
ak $ q=0 $ , stačí položiť $ n_{0}=1 $ pre každé $ \varepsilon  > 0 $, ak $ q\neq 0 $, $ \vert q \vert < 1 $, sú riešeniami nerovnice $ \vert q^{n} \vert < \varepsilon $ všetky tie $ n \in \mathbb{N} $, pre ktoré platí $ n  > \frac{\ln \varepsilon}{\ln \vert q \vert} $, stačí teda položiť $ n_{0}=1 $ pre $ \varepsilon\geq 1 $, $ n_{0}= [  \frac{\ln \varepsilon}{\ln \vert q \vert} ]+1 $ pre $ \varepsilon \in (0,1) $;
$\boxed{106.}$
$\boldsymbol{1.}$
$ \lim_{n \to \infty}a_n=1$; nerovnosti $ \vert a_{n} -1 \vert < \varepsilon $ vyhovujú všetky tie párne $ n \in \mathbb{N} $, pre ktoré $ n > \frac{1}{\varepsilon} $ a všetky tie nepárne $ n \in \mathbb{N} $, pre ktoré $ n > \frac{1}{\sqrt{\varepsilon}} $, pre dané $ \varepsilon > 0 $ stačí teda položiť napr. $ n_{0}= [ \max \lbrace \frac{1}{\varepsilon} , \frac{1}{\sqrt{\varepsilon}} \rbrace ]+1 $;
$\boldsymbol{2.}$
$ \lim_{n \to \infty}a_n=0$;
$\boxed{108.}$
$\boldsymbol{1.}$
neexistujú, nerovnica $\vert a_{n}  -b \vert < \varepsilon   $ nemá totiž riešenie pre $ \varepsilon < 0 $;
$\boldsymbol{2.}$
len konštantná postupnosť $ b,b,b,b,...(b \in \mathbb{R}) $;
$\boldsymbol{3.}$
na overenie pravdivosti tvrdenia  $ \exists  \, \varepsilon > 0 \quad \exists \, n_{0 } \in \mathbb{N} \quad \forall n \in \mathbb{N} , \, n> n_{0}: \vert 2-7 \vert < \varepsilon $ stačí položiť $ \varepsilon = 6 $;
$\boldsymbol{4.}$
len konštantná postupnosť $ b,b,b,b,...(b \in \mathbb{R}) $;
$\boxed{109.}$
áno; z tvrdenia $ \exists N^{*}\in \mathbb{N} \quad \forall \, \varepsilon > 0 $ $ \forall n \in \mathbb{N} , \, n> N^{*} $: $ \vert a_{n} -b \vert < \varepsilon $ vyplýva totiž tvrdenie $ \forall \varepsilon > 0 $, $ \exists n_{0 }\in \mathbb{N} $, $ \forall n \in \mathbb{N}, \, n>n_{0} $: $ \vert a_{n} -b \vert < \varepsilon $; pre každé $ \varepsilon > 0  $ stačí položiť $ n_{0} = N^{*} $ ( z tvrdenia   $ \exists N^{*}\in \mathbb{N} \quad \forall \, \varepsilon > 0 $ $ \forall n \in \mathbb{N} , \, n> N^{*} $: $ \vert a_{n} -b \vert < \varepsilon $ vyplýva aj to, že pre  $ n> N^{*}$  je $ a_{n}=b $);
$\boxed{110.}$
$\boldsymbol{1.}$
všetky ohraničené postupnosti; 
$\boldsymbol{2.}$
len postupnosť   $ 0,0,0,... $;
$\boldsymbol{3.}$
všetky  ohraničené  postupnosti (každú z uvedených odpovedí treba samozrejme podrobne zdôvodniť);
$\boxed{111.}$
$\boldsymbol{2.}$
$\forall K \in \mathbb{R} \quad \exists \delta > 0  \,  \forall x \in D(f), \, x \ne a: \, \vert x-a \vert < \delta  $  $\Rightarrow   \, f(x)  > K  $; 
$\boldsymbol{3.}$
$\forall K \in \mathbb{R} \quad \exists \delta  >0  \, \forall x \in D(f), \, x \ne a: \, \vert x-a \vert < \delta  $  $\Rightarrow   \, f(x)  < K  $;
$\boldsymbol{4.}$
$ \forall \varepsilon >0  $  $\exists L \in \mathbb{R} \quad \forall x \in D(f): \, x>L $ $\Rightarrow \, \vert f(x) - b \vert  < \varepsilon   $;
$\boldsymbol{5.}$
$ \forall K \in \mathbb{R} \quad \exists L \in \mathbb{R} \quad \forall x \in D(f): \, x>L \, \Rightarrow \, f(x)>K $;
$\boldsymbol{6.}$
$ \forall K \in \mathbb{R} \quad \exists L \in \mathbb{R} \quad \forall x \in D(f): \, x>L \, \Rightarrow \, f(x)<K $;
$\boldsymbol{7.}$
$ \forall \varepsilon >0  $  $\exists L \in \mathbb{R} \quad \forall x \in D(f), \, x<L : \, \vert f(x) - b \vert  < \varepsilon   $;
$\boldsymbol{8.}$
$ \forall K \in \mathbb{R} \quad \exists L \in \mathbb{R} \quad \forall x \in D(f), \, x<L \, : \, f(x)>K $;
$\boldsymbol{9.}$
$ \forall K \in \mathbb{R} \quad \exists L \in \mathbb{R} \quad \forall x \in D(f), \, x<L \, : \, f(x)<K $;
$\boxed{112.}$
$\boldsymbol{2.}$
ak $ \vert x-1 \vert<\delta, \, x\ne 1 $, tak  $\frac{1}{(1-x)^{2}}> \frac{1}{ \delta ^{2}}$; pre dané $ K \in \mathbb{R} $ stačí položiť napr.  $ \delta = \sqrt{1+ \vert K \vert} $ (definíciu limity používame v takej podobe, ako je v pr. 111.2);
$\boldsymbol{4.}$
ak $x\geq 0 $ a $\vert x-8 \vert < \delta $, tak $\vert \sqrt[3]{x}-2\vert = \frac{\vert x-8 \vert}{\vert \sqrt[3]{x^{2}}+2 \sqrt[3]{x}+4 \vert}< \frac{\delta}{4}  $ (pre $x\geq 0  $ je $ \sqrt[3]{x^{2}}+2 \sqrt[3]{x}+4 \geq 4 $, pre dané $\varepsilon > 0  $ stačí položiť  $\delta = \min \lbrace 4 \varepsilon, 8 \rbrace   $ (podmienka  $\delta \leq 8 $ zaručuje, že pre všetky  $x \in O_{\delta}(8) $ platí nerovnosť  $  x > 0$, ktorú sme potrebovali pri odhadovaní výrazu $ \vert \sqrt[3]{x}-2 \vert $);
$\boldsymbol{5.}$
ak $ x<0 $ a $\vert x+2\vert< \delta $ , tak  $ \vert x^{2}-4 \vert = \vert x-2 \vert \vert x+2 \vert  <2 \delta $; pre dané $\varepsilon > 0$ stačí položiť  $\delta= \min \lbrace \frac{\varepsilon}{2},2 \rbrace $;
$\boxed{113.}$
$\boldsymbol{1.}$
$\exists \varepsilon  > 0 \quad \forall \delta > 0 \quad \exists x \in D(f): \, x \ne 0 \land \vert x \vert < \delta \land \, \vert f(x)-4 \vert > \varepsilon $;
$\boldsymbol{2.}$
$\forall a \in \mathbb {R^{*}} \quad \exists O(a) \quad \forall O_{\delta} (O) \quad \exists x \in D(f), \, x\in O_{ \delta}(O), \, x \ne 0: \, f(x)\notin O(a) $;
$\boxed{114.}$
z nerovnosti $ \vert \vert f(x) \vert - \vert b \vert \vert \leq \vert f(x) - b \vert$ vyplýva: ak pre všetky $ x \in O(a)$ platí $ \vert f(x) - b \vert < \varepsilon $ (pre dané $\varepsilon > 0  $ existenciu takého $O(a) $ zaručuje predpoklad $\lim_{x \to a}f(x)=b$), tak pre všetky $ x \in O(a) $ platí $ \vert \vert f(x) \vert - \vert b \vert \vert < \varepsilon  $; opačná implikácia neplatí (napr. $f(x)= \sgn (x-a) $; neexistenciu limity v bode $a$ môžeme dokázať rovnako ako v pr. 115.2);
$\boxed{115.}$
$\boldsymbol{1.}$
ak pre všetky $ z\in O(a )$ platí $ \vert f(z) - b \vert < \varepsilon $ a ak $ x, y \in O(a ) $, tak $ \vert f(x) - f(y) \vert \leq \vert f(x) - b \vert + \vert b - f(y) \vert < 2\varepsilon $;
$\boxed{116.}$
$\boldsymbol{1.}$
$ \frac{1}{2}$;
$\boldsymbol{2.}$
$0$;
$\boldsymbol{3.}$
$5^{-5}$ (uvedenú limitu  možno napr. zapísať ako súčin piatich limít, z ktorých každá je rovná $ \frac{1}{5}$);
$\boldsymbol{4.}$
$\frac{2}{9}$ (najprv dať na spoločného menovateľa;
$\boldsymbol{5.}$
$x+ \frac{a}{2}$ (limitovaný výraz má tvar $ \frac{1}{n} \quad [(n-1)x+ \frac{a}{n} (1+2+...+(n-1))]= \frac{n-1}{n}x+\frac{a}{n^{2}}. \frac{(n-1)n}{2}$);
$\boldsymbol{6.}$
$\frac{1}{3}$ (zlomok rozšíriť $3^{-n}$ a využiť pr. 105.5);
$\boxed{117.}$
$\boldsymbol{2.}$
$\frac{2}{3}$;	
$\boldsymbol{3.}$
$1$ (stačí dosadiť, daná funkcia je elementárna a definovaná v bode 0);
$\boldsymbol{4.}$
$1$;
$\boldsymbol{5.}$
$0$;
$\boldsymbol{6.}$
$(\frac{3}{2})^{10}$;
$\boldsymbol{7.}$
$\frac{m}{n}$;
$\boldsymbol{8.}$
$-\frac{1}{2}$ (spoločný menovateľ je $ x(x-1)(x-2))$;
$\boldsymbol{9.}$
$\frac{n(n+1)}{2}$ (delenie výrazom $ x^{-1} $ si uľahčíme, ak čitateľ napíšeme v tvare $ (x-1)+(x^{2}-1)+...(x^{n}-1))$;
$\boldsymbol{10.}$
$\frac{49}{24}(x^{100}-2x+1=(x^{100}-1)-2(x-1)$; $x^{50}-2x+1=(x^{50}-1)-2(x-1)$; tento prepis uľahčí delenie výrazom $x-1$);
$\boxed{118.}$
$\boldsymbol{1.,2.}$
napr. $f(x)=\chi (x), \, g(x)=1-\chi (x) $ (tým sme vlastne ukázali, že implikácie  vo vetách o limite súčtu a limite súčinu nemožno obrátiť);
$\boxed{119.}$
áno; napr. $ s_{n}=(-1)^{n}\frac{1}{n}$ možno dokázať, že pre každú postupnosť  $\lbrace a_{n}\rbrace_{n=1}^\infty $ vyhovujúcu pr. 119 platí $\lim_{n \to \infty}a_n =0$);
$\boxed{120.}$
$\boldsymbol{2.}$
$\frac{4}{3}$;
$\boldsymbol{3.}$
$\frac{1}{\sqrt{2a}}$ (= $\lim_{x \to a}( \frac{\sqrt{x-a}}{\sqrt{x+a}(x+a)}+\frac{1}{\sqrt{x+a}})$) ;
$\boldsymbol{4.}$
$-\frac{1}{16}$;
$\boldsymbol{5.}$
$\frac{2}{27}$ na úpravu čitateľa sa použije vzťah $A^{3}-B^{3}=(A-B)(A^{2}+AB+B^{2}) $, v menovateli stačí vyňať pred zátvorku premennú $ x $);
$\boldsymbol{6.}$
$\frac{1}{n!}$;
$\boxed{121.}$
$\boldsymbol{2.}$
$\frac{112}{27}$;
$\boldsymbol{3.}$
$\frac{7}{36}$;
$\boxed{122.}$
$\boldsymbol{1.}$
$\frac{5}{3}$ (subst. $t=\sqrt[15]{x} $);
$\boldsymbol{3.}$
$\frac{n}{m}$ (subst. $t=\sqrt[m.n]{x} $);
$\boldsymbol{4.}$
$\frac{1}{2}$;
$\boxed{122.}$
$\boldsymbol{1.}$
1 (zlomok rozšíriť výrazom $\frac{1}{\sqrt{x}}$; uvedomte si, že pri výpočte limity v čitateli aj v menovateli takto získaného zlomku sa používa veta o limite zloženej funkcie);
$\boldsymbol{2.}$
$\frac{1}{\sqrt{2}}$ (zlomok rozšíriť výrazom $\frac{1}{\sqrt{x}}$;
$\boldsymbol{3.}$
-2 (pre $x>0$ je $\vert x \vert=x $);
$\boldsymbol{4.}$
$\frac{a+b}{2}$ (rozšíriť výrazom $\sqrt{(x+a)(x+b)}+x $);
$\boldsymbol{5.}$
$\frac{2}{3}$;
$\boldsymbol{6.}$
$\frac{a_{1}+...+a_{n}}{n}$;
$\boldsymbol{7.}$
$-\frac{1}{4}$ (rozšíriť výrazom $ x+\sqrt{x^{2}+2x}+2\sqrt{x^{2}+x}$, potom napísať v tvare súčinu tak, aby jeden zo súčiniteľov bol $ \frac{2x}{x+\sqrt{x^{2}+2x}+2\sqrt{x^{2}+x}}$, jeho limitu už vieme vypočítať; z druhého súčiniteľa vyňať pred zátvorku $ x $ potom rozšíriť výrazom $\sqrt{x^{2}+2x}+x+1 $);
$\boldsymbol{8.}$
$\frac{1}{3}$;
$\boxed{125.}$
áno, podkladom pre konštrukciu funkcií $f,g $ môže byť rozbor dôkazov vety o limite zloženej funkcie a tvrdenia z pr. 124; $g $ nesmie spĺňať podmienku (*) z uvedenej vety, pre $f$ nesmie platiť $\lim_{x \to 2}f(x) = f(2)$);
$\boxed{126.}$
$\boldsymbol{1.}$ 5;
$\boldsymbol{2.}$
$\sin 1$ (stačí dosadiť);
$\boldsymbol{3.}$
$\frac{m}{n}$ (rozšíriť výrazom $\frac{1}{x}$);
$\boldsymbol{4.}$ 2 (rozšírte výrazom $x^{3}+2x$ a využite, že  $\lim_{x \to 0}\frac{\sin (x^{3}+2x)}{x^{3}+2x} = 1$ podľa vety o limite zloženej funkcie);
$\boldsymbol{5.}$
$\frac{1}{3}$;
$\boldsymbol{6.}$ $x$ (subst. $t= \frac{x}{2^{n}} $);
$\boxed{127.}$
$\boldsymbol{1.}$
$\frac{1}{2} $ (rozšíriť výrazom $1+ \cos x $ alebo použiť vzorec $1-\cos x =2 \sin ^{2}\frac{x}{2}$);
$\boldsymbol{2.}$
$\frac{1}{2} $;
$\boldsymbol{3.}$ 2 (stačí rozšíriť  výrazom $\frac{1}{x} $);
$\boldsymbol{4.}$
$\frac{1}{2} $;
$\boldsymbol{5.}$ 4 (v čitateli stačí odpočítať a pripočítať číslo 1 a použiť výsledok pr. 127.1);
$\boldsymbol{6.}$
$\frac{1}{p} $ (rozšíriť výrazom $\frac{1}{x} $; využiť, že $\lim_{x \to 0}\frac{1-\cos x}{x} = \lim_{x \to 0}\frac{1-\cos x}{x^{2}}x=0$, pozri pr. 127.1);
$\boxed{128.}$
$\boldsymbol{1.}$ $\cos a $ (použite vzorec $\sin x- \sin a = 2 \sin \frac{x-a}{2} \cos \frac{x+a}{2} $);
$\boldsymbol{2.}$
$-\frac{1}{\sin ^{2}a} $;
$\boldsymbol{3.}$
$- \cos a$;
$\boldsymbol{4.}$
$\frac{3}{2} \sin 2a $ (v čitateli napr. odpočítať a pripočítať výraz $\sin ( a+x) \sin a$);
$\boxed{129.}$
$\boldsymbol{1.}$ 
$\frac{1}{2} $;
$\boldsymbol{2.}$ -3 (subst. $ t= \sin x$);
$\boldsymbol{3.}$
$\frac{3}{4} $;
$\boldsymbol{5.}$
$\frac{2}{\pi} $;
$\boldsymbol{6.}$
1 (subst. $t = \arcsin x $;
$\boxed{130.}$
$\boldsymbol{1.}$ 
$\frac{1}{4} $;
$\boldsymbol{2.}$ 
$\frac{4}{3} $;
$\boldsymbol{3.}$ 
$-\frac{1}{12} $;
$\boldsymbol{4.}$ 
$\sqrt{2} $ (limitovaná funkcia je definovaná len pre $ x>0$, na jej definičnom obore teda platí $x= \sqrt{x^{2}} $);
$\boldsymbol{5.}$ 
$-\frac{1}{20} $;
$\boxed{131.}$
$\boldsymbol{1.}$ 
0 ($\lim_{x \to \infty}  \frac{1}{x}=0$, $\sin  $ je ohroničená funkcia);
$\boldsymbol{2.}$ 
$\frac{1}{2} $ (rozšíriť výrazom $\frac{1}{x^{2}} $);
$\boldsymbol{3.}$ 0;
$\boldsymbol{4.}$ 3;
$\boldsymbol{5.}$
$\frac{1}{2}$;
$\boxed{132.}$
$\boldsymbol{2.}$ 
$\frac{n}{2^{n}}= \frac{n}{(3/2)^{n}}.(3/4)^{n}\leq (3/4)^{n}$;
$\boldsymbol{3.}$
$( \frac{n}{\sqrt[3]{25}^{n}})^{3/2}\leq (\frac{n}{2^{n}})^{3/2}$ (rovnosti z pr. 132 možno dokázať aj na základe pr. 133.1;
$\boxed{133.}$
$\boldsymbol{1.}$ 
 $ 0 \leq a_{n} \leq \frac{a_{1}}{q}.q^{n}  $;
$\boldsymbol{2.}$ neplatí, stačí zvoliť  $ a_{n}= \frac{n+1}{2n}$;
$\boldsymbol{3a/.}$ $0$ (použiť pr. $133.1$ a nerovnosť $(1+ \frac{1}{n})^{n}\geq 2) $;$\boldsymbol{3b/.}$ $0$ (treba si uvedomiť, že tvrdenie z príkladu $133.1$ zostane v platnosti, aj keď $\frac{a_{n+1}}{a_{n}}\leq q$ bude platiť len pre všetky $n\geq K$, kde $K\in \mathbb{N}$ je dané);
$\boxed{134.}$
$\boldsymbol{1.}$ $0$ $((\frac{7}{n})^{n}\leq (\frac{7}{8})^{n}$ pre $n\geq 8$);$\boldsymbol{2.}$ $0$;
$\boxed{135.}$
$\boldsymbol{2.}$
pre $\omega(n):=\sqrt[n]{n}-1$ platí $0\leq \omega(n) \leq \sqrt{\frac{2}{n}}$;
$\boxed{136.}$
$\boldsymbol{1.}$  $0$ (rozšíriť výrazom $\frac{1}{n!}$);
$\boldsymbol{2.}$  $1$ (pre $n\geq 19 $ je $\sqrt[n]{4}\leq \sqrt[n]{\frac{5n+1}{n+5}} \leq \sqrt[n]{6} $; alebo všeobecnejšie: $\lim_{n\rightarrow\infty} \frac{5n+1}{n+5} =5$, preto pre dané $\varepsilon \in (0,5) $ existuje $O (\infty )$, v ktorom platí $5- \varepsilon \leq \frac{5n+1}{n+5} \leq 5+ \varepsilon $, a teda aj $\sqrt[n]{5- \varepsilon} \leq  \sqrt[n]{\frac{5n+1}{n+5}}\leq \sqrt[n]{5+ \varepsilon} $)
$\boldsymbol{3.}$  
$3$ $ (3^{n}-2^{n} = (3-2).(3^{n-1}+3^{n-2}.2+...+2^{n-1}) 3^{n-1}$; iná možnosť: $ \sqrt[n]{3^{n}-2^{n}}= 3 \sqrt[n]{1-(\frac{2}{3})^{n}}$, pritom $\lim_{n\rightarrow\infty} (1- ( \frac{2}{3})^{n})=1$, preto pre dané $\varepsilon \in (0,1) $  existuje $O( \infty )$, v ktorom platí $1- \varepsilon \leq 1- (\frac{2}{3})^{n} \leq 1+ \varepsilon $, a teda aj $ \sqrt[n]{1- \varepsilon } \leq \sqrt[n ]{1- (\frac{2}{3})^{n}} \leq \sqrt[n]{1+ \varepsilon }$; ďalšou možnosťou výpočtu limít z pr. 136.2-4. je použitie postupov uvedených v odseku 2.5);
$\boldsymbol{4.}$  $1 $  ($\frac{1}{2^{n}}\leq \frac{1}{2n} $, iná možnosť:  $\sqrt[n]{\frac{1}{n}-\frac{1}{2^{n}}}= \sqrt[n]{\frac{1}{n}}\sqrt[n]{1-\frac{n}{2^{n}}} $, ďalej využite analogicky ako v pr. 136.2,3 fakt, že $\lim_{n\rightarrow\infty} (1-\frac{n}{2^{n}})=1$);
$\boxed{137.}$
$\boldsymbol{1.}$  $-\infty$  ($\lim_{x\rightarrow-\infty} \sqrt{(x+a)(x+b)}=+\infty$, $\lim_{x\rightarrow-\infty} -x = + \infty$, preto $\lim_{x\rightarrow-\infty}( \sqrt{(x+a)(x+b)}-x)=+\infty$;
$\boldsymbol{2.}$  $-\infty$;
$\boldsymbol{3.}$  $+\infty$;
$\boldsymbol{4.}$  $+\infty$ (rozšíriť výrazom $\frac{1}{x^{2}} $; $x^{2}-\frac{5}{x}\rightarrow \, +\infty $, $1/(1-\frac{3}{x}+\frac{1}{x^{2}}) \rightarrow \, 1 $ pre  $x\rightarrow -\infty $);
$\boldsymbol{5.}$  $+\infty$;
$\boldsymbol{6.}$  $+\infty$;
$\boldsymbol{7.}$  $+\infty$ ($2+ \sin x \geq 1 $ pre $x \in \mathbb{R} $, preto $(2+ \sin x).x\geq x $ pre $x\geq0 $);
$\boxed{138.}$
$\boldsymbol{1.}$  $+\infty$;
$\boldsymbol{2.}$  $+\infty$ ; ($\frac{\sin x}{x}\rightarrow 1$,  $\frac{1}{x^{2}}\rightarrow +\infty$ pre $x\rightarrow 0 $);
$\boldsymbol{3.}$  $+\infty$ (rozšíriť výrazom $\frac{1}{x^{2}} $  alebo $\frac{1}{x^{4}} $);
$\boldsymbol{4.}$  $+\infty$;
$\boldsymbol{5.}$  $+\infty$;
$\boldsymbol{6.}$  $+\infty$;
$\boxed{139.}$  $\boldsymbol{3.}$ napr.  $f(x)= \frac{1}{(x-1)^{2}}+ \chi (x) $, $g(x)= -\frac{1}{(x-1)^{2}} $;
$\boxed{140.}$
$\boldsymbol{4.}$ napr. $a_{n}=n, \, b_{n}=(-1)^{n}. \frac{1}{n} $ ( príklady 139 a 140 ukazujú, že možno sformulovať všeobecné tvrdenia, ktoré by umožňovali výpočet limít neurčitých výrazov $ +\infty$ $ -\infty $, resp. $0.(+\infty) $);
$\boxed{141.}$ 
$\lim_{x\rightarrow\infty} R(x) =$ $ f(n)=\left\{\begin{matrix} +\infty, & \mbox{ak }n>m, \frac{a_{0}}{b_{0}}>0 \\ -\infty, & \mbox{ak }n>m, \frac{a_{0}}{b_{0}}<0 \\ \frac{a_{0}}{b_{0}} & \mbox {ak }n=m, \\ 0 & \mbox {ak }n<m\end{matrix}\right.$
najprv dokážte - napr. matematickou indukciou - že $\lim_{x\rightarrow-\infty}(c_{0}x^{s}+c_{1}x^{s-1}+...+c_{s})=$ 
$\left\{\begin{matrix} +\infty, & \mbox{ak } c_{0}>0 \\ - \infty, & \mbox{ak} c_{0}<0\end{matrix}\right. $);
$\boxed{142.}$
$\boldsymbol{1.}$
$\forall \varepsilon >0 \quad \exists \delta >0 \quad \forall x \in D(f): \, 0< x-a < \delta \, \Rightarrow \vert f(x) - b \vert < \varepsilon$;
$\boldsymbol{2.}$
$\forall \varepsilon >0 \quad \exists \delta >0 \quad \forall x \in D(f): \, -\delta< x-a < 0 \, \Rightarrow \vert f(x) - b \vert < \varepsilon$ (pritom predpokladáme, že $a$ je hromadný bod množiny $ D(f) \cap (a, \infty )$ resp. množiny $D(f) \cap (-\infty, a)$);
$\boxed{143.}$
$\boldsymbol{1.}$
$\lim_{x\rightarrow 1^{+}} f(x) =2$, $\lim_{x\rightarrow 1^{-}} f(x) =-2$;
$\boldsymbol{2.}$
$\lim_{x\rightarrow 0^{+}} f(x) =\sqrt{2}$, $\lim_{x\rightarrow 0^{-}} f(x) =-\sqrt{2}$ ($1-\cos 2x = 2\sin ^{2}x $);
$\boldsymbol{3.}$
$\lim_{x\rightarrow 2^{+}} f(x) =\infty$, $\lim_{x\rightarrow 2^{-}} f(x) =-\infty$;
$\boldsymbol{4.}$
$\lim_{x\rightarrow 0^{+}} f(x) =0$, $\lim_{x\rightarrow 0^{-}} f(x) =\frac{1}{2}$;
$\boxed{144.}$
$\boldsymbol{1.}$
$\lim_{x\rightarrow \pi/2^{+}} tg x =-\infty \ne +\infty $ $=\lim_{x\rightarrow \pi/2^{-}} tg x$, preto $\lim_{x\rightarrow \pi/2} tg x$ neexistuje;
$\boldsymbol{2.}$
$\lim_{x\rightarrow 0}x \sgn x = 0$ (možno dokazovať pomocou jednostranných limít alebo využiť fakt, že $sgn$ je ohraničená funkcia a $\lim_{x\rightarrow 0}x  = 0$;
$\boldsymbol{3.}$ neexistuje ( $ \frac{\sin x}{x^{2}} = \frac{\sin x}{x} . \frac{1}{x}$);
$\boxed{146.}$
$\boldsymbol{2.}$
za daných predpokladov $\lim_{x\rightarrow 0} f(x)$ existuje práve vtedy, keď $\lim_{x\rightarrow 0^{+}} f(x)=0$;
$\boxed{149.}$
Nech sú dané funkcie $f,g $, nech $b \in \mathbb{R^{*}} $ je hromadný bod množiny $D(f) \cap D(g) $, nech $\lim_{x\rightarrow b} f(x)=a$, $\lim_{x\rightarrow b} g(x)=A$. Potom: $\boldsymbol{1.}$ ak $a \in (0,1),\, A= +\infty  $, tak 
$\lim_{x\rightarrow b} f(x)^{g(x)}=0$;
$\boldsymbol{2.}$ ak $a \in (0,1),\, A= -\infty  $, tak $\lim_{x\rightarrow b} f(x)^{g(x)}=+\infty$;
$\boldsymbol{3.}$ ak $a \in (1,+\infty)\cup \lbrace +\infty \rbrace,\, A= +\infty  $, tak $\lim_{x\rightarrow b} f(x)^{g(x)}=+\infty$; 
$\boldsymbol{4.}$ ak $a \in (1,+\infty)\cup \lbrace +\infty \rbrace,\, A= -\infty  $, tak $\lim_{x\rightarrow b} f(x)^{g(x)}=0$ (pri dôkazoch sa stačí obmedziť na také prstencové okolie bodu $b$, v ktorom je $f$ kladná (také okolie existuje, pretože $\lim_{x\rightarrow b} f(x)$ je kladné číslo alebo bod $+ \infty$)); 
$\boxed{150.}$
$\boldsymbol{1.}$  $\sqrt{\frac{2}{3}} $;
$\boldsymbol{2.}$  $\frac{1}{4} $;
$\boldsymbol{3.}$  $+\infty $;
$\boldsymbol{4.}$  $\frac{1}{2} $;
$\boldsymbol{5.}$  $0 $;
$\boldsymbol{6.}$  $0 $;
$\boldsymbol{7.}$  $+ \infty $;
$\boldsymbol{8.}$ neexistuje (limita sprava (zľava) je $+\infty, (0) $);
$\boxed{151.}$
$\boldsymbol{1.}$  $e^{3} $;
$\boldsymbol{2.}$  $e^{-2} $;
$\boldsymbol{3.}$  $e^{3} $;
$\boldsymbol{4.}$  $e^{-1} $;
$\boldsymbol{5.}$  $e^{ctg a} $;
$\boldsymbol{6.}$  $1 $;
$\boldsymbol{7.}$  $e $ (možno použiť substitúciu $\frac{1}{x}=t $, potom $t\rightarrow 0^{+} $ pre $x \rightarrow +\infty $);
$\boldsymbol{8.}$  $e^{-x^{2}/2} $;
$\boxed{152.}$
$\boldsymbol{1.}$  $1 $; (= $\lim_{x\rightarrow 0} ln[(1+x)^{1/x}  ] $; použije sa veta z pr. 124);
$\boldsymbol{2.}$  $1 $;
$\boldsymbol{3.}$  $\frac{1}{a} $; 
$\boldsymbol{5.}$  $0 $;
$\boldsymbol{6.}$  $n $ (využite, že $\lim_{x\rightarrow 0} \frac {ln [1+(x+ \sqrt{1-x^{2}}-1)]}{x+\sqrt{1-x^{2}}-1} =1$, analogický vzťah použite pre $ln (nx+ \sqrt{1-n^{2}x^{2}}) $);   
$\boxed{153.}$
$\boldsymbol{1.}$  $1 $ (možno použiť substitúciu  $e^{x}-1=t  $);
$\boldsymbol{2.}$  $1 $;
$\boldsymbol{3.}$  $4. ln \frac{2}{e} $ (v čitateli možno pripočítať a odpočítať číslo $4$, potom previesť na súčet dvoch limít);
$\boldsymbol{5.}$  $\frac{25}{4} $;
$\boldsymbol{6.}$  $1 $ (v čitateli možno pripočítať a odpočítať číslo $1$, potom zlomok rozšíriť výrazom $\frac{1}{x}$);
$\boldsymbol{7.}$  $ln x $;
$\boldsymbol{8.}$  $e^{2} $;
$\boxed{154.}$
$\boldsymbol{1.}$  $\frac{1}{3} $;
$\boldsymbol{3.}$  $\frac{1}{2} $;
$\boxed{153.}$
$\boldsymbol{1.}$  $a_{n+1}=(1-\frac{1}{2^{n+1}}) a_{n} < a_{n} \, (n \in \mathbb{N}) $, teda $ \lbrace a_{n} \rbrace_{n=1}^\infty $ je klesaajúca postupnosť; $ \forall n \in \mathbb{N}: \, a_{n} 	>0 $, teda $ \lbrace a_{n} \rbrace_{n=1}^\infty $  je zdola ohraničená; preto existuje konečná  $\lim_{n\rightarrow \infty} a_{n}  $(= $\inf \lbrace a_{n} ; n \in \mathbb{N}\rbrace $);
$\boldsymbol{2.}$ $ \lbrace a_{n} \rbrace_{n=1}^\infty $ je rastúca; z nerovnosti  $n!\geq 2^{n-1}  $ ($n \in \mathbb{N}  $) vyplýva  $a_{n}\leq 1+1+\frac{1}{2}+...+\frac{1}{2^{n-1}} = 2+1-\frac{1}{2^{n-1}} < 3$;
$\boldsymbol{3.}$ $\lbrace a_{n}\rbrace_{n=1}^{\infty}  $  je rastúca $a_{n+1}-a_{n}>0  $) a zhora ohraničená ( $a_{n} < 1,\, n \in \mathbb{N} $);
$\boldsymbol{4.}$ $\lbrace a_{n}\rbrace_{n=1}^{\infty}  $  je zhora ohraničená ( $a_{n} < \frac{1}{5}+...+ \frac{1}{5^{n}} = \frac{1}{4}(1- \frac{1}{5^{n}} < \frac{1}{4}$) a rastúca;
$\boldsymbol{5.}$ pre $n>10 $ je $a_{n}< a_{n-1} $, $\lbrace a_{n}\rbrace_{n=10}^{\infty} $ je klesajúca a zdola ohraničená  ($ \forall n \in \mathbb{N}: \, a_{n} 	>0 $);
$\boxed{156.}$
$\boldsymbol{2.}$ ak si najprv pomôžeme grafmi funkcií $y=\frac{1}{2}(x^{2}+1) $  a   $y=x $ nakreslenými do jedného obrázka (podobne ako obr. 2), zistíme, že postupnosť $\lbrace a_{n}\rbrace_{n=10}^{\infty} $ je rastúca a zhora ohraničená (napr. číslom  $ 1 $); pretože však i najlepšie nakreslený obrázok nieje dôkazom, treba obidve uvedené tvrdenia ešte dokázať; hľadaná limita vyhovuje rovnici  $ 2a=a^{2}+1 $, preto $\lim_{n\rightarrow \infty} a_{n} =0$ (druhý z koreňov rovnice $x= \frac{x}{2+x} $, t.j. číslo $-1$ nemôže byť hľadanou limitou, pretože platí: $\lim_{k\rightarrow \infty} a_{2k} = \lim_{k\rightarrow \infty} a_{2k-1}= \frac{1}{2}(1+\sqrt{5})$, preto $\lim_{n\rightarrow \infty} a_{n} = \frac{1}{2}(1+\sqrt{5})$;
$\boldsymbol{5.}$ pre $x=k \pi $  ($k \in \mathbb{Z}$) je postupnosť  $\lbrace a_{n}\rbrace_{n=1}^{\infty}$ konštantná, pre  $ x \in \bigcup\limits_{k \in \mathbb{Z}}2k \pi, (2k+1)\pi) $ klesajúca a zdola ohraničená, pre  $ x \in \bigcup\limits_{k \in \mathbb{Z}}(2k-1) \pi, 2k\pi) $ rastúca  a zhora ohraničená (pri dôkazoch monotónnosti využite platnosť nerovností  $0 < \sin x < x  $ pre  $x \in (0,\pi)  $, $x < \sin x < 0  $ pre  $x \in (-\pi,0)  $), $\lim_{n\rightarrow \infty}a_{n}=0$;
$\boxed{157.}$  $\frac{a_{n+1}}{a_{n}}= \frac{n+1}{n}. \frac{1}{(1+\frac{1}{n(n+2)})^{n+2}}<1$ (na odhad menovateľa v druhom súčiniteli použite nerovnosť z pr. 10.1), teda $\lbrace a_{n}\rbrace_{n=1}^{\infty }$ je klesajúca postupnosť (fakt, že $\lbrace (1+\frac{1}{n})^{n}\rbrace_{n=1}^{\infty }$ je rastúca postupnosť - ktorý pokladáme za známy z prednášok - možno dokázať podobne); $e=\lim_{n\rightarrow \infty} a_{n} = \inf \lbrace a_{n}; n\in \mathbb{N}\rbrace$;
$\boxed{159.}$ $\boldsymbol{1.}$ nie je (platí len implikácia "$\Rightarrow $ "; protipríkladom, dokumentujúcim neplatnosť implikácie "$\Leftarrow $ ", je napr. Dirichletova funkcia $\chi (x) $);
$\boldsymbol{2.}$ je;
$\boxed{160.}$ (H označuje množinu všetkých hromadných hodnôt postupnosti $\lbrace a_{n}\rbrace_{n=1}$; $a$ resp. $b$ označujú $\lim_{n\rightarrow \infty} \inf a_{n}$, resp. $\lim_{n\rightarrow \infty} \sup a_{n}$; v pr. 160.1-6 použijeme túto únahu: ak $\lbrace a_{k_{1}}\rbrace_{n=1}^{\infty}$,..., $\lbrace a_{k_{1}}\rbrace_{n=1}^{\infty}$ sú podpostupnosti $\lbrace a_{n}\rbrace_{n=1}^{\infty}$ také, že: a/ každá z nich má limitu; b/ množiny  $N_{1}:=\lbrace k_{1} (n); \, n \in \mathbb{N} \rbrace ,...,N_{1}:=\lbrace k_{1} (n); \, n \in \mathbb{N} \rbrace  $ sú po dvoch disjunktné a $N_{1}\cup ... \cup N_{1} = N $ (t.j. ak sa podarí postupnosť $\lbrace a_{n}\rbrace_{n=1}^{\infty}$ "rozdeliť" na konečný počet podpostupností, z ktorých každá má limitu), tak H pozostáva z limít postupností $\lbrace a_{k_{1}}\rbrace_{n=1}^{\infty}$,..., $\lbrace a_{k_{1}}\rbrace_{n=1}^{\infty}$;
$\boldsymbol{1.}$  $\lim_{k\rightarrow \infty} a_{2k}=-2$, $\lim_{k\rightarrow \infty}a_{2k-1}=2$, $ H= \lbrace -2,2\rbrace $, $a=-2, \, b=2 $;
$\boldsymbol{2.}$   $ H= \lbrace -4,0,2,6\rbrace $, $a=-4, \, b=6 $ ( $\lim_{k\rightarrow \infty} a_{4k} = 2$, $\lim_{k\rightarrow \infty} a_{4k+1} = 6$, $\lim_{k\rightarrow \infty} a_{4k+2} = -4$, $\lim_{k\rightarrow \infty} a_{4k+3} = 0$,  
$\boldsymbol{3.}$   $ H= \lbrace -1,-\frac{1}{2},0,\frac{1}{2},1\rbrace $, $a=-1, \, b=1 $;
$\boldsymbol{4.}$   $ H= \lbrace -\frac{1}{\sqrt{2}}-e,\frac{1}{\sqrt{2}}-e,e-1,e,e+1\rbrace $, $a=-\frac{1}{\sqrt{2}}-e, \, b=e+1 $;
$\boldsymbol{5.}$   $ H= \lbrace 0,\infty \rbrace $, $a=0, \, b=+\infty $;
$\boldsymbol{6.}$   $ H= \lbrace 1,2\rbrace $, $a=1, \, b=2 $;
 $\boldsymbol{7.}$ , $\boldsymbol{8.}$ $a=0,\, b=1$; o čísle $0$, resp. $1$ možno totiž dokázať, že vyhovuje obidvom podmienkam z vety 18 (v pr. 160.7,8 možno dokázať, že $ H= \langle0,1\rangle $);
$\boxed{161.}$ $\boldsymbol{4.}$  napr. $\lbrace 1,2,1,2,3,1,2,3,4,1,2,3,4,5,... \rbrace $;
$\boxed{163.}$ (skôr než začnete dokazovať, by ste si mali uvedomiť, že platí  $a,b \in \mathbb{R}, \, a<b$) sporom: keby $N_{\varepsilon}^{*} $ bola nekonečná pre niektoré $\varepsilon^{*} \in (0, \frac{b-a}{2} $), existovala by hromadná hodnota $c$ postupnosti $\lbrace a _{n} \rbrace_{n=1^{\infty}}  $, $c \in \langle a+\varepsilon^{*},b-\varepsilon^{*} \rangle $ (teda $c\ne a, \, c\ne b $);
$\boxed{164.}$ stačí dokázať: ak $a,b \in \mathbb{R}, \, a\leq b $, pričom platí $\forall \varepsilon >0: \, b< a+\varepsilon $, tak   $a=b $ (dôkaz možno vykonať sporom), tvrdenie pr. 164 je potom dôsledkom tohoto a vety 19	
$\boxed{165.}$ $\boldsymbol{1.}$ v $O(a) $ ( $O(a)$ je ľubovolné okolie bodu $a$) leží aspoň jeden prvok $z \in A $; existuje také okolie $O(z) $ prvku $z$, že $O(z) \subset  O(a) $ v $O(z) $ leží (pretože $z \in A $) nekonečne veľa členov postupnosti $\lbrace a_{n}\rbrace_{n=1}^{\infty} $, teda pre každé  $O(a) $ je množina $\lbrace n \in \mathbb{N}; \, a_{n}\in O(a)  \rbrace $ nekonečná;
$\boldsymbol{2.}$ nie; $A$ neobsahuje prvok $0$, čo je v spore s pr. 165.1;
$\boxed{166.}$ $\boldsymbol{1a/}$ postupnosť $\lbrace a _{n(k)}+b_{n(k)} \rbrace_{k=1}^{\infty}  $ vybraná z postupnosti  $\lbrace a _{n} + b_{n} \rbrace_{n=1}^{\infty}$ konverguje práve vtedy, keď konverguje postupnosť $\lbrace b _{n(k)} \rbrace_{k=1}^{\infty}$; preto množinu  $ H $ hromadných hodnôt postupnosti $\lbrace a _{n} + b_{n} \rbrace_{n=1}^{\infty}$ možno písať v tvare $H = a+h;\, h \in H_{1}$, kde  $a:=\lim_{n\rightarrow \infty} a_{n}$, $H_{1}$ je množina hromadných honôt postupnosti  $\lbrace  b_{n} \rbrace_{n=1}^{\infty}$; 
$\boldsymbol{1b/}$ dôkaz prvej nerovnosti možno založiť na tejto úvahe: ak pre postupnosť  $\lbrace  c_{n} \rbrace_{n=1}^{\infty}$ a číslo $c\in \mathbb{R}$ platí: pre každé $\varepsilon > 0 $ je množina  $\lbrace n\in \mathbb{N}; \, c_{n}\leq c- \varepsilon \rbrace$ konečná, tak  $ c \leq \lim_{n\rightarrow \infty} c_{n}$ (porovnaj s vetou 18); stačí teda ukázať, že číslo $ \lim_{n\rightarrow \infty} a_{n + \lim_{n\rightarrow \infty} b_{n}}$ vyhovuje tejto podmienke; posledná nerovnosť sa dokazuje analogicky;
$\boldsymbol{2.}$ napr. $\lbrace a_{n} \rbrace_{n=1}^{\infty} = \lbrace 0,1,2,3,0,1,2,3,...\rbrace $, $\lbrace a_{n} \rbrace_{n=1}^{\infty} = \lbrace 1,2,0,1,1,2,0,1,1,2,0,1...\rbrace $;
$\boxed{168.}$ nech je dané $O(a)$; ak $b \in (O(a)-\lbrace a \rbrace )\cap A^{)} $, tak existuje také $O(b)$, že $O(b) \subset O(a)-\lbrace a \rbrace$; v $O(b)$ leží aspoň jeden prvok z $A$;
$\boxed{169.}$ neexistuje, vyplýva to z pr. 168;
$\boxed{170.}$ doporučujeme dokazovať nepriamo, pokusy o priamy dôkaz radšej konzultujte s privolaným odborníkom;
$\boxed{171.}$ neplatnosť obrátenej implikácie dokumentuje postupnosť $a_{n}=(-1)^{n}$;
$\boxed{172.}$ (využite pr. 171); $\boldsymbol{1.}$ $0$;
$\boldsymbol{2.}$ $0$;
$\boxed{173.}$ $\boldsymbol{1.}$  ak $N= \max \lbrace m \in \mathbb{N}; \, a_{m} \leq K\rbrace$; ( $K > 1$ je dané), tak pre všetky $n \in \mathbb{N}, \, n>N $ platí $a_{n} \leq K $;
$\boldsymbol{2.}$ vyplýva to z pr. 173.1 a vety o limite zloženej funkcie;
$\boxed{174.}$ $\boldsymbol{1.}$ nie; $\boldsymbol{2.}$ áno;
$\boxed{175.}$ $\boldsymbol{1.}$ ak $\lim_{n \to \infty} \vert a_n \vert =0 $  tak aj $\lim_{n \to \infty}  a_n =0 $;
$\boldsymbol{2.}$ nepriamo, keby napr. množina $N^{+}$ bola konečná, vyplývala by z existencie $\lim_{n \to \infty} \vert a_n \vert =b $ existencia $\lim_{n \to \infty}  a_n =-b $;
$\boldsymbol{3.}$ nepriamo, keby $\mathbb{N}-(N^{+}\cup N^{-}) $ bola nekonečná, platilo by $\lim_{n \to \infty} \vert a_n \vert =0 $;
$\boxed{177.}$ napr. $f(x)=x $ pre $x \in \mathbb{Q} $,  $f(x)=-x $ pre $x \notin \mathbb{Q} $;
$\boxed{178.}$ sporom, ak $\sin n $ konverguje, tak $\lim_{n \to \infty} (\sin (n+2)-\sin n) =0 $, odtiaľ $\lim_{n \to \infty} \cos (n+1) =0 $; ak zo súčtového vzorca pre $\cos (n+1)$ vyjadríme $\sin n$, dostaneme $\lim_{n \to \infty} \sin n = \lim_{n \to \infty} \cos n =0 $, čo je v spore zo vzťahom $\sin^{2} n + \cos ^{2} n =1$;
neexistencia nevlastnej limity vyplýva z ohraničenosti postupnosti $\lbrace \sin n \rbrace _{n=1} ^{\infty}$;
$\boxed{179.}$ $\boldsymbol{1.}$ $\frac{1}{2} m.n.(n-m)$ (skontrolujte si, či Vami použitý postup je správny, keď $m=1$ alebo $n=1$);
$\boldsymbol{2.}$ $\frac{n.(n-1)}{2} a^{n-2}$ (možno použiť subst. $x-a=t$ ; je použitý postup správny pre $n=1$?);
$\boldsymbol{3.}$ $1 $;
$\boldsymbol{4.}$ $\frac{m-n}{2} $ (skontrolujte si správnosť použitého postupu, ak  $m=1$ alebo $n=1$);
$\boldsymbol{5.}$ $x^{2}+ax+a^{2}/3 $;
$\boldsymbol{6.}$ $n^{-n(n+1)/2}$ (skontrolujte si správnosť použitého postupu pre párne $n$, vtedy totiž exponent v menovateli nie je celočíselný);
$\boldsymbol{7.}$ $+\infty$;
$\boldsymbol{8.}$ $ \frac{2}{3}$( $\lim_{n \to \infty} \frac{2\frac{n(n+1)}{2}-n}{3\frac{n(n+1)}{2}-2n} $); 
$\boldsymbol{9.}$ neexistuje ($a_{2k}=-\frac{1}{2}$, $a_{2k-1}=1-\frac{k-1}{2k-1}$, $k\in \mathbb{N}$, kde $a_{n}=\frac{1-2+3-4+...+(-1)^{n-1}.n}{n}$);
$\boldsymbol{10.}$ $ \frac{3}{2}$;
$\boxed{180.}$ ak $\lim_{n \to \infty} a_{n}=a \in \mathbb{R} $, tak $\lbrace a_{n} \rbrace_{n=1} ^{\infty}$ je ohraničená, existuje teda $ B:= \sup_{n\in \mathbb{N}}a_{n} $, $A:= \inf_{n\in \mathbb{N}}a_{n} $, pritom $A \leq a\leq B $; predpokladajme $ A < a $, potom $ \exists N \in \mathbb{N} $  $ \forall n>N: \, a_{n}>A $ (stačí zvoliť $\varepsilon= \frac{a-A}{2} $ a použiť definíciu limity), potom najmenšie z čísel $a_{1},...,a_{N} $ musí byť rovné $A $; podrobné vypracovanie ako aj ostatné prípady prenechávame čitateľovi;
$\boxed{181.}$ R má limity v bodoch $+\infty$, $-\infty$; ak $\lim_{x \to \infty} R(x)=+\infty$,   $\lim_{x \to -\infty} R(x)=-\infty$ (ostatné prípady urobte sami), tak  $\exists a_{1} \in \mathbb{R} \, \forall x < a_{1}$, $x\in D(R):\, R(x)>1$, pritom $a_{1}< a_{2} $; ak má platiť  $ R(\mathbb{Z})=\mathbb{Q} $, musí byť $R( \mathbb{Z} \cap \langle a_{1}, a_{2}\rangle$) = $ \mathbb{Q} \cap \langle-1,1 \rangle$, čo je nemožné (množina vľavo je konečná, množina vpravo nekonečná);
$\boxed{182.}$ $\boldsymbol{1.}$ $ \frac{\alpha}{m} -\frac{\beta}{n}$ (stačí nájsť "vzorec" pre $ \lim_{x \to 0}\frac{\sqrt[k]{1+\vartheta x}-1}{x}$, $k\in \mathbb{Z}- \lbrace 0 \rbrace $);
$\boldsymbol{2.}$ $ \frac{\alpha}{m} -\frac{\beta}{n}$ ( v čitateli stačí pripočítať a odpočítať výraz  $\sqrt[m]{1+\alpha_{x}}$);
$\boldsymbol{3.}$ $1 $;
$\boldsymbol{4.}$ $-\frac{1}{2} $; 
$\boldsymbol{5.}$ $-\frac{\pi}{6} $ (je vhodné použiť modifikáciu vety o limite zloženej funkcie z pr. 124);
$\boldsymbol{6.}$ $5$;
$\boldsymbol{7.}$ $-\frac{1}{4} $;
$\boldsymbol{8.}$ $2 $ (stačí pripočítať a odpočítať $x$);
$\boldsymbol{9.}$ $2^{n}$ ($= \lim_{n \to \infty}((\frac{x-\sqrt{x^{2}-1}}{x})^{n} + (\frac{x+\sqrt{x^{2}-1}}{x})^{n}) $;
$\boxed{184.}$ $\boldsymbol{1.}$ $-\frac{9}{2^{7}} $; 
$\boldsymbol{2.}$ $\sin a $; 
$\boldsymbol{3.}$ $\sqrt{3} $ (pozri pr. 128.4); 
$\boldsymbol{4.}$ $\frac{2}{3\pi} $ (využite, že   $\lim_{u \to 0}\frac{\arctan u}{u}=1$; 
$\boldsymbol{5.}$ $+\infty $;
$\boldsymbol{6.}$ $3 $;
$\boldsymbol{7.}$ $-24 $;
$\boldsymbol{8.}$ $\sqrt{2} $ (ak použijeme subs. $\arccos (1-x)=t$, musíme si uvedomiť, že pre $x\rightarrow 0$ platí $t\rightarrow 0^{+}$, stačí teda hľadať $\lim_{t \to 0^{+}}\frac{t}{\sqrt{1-\cos t}}$;  $\lim_{t \to 0}\frac{t}{\sqrt{1-\cos t}}$ totiž neexistuje);
$\boldsymbol{9.}$  $+\infty $
$\boldsymbol{10.}$  $\frac{3}{2} $;
$\boldsymbol{11.}$  $\frac{5}{8}\sqrt{2}\pi ^{2}$;
$\boldsymbol{12.}$  $-\infty $;
$\boldsymbol{13.}$  $0 $;
$\boldsymbol{14.}$  $0 $;
$\boldsymbol{15.}$  $0 $;
$\boxed{185.}$ $\boldsymbol{2.}$  $0 $ (  $\lim_{n \to \infty}\sqrt[n]{a_{n}}= \frac{1}{\sqrt{2}} $, pre vhodne zvolené  $\varepsilon > 0 $ preto platí $\exists n_{0} \in \mathbb{N}, \, n>n_{0} $: $\sqrt[n]{a_{n}}<\frac{1}{\sqrt{2}}+\varepsilon < 1$); 
$\boxed{186.}$ $\boldsymbol{1,2.}$ možno využiť pr. 133.1;
$\boldsymbol{3.}$ možno použiť nerovnosť $n!>(\frac{n}{3})^{n} $ iná možnosť: pre $i=1,....,n $ platí $i.(n-i+1)\geq n $, odtiaľ $(n!)^{2}\geq n^{n} $;
$\boxed{187.}$ $\boldsymbol{1.}$  $1 $ (použite výsledok pr. 135);
$\boldsymbol{2.}$  $\max \lbrace a,b \rbrace $;
$\boldsymbol{3.}$  $3 $;
$\boldsymbol{4.}$  $11 $;
$\boldsymbol{5.}$  $0 $ (možno použiť pr. 133 a fakt, že $\lim_{n \to \infty}(1+ \frac{1}{n})^{n}=e $);
$\boldsymbol{6.}$  $+\infty $;
$\boldsymbol{7.}$  $+\infty $;
$\boldsymbol{8.}$  $\frac{1}{3} $ (využite pr. 186.1);
$\boldsymbol{9.}$  $1$;
$\boldsymbol{10.}$  $-\infty $;
$\boxed{188.}$ správny je postup v bode a/;
$\boxed{189.}$ Nech je daná funkcia $g$ a kladná funkcia $f$, nech $b$ je hromadný bod množiny $D(f)\cap D(g)$, nech $\lim_{x \to b} f(x)=0$. Potom: a/ ak $\lim_{x \to b} g(x)=+\infty$, tak $\lim_{x \to b} f(x)^{g(x)}=0$; b/ ak $\lim_{x \to b} g(x)=-\infty$, tak $\lim_{x \to b} f(x)^{g(x)}=+\infty$;
$\boxed{190.}$ $\boldsymbol{1.}$  $0$;
$\boldsymbol{2.}$  $e^{2}$;
$\boldsymbol{3.}$  $+\infty$;
$\boldsymbol{4.}$  $e^{-1/2}$;
$\boldsymbol{5.}$  $0$;
$\boldsymbol{6.}$  $0$;
$\boldsymbol{7.}$  $e^{(\beta^{2}-\alpha^{2}})/2$; (prípad $\alpha=\beta$ treba robiť samostatne;
$\boldsymbol{8.}$ neexistuje;
$\boldsymbol{9.}$  $0$ pre $a_{1}>a_{2}$, $e^({b_{1}-b_{2})/a_{1}}$ pre $a_{1}=a_{2}$ (prípad $a_{1}=a_{2}$, $b_{1}=b_{2}$ treba robiť samostatne);
$\boldsymbol{10.}$  $e^{-1}$;
$\boldsymbol{11.}$  $e^{-(a+b)}$ (prípady $a=0,\, b=0$ sa robia zvlášť);
$\boxed{191.}$ $\boldsymbol{1.}$  $- \ln 2$ (pre $x>0 \, x\ne 1$ je $\log_{x}2= \frac{1}{\log_{2}x}$;
$\boldsymbol{2.}$  $\ln 8$;
$\boldsymbol{3.}$  $\frac{1}{8}$;
$\boldsymbol{4.}$  $\alpha a^{\alpha}-1$ pre $\alpha\ne 0$, $0$ pre $\alpha=0$;
$\boldsymbol{5.}$  $a^{a}(1+\ln a)$;
$\boldsymbol{6.}$  $\sqrt{a.b}$ (prípad $a=b=1$ treba robiť samostatne);
$\boldsymbol{7.}$  $\ln x$ (= $\lim_{n \to \infty}\sqrt[n+1]{x}n^{2}(x^{1/n-1/(n+1)}-1)$; prípad $x=1$ treba urobiť samostatne);
$\boldsymbol{8.}$  $\frac{1}{\ln a - \ln b}$;
$\boldsymbol{9.}$  $\frac{1}{sqrt{a.b}}$ (prípad $a=b=1$ treba robiť samostatne);
$\boldsymbol{10.}$  $-2$;
$\boldsymbol{11.}$  $\frac{\alpha^{2}}{\beta ^{2}}$;
$\boldsymbol{12.}$  $-2$;
$\boxed{192.}$ $\boldsymbol{1a/}$  použite nerovnosti $\frac{[x]^{\alpha}}{a^{[x]+1}}$ $\leq \frac{x^{\alpha}}{a^{x}}\leq$ $\frac{([x]+1)^{\alpha}}{a^{[x]}}$ $x\geq0, \, a>1, \, \alpha >0$) a pr. 186.1;
$\boldsymbol{1b/}$  možno použiť substitúciu  $\log_{a}x=t$;
$\boldsymbol{2a/}$   $0$;
$\boldsymbol{2b/}$   $0$; 
$\boldsymbol{2c/}$   $1$;  
$\boxed{194.}$ (je zrejmé, že $g$ musí mať nevlastnú limitu); napr. $g(x)=1$ pre ($x<N, \, g(x)=f(x)$ pre $x \in \langle N+2k, N+2k+1 )$, $g(x)=2f(x)$ pre $x \in \langle N+2k+1, N+2k+2 )$ $k=1,2,...;N \in \mathbb{N}$ je číslo, pre ktoré platí $\forall x\geq N :\, f(x)>0$);
$\boxed{194.}$  $\boldsymbol{1.}$ $\lbrace a_{n} \rbrace _{n=1} ^{\infty}$ je klesajúca zdola ohraničená;
$\boldsymbol{2.}$ $\lbrace a_{n} \rbrace _{n=1} ^{\infty}$ je rastúca  a zhora ohraničená (použite nerovnosť  $\frac{1}{n^{2}}\leq \frac{1}{n-1}- \frac{1}{n}$  ($n>1$) alebo  $\frac{1}{(2^{k})^{2}}+\frac{1}{(2^{k}+1)^{2}}+...+\frac{1}{(2^{k+1}-1)^{2}}\leq \frac{1}{2^{k}}$;
$\boldsymbol{3.}$ $\lbrace a_{n} \rbrace _{n=1} ^{\infty}$ je rastúca  a zhora ohraničená (prísť na  odhad $a_{n}\leq (2+ \frac{1}{2}+ \frac{1}{4}+...+ \frac{1}{2^{n-2}} $) pre $n\geq 3 $ dá asi trocha námahy; iná možnosť: využiť nerovnosť $\ln (1+ \frac{1}{n})< \frac{1}{n} $ ( $ n \in \mathbb{N} $), ktorú dostaneme logaritmovaním nerovnosti z pr. 157, potom $(1+\frac{1}{2^{k}})=e^{ln(1+1/2^{k}})< e^{1/2^{k}}  $ );
$\boxed{196.}$   $\lbrace a_{n} \rbrace _{n=1} ^{\infty}$ je rastúca  a zhora ohraničená (ak postupne pre $k=n,n-1,...,1 $ použijeme nerovnosť $ k< 2^{2^{k}}$, dostaneme  $b_{n}< 2\sqrt{1+\sqrt{1+...+\sqrt{1+\sqrt{2}}}}   $  $\rightarrow n-1 $ odmocnín);
$\boxed{197.}$  $\boldsymbol{1.}$ ak zobrazíme na číselnej osi niekoľko prvých členov, zistíme, že by malo platiť $a_{2k}=1-(\frac{1}{4} +\frac{1}{4^{2}}+...+\frac{1}{4^{k-1}})$,  $ k>1  $,  $a_{2k+1}=\frac{1}{2}(1+\frac{1}{4} +\frac{1}{4^{2}}+...+\frac{1}{4^{k-1}}) $,  $ k\geq 1  $, čo možno dokázať indukciou; potom  $ \lim_{k \to \infty}a_{2k}  $ $= \lim_{k \to \infty}a_{2k+1}  = \frac{2}{3}$ $= \lim_{n \to \infty}a_n  $;
$\boldsymbol{2.}$ $\lim_{n \to \infty}a_n =a$, ak $a=b$; $ \lim_{n \to \infty}a_n =\frac{a}{3}+ \frac{2b}{3} $, ak $a \ne b $ (ak definujeme postupnosť $\lbrace t_n \rbrace_{n=1}^{\infty}$ vzťahom $a+(b-a)t_{n}=a_{n}$, dostaneme postupnosť z pr. 197.1);
$\boxed{198.}$  $\boldsymbol{1,2.}$ platí; obidva dôkazy sú analogické dôkazu vety o Heineho definícii limity;
$\boxed{200.}$ stačí dokázať, že čísla na pravých stranách rovností vyhovujú obidvom podmienkam z vety 18;
$\boxed{201.}$  $\boldsymbol{1.}$ pozri napr. pr.160.7,8;
$\boldsymbol{2.}$ nie (pozri pr.165.1);
$\boxed{202.}$ sporom; označme $a=\lim_{n \to \infty}a_n$, $b=\lim_{n \to \infty}a_n$, ak $c \in (a,b)$, ale $c \notin H$, tak existuje $\eta > 0 $ tak, že $(c-\eta ,c+\eta) \cap H=\varnothing $ (pozri pr. 165.1, pritom $a\notin (c-\eta ,c+\eta)$, $b\notin (c-\eta ,c+\eta)$; potom $\lbrace n \in \mathbb{N};\, a_{n }\in (c-\eta /2 ,c+\eta /2)) \rbrace$ je konečná, preto $\exists N_{1} \in \mathbb{N}\, \forall n\in \mathbb{N}\,  n > N_{1}: \, (a_{n} > c+ \eta /2 \, \lor \, a_{n} < c- \eta /2 )$; z podmienky  $\lim_{n \to \infty} (a_{n+1} - a_{n})=0$ vyplýva $\exists N_{2} \in \mathbb{N}\,  \forall n \in \mathbb{N},\, n >N_{2}: \quad \vert a_{n+1} - a_{n} \vert < \eta $; nech $N= \max \lbrace N_{1}, N_{2} \rbrace $, potom platí ( $\forall n \in \mathbb{N},\, n >N: \quad a_{n} < c-\eta /2 $) $\lor $ ( $\forall n \in \mathbb{N},\, n >N: \quad a_{n} > c+ \eta /2 $); to je ale spor s faktom, že $a$  aj   $b$ sú hromadné hodnoty postupnosti $\lbrace a_{n} \rbrace _{n=1} ^{\infty} $;
$\boxed{204.}$  $\boldsymbol{1.}$  všetky postupnosti, ktorých hromadnou hodnotou je číslo $0 $;
$\boldsymbol{2.}$  všetky postupnosti z pr. 204.1 a všetky postupnosti $\lbrace a_{n} \rbrace _{n=1} ^{\infty} $, pre ktoré platí $\exists  n \in \mathbb{N}, n\geq 2: \, a_{n} =0$;
$\boxed{204.}$  $\boldsymbol{1.}$ vyplýva to z vety 19 a z tvrdenia: ak $\lbrace c_{n} \rbrace _{n=1} ^{\infty} $ je postupnosť kladných čísel a $\limsup_{n \to \infty}c_n=L \in (0,\infty) $, tak $\liminf_{n \to \infty} \frac{1}{c_n}=\frac{1}{L} $;
$\boldsymbol{2.}$ pre všetky, ktoré vyhovujú podmienke $0 <\liminf_{n \to \infty}a_n < +\infty$;
$\boxed{206.}$  $\boldsymbol{1.}$ v prípade  $\lim_{n \to \infty}a_n =0 $ uvedená rovnosť iste platí, pretože vtedy  $\lim_{n \to \infty}a_n . b_{n} =0 $; v prípade  $\lim_{n \to \infty}a_n \ne 0 $ sa postupuje analogicky ako v pr. 166;
$\boxed{207.}$ sporom; nech teda ($*$) $\exists N \in \mathbb{N} \, \forall n \in \mathbb{N}, n>N $ $\exists k \in \mathbb{N} \, k<n:\, a_{k} \leq a_{n} $; nech $m= \min \lbrace a_{n}; n=1,...,N \rbrace$ (zrejme $m>0 $), potom $a_{N+1}\geq m $  (podľa ($*$) niektoré z čísel $a_{1}, ... , a_{N}$  je menšie alebo rovné číslu $a_{N+1}$) atď, teda $\forall n\in \mathbb{N}: \, a_{n}\geq m $, čo je spor s predpokladom $\liminf_{n \to \infty}a_n =0 $;











\section{Dodatok 2}
$\boxed{208.}$  $\boldsymbol{1.}$ nie je otvorená (2 a 3 nie sú vnútorné body), nie je uzavretá (neobsahuje svoje hromadné body 0, 1);
$\boldsymbol{2.}$ otvorená , nie je uzavretá (neobsahuje napr. svoj hromadný bod 0);
$\boldsymbol{3.}$ nie je otvorená (bod 0),  nie je uzavretá (napr.  bod 1);
$\boldsymbol{4.}$ uzavretá, nie je otvorená;
$\boldsymbol{5.}$  nie je otvorená, nie je uzavretá;
$\boldsymbol{6.}$  $\chi (R)=\lbrace 0,1\rbrace $,  nie je otvorená, je uzavretá;
$\boldsymbol{7.}$ nie je otvorená,  nie je uzavretá (neobsahuje svoj hromadný bod 0);
$\boldsymbol{8.}$ $f(R)=\lbrace \frac{1}{n}; n \in \mathbb{N} \rbrace  \cup \lbrace 0 \rbrace$; je  uzavretá, nie je otvorená;
$\boldsymbol{9.}$  otvorená , nie je uzavretá;
$\boxed{209.}$  $\boldsymbol{2.}$ stačí dokázať, že doplnky sú otvorené; využite de Morganove vzorce $\mathbb{R}-(A \cup B)= (\mathbb{R} - A) \cap ( \mathbb{R} - B) $, $\mathbb{R}-(A \cap B)= (\mathbb{R} - A) \cup ( \mathbb{R} - B) $ a pr. 209.1;
$\boldsymbol{3.}$ $ A-B=A \cap (\mathbb{R} - B)$, ďalej využite pr. 209.1,2;
$\boldsymbol{4.}$ nech $x \in \bigcup\limits_{\alpha \in I }A_{\alpha} $, potom  $x \in A_{\alpha_{0}}  $ pre  niektoré $\alpha_{0}\in I $; $ A_{\alpha_{0}} $ je otvorená množina, teda pre niektoré okolie $O(x)$ bodu $x$ platí $O(x) \subset A_{\alpha_{0}}$; pretože $A_{\alpha_{0}}\subset \bigcup\limits_{\alpha \in I}A_{\alpha} $, platí aj $O(x) \subset \bigcup\limits_{\alpha \in I}A_{\alpha}$; teda každý prvok množiny $\bigcup\limits_{\alpha \in I}A_{\alpha}$ je jej vnútorným bodom;
$\boxed{210.}$ $A+B=\bigcup\limits_{b \in B} ( \lbrace b \rbrace + A)$, každá z množín $\lbrace b \rbrace +A $ je otvorená, zjednotenie ľubovoľného systému otvorených množín je otvorená množina (pozri pr. 209.4);
$\boxed{211.}$ $\boldsymbol{1.}$ $(a,b)=\bigcup\limits_{n \in \mathbb{N}} \langle a+ \frac{b-a}{n+2}, b-\frac{b-a}{n+2} \rangle $;
$\boldsymbol{2.}$ nepriamo, ak $\langle a,b \rangle =\bigcup\limits_{t \in I}A_{t} $, kde $ A_{t} $ je otvorený interval pre každé $t \in I $, tak $\exists t_{0} \in I: \, b \in A_{t_{0}} \subset \langle a,b \rangle$, teda $b$ je vnútorný bod množiny $\langle a,b \rangle$, čo neplatí;
$\boxed{212.}$ $C =\bigcup\limits_{b \in B}C_{b} $, kde $C_{b}= \lbrace \vert x-b \vert; x \in A \rbrace $; pre $b \notin A $  je množina $ C_{b}$ otvorená (ak $ \varepsilon  $-okolie bodu  $x$ patrí do $A$, tak $ \varepsilon  $-okolie bodu $\vert x-b \vert $ patrí do $ C_{b}$); pre $b \in A $ možno $ C_{b}$ písať v tvare $\lbrace 0 \rbrace \cup \lbrace \vert x-b \vert ; x \in A- \lbrace b \rbrace \rbrace $, pritom $A-\lbrace b\rbrace  $ je otvorená podľa pr. 209.3 a $b \notin A-\lbrace b \rbrace $, preto - v predchádzajúcej úvahe stačí nahradiť množinu $A$   množinou $  A- \lbrace b \rbrace $  - je $\lbrace \vert x-b \vert; x \in A - \lbrace b \rbrace  \rbrace$ otvorená množina;
dokončenie dôkazu ponechávame čitateľovi;
$\boxed{213.}$ sporom, ak $\varnothing \ne A \ne \mathbb{R}$ je súčasne uzavretá a otvorená, tak $ A $ aj $ \mathbb{R} -A $ sú otvorené a neprázdne; nech $ a \in A $, $b \in \mathbb{R} - A$, predpokladajme $ a<b $ (pre $ a>b $ je postup analogický); nech $ c:= \sup \lbrace x \in (a,b); \, x \in A \rbrace $, potom $ c \ne a $  (lebo  $ A $  je otvorená), $ c \ne b$  (lebo  $\mathbb{R} - A $  je otvorená); pritom  $ c  $ nemôže ležať v $ A $ ($ c $ by muselo patriť do $ A $  spolu s niektorým svojim okolím, teda existovali by prvky patriace do $ A $  a ležiace v intervale $ (c,b) $, čo odporuje vlastnostiam $ c $  ako suprema) a nemôže patriť ani do $\mathbb{R} - A $, čo je spor; 
$\boxed{214.}$ využite, že množina $A \subset \mathbb{R}$ je kompaktná práve vtedy, keď je uzavretá a ohraničená (tvrdenie pr. 209.2 možno matematickou indukciou rozšíriť na konečný počet množín);
$\boxed{215.}$ ak $N:= \max \lbrace n \in \mathbb{N};\, \frac{1}{n} \geq \epsilon \rbrace $, tak konečné podpokrytie je $\lbrace (-\varepsilon , \varepsilon),  ( 1- \varepsilon , 1+ \varepsilon),..., (\frac{1-\varepsilon}{2^{N}},\frac{1+\varepsilon}{2^{N}}) \rbrace$;
$\boxed{216.}$ nie, každý y uvedených intervalov obsahuje práve jeden prvok nekonečnej množiny $E$;
$\boxed{217.}$ napr. $\lbrace ( \frac{1}{n},2);\, n \in \mathbb{N} \rbrace $;
$\boxed{218.}$ $\boldsymbol{1.}$ z množiny $A_{n}$ vyberme prvok $a_{n}$ ( potom $\lbrace a_{n} \rbrace _{n=k}  ^{\infty}$ je postupnosť prvkov z  $A_{k}$), z $\lbrace a_{n} \rbrace _{n=1}  ^{\infty}$ možno vybrať konvergentnú postupnosť, jej limita leží v každej z množín  $A_{n}$, $ n \in \mathbb{N}$, teda leží aj v $ A $;
$\boldsymbol{2.}$ $\bigcap\limits_{n \in \mathbb{N}} A_{n} = \bigcap\limits_{n \in \mathbb{N}} B_{n}$, kde $B_{n}:= A_{1}\cap A_{2}\cap ... \cap A_{n}$;  $\lbrace B_{n} \rbrace _{n=1}  ^{\infty}$ vyhovuje podmienkam z pr. 218.1;
$\boxed{219.}$ $\boldsymbol{1.}$ áno (využite, že každú otvorenú množinu možno písať v tvare zjedotenia otvorených intervalov);
$\boldsymbol{2.}$ nie (túto vlastnosť majú len konečné množiny; využite, že $\lbrace \lbrace a \rbrace;a \in A \rbrace$  je pokrytie množiny $A \ne \varnothing $  uzavretými množinami);
$\boldsymbol{3.}$ nie (uvažujte napr. množinu $\langle 0,1\rangle $ a jej pokrytie $\lbrace \langle -\frac{1}{4}, \frac{1}{4} \rangle$, $\langle \frac{3}{4}, \frac{5}{4} \rangle \rbrace \cup $ $\lbrace \langle \frac{1}{4} +\frac{1}{n} , \frac{3}{4} - \frac{1}{n} \rangle$; $n\geq 3 \rbrace $; platí ale, že každá množina s uvedenou vlastnosťou je kompaktná, využite pr. 219.1);
$\boldsymbol{4.}$ nie (uvažujte napr. množinu $\langle 0,1\rangle$ a jej  pokrytie $\lbrace \langle 1, 2 ) \rbrace \cup \lbrace\langle 0, \frac{n-1}{n} ); \, n \geq 2 \rbrace $; platí ale,  že každá množina s uvedenou vlastnosťou je kompaktná);
$\boxed{220.}$ využite ekvivalenciu $a/ \iff c/$ z vety 21; ak $c_{n} \in A+B$, $n \in \mathbb{N} $, tak existujú $a_{n} \in A$, $b_{n} \in B$ ($n \in \mathbb{N} $) tak, že $\lbrace c_{n} \rbrace _{n=1}  ^{\infty}$ $=\lbrace a _{n}+b_{n} \rbrace _{n=1}  ^{\infty}$; z $\lbrace a_{n} \rbrace _{n=1}  ^{\infty}$ možno vybrať konvergentnú $\lbrace a_{n_{k}} \rbrace _{k=1}  ^{\infty}$, z $\lbrace b_{n} \rbrace _{n=1}  ^{\infty}$ možno vybrať konvergentnú $\lbrace b_{n_{k}} \rbrace _{k=1}  ^{\infty}$, potom  $\lbrace c_{n_{k_{l}}} \rbrace _{l=1}  ^{\infty}$ je konvergentná a jej limita leží v množine  $A+B $; podobne pre $A.B$;
 





\section{Spojitosť funkcie}
$\boxed{221.}$ $\exists \varepsilon > 0 \, \forall \delta > 0 \, \exists  x_{\delta} \in D(f): \vert x_{\delta} -a \vert < \delta \land \vert f(x_{\delta}) - f(a) \vert \geq \varepsilon $;
$\boxed{222.}$ $\boldsymbol{1.}$ $0$ je hromadný bod množiny $D(f)$, $\lim_{x \to 0} f(x) = 1 = f(0)$, teda $f$ je spojitá v bode $0$;
$\boldsymbol{2.}$  nie je spojitá v bode $0$ ($0$  je hromadný bod $D(f)$  a   $\lim_{x \to 0} f(x) $ neexistuje);
$\boldsymbol{3.}$ nie je spojitá v bode $-1$;
$\boldsymbol{4.}$ nie je spojitá v bode $0$;
$\boldsymbol{5.}$ je spojitá v bode $0$;
$\boxed{223.}$ $\boldsymbol{1.}$ $f(0)$ treba položiť rovné $\frac{3}{2} $ ($= \lim_{x \to 0} f(x)$);
$\boldsymbol{2.}$ $f(0):=\sqrt{e}$;
$\boldsymbol{3.}$ $f(0):=0$;
$\boldsymbol{4.}$ $f(1):=1$;
$\boldsymbol{5.}$ $f(-3):=0$;
$\boxed{224.}$ $\boldsymbol{1.}$ spojitá (vyplýva to z viet 1 a 3);
$\boldsymbol{3.}$ spojitá  v každom bode množiny $\mathbb{R}- \lbrace -1,0,1 \rbrace $, nespojitá v bodoch $-1,0,1 $;
$\boldsymbol{4.}$ spojitá len  v bodoch množiny $(\mathbb{R}- \mathbb{Z}) \cup \lbrace 0 \rbrace $;
$\boldsymbol{5.}$ $f(x)=1$ pre $x \in \langle0,1)$, $f(x):=\frac{1}{2}$  pre $x =1$, $f(x):=0$  pre $x > 1$; $f$ je spojitá práve v bodoch množiny $\langle 0, +\infty) -\lbrace 1 \rbrace $;
$\boldsymbol{6.}$ $f(x)=1$ pre $ \vert x \vert \leq 1$, $f(x)=x^{2}$  pre $ \vert x \vert >1$; $f$ je spojitá;
$\boldsymbol{7.}$ (množina $D(f)= \mathbb{Z}$ pozostáva len z izolovaných bodov);
$\boxed{227.}$ nemožno; napr. 1. ak $x_{0}$ je izolovaný bod množiny $D(f)$, je funkcia $g\circ f $ spojitá v bode $x_{0}$ ($x_{0}$ je potom totiž aj izolovaný bod množiny  $D(g\circ f)$); alebo 2. $f(x)= const$, $x \in  \mathbb{R}$  ; $ g $ je ľubovoľná funkcia nespojitá v bode $ f(a)$;
$\boxed{228.}$ pri dôkaze spojitosti funkcie $\vert f \vert$ v bode $a$ využite nerovnosť $\vert \vert r \vert - \vert s \vert \vert \leq \vert r-s \vert$ ($r,s \in \mathbb{R}$) a spojitosť funkcie  $f$  v bode $a$; ďalej platí $\max \lbrace f,g \rbrace =\frac{f+g+\vert f-g\vert}{2}$, analogicky možno vyjadriť $\min \lbrace f,g \rbrace $;
$\boxed{229.}$ funkcia $\chi $ nie je spojitá v žiadnom bode  množiny $\mathbb{R}$;
$\boxed{230.}$ nech $a \in \mathbb{R} - \mathbb{Q} $, nech $\varepsilon  >0$ je dané; nech $\overline{n} \in \mathbb{N} $ je také, že $\frac{1}{\overline{n}}< \varepsilon $; pre každé $n \in \lbrace1,...,\overline{n} \rbrace $ existuje $\delta_{n}>0 $ tak, že v $O_{\delta_{n}}(a)$ neležia body s funkčnými hodnotami  $\frac{1}{n} $; pre dané  $\varepsilon $ potom stačí položiť  $\delta = \min \lbrace \delta_{1},...,\delta_{\overline{n}} \rbrace$; dôkaz druhého tvrdenia z pr. 230 prenechávame čitateľovi;
$\boxed{232.}$ $\boldsymbol{1.}$ $-1$ je bod neodstrániteľnej nespojitosti, bod nespojitosti 2. druhu;
$\boldsymbol{2.}$ $1$ a $-2$ sú body neodstrániteľnej nespojitosti, body nespojitosti 2. druhu;
$\boldsymbol{3.}$ $0$ je bod odstrániteľnej nespojitosti;
$\boldsymbol{4.}$ $1$ a $0$ sú body odstrániteľnej nespojitosti, $-1$ je bod neodstrániteľnej nespojitosti, bod nespojitosti 2. druhu;
$\boldsymbol{5.}$ $2$ a $-2$ sú body odstrániteľnej nespojitosti;
$\boldsymbol{6.}$ $0$ je bod neodstrániteľnej nespojitosti, bod nespojitosti 2. druhu;
$\boldsymbol{7.}$ $0$ je bod odstrániteľnej nespojitost;
$\boldsymbol{8.}$ $0$ je bod odstrániteľnej nespojitost, $1$ je bod neodstrániteľnej nespojitosti, bod nespojitosti 2. druhu;
$\boldsymbol{9.}$ $-1$ je bod neodstrániteľnej nespojitost, bod nespojitosti 1. druhu;
$\boldsymbol{10.}$ každý prvok množiny $\lbrace \frac{1}{n}; \, n \in \mathbb{Z} -\lbrace 0 \rbrace \rbrace$ je bod neodstrániteľnej nespojitosti, bod nespojitosti 1. druhu;  $0$ je bod odstrániteľnej nespojitost  $\lim_{x \to 0} f(x) =1$; platí totiž $f \langle ( \frac{1}{n+1}, \frac{1}{n} \rangle ) $ pre  $n \in \mathbb{N}$,  $f (( \frac{1}{n}, -\frac{1}{n+1} \rangle ) = \langle 1, \frac{n+1}{n})$ pre  $n \in \mathbb{N}$;
$\boldsymbol{11.}$ $f(x) = x$, $D(f) = \mathbb{R} - \lbrace 0,1 \rbrace$; $0$ a $1$ sú body odstrániteľnej nespojitosti;
$\boldsymbol{12.}$ každý prvok  $a \in\mathbb{R}- \mathbb{Z} $ je bod neodstrániteľnej nespojitosti, bod nespojitosti 2. druhu;
$\boldsymbol{13.} $ platí $g(\varphi (x)) = g(x)$;  $0$ a  $2$ sú body odstrániteľnej nespojitost (všimnite si napriek tomu, že ;$\varphi $ je spojitá len v bode $1$, je $g \circ \varphi $ spojitá fubkcia);
$\boxed{234.}$ $\boldsymbol{2.}$ pre spojitú funkciu  $f(x) = x- \cos x$ platí  $f(0)<0$, $f(\frac{\pi}{2})>0$;
$\boldsymbol{3.}$ $\lim_{x \to \infty} P(x) $ a $\lim_{x \to -\infty} P(x) $ sú nevlastné a opačných znamienok, preto musia existovať $a,b \in \mathbb{R}$ tak, že $P(a)> 0$, $P(b)< 0 $ ;
$\boxed{235.}$ využite nerovnosti $\min \lbrace f(x_{1}),..., f(x_{n}) \rbrace \leq \frac{1}{n}(f(x_{1})+...+f(x_{n}))\leq \max \lbrace f(x_{1}),..., f(x_{n}) \rbrace$;
$\boxed{236.}$ ak $f(0)=0$ alebo   $f(1)=1$, je dôkaz skončený; ak  $f(0)\ne 0$, $f(1)\ne 1$, tak  pre  spojitú  funkciu  $g(x) = f(x)-x $ platí  $fg(0)>0$, $fg(1)< 0$;
$\boxed{238.}$ stačí dokázať tvrdenie: každá nekonštantná funkcia s uvedenou vlastnosťou je spojitá;
$\boxed{240.}$ $f$ nadobúda na $\langle a,b \rangle$ minimum, ktoré (pretože je funkčnou hodnotou), je kladné;
$\boxed{241.}$ ak $a,b \in \mathbb{R}$, tak funkcia $g$ určená predpisom $g(a) = \lim_{x \to a^{+} }f(x), \, g(b) = \lim_{x \to b^{-}} f(x), \, g(x)=f(x)$ pre $x \in (a,b)$ je spojitá, a teda ohraničená na kompakte $\langle a,b \rangle$; $f$ je potom ohraničená ako zúženie ohraničenej funkcie; vo všeobecnosti (tj. ak $a, b \in \mathbb{R^{*}}$) možno dokazovať nasledovne: $f$ je ohraničená na niektorom okolí $O(a)$ (to vyplýva z existencie vlastnej  $\lim_{x \to a} f(x)$)  a na niektorom okolí $O(b)$; množina $A:=(a,b)-(O(a) \cup O(b)) $ môže byť prázdna, jednoprvková (v týchto prípadoch je dôkaz skončený) alebo kompaktný interval, na tom je $f$ ohraničená;
$\boxed{242.}$ ak  $\lim_{x \to \infty }P(x) = \lim_{x \to -\infty }P(x)= +\infty $, tak existuje $a>0$ tak, že $\forall x \in \mathbb{R}, \, \vert x \vert > a: \, f(x) > f(0) $; na kompakte $\langle -a,a \rangle$ nadobúda $f$ minimum, ktoré je globálnym minimom na $ \mathbb{R} $;
$\boxed{245.}$ dôkaz je analogický dôkazu veta 5;
$\boxed{246.}$ napr. $f(x)=e^{x} \sin \frac{1}{x}$, $ x \in (0,1 \rangle$; 
$\boxed{247.}$ ($\forall x \in (a,b) \quad \forall\varepsilon  >0 \quad   \exists \delta >0 \,  \forall y \in (a,b): \, \vert x-y \vert < \delta \, \Rightarrow \, \vert f(x)-f(y) \vert < \varepsilon$)  $\land$ ( $\exists \varepsilon _{0} >0 \quad \forall \delta  >0 \, \exists x_{\delta}, y_{\delta}  \in (a,b): \vert  x_{\delta}- y_{\delta}\vert <  \delta $  $\land$  $\vert f(x_{\delta})-f(y_{\delta}) \vert \geq \varepsilon_{0} $);
$\boxed{248.}$  $\boldsymbol{1.}$ áno;
$\boldsymbol{3.}$  nie;
$\boldsymbol{4.}$ áno;
$\boldsymbol{6.}$ áno (Ťažkosti spôsobuje bod $0$, keby sme boli na intervale $\langle \eta , +\infty)$, kde $\eta >0$ , mohli by sme použiť odhad $ \vert \sqrt{x} - \sqrt{y} \vert =\vert \frac{x-y}{\sqrt{x} + \sqrt{y}}\vert \leq \frac{\vert x-y \vert}{2 \sqrt{\eta}} $. Možmo uvažovať nasledovne: nech je da né $\varepsilon  >0 $; ak $0 <  x <  \varepsilon ^{2} $, $0 <  y <  \varepsilon ^{2} $, tak  $ \vert \sqrt{x} - \sqrt{y} \vert < \varepsilon$. Zostáva sa ešte zaoberať prípadom $x\geq \frac{\varepsilon^{2}}{2} $, $y\geq \frac{\varepsilon^{2}}{2} $; potom $ \vert \sqrt{x} - \sqrt{y} \vert \leq \frac{\vert x-y \vert}{2 \sqrt{\eta}}$; teda pre dané $\varepsilon$ a $x,y \geq \frac{\varepsilon^{2}}{2}$ stačí položiť $\delta <  \sqrt{2}\varepsilon^{2} $, tj. napr. $\delta =\frac{\varepsilon^{2}}{2}$. Celkovo sme dokázali: ak  $ \vert x-y  \vert < \frac{\varepsilon^{2}}{2}$, tak $ \vert \sqrt{x} - \sqrt{y} \vert < \varepsilon$.);
 $\boldsymbol{7.}$ áno  (využite vzorec $ \sin x - \sin y = 2 \sin \frac{x-y}{2} \cos \frac{x+y}{2}$  a nerovnosť $\sin \alpha < \vert \alpha\vert$);
$\boxed{249.}$ "$\Rightarrow$ ": k číslu $\frac{\varepsilon}{2}$	nájdime $\delta$ z definície rovnomernej spojitosti, nech pre  body $a=x_{0 }< x_{1}<...< x_{n}=b$ platí $\vert x_{i} -x_{i-1}\vert < \delta$ ($i=1,...,n $), potom graf funkcie $\varphi $ je lomená čiara s vrcholmi  $(x_{0 }, f(x_{0})),(x_{1 }, f(x_{1})), ..., (x_{n }, f(x_{n}))$; pri dôkaze implikácie "$\Leftarrow$" využite, že každá spojitá po častiach lineárna funkcia definovaná na ohraničenom intervale $I$ je rovnomerne spojitá (to možno dokázať priamo z definície), $\vert f(x)-f(y) \vert \leq  \vert f(x)-\varphi(x) \vert + \vert \varphi(x)-\varphi(y) \vert + \vert \varphi(y)-f(y) \vert$;
$\boxed{250.}$ návodom je dôkaz implikácie "$\Rightarrow$ " v pr. 249;
$\boxed{251.}$ $\boldsymbol{1.}$ áno ($f$ je spojitá, a teda podľa vety 6 aj rovnomerne spojitá na kompakte $\langle -1,1 \rangle $);
$\boldsymbol{3.}$ áno (zvoľme $a>0$ pevné; rovnomernú spojitosť na $\langle a,+ \infty) $ možno dokázať z definície, rovnomerná spojitosť na $\langle 0,a \rangle $ vyplýva z vety 6; ak je funkcia rovnomerne spojitá na $\langle 0,a \rangle $ a na  $\langle a,+ \infty) $, tak je rovnomerne spojitá na ich zjednotení);
$\boldsymbol{4.}$ áno ( $f$ je spojitá na kompakte $\langle 0,\pi \rangle$);
$\boldsymbol{5.}$ nie (rodeľme každý z intervalov $\langle \frac{\pi}{2}+k \pi,\frac {3 \pi}{2}+ k\pi \rangle $, $k=0,1,2,...,$ na $n$ rovnakých častí bodmi $x_{0}^{(k)}:= \frac{\pi}{2}+k\pi$,  $x_{1}^{(k)}:= \frac{\pi}{2}+k\pi + \frac{\pi}{n}$,...,$x_{n-1}^{(k)}:= \frac{\pi}{2}+k\pi + \frac{(n-1)\pi}{n}$, $x_{n}^{(k)}:= \frac{3\pi}{2}+k\pi $. Pretože $2\pi({1+k})=$ $\vert f(x_{0}^{(k)})-f(x_{n}^{(k)})  \vert\leq $ $\vert f(x_{0}^{(k)})-f(x_{1}^{(k)})  \vert+...+\vert f(x_{n-1}^{(k)})-f(x_{n}^{(k)})  \vert $, musí byť aspoň jedno z čísel na pravej strane tejto  nerovnosti väčšie ako $\frac{2\pi(1+k)}{n} $. Teda pre ľubovoľné $n \in \mathbb{N}$ existuje v každom z intervalov $\langle \frac{\pi}{2}+k\pi,\frac{3\pi}{2}+k\pi\rangle $ dvojica ($x_{k},y_{k} $) taká, že $\vert x_{k}-y_{k}\vert \leq  \frac{\pi}{n}$ a $\vert f(x_{k})-f(y_{k})\vert > \frac{2\pi(1+k)}{n}$, pritom $\frac{2\pi(1+k)}{n} \rightarrow +\infty$ pre $k \rightarrow +\infty$);
$\boldsymbol{6.}$ áno ($\lim_{x \to 1}f(x)=\lim_{x \to -1}f(x)=0 $,ďalej pozri riešenie pr. 251.2);
$\boldsymbol{7.}$ áno (Ak $x,y\geq 1 $, $0< x-y < \eta  $, tak  $\frac{x}{y}=\frac{x-y+y}{y} < 1+\eta $.
Pretože $x>y\geq 1 $ je nerovnosť  $\vert\ln x - \ln y \vert < \varepsilon $ ekvivalentná s nerovnosťou $\frac{x}{y}< e^{\varepsilon}=1+(e^{\varepsilon}-1) $, stačí pre dané $\varepsilon>0$ položiť $\delta= \varepsilon-1 $.);
$\boxed{252.}$ ak $a:=\lim_{x \to \infty }f(x)$, tak k danému  $\varepsilon >0$ existuje $O(+\infty) $ tak, že $f(O(+\infty)) \subset (a-\frac{\varepsilon}{2}, a+\frac{\varepsilon}{2}) $; pre $ x,y \in O(+\infty)$ potom platí $\vert  f(x)-f(y) \vert < \varepsilon$; existuje kompaktný interval $I$ obsahujúci množinu $ A:= \langle 0,+\infty)-O(+\infty) $ (pozn.: v prípade $A=\varnothing$ netreba už vlastne postupovať ďalej, pre dané $\varepsilon>0$ stačí za $\delta$
 zvoliť ľubovoľné kladné číslo; úvahu s intervalom   $I$  možno ale použiť aj na tento prípad, vyhneme sa tým predlžovaniu dôkazu spôsobenému rozlišovaním jednotlivých možmostí) na  $I$ je  $f$ rovnomerne spojitá, a preto k číslu  $\varepsilon$ možno nájsť príslušné  $\delta$ z definície rovnomernej spojitosti (na  $I$; zvyšok úvah, ako nájsť  $\delta$ vyhovujúce definícii rovnomernej spojitosti na celom intervale  $\langle 0, +\infty)$, prenecháme čitateľovi; 
$\boxed{254.}$ $\boldsymbol{1.}$ sporom; nech $f$ je rovnomerne spojitá na $(a,b)$, zvoľme $\varepsilon >0$ pevné, napr.   $ \varepsilon=1$ a nájdime k nemu príslušné $\delta$ z definície rovnomernej spojitosti; potom pre ľubovoľné $x \in (b-\delta, b) $ je $\vert x-(b-\delta) \vert < \delta$, preto $\vert f(x)-f(b-\delta)) \vert < \varepsilon$; funkcia $f$ je teda ohraničená na $(b-\delta, b) $, čo je spor s tým, že $\lim_{x \to b^{-} }f(x)=+\infty$;
$\boldsymbol{2.}$ ak neexistuje  $\lim_{x \to b^{-} }f(x)$, tak existujú postupnosti $\lbrace a_{n} \rbrace _{n=1}^{\infty} $, $\lbrace b_{n} \rbrace _{n=1}^{\infty} $ ( $a_{n}, b_{n} \in (a,b))$ tak, že $\lim_{n \to \infty }a_{n}=\lim_{n \to \infty }b_{n}=b$, ($A:=)$  $\lim_{n \to \infty }f(a_{n}) \ne \lim_{n \to \infty }f(b_{n})$ ($B:=)$ (pozri pr. 199); ak $A, B \in \mathbb{R}, \, A > B $ (ostatné prípady sú analogické), tak v ľubovoľne malom okolí bodu  $b$ existujú body s funkčnými hodnotami väčšími ako  $ A-\frac{A-B}{3} $ aj body s funkčnými hodnotami menšími ako  $ B+\frac{A-B}{3} $; (iný možný dôkaz tvrdení z pr.254: ak  $f $ je rovnomerne spojitá na  $(a,b) $, tak v bode  $b $ je splnené Cauchyho-Bolzanovo kritérium konvergencie);
$\boxed{255.}$ inšpiráciou pre dôkaz  $" \Leftarrow "$ je pr. 251.2, pri dôkaze $" \Rightarrow "$ využite pr. 254;
$\boxed{256.}$ $\boldsymbol{1.}$ využite pr. 255 a 241;
$\boldsymbol{2.}$ pozri pr. 248.3;
$\boxed{257.}$ $\boldsymbol{1.}$ každý prvok množiny $\lbrace (-1)^{k}\sqrt{n}; \, k \in \lbrace -1,1 \rbrace, n \in \mathbb{N}\rbrace $ je bod neodstrániteľnej nespojitosti, bod  nespojitosti 1. druhu; $0$ je bod neodstrániteľnej nespojitosti, bod  nespojitosti 2. druhu;
$\boldsymbol{2.}$ každý prvok množiny $\lbrace \frac{(-1)^{k}}{n}; \, k \in \lbrace -1,1 \rbrace, n \in \mathbb{N}\rbrace $ je bod neodstrániteľnej nespojitosti, bod  nespojitosti 1. druhu; $0$ je bod neodstrániteľnej nespojitosti, bod  nespojitosti 2. druhu;
$\boldsymbol{3.}$  $0$ je bod neodstrániteľnej nespojitosti, bod  nespojitosti 2. druhu;
$\boldsymbol{4.}$ každý prvok množiny $\lbrace \frac{k \pi}{2}; \, k \in  \mathbb{Z}- \lbrace 0 \rbrace\rbrace $ je bod neodstrániteľnej nespojitosti, bod  nespojitosti 1. druhu; 
$\boldsymbol{5.}$ $f$ je spojitá funkcia;
$\boldsymbol{6.}$ $f$ je spojitá funkcia; 
$\boxed{258.}$ (pozri pr. 192) $\boldsymbol{1.}$ $0$;
$\boldsymbol{2.}$ $1$;
$\boldsymbol{3.}$ $0$; 
$\boxed{259.}$ (možno využiť pr. 228) ($F(x)= \min \lbrace c, \max \lbrace f(x), -c \rbrace \rbrace $); 
$\boxed{260.}$ obrátená implikácia neplatí (položme napr. $ a_{n}=f(q_{n})$, kde $f(x)=0$ pre $x \in \mathbb{Q}-\mathbb{N}, \, f(x)=1 $ pre $x \in \mathbb{N}$ ;
$\boxed{261.}$ $f$ je spojitá v každom bode množiny $A= \lbrace (\mathbb{R}-\mathbb{Q} \cap (0+\infty))$ a nespojitá v bodoch množiny $\mathbb{R} -A$;
$\boxed{262.}$  $f$ je spojitá;
$\boxed{263.}$ $\boldsymbol{1.}$ nevyplýva (vhodný protipríklad už nájdite sami); uvedená podmienka je ekvivalentná s podmienkou:  $f$ je ohraničená v každom okolí bodu  $x_{0}$;
$\boldsymbol{2.}$ nevyplýva (uvažujte napr. $f(x) = x$ pre $x<0$,   $f(x) = x+1$ pre $x\geq 0$); uvedená podmienka je ekvivalentná s podmienkou:  $f$ má spojitú inverznú funkciu;
$\boxed{264.}$ $\boldsymbol{1.}$ nie (možno využiť pr. 232.13);
$\boldsymbol{2.}$ áno (pri dôkaze možno využiť pr.239);
$\boxed{266.}$ napr. $f(x)=x$ pre $x \in \mathbb{Q}- \lbrace 0,1 \rbrace $, $f(x)=-x$ pre $x \in \mathbb{R}- \mathbb{Q} $, $f(0)=1$, $f(1)=0$;
$\boxed{267.}$ Sporom. Zvoľme $x_{0}$, platí $1_{0}:=\lim_{x \to x_{0} }f(x) \ne f(x_{0})$, preto $\exists \varepsilon_{0} > 0 $  $ \exists\delta_{0 } > 0: f(x_{0}) \notin \langle 1_{0}-\varepsilon_{0}, 1_{0}+\varepsilon_{0} \rangle \, (=:B_{0}) \, \land f(\langle x_{0}+\delta_{0}/3, x_{0}+\delta_{0}\rangle ) \subset B_{0}$. Pre bod $x_{1} :=x_{0}+2\delta_{0}/3$ platí $l_{1}:= \lim_{x \to x_{1} }f(x) \ne f(x_{1}) $, preto $\exists 0< \varepsilon_{1} < \varepsilon_{0}/2 $  $ \exists 0<\delta_{1 } < \delta_{0}/3: f(x_{1}) \notin \langle 1_{1}-\varepsilon_{1}, 1_{1}+\varepsilon_{1} \rangle \, (=:B_{0}) \, \land f(\langle x_{1}+\delta_{1}/3, x_{1}+\delta_{1}\rangle ) \subset \langle l_{1}-\varepsilon_{1},l_{1}+\varepsilon_{1} \rangle$; teraz položme  $x_{2} :=x_{1}+2\delta_{1}/3$ atď. Tak dostaneme postupnosť prvkov $\lbrace x_{n} \rbrace _{n=0}^{\infty}$, postupnosť uzavretých intervalov  $\lbrace B_{n} \rbrace _{n=0}^{\infty} = \lbrace \langle l_{n}-\varepsilon_{n} , l_{n}+\varepsilon_{n}\rangle \rbrace _{n=0}^{\infty}$ a postupnosť uzavretých intervalov  $\lbrace A_{n} \rbrace _{n=0}^{\infty}= \lbrace \langle x_{n}+ \delta_{n}/3, x_{n}+\delta_{n} \rangle \rbrace _{n=0}^{\infty}$ s vlastnosťami: $A_{0}\supset A_{1} \supset  A_{2} \supset A_{3} \supset ..., B_{0}\supset B_{1} \supset B_{2} \supset...,$  $\delta_{n+1}< \delta_{n}/3 $; $\varepsilon_{n+1}< \varepsilon_{n}/2$,  $x_{n+1}\in A_{n}- A_{n+1} $, $f(x_{n+1})\in B_{n} $, $f(A_{n})\subset B_{n} $. Označme $a,b$ tie prvky, pre ktoré platí $\lbrace a\rbrace = \bigcap\limits_{n=0}^\infty A_{n} $, $\lbrace b\rbrace = \bigcap\limits_{n=0}^\infty B_{n} $; potom platí  $ a= \lim_{n \to \infty}x_n , \,  b= \lim_{n \to \infty}f(x_n)$, preto $\lim_{n \to \infty}f(x=b)$. Súčasne ale (pretože $a \in A_{n} $ pre každé $n=0,1,...) $ je $f(a) \in B_{n} (n=0,1,...)$, preto $f(a)=b $, čo je spor.
$\boxed{268.}$ $\boldsymbol{1.}$ každá z množín $ A_{n}:=\lbrace x \in \mathbb{R}; \, f(x)>\frac{1}{n}\rbrace $,  $ B_{n}:=\lbrace x \in \mathbb{R}; \, f(x)<\frac{1}{n}\rbrace $ ($ n \in \mathbb{N}$ je konečná a $A = \bigcup\limits_{n=1}^\infty (A_{n} \cup  B_{n} $;
$\boldsymbol{2.}$ množina  $ M=\lbrace x \in \mathbb{R};f(x)=0 \rbrace  $ je konečná, funkcia  $g(x) = \frac{1}{f(x)}, x \in \mathbb{R}-M $, vhodne dodefinovaná v bodoch množiny $M$, vyhovuje predpokladom pr. 268.1;
$\boxed{271.}$ existujú postupnosti $\lbrace a_{n}\rbrace_{n=1}^{\infty} $, $\lbrace b_{n}\rbrace_{n=1}^{\infty}$ tak, že $  \lim_{n \to \infty}a_n =\lim_{n \to \infty}b_n = +\infty$, $(a:=\lim_{n \to \infty}f(a_n)<\lim_{n \to \infty}f(b_n) = +\infty (=:b$, $a,b \in \mathbb{R}$ (pozri pr. 199); potom stačí položiť napr. $A=\frac{a+b}{2} $;
$\boxed{273.}$ sporom; ak $g(x):=f(x+ \frac{T}{2})-f(x) \ne 0$ pre každé $x\in \mathbb{R}$, tak $g$ nemení znamienko na $ \mathbb{R}$, potom ale (ak napr. $g(x) >0 $ na  $\mathbb{R} $ je  $f(x+T) - f(x)= (f(x+T)-f(x+T/2))+(f(x+T/2)-f(x)) >0$;
$\boxed{274.}$ $\boldsymbol{2.}$ napr. $f(x)=0$ pre  $x\in \mathbb{R}-(\mathbb{N} \cup \lbrace \frac{1}{n}; n\in \mathbb{N} \rbrace$), $f(n)=\frac{1}{n} \in \mathbb{R}$;
$\boldsymbol{3.}$ napr. $f(x)=x$ pre  $x\in \langle 0,1\rangle $, $f(x)=x-1$ pre  $x\in (2,3\rangle $;
$\boxed{276.}$ $\boldsymbol{1.}$ pri dôkaze rovnomernej spojitosti funkcie $f.g$ využite, že $f$ aj $g$ sú na $(a,b)$ ohraničené (pozri pr. 256.1);
$\boxed{277.}$ tvrdenie je nepravdivé;
$\boxed{278.}$ $\boldsymbol{1,2.}$ každá funkcia rovnomerne spojitá na $\mathbb{R}$;
$\boldsymbol{3.}$ každá konštantná funkcia;


\section{Diferenciálny počet funkcií jednej premennej}
$\boxed{279.}$
$\boldsymbol{1.}\ 0,2;\boldsymbol{2.}\ 0;\boldsymbol{3.}\ 2;\boldsymbol{4.}\ 1;\boldsymbol{5.}\ 2;$
$\boxed{280.}$
$\boldsymbol{1.}\ f'(1)=+\infty;\boldsymbol{2.}\ \lim_{x\rightarrow 0+}\frac{\sgn x}{x}=\lim_{x\rightarrow 0+}\frac{1}{x}=\infty=\lim_{x\rightarrow 0-}-\frac{1}{x}=\lim_{x\rightarrow 0-}\frac{\sgn x}{x};$
$\boxed{281.}$
$\boldsymbol{1.}\ f'(x)=3x^2+2,D(f')=\mathbb{R};\boldsymbol{2.}\ f'(x)=\frac{4}{3}\sqrt[3]{x},D(f')=\mathbb{R};\boldsymbol{3.}\ f'(x)=-\frac{2x}{(1+x^2)^2},D(f')=\mathbb{R};\boldsymbol{4.}\ f'(x)=2^{x+1} \ln 2,D(f')=\mathbb{R};\boldsymbol{5.}\ f'(x)=\frac{1}{x},D(f')=(0,\infty);\boldsymbol{6.}\ f'(x)=-\frac{1}{{\sin}^2 n},D(f')=\mathbb{R}\setminus \{k\pi; k\in\mathbb{Z}\};\boldsymbol{7.}\ f'(x)=\frac{2}{3(x+1)^{\frac{1}{3}}},D(f')=\mathbb{R}\setminus {\ -1}\,$ pre $x=-1$ platí $\lim_{h\rightarrow 0+}\frac{f(-1+h)-f(-1)}{h}=\lim_{h\rightarrow 0+}\frac{1}{\sqrt[3]{h}}=\infty \neq -\infty=\lim_{h\rightarrow 0-}\frac{1}{\sqrt[3]{h}}=\lim_{h\rightarrow 0-}\frac{f(-1+h)-f(-1)}{h}$, teda $f'(-1)$ neexistuje; $\boldsymbol{8.}\ f'(x)=1,D(f')=\mathbb{R}\setminus \{0\}$; pre $x=0$ platí $f'(0)=\infty$ (pozri pr. $280.2$), preto $0\notin D(f')$;
$\boxed{282.}$
$\boldsymbol{1.}\ f'_+(\frac{\pi}{2})=1,f'_-(\frac{\pi}{2})=-1;\boldsymbol{2.}\ f'_+(1)=-\pi,f'_-(1)=0;\boldsymbol{3.}\ f'_+(2)=1,f'_-(2)=-1;\boldsymbol{4.}\ f'_+(3)=1,f'_-(3)=-1;$
$\boxed{283.}$
$\boldsymbol{1.}\ f'_+(0)=1=f'_-(0)$, teda $f'(0)=1;\boldsymbol{2.}\ f'_-(-1)=-2,f'_+(-1)=\lim_{x\rightarrow -1_+}\frac{-2x-0}{x+1}$ ($f(-1)$ sa vypočíta dosadením čísla $-1$ do výrazu $x^2-1$)$=\lim_{x\rightarrow -1_+}-\frac{2x}{x+1}=\infty;f'_+(-1)\notin f'_-(-1)$, preto $f'(-1)$ neexistuje; $\boldsymbol{3.}\ f'_+(0)=0=f'_-(0)$, teda $f'(0)=0;\boldsymbol{4.}\ f'_+(\pi)=f'_-(\pi)=0,$ teda $f'(\pi)=0$;
$\boxed{284.}$
$\boldsymbol{1.}\ \lim_{h\rightarrow 0}\frac{f(h)-f(-h)}{2h}=\lim_{h\rightarrow 0}(\frac{1}{2}(\frac{f(h)-f(0)}{h}+\frac{f(0)-f(-h)}{h})),$ pri výpočte limity druhého sčítanca použite substitúciu $t=-h$; $\boldsymbol{2.}\ $ nevyplýva, stačí uvažovať $f(x)=|x|$; 
$\boxed{286.}$
$\boldsymbol{1.}\ y'=1-2x;\boldsymbol{2.}\ y'=x^{\frac{2}{3}}+\sqrt{5}x^{-\sqrt{5}-1}$; $\boldsymbol{3.}\ y'=\frac{2(1-2x)}{(1-x-x^2)^2};\boldsymbol{4.}\ y'=\frac{6-\sqrt[3]{x^2}}{6\sqrt{x}(2+\sqrt[3]{x^2})^2};\boldsymbol{5.}\ y'=30(3x-7)^9;\boldsymbol{6.}\ y'=\frac{1+2x^2}{\sqrt{1+x^2}};\boldsymbol{7.}\ y'=-(p+q+(p-q)x)*\frac{(1-x)^{p-1}}{(1+x)^{q+1}};\boldsymbol{8.}\ y'=-\frac{1}{(1+x^2)^{\frac{3}{2}}};\boldsymbol{9.}\ y'=\frac{7\sqrt[5]{2}}{65\sqrt[5]{x^4}\sqrt[13]{(9+7\sqrt[5]{2x})^12}};\boldsymbol{10.}\ y'=\frac{4\sqrt{x+\sqrt{x}}*\sqrt{x}+2\sqrt{x}+1}{8\sqrt{x}*\sqrt{x+\sqrt{x}}*\sqrt{x+\sqrt{x+\sqrt{x}}}};$
$\boxed{287.}$
$f'(a)=\lim_{x\rightarrow a}\frac{(x-a)\varphi (x)}{x-a}=\varphi (a);$
$\boxed{288.}$
$\boldsymbol{1.}\ f'_+(a)=\varphi (a)\neq -\varphi (a)=f'_-(a);\boldsymbol{2.}\ f'(a)=\lim_{x\rightarrow a}(\varphi (x) \sgn (x-a))*|x-a|^{\varepsilon}=0$ (funkcia $\varphi (x) sgn (x-a)$ je ohranižená, $\lim_{x\rightarrow a}|x-a|^{\varepsilon}=0$);
$\boxed{289.}$
$\boldsymbol{1.}\ y'=\cos x - x\sin x;\boldsymbol{2.}\ y'=\frac{x^2}{(\cos x + x\sin x)^2};\boldsymbol{3.}\ y'=\frac{\sin x \cos x - 2x}{2\sqrt{x}\sin ^2 x},x>0,x\neq \frac{k\pi}{2},k\in \mathbb{N}$ (Doteraz sme pri zápisoch využívali dohodu, podľa ktorej definičný obor funkcie - pokiaľ nie je výslovne určený - je množina všetkých tých $x\in \mathbb{R}$, pre ktorý má daný predpis 'zmysel'. Funkcia $f(x)=\frac{\sqrt{x}}{\tan x}$ je definovaná na $D(f)=\{x\in \mathbb{R};x>0 \wedge x\neq \frac{k\pi}{2},k\in \mathbb{N}\}$ a podľa vety o derivácii podielu má v každom jej bode konečnú deriváciu, preto $D(f')=D(f)$; množina $D(f)$ sa v tomto prípade nezhoduje s množinou $\{x\in \mathbb{R}; x>0 \wedge x\neq k\pi, k\in \mathbb{N} \}$ všetkých tých reálnych čísel, pre ktoré má výraz $f'(x)=\frac{(\sin x \cos x - 2x)}{(2\sqrt{x}\sin ^2 x)}$ 'zmysel');$\boldsymbol{4.}\ y'=n\sin^{n-1}x \cos (n+1)x; \boldsymbol{5.}\ y'=-\sin 2x*\cos (\cos 2x);\boldsymbol{6.}\ y'=2\sin x * \frac{\cos x \sin x^2 - x\sin x \cos x^2}{\sin^2 x^2};\boldsymbol{7.}\ y'=-\cos 2x,x\neq \frac{k\pi}{2}, x\neq -\frac{\pi}{4}+k\pi,k\in \mathbb{Z}$ (pozri poznámku k pr. $289.3$); $\boldsymbol{8.}\ y'=\frac{x^4-1}{x^3\cos^2 (x^2+x^{-2})\sqrt{1+\tan^2(x^2+x^{-2})}}$
$\boxed{290.}$ 
$\boldsymbol{1.}$ $\ y'= \frac{1}{2}(\sqrt{2^{x}}\ln 2-\sqrt{5^{-x}} \ln 5)$;
$\boldsymbol{2.}$ $\ y'=-2x e^{-x^{2}}$;
$\boldsymbol{3.}$ $\ y'= e^{x}x^{2}$;
$\boldsymbol{4.}$ $\ y'= 2^{1+ \sin x^{2}}.x \cos x^{2}\ln 2$; 
$\boldsymbol{5.}$ $\ y'= -e^{\sqrt{\frac{1-x}{1+x}}}(1+x)^{-3/2}(1-x)^{-1/2}$;
$\boldsymbol{6.}$ $\ y'= e^{ax}\sqrt{a^{2}+b^{2}}\sin bx$;
$\boldsymbol{7.}$ $\ y'=(\frac{a}{b})^{x}(\frac{b}{x})^{a}(\frac{x}{a})^{b}(\ln a-\ln b-\frac{a}{x}+\frac{b}{x})$;
$\boldsymbol{8.}$ $\ y'= x^{2}e^{-x}\sin x$;
$\boxed{291.}$ 
$\boldsymbol{1.}$ možno (v opačnom prípade by existovala vlastná alebo nevlastná $\ g'(a) $, pretože  $\ g=(f+g)-f)$;
$\boldsymbol{2.}$ nemožno (ak zvolíme $f(x)=\sqrt[3]{x}$,  $g_{1}(x)=\vert x \vert$, $g_{2}(x)=\sgn x .\sqrt[3]{x}$, tak existuje nevlastná $(f+g_{1})'(0)$, a neexistuje $(f+g_{2})'(0)$; za uvedených predpokladov ale nemôže existovať vlastná $(f+g)'(a)$;
$\boldsymbol{3.}$ nemožno (ak zvolíme $f_{1}(x)=g_{1}(x)=\vert x \vert$, $f_{2}(x)=x+\vert x \vert$, $g_{2}(x)=x-\vert x \vert$, $f_{3}(x)=(1+ \sgn x)\sqrt[3]{x}$, $g_{3}(x)=(1- \sgn x)\sqrt[3]{x}$, tak $(f_{1}+g_{1})'(0)$ neexistuje, $(f_{2}+g_{2})'(0)$ je vlastná a $(f_{3}+g_{3})'(0)$ je nevlastná);
$\boxed{292.}$ dôkaz prebieha analogicky ako dôkaz vety o derivácii zloženej funkcie pre prípad vlastvých derivácií;
$\boxed{293.}$ 
$\boldsymbol{1.}$ $\ y'= \frac{1}{\ln 3 . \sin x}, \, x \in \bigcup\limits_{k \in \mathbb{Z}}(2k\pi, (2k+1)\pi)$ (pozri poznámku k pr. 289.3);
$\boldsymbol{2.}$ $\ y'= \frac{x}{x^{4}-1}$, $\vert x\vert >1 $;
$\boldsymbol{3.}$ $\ y'= \frac{1}{\sqrt{x^{2}+1}}$;
$\boldsymbol{4.}$ $\ y'= \frac{1}{(1+x^{2})(1+x)^{2}}$, $ x >-1 $;
$\boldsymbol{5.}$ $\ y'= \ln (x+\sqrt{1+x^{2}})$;
$\boldsymbol{6.}$ $\ y'= \frac{(2x+1) e^{\sqrt{\log _{2}(x^{2}+x+1)}}}{\ln 4(x^{2}+x+1)\sqrt{\log_{2}(x^{2}+x+1)}}$;
$\boldsymbol{7.}$ $\ y'= \frac{2\sin a}{(1-x^{2})(1-x^{2}\cos ^{2}a)}$,$\vert x\vert <1 $ ;
$\boldsymbol{8.}$ $\ y'= 2\sin \ln x $;
$\boldsymbol{9.}$ $\ y'= \frac{x^{2}+\sgn x}{x\sqrt{x^{2}+1}}$;
$\boldsymbol{10.}$ $\ y'= \frac{1}{x}(\frac{\ln x. \log _{3}x}{\ln 2}+\log _{2}x.\log _{3}x+\frac{\log _{2}x.\ln x}{\ln 3})$;
$\boxed{294.}$ 
$\boldsymbol{1.}$ $\ y'=\frac{x+2}{2\sqrt{x}(x+1)} $;
$\boldsymbol{2.}$ $\ y'=\arcsin x \frac{x}{\sqrt{(1-x^{2})}} $;
$\boldsymbol{3.}$ $\ y'=\frac{\pi}{2\sqrt{(1-x^{2})}\arcsin ^{2}x} $; využite pr. 62.1);
$\boldsymbol{4.}$ $\ y'=\frac{1}{2\sqrt{x}(1-x^{2})} $;
$\boldsymbol{5.}$ $\ y'=\arcsin \sqrt {\frac{x}{1+x}}, \, x\geq 0 $ (tabuľkovým derivovaním možno nájsť hodnoty $\ y'$ len pre $x>0$ pretože v bode $0$ má funkcia $\sqrt{x}$ nevlastnú deriváciu; derivácia funkcie $x \arcsin \sqrt{\frac{x}{1+x}} + \arctan \sqrt{x}-\sqrt{x}$ v bode $0$ sa vypočíta na základe definície derivácie;
$\boldsymbol{6.}$ $\ y'=\frac{1}{2x\sqrt{x-1}\arccos \frac{1}{\sqrt{x}}} $;
$\boldsymbol{7.}$ $\ y'=\frac{3^{\arctan (2x+\pi)}.\ln 9}{1+(2x+\pi)^{2}} $;
$\boldsymbol{8.}$ $\ y'=\frac{\sin x. \sgn (\cos x-\cos a)}{1-\cos a. \cos x} $, $\cos x\ne \cos a $ (využite rovnosť $(\cos x- \cos a )^{2}=(1-\cos x. \cos a )^{2}-(\sin x. \sin a )^{2}$ ;
$\boldsymbol{9.}$ $\ y'=(\arcsin x)^{2} $; (tabuľkovým derivovaním možno $\ y'$ nájsť len pre $ \vert x\vert <1$, pretože $arcsin  $ má v bodoch $1,-1 $ nevlastnú deriváciu; platnosť uvedeného vzťahu pre body $1,-1 $ by sa nemusela overovať z definície; neskôr - pozri pr. 384 - uvedieme efektívnejší spôsob nájdenia derivácie v bodoch $1,-1 $);
$\boldsymbol{10.}$ $\ y'=\frac{\cos x}{\sqrt{3-2\sin ^{2}x}} $;
$\boxed{295.}$ 
$\boldsymbol{2.}$ $\ y'=\frac{24-x-5x^{2}}{3\sqrt{x-1}\sqrt[3]{(x+2)^{5}}\sqrt{(x+3)^{5}}} $; (uvedený postup možno použiť pre $x>1$; derivácia v bode $1$  je nevlastná;
$\boldsymbol{3.}$ $\ y'=\frac{ctg 3x}{1-\sin 3x}.\sqrt[3]{\frac{\sin 3x}{1-\sin 3x}} $, $x\in (0,\frac{\pi}{6} )$;
$\boldsymbol{5.}$ $\ y'=\frac{1}{2}.(1-\ln x) \sqrt[x]{x} $;
$\boldsymbol{6.}$ $\ y'=x(2\ln x+1)x^{x^{2}} $;
$\boldsymbol{7.}$ $\ y'=(\cos x)^{\sin x}(\cos x.\ln\cos x-\tan x.\sin x)^{\cos x}(\sin x.\ln \sin x-\cos x. ctg x)$, $x\in (0,\frac{\pi}{2} ) $ ;
$\boxed{296.}$ 
$\boldsymbol{1.}$ napr. $f(x)= g(x)=\vert x-a\vert $;
$\boldsymbol{2.}$ napr. $f(x)=x-a, \, g(x)=\vert x-a\vert $;
$\boxed{297.}$ dôkaz sa robí analogicky ako v prípade vlastnej $\ f'(a)$;
$\boxed{299.}$ 
$\boldsymbol{1.}$ $\ y'=\frac{\sqrt{b^{2}-a^{2}}}{a+b \cos x},\,x\in (0,\frac{\pi}{2} )  $;
$\boldsymbol{2.}$ $\ y'=-2\cos x.\arctan(\sin x) $;
$\boldsymbol{3.}$ $\ y'=-\frac{\arccos x}{x^{2}},\,0<\vert x\vert<1 )  $;
$\boldsymbol{4.}$ $\ y'=\ln^{2}(x+\sqrt{1+x^{2})}$;
$\boldsymbol{5.}$ $\ y'=\frac{a^{2}+b^{2}}{(x^{2}+b^{2})(x+a)}$, $x > -a$;
$\boldsymbol{6.}$ $\ y'=\frac{1}{x^{4}+1}$, $ \vert x \vert\ne 1$;
$\boldsymbol{7.}$ $\ y'=-e^{-x}arcctg\, e^{x}$;
$\boldsymbol{8.}$ $\ y'=\frac{\cos x.\sin ^{2}x}{\sqrt{\cos ^{2}x-2\sin x}}$, $ \sin x \leq\sqrt{2}- 1$ (bolo možné použiť substitúciu $t=1+\sin x $);
$\boldsymbol{9.}$ $\ y'=\frac{m}{\sqrt{1-x^{2}} }.2\cos (m\arcsin x)e^{m \arcsin x},\,\vert x\vert<1 )$;
$\boldsymbol{10.}$ $\ y'=\frac{x}{x^{2}-2x\cos a+1}$;
$\boldsymbol{11.}$ $\ y'=(5x-1)(x-1)(x+1)^{2}$ pre $x > -1$, $\ y'=-(5x-1)(x-1)(x+1)^{2}$ pre $x < -1$, $y_{+}'(-1)=y_{-}'(-1)=0$, teda $y'(-1)=0$; to možno naraz zapísať v tvare $\ y'=(5x-1)(x-1)(x+1)^{2}\sgn (x+1)$;
$\boldsymbol{12.}$ $\ y'=\pi n \sin 2\pi x$ pre ; $x \in (n,n+1),\, y_{+}'(n)=y_{-}'(n)=0$, teda $y'(n)=0$ ($n\in \mathbb{Z}$); to možno naraz zapísať  $y'=\pi[x] \sin 2\pi x$;
$\boldsymbol{13.}$ $\ y'= \left\{\begin{matrix} -1, & \mbox{ak }x <1 \\ 2x-3, & \mbox{ak }x \in \langle 1,2\rangle \\ 1, & \mbox {ak  }x >2\end{matrix}\right. $;
$\boldsymbol{14.}$ $\ y'= \left\{\begin{matrix} \frac{1}{1+x^{2}}, & \mbox{ak } -1<x\leq 1 \\ \frac{1}{2}, & 
         \mbox{ak   } \vert x \vert > 1\end{matrix}\right.$; $y_{+}'(-1)=\frac{1}{2}$, $y_{-}'(-1)=+\infty$, teda $y'(-1)$ neexistuje;
$\boldsymbol{15.}$ $\ y'=th^{3}x$, pritom sa použije vzťah $\frac{1}{ch^{2}x}=1-th^{2}x$;  
$\boldsymbol{16.}$ $\ y'=(\ln x)^{x}.(\ln(\ln x)+\frac{1}{\ln x}-\frac{2\ln x}{x}/(x^{\ln x})$, $ x >1$;
$\boldsymbol{17.}$ $\ y'=(\frac{\sin 2x . \arctan x}{\sqrt{1-\sin ^{4}x}\arcsin (\sin ^{2}x}+\frac{\ln \arcsin (\sin ^{2}x)}{1+x^{2}}).(\arcsin (\sin ^{2}x))^{\arctan x}$;
$\boldsymbol{18.}$ $\ y'=x^{x^{x}}.x^{x}(\ln ^{2}x+\ln x+\frac{1}{x})$;
$\boldsymbol{19.}$  použitím vzťahu $\log _{x }e=\frac{1}{\ln x}$   možno daný výraz upraviť na $y=e^{2}-2xe^{2}+ex^{2}$  $x>0, \, x\ne 1)$; teda $\ y'=ex(x-e)$,  $x>0, \, x\ne 1)$;
$\boxed{300.}$ $a=2x_{0}$, $b=-x_{0}^{2}$;
$\boxed{301.}$ dôkaz sa robí matematickou indukciou, pričom sa použije veta o derivácii súčinu; $f'(0)=10001$;
$\boxed{302.}$ $\boldsymbol{1.}$ $\ y'=\frac{ff'+gg'}{\sqrt{f^{2}+g^{2}}}$;
$\boldsymbol{2.}$ $\ y'=\frac{f'g+fg'}{f^{2}+g^{2}}$;
$\boldsymbol{3.}$ $\ y'=\frac{fg'.\ln f+gf'.\ln g}{f.g.\ln^{2}f}$ (treba použiť vzťah  $\log _{f}g=\frac{\ln g}{\ln f}$);
$\boldsymbol{4.}$ $\ y'=2x f'(x^{2})$;
$\boldsymbol{5.}$ $\ y'=\sin 2x f'(\sin ^{2}x)-g'(\cos ^{2}x))$;
$\boldsymbol{6.}$ $\ y'=e^{f(x)}(e^{x}.f'(e^{x})+f(e^{x}).f'(x))$;
$\boxed{303.}$ $\boldsymbol{1.}$ $P_{n}(x)=(x+x^{2}+...+x^{n})'=(\frac{x^{n+1}-x}{x-1})'=\frac{nx^{n+1}-(n+1)x^{n}+1}{(x-1)^{2}} $ pre $x \ne 1$,  $P_{n}(1)=\frac{n(n+1)}{2} $;
$\boldsymbol{2.}$ $Q_{n}(x)=(\frac{x^{2n+1}-x}{x^{2}-1})'=\frac{(2n-1)x^{2n+2}-(2n+1)x^{2n}+x^{2}+1}{(x^{2}-1)^{2}}$; 
$\boldsymbol{3.}$ $R_{n}(x)=(xP_{n}(x))'=\frac{n^{2}x^{n+2}-(2n^{2}+2n-1)x^{n+1}+(n+1)^{2}x^{n}-x-1}{(x-1)^{3}}$; 
$\boldsymbol{4.}$ $T_{n}(x)=(\sum_{k=1}^n \sin kx)'=(\frac{\cos \frac{x}{2}- \cos \frac{(2n+1)\pi}{2}}{2\sin \frac{x}{2} })'=\frac{-1+(n+1)\cos nx -n \cos (n+1)x}{4 \sin ^{2}\frac{x}{2}}$;
$\boxed{304.}$ $\boldsymbol{1.}$ $\alpha >0 $;
$\boldsymbol{2.}$ $\alpha >1 $;
$\boldsymbol{3.}$ $\alpha >2 $ (pre $x \ne 0$ možno  $f'(x)$ nájsť tabuľkovým derivovaním, $f'(0)$ sa musí hľadať priamo z definície; treba si uvedomiť, že $\lim_{x \to 0}\cos \frac{1}{x} $ neexistuje, $\lim_{x \to 0}x^{\beta} \sin \frac{1}{x}=0 $ pre $\beta >0 $; pre $\beta \leq 0 $ táto limita neexistuje);
$\boxed{305.}$ nie (funkcia $\sqrt[3]{x^{2}}$ nemá vlastnú ani nevlastnú deriváciu v bode $0$; jednostranné derivácie sú nevlastné a opačných znamienok); $y'(0)=0$ (vypočíta sa z definície);
$\boxed{306.}$ $\boldsymbol{1.}$ napr. $f(x)=(x-1)^{2}(x-2)^{2} $ pre $x \in \mathbb{Q}, \, f(x)=0 $ pre $x \in \mathbb{R} -\mathbb{Q}$;
$\boldsymbol{2.}$ napr. funkcia určená podmienkami $f(x)=\sqrt[3]{\frac{x}{4}}$ pre $x \in \langle -\frac{1}{2},\frac{1}{2}\rangle$, $f(x+1)=f(x)+1 $ ($x \in \mathbb{R}$);
$\boldsymbol{3.}$ napr. $f(x)=ax+g(x)$, kde $g$ je určená podmienkami $g(x)=0$ pre $x \leq 0, \, g(x)=x^{2}(x-1)^{2}\sin \frac{1}{x(x-1)} $  pre $x \in (0,1), \,g(x+1)=g(x) $ pre $x \geq 0 $;
$\boxed{308.}$ napr. $f(x)=(1+\sgn (x-1)). \frac{\arccos(\cos \pi(x-1)}{\pi}$, tj. $f(x)=0 $ pre $x<1 $, $f(x)=x-k $ pre $x\in \langle k, k+1), \, k \in \mathbb{N} $, $k$ nepárne, $f(x)=k+1-x$ pre $x\in \langle k, k+1), \, k \in \mathbb{N} $, $k$ párne;
$\boxed{309.}$ $\boldsymbol{1.}$ napr. $f(x)=-\sqrt[3]{x-a}$;
$\boldsymbol{2.}$ napr. $f(x)=\sgn (x-a)$;
$\boldsymbol{3.}$ napr. $f(x)=\sgn (x-a).(\sin \frac{1}{x-a}-2)$ pre  $x \ne a $, $f(a)=0$;
$\boxed{310.}$ $\boldsymbol{1.}$ $\beta /b=f(\sqrt{2}), \, (f^{-1})'(\frac{6}{5})=\frac{1}{2}$;
$\boldsymbol{2.}$ $b=f(0), \, (f^{-1})'(-\frac{1}{2})=\frac{1}{2}$;
$\boldsymbol{3.}$ $b=f(0), \, (f^{-1})'(1)=5$; 
$\boldsymbol{4.}$ $b=f(\sqrt{2}), \, (f^{-1})'(0)= - \frac{1}{4\sqrt{2}}$;
$\boldsymbol{5.}$ $b=f(\frac{1}{\sqrt{2}}), \, (f^{-1})'(\frac{3}{4})= \frac{1}{\sqrt{2}}$;
$\boxed{311.}$ pretože $f^{-1}(f(x))=x $ v niektorom okolí bodu $a$, je $f^{-1}(f(a)).f'(a)=1 $; treba predpokladať existenciu vlasnej  $(f^{-1})'(f(a)) $ a vlasnej  $f'(a)\ne 0 $;
$\boxed{312.}$ $\boldsymbol{1.}$ $f'(x)= (1-x^{2})^{-1/2}$;
$\boldsymbol{2.}$ $f'(x)= -(1-x^{2})^{-1/2}$;
$\boldsymbol{3.}$ $f'(x)= \frac{1}{1+x^{2}}$;
$\boldsymbol{4.}$ $f'(x)= -\frac{1}{1+x^{2}}$;
$\boldsymbol{5.}$ $f'(x)= \frac{1}{x}, \, x>0$;
$\boxed{313.}$ dôkaz je analogický dôkazu vety o derivácii inverznej funkcie;
$\boxed{314.}$ $\boldsymbol{1\alpha/}$ $y= \sqrt[3]{4}(x+1)$; $\boldsymbol{1\beta/}$ $y= 3$;
$\boldsymbol{2.}$ $y= x+1$;
$\boldsymbol{3.}$ $x=1$;
$\boldsymbol{4.}$ $A=(0,1), y=1$;
$\boldsymbol{5.}$ $A=(1,\frac{1}{2}), x+2y-2=0$;
$\boxed{315.}$ $\boldsymbol{1.}$ $12x-4y-13=0$;
$\boldsymbol{2.}$ $4x-4y+3=0$;
$\boxed{316.}$ $\boldsymbol{1.}$ $b^{2}-4ac=0$;
$\boldsymbol{2.}$ $a=\frac{1}{2e}$;
$\boxed{317.}$ $\boldsymbol{1.}$ $d(\frac{\ln x}{\sqrt{x}})(a)=\frac{2- \ln a}{2a\sqrt{a}}\mathrm{d}x (a)$; (teda  $d(\frac{\ln x}{\sqrt{x}})(a)$ je funkcia daná predpisom $\frac{2- \ln a}{2a\sqrt{a}}(x-a)$);,
$\boldsymbol{2.}$ $\frac{a}{\sqrt{A^{2}+a^{2}}}\mathrm{d}x (a)$;
$\boldsymbol{3.}$ $\frac{2a}{a^{2}-1}\mathrm{d}x (a)$ ($\vert a \vert <1) $;
$\boldsymbol{4.}$ $\frac{1}{\cos ^{3}a}\mathrm{d}x (a)$;
$\boxed{318.}$ $\boldsymbol{1.}$ $dy(a)=(u(a).v(a).w'(a)+u(a).v'(a).w(a)+u'(a).v(a).w(a) )\mathrm{d}x (a)$, čo sa stručne zapisuje v tvare $dy=u.v.dw+u.dv.w+du.v.w$;
$\boldsymbol{2.}$ $dy=\frac{v^{2}.du-2u.v.dv}{v^{4}}$;
$\boldsymbol{3.}$ $dy=\frac{w.du-u.dw}{u^{2}+w^{2}}$;
$\boldsymbol{4.}$ $dy=\frac{u.du-v.dv}{u^{2}+v^{2}}$;
$\boxed{319.}$ $\boldsymbol{1.}$ $\sqrt{a^{2}+x^{2}}\approx a+\frac{x}{2a}$ pre $x$ blízke $0$;
$\boldsymbol{2.}$ $\ln{x}+\sqrt{1+x^{2}} \approx x$ pre $x$ blízke $0$;
$\boldsymbol{3.}$ $\arctan (1+x^{2}) \approx \frac{\pi}{4}+\frac{1}{2} x$ pre $x$ blízke $0$;
$\boldsymbol{4.}$ $\sqrt[n]{a^{n}+x}\approx a+\frac{x}{na^{n-1}}$ pre $x$ blízke $0$;
$\boxed{320.}$ "$\Rightarrow $": stačí položiť $\varphi (x)=A+\frac{\omega(x)}{x-a}$ pre $x \ne a$, $\varphi(a)=0$; "$\Leftarrow $": $\varphi(x)=0$ možno písať v tvare $A+(\varphi(x)-A)$, kde $A=\lim_{x \to a}\varphi(x)$; potom $\varphi(x)=(\varphi(x)-A)(x-a)$;
$\boxed{321.}$ $\boldsymbol{1.}$  $y''=\frac{2x^{3}+3x}{(1+x^{2})^{3/2}} $;
$\boldsymbol{2.}$  $y''=3x(1-x^{2})^{-5/2} $;   
$\boldsymbol{3.}$  $y''=2(2x^{2}-1)e^{-x^{2}} $;   
$\boldsymbol{4.}$  $y''=2\arctan x+\frac{2x}{1+x^{2}} $;
$\boldsymbol{5.}$  $y''=\frac{1}{x}, \, x>0 $;
c
$\boxed{322.}$ $\boldsymbol{2.}$  $y^{V}=2^{2}.3^{3}.5! $; $y^{VI}=0$;
$\boldsymbol{3.}$  $y^{(10)}=-(3.5.7.....17)2^{-10}.x^{-19/2}$;
$\boldsymbol{4.}$  $y=-1+\frac{2}{1-x}, \, y^{(22)}=2.221.(1-x)^{-23} $;
$\boldsymbol{6.}$  $y=\frac{1}{2}(\sin 6x-\sin 2x), \, y^{(15)}=\frac{1}{2}(2^{15}\cos 2x -6^{15}\cos 6x) $;
$\boldsymbol{7.}$  $y=\frac{1}{2}\sin 2x(\cos 2x-\cos 4x)= \frac{1}{4}\sin 4x -\frac{1}{4}(\sin 6x-\sin 2x) , \, y^{(10)}=\frac{1}{4}(-4^{10}\sin 4x +6^{10}\sin 6x-2^{10}\sin 2x) $;
$\boxed{323.}$ $\boldsymbol{1.}$  $((1-x)^{-1/2})^{(100)}=\frac{3.5.....199}{2^{100}}(1-x)^{-201/2} $, $((1-x)^{-1/2})^{(99)}=\frac{3.5.....197}{2^{99}}(1-x)^{-199/2} $, $y^{(100)}=\frac{3.5.....197}{2^{100}}\sqrt{(1-x)^{201}}(399-x) $;
$\boldsymbol{2.}$  $y^{V}=\frac{-6}{x^{4}}, x>0$;
$\boldsymbol{3.}$  $y^{(50)}=2^{49}(1225-2x^{2})\sin 2x + 50.2^{50}x \cos 2x$;
$\boldsymbol{4.}$  $y'''=-\frac{9\sin 3x}{\sqrt[7]{(1-3x)^{3}}}(27x^{2}-18x-1)-\frac{\cos 3x}{\sqrt[3]{(1-3x)^{10}}}(243x^{2}-162x-1)$;
$\boldsymbol{5.}$  $y^{(101)}=3^{99}(10100+18x-9x^{2})\sin 3x + 202.3^{100}(x-1). \cos 3x$;
$\boldsymbol{6.}$ Leibnizovým vzorcom sa vypočíta  $15$. derivácia funkcie; $y'=2x+\sin 2x - 2 \sin x-2x \cos x$,   $y^{(16)}=-2^{15}\cos 2x+32 \cos x - 2x \sin x$; 
$\boldsymbol{7.}$  $y=\frac{1}{2}x(\sin 3x -\sin x)$, $y^{(100)}=\frac{1}{2}(x(3^{100}\sin 3x \sin x)+100\cos x - 100.3^{99}\cos 3x)$;
$\boxed{325.}$ $\boldsymbol{1.}$ $y^{n}=\frac{(-1)^{n-1}(n-1)!a^{n}}{(ax+b)^{n}}, ax+b>0$;
$\boldsymbol{2.}$ $y^{(n)}=n![\frac{(-1)^{n}}{x^{n+1}}+\frac{1}{(1-x)^{n+1}}](y=\frac{1}{x}+\frac{1}{1-x}) $;
$\boldsymbol{3.}$ $y^{(n)}=\frac{(-1)^{n}n!}{(x-2)^{n+1}}+\frac{n!}{(1-x)^{n+1}}(y=\frac{1}{x-2}+\frac{1}{1-x}) $;
$\boldsymbol{4.}$ $y^{(n)}=\frac{1.3.....(2n-1)}{(1-2x)^{2n+1}}$;
$\boldsymbol{5.}$  $\ y'=\sin 2x$;  $y^{(n)}=(\sin 2x)^{(n-1)}=2^{n-1}\sin (2x+\frac{(n-1)\pi}{2})$;
$\boldsymbol{6.}$  $\ y^{(n)}=\frac{1}{2}(8^{n}\sin (8x+\frac{n\pi}{2})-2^{n}\sin (2x+\frac{n\pi}{2}))$  $y=\frac{1}{2}(\sin 8x- \sin 2x))$;
$\boldsymbol{7.}$  $\ y^{(n)}=4^{n-1}\cos (4x+\frac{n\pi}{2})$ (pozri návod k pr. 40.2);
$\boldsymbol{8.}$  $\ y^{(n)}=\frac{1}{2} \sin ax.(\sin (a+b)x +\sin (a-b)x)=\frac{1}{4}(2\cos bx-\cos(2a+b)x -\cos (2a-b)x)$, $\ y^{(n)}=\frac{1}{4} (2b^{n} \cos(bx+\frac{n\pi}{2})-(2a+b)^{n}\cos ((2a+b)x+\frac {n\pi}{2})-(2a-b)^{n}\cos ((2a-b)x+\frac{n\pi}{2}))$;
$\boxed{326.}$ $\boldsymbol{1.}$ $y^{(n)}=2^{x-1}.\ln^{n}2 (x-1+\frac{n}{\ln 2}) $ (použil sa Leibnizov vzorec);
$\boldsymbol{2.}$ $((1+x)^{-1/3})^{(n)} =(-1)^{n}\frac {1.4.7.....(3n-2)}{3^{n}}(1+x)^{-(3n+1)/3}$,  $y^{(n)}=(-1)^{n+1}\frac{1.4.7.....(3n-5)}{3^{n}(1+x)^{n+1/3}}.(2x+3n)$ pre $n\geq 2$, $\ y'=\frac{3+2x}{\sqrt[3]{(1+x)^{4}}}$ (výsledok pre $\ y'$ sa musí uviesť samostatne, pretože vzťah pre $((1+x)^{-1/3})^{(n-1)}$, ktorý dosádzame do Leibnizovho vzorca, neplatí pre $n=1$);
$\boldsymbol{3.}$ $y^{n}=b^{n-2}.\sin (bx+\frac{n\pi}{2})(b^{2}x^{2}-n(n-1))+2nb^{n-1} x \sin (bx+\frac{n-1}{2}\pi )$ (pri úprave sa použil vzťah $\sin (A-\pi)=-\sin A$; pre $n\geq 2$ má Leibnizov vzorec tvar  $y^{(n)}=x^{2}(\sin bx)^{(n)}+ \begin{pmatrix}n \\1 \end{pmatrix} 2x.(\sin bx)^{(n-1)}+ \begin{pmatrix}n \\2 \end{pmatrix} 2.(\sin bx)^{(n-2)}$, $\ y'$ treba vypočítať samostatne, zhodou okolností možno výsledky písať v uvedenej spoločnej podobe);
$\boldsymbol{4.}$ $y^{(n)}=(-1)^{n}e^{-x}(x^{2}+(n-1)(n-2))$;
$\boldsymbol{5.}$ $(\log _{2}(1-3x))^{n}=-\frac{(n-1)!3^{n}}{(1-3x)^{n}\ln 2}$,  $y^{(n)}=-\frac{(n-2)!3^{n-1}}{(1-3x)^{n}\ln 2}(3x-n)$ pre $n>1$, $\ y'=\log _{2}(1-3x)-\frac{3x}{(1-3x)\ln 2}$ ($\ y'$ uvádzame samostatne z toho istého dôvodu ako v pr. 326.1); 
$\boldsymbol{6.}$ $y^{(n)}=(-3)^{n-2}e^{2-3x}(36x^{2}-12(9+2n)x+4n^{2}+32n+81)$ pre $n\geq 2$, $\ y'=e^{2-3x}(3-2x)(6x-7)$ ($\ y'$ treba vypočítať samostatne, pretože Leibnizov vzorec pre $n\geq 2$ má iný tvar ako pre $n=1$);
$\boxed{327.}$ $\boldsymbol{1.}$ $f(x)=\frac{2}{3}.\frac{1}{1-2x}+\frac{1}{3}.\frac{1}{1+x}$, $f^{(n)}(0)=\frac{n!}{3}(2^{n+1}+(-1)^{n})$;   
$\boldsymbol{2.}$ $f^{(n)}(x)=2^{n}x^{2}e^{2x}+2^{n}nxe^{2x}+n(n-1)2^{n-2}e^{2x}$ (použili sme Leibnizov vzorec pre $n\geq 2$, samostatne vypočítali $\ y'$ a presvedčili sa, že sa dá vyjadriť tým istým vzťahom), $f^{(n)}(0)=n(n-1)2^{n-2}$; 
$\boldsymbol{3.}$ $f^{(n)}(x)=\frac {(2n-x)}{2^{n}(1-x)^{n+1/2}}\prod_{k=1}^{n-1} (2k-1)$ pre $n\geq 2$, $ f'(x)=\frac{2-x}{2(1-x)^{3/2}}$; $f^{(n)}(0)=\frac{n}{2^{n-1}}\prod_{k=1}^{n-1} (2k-1)$ pre $n\geq 2$, $f^{(1)}(0)=1$;
$\boldsymbol{4.}$ $f(x)=\frac{1}{2}(\frac{1}{1-x)}-\frac{1}{1+x})$, $f^{(n)}(0)=\frac{1}{2}n!(1-(-1)^{n})$;$\boldsymbol{5.}$ platí $(1-x^{2}).f''(x)-x.f'(x)=0$, odtiaľ $f^{(n)}(0)=(n-2)^{2}(0)$ (odvodí sa Leibnizovým vzorcom za predpokladu $n\geq 4$, ale zhodou okolností platí tento rekurentný vzťah aj pre $n=3$; vzťah pre  $n=3$ dostaneme, ak prvú uvedenú identitu raz zderivujeme);  $f'(0)=1$, $f''(0)=0$, $f^{(2k+2)}(0)=0$, $f^{(2k+1)}(0)=1^{2}.3^{2}.....(2k-1)^{2}$  ($k=0,1,...$);
$\boldsymbol{7.}$ $f(x)=g(x).g(x)$, kde $g(x)=\frac{1}{2}(\frac{1}{1-x}+\frac{1}{1+x})$; $g^{(k)}(0)=0$    pre $k$ nepárne, $g^{(k)}(0)=k!$ pre $k$ párne; na súčin  $g.g$  použijeme Leibnizov vzorec;  $f^{(n)}(0)=0$   pre $n$ nepárne,  $f^{(2n)}(0)=\sum_{k=0}^m(2 k)!$;
$\boxed{329.}$  nepriamo, ak nie je prostá, tak vyhovuje predpokladom Rolleho vety na niektorom intervale;
$\boxed{330.}$ $n+1$ nulových bodov vytvorí $n$ intervalov, na každom z nich sú splnené predpoklady Rolleho vety; tak dostaneme $n$ nulových bodov funkcie $f'$, tie vytvoria $n-1$ intervalov, naktorých $f'$ vyhovuje predpokladom Rolleho vety atď.;
$\boxed{331.}$ Ak $x_{1}<x_{2}<...<x_{k}$ sú korene polynómu $P_{n}$, tak podľa návodu k pr. 330 existuje $k-1$ koreňov polynómu $P'_{n}$, tak, že $x_{1}<c_{1}<x_{2}<c_{2}...<c_{k-1}<x_{k}$. Ďalej: ak $x_{i}$ je m-násobný koreň $P_{n}$ ($m\geq 2$), tak $x_{i}$ je $m+1$-násobný koreň $P'_{n}$. Súčet násobností takto získaných koreňov je $n-1$ (podľa základnej vety algebry nemôže byť väčší než $n-1$; ak násobnosť koreňa $c_{i}$ ($i=1,...,k-1$) zdola odhadneme číslom $1$, je číslo $n-1$ aj dolným odhadom súčtu všetkých násobností.
$\boxed{332.}$ z návodu k pr.331 vyplýva: ak súčet násobností reálnych koreňov polynómu $P_{n}$ je  $k$ tak súčet násobností reálnych koreňov $P'_{n}$ je aspoň $k-1$ ($n\geq 2, k\geq 2 $); derivácia polynómu $x^{3}-3x^{2}+6x-1$ nemá reálne korene;
$\boxed{333.}$ treba dokázať, že existujú body $c,d \in \mathbb{R}$, $a<c<d<b$, v ktorých $f(c)=f(d)$; na intervale $\langle c,d\rangle$  potom $f$ vyhovuje predpokladom  Rolleho vety; 
$\boxed{334.}$ $f$ musí byť spojitá na $(a,b)$; keby bola naviac spojitá aj v bodoch $a,b$, vyhovovala by predpokladom  Rolleho vety, čo by bolo v spore z podmienkou $f'(x)\ne 0$ pre všetky  $x \in(a,b)$; zadaniu vyhovuje napr. funkcia  $f$ daná predpisom  $f(x)=x$ pre  $x \in(a,b\rangle $,  $f(x)=b$;
$\boxed{335.}$  nie je;  $y$ totiž nemá vlastnú ani nevlastnú deriváciu v bode $0$;
$\boxed{336.}$ $\boldsymbol{1.}$ zvoľme pevne $a\in I$; ak $x\in I$, tak podľa Lagrangeovej vety  $f(x)-f(a)=0.(x-a)$, tj. $f(x)=f(a)$;
$\boldsymbol{2.}$ z pr. 336.1 vyplýva, že $M$ nemôže byť otvorený interval; zadaniu vyhovuje napr. $M=\mathbb{R}-\lbrace 0\rbrace$, $f(x)=\sgn x/M$;
$\boxed{338.}$ postup je rovnaký ako v pr. 337, v pr. 338.3 treba robiť úvahy zvlášť pre $(-\infty, -1)$ a zvlášť pre $(1,\infty)$ (množina $(-\infty, -1)\cup (1,\infty)$ nie je interval);
$\boxed{339.}$ uvedieme dva návody: 1. zvoľme pevne $a \in I $; ak $x \in I $, tak podľa Lagrangeovej vety  $f(x)-f(a)=k.(x-a)$; 2. funkcia $g(x)=kx+b $ má požadovanú vlastnosť ($g'(x) * k$ na $I$), ak ju má aj funkcia $f$, je $(g-f)'=0$ na $I$, podľa pr. 336 potom $g-f$ $* const$ na $I$;
$\boxed{340.}$ sporom; ak $f'(x)=\sgn x$, tak podľa pr. 339 $f(x)=\vert x\vert +c$, ale  $\vert x\vert +c$ nemá deriváciu v bode $0$;
$\boxed{341.}$ $\boldsymbol{2.}$ treba využiť, že z nerovnosti $0<y<c<x$ vyplýva pre $p>1$ nerovnosť $y^{p-1}<c^{p-1}<x^{p-1}$;
$\boldsymbol{4.}$  $\ln \frac{a}{b}= \ln a-\ln b$; ak $0<b<c<a$, tak $\frac{1}{a}<\frac{1}{c}<\frac{1}{b}$;
$\boxed{342.}$ $\boldsymbol{1.}$ $f$ je na $(a,b)$ spojitá (má totiž v každom bode deriváciu), na každom intervale $\langle x,y \rangle \subset (a,b)$ vyhovuje predpokladom Lagrangeovej vety; ak $\vert f'(x)\vert \leq K$ pre všetky $x \in (a,b)$, tak $\vert f(x)-f(y)\vert \leq K \vert x-y\vert\quad (x,y \in (a,b)$), odtiaľ vyplýva naše tvrdenie na základe definície rovnomernej spojitosti;
$\boxed{343.}$ nepriamo;  ak $\vert f'(x)\vert \leq K$   $x \in (a,b)$, tak použitím Lagrangeovej vety $\vert f(x)-f(A)\vert \leq K \vert x-A\vert \leq K(b-a)$  ($A \in (a,b)$ je pevne zvolené), teda $ \vert f(x) \vert \leq \vert f(A) \vert + \vert f(x)-f(A)  \vert \leq  \vert f(A) \vert+K(b-a)$, tj. $f$ je uhraničená na $(a,b)$;
$\boxed{344.}$ zvoľme $a \in I$ pevné; pre $a \in I, \, x \ne a$ označme $z(x)$ to číslo $z$ ležiace medzi $x,a $ pre ktoré $F(x) - F(a) = f(z)(x-a)$ (ak je takých čísel viacej, vyberieme jedno z nich); dokážte implikáciu: ak $x\rightarrow a$, tak $z(x) \rightarrow a$; potom $F'(a)=\lim_{x \to a}\frac{F(x) - F(a)}{x-a}=\lim_{x \to a}f(z(x))=\lim_{z \to a}f(z)=f(a)$ ($f$ je na $I$ spojitá);
$\boxed{345.}$ nie; stačí uvážiť $f(x)=x^{3}$, interval $\langle -1,1 \rangle$, $c=0$;
$\boxed{346.}$ číslo $c_{1}$, pre ktoré $f(b)-f(a)=f'(c_{1})(b-a)$, nemusí byť totožné s číslom $c_{2}$, pre ktoré $g(b)-g(a)=g'(c_{2})(b-a)$; dokumentujeme to na príklade $f(x)=x^{3}$, $g(x)=x^{2}$, $a=0$, $b=1$; potom $f(1)-f(0)=f'(\frac{1}{\sqrt{3}})(1-0)$,  $g(1)-g(0)=g'(\frac{1}{2})(1-0)$, ale $\frac{f(1)-f(0)}{g(1)-g(0)}=\frac{f'(\frac{2}{3}}{g'(\frac{2}{3}})$;
$\boxed{347.}$ stačí použiť Cauchyho vetu pre funkcie $f(x)$ a $\frac{1}{x}$ na intervale $\langle 1,2 \rangle$ ( a samozrejme predtým preveriť splnenie predpokladov (i), (ii));
$\boxed{348.}$ ak rozšírime zlomok na ľavej strane rovnosti výrazom $\frac{1}{a.b}$, vidíme, že treba použiť Cauchyho vetu pre funkcie $\frac{f(x)}{x}$  a  $\frac{1}{x}$ na $\langle a,b \rangle$;
$\boxed{349.}$ $\boldsymbol{1.}$ klesajúca na $(-\infty,-1 \rangle$ a na $\langle 1,\infty) $,rastúca na $\langle -1,1\rangle $;
$\boldsymbol{2.}$ klesajúca na $(-\infty,-1 \rangle$ a na $\langle 1,\infty) $,rastúca na $\langle -1,1\rangle $;
$\boldsymbol{3.}$  na každom z intervalov $((2k-1)\pi,(2k+1)\pi)$ je $y'$ kladná, teda  $y$ rastie na každom z intervalov $\langle (2k-1)\pi,(2k+1)\pi\rangle, \, k\in \mathbb{Z}$, odtiaľ vyplýva, že $y$ rastie na $\mathbb{R}= \bigcup_{k\in \mathbb{Z}}\langle (2k-1)\pi,(2k+1)\pi\rangle $;
$\boldsymbol{4.}$ rastie na každom z intervalov $\langle 2k\pi, \frac {\pi}{3}+2k\pi \rangle$, $\langle \frac{2\pi}{3}+2k\pi,(2k+1)\pi \rangle$, $\langle \frac{7\pi}{6}+2k\pi,\frac{11\pi}{6}+2k\pi \rangle$, klesá na každom z intervalov $\langle \frac{\pi}{3}+2k\pi,\frac{2\pi}{3}+2k\pi \rangle$, $\langle (2k+1)\pi,\frac{7\pi}{6}+2k\pi \rangle$, $\langle \frac{11\pi}{6}+2k\pi,(2k+2)\pi \rangle$, $k \in \mathbb{Z}$;
$\boldsymbol{5.}$ rastie na každom z intervalov $\langle \frac{1}{2k+1}, \frac{1}{2k}\rangle$ ($k=\pm 1,\pm 2,\pm 3,...)$ a na $\langle 1,\infty$, klesá na $\langle \frac{1}{2k+2}, \frac{1}{2k+1}\rangle$ ($k=\pm 1,\pm 2,\pm 3,...)$ a na $(-\infty , -1\rangle$;
$\boldsymbol{6.}$ klesá na $(-\infty , 0\rangle$ a na $\langle 2\log _{2}e,\infty)$, rastie na $\langle 0, 2\log _{2}e\rangle$ ($\log _{2}e=\frac{1}{\ln 2} $);
$\boldsymbol{7.}$ klesá na $(-\infty , -1\rangle$ a na $(0,1\rangle$, rastie na $\langle -1, 0)$ a na $\langle 1,\infty)$;
$\boldsymbol{8.}$ rastie na $\langle e^{-7\pi/12+2k\pi},e^{13\pi/12+2k\pi}\rangle$,   klesá na $\langle e^{13\pi/12+2k\pi},e^{17\pi/12+2k\pi}\rangle$, $k\in \mathbb{Z}$ (všimnite si, že číslo $0$ neleží v žiadnom z týchto intervalov);
$\boxed{350.}$ pre dané $a\in \mathbb{R}$ je funkcia $F(x):= \frac{f(x)-f(a)}{x-a}$ kladná na $\mathbb{R}- \lbrace a \rbrace$, teda $\lim_{x \to a}F(x)\geq 0$;
$\boxed{351.}$ $a$ stačí zvoliť tak, aby všetky korene polynómu $P'_{n}$ ležali v $\langle -a,a \rangle$, potom $P'_{n}$ nemení znamienko na $(-\infty, -a)$ a na $(a, \infty)$;
$\boxed{352.}$ $\boldsymbol{1.}$ označme $f(x)=e^{x}$, $g(x)=1+x$; potom $f(0)=g(0)$, $f'(x)>g'(x)$ pre všetky $x>0$, $f'(x)<g'(x)$ pte všetky $x<0$; odtiaľ už vyplýva na základe vety 12 uvedená nerovnosť;
$\boldsymbol{2.}$ nerovnosť $\sin x< x$ iste platí pre $x>1$, petože vtedy $\sin x\leq 1<x$; ak $f(x):=x$, tak $f(0)=g(0)$, $f'(x)>g'(x)$ pre všetky $x \in (0,1 \rangle$, teda $\sin x<x$ platí aj pre $x \in (0,1 \rangle$ (nerovnosť $f'(x)>g'(x)$ neplatí pre všetky $x>0$ - pre  $x>0$ platí len nerovnosť $f'(x)\geq g'(x)$ - preto vetu 12 nemôžeme použiť na dôkaz nerovnosti $\sin x<x$ na celom intervale $(0, \infty)$); ak $f(x):= \sin x$, $g(x):=x-x^{3}/6 $,tak $f(0)=g(0)$, $f'(0)=g'(0)$ a  pre všetky $x>0 $ platí  $f''(x)>g''(x)$  (to vyplýva z predtým dokázanej nerovnosti $\sin x<x$ pre  $x>0 $);
$\boldsymbol{4.}$ ak $f(x):=\tan x$, $g(x):=x+\frac{x^{3}}{3}$, tak $f(0)=g(0)$, $f'(0)=g'(0)$, $f''(x)=2 \cos ^{-3}x.\sin x$, $g''(x)=2x$; pre  $x \in (0,\frac{\pi}{2})$ je $0<\cos x<1$, teda $\cos^{-2}x>1 $, súčasne  $\tan x>x $ (pozri pr. 352.3), teda $\tan x. \cos^{-2}x=\cos^{-3}x.\sin x>x$; odtiaľ vyplýva $f''(x)>g''(x)$ pre  $x \in (0,\frac{\pi}{2})$;
$\boldsymbol{7.}$ ak  $f(x):=1+2\ln x$, $g(x):=x^{2}$, tak $f(1)=g(1)$, $f'(x)<g'(x)$ pre všetky $x>1$, $f'(x)>g'(x)$ pre všetky $x\in(0,1)$;
$\boxed{353.}$ ak do uvedenej nerovnosti dosadíme $x=\frac{1}{n-1},...,x=\frac{1}{2n}$ a získané nerovnosti sčítame, dostaneme odhad $\ln (\frac{2n+1}{n})<\frac{1}{n}+...+\frac{1}{2n}<\ln (\frac{2n}{n-1}) $, $n>1$;
$\boxed{354.}$ nech $M=\lbrace a_{1},...,a_{m} \rbrace$, $a_{1}<a_{2}<...<a_{m}$; potom $f$ rastie na každom z intervalov  $(a,a_{1} \rangle$, $\langle a_{1},a_{2} \rangle$,..., $\langle a_{m},b)$; odtiaľ vyplýva, že $f$ rastie na $(a,b)$;
$\boxed{355.}$ funkcia $f-g$ je neklesajúca na $\langle a, \infty)$ a platí $f(a)-g(a)>0$; odtiaľ vyplýva, že $f(x)-g(x)\geq f(a)-g(a)>0$ pre všetky $x\geq a$;
$\boxed{356.}$ nie (príslušný protiklad už musíte nájsť sami);
$\boxed{357.}$ $\boldsymbol{1.}$ rýdzo konvexná na $(-\infty ,1\rangle$, rýdzo konkávna na $\langle 1, \infty)$;
$\boldsymbol{2.}$ rýdzo konvexná na $\langle 0, \infty)$, rýdzo konkávna na $(-\infty ,0\rangle$;
$\boldsymbol{3.}$ rýdzo konkávna na každom z intervalov $\langle 2k\pi, (2k+1)\pi \rangle$, rýdzo konvexná na  každom z intervalov $\langle (2k-1)\pi, 2k\pi \rangle$, $k\in \mathbb{Z}$;
$\boldsymbol{4.}$ rýdzo konvexná na $\langle -1,1 \rangle$, rýdzo konkávna na $(-\infty ,-1\rangle$ a na $\langle  1,\infty )$; 
$\boldsymbol{5.}$ rýdzo konvexná na $\langle e^{\pi/4-(2k+1)\pi},e^{\pi/4-2k\pi} \rangle$, rýdzo konkávna na $\langle e^{\pi/4-2k\pi},e^{\pi/4-(2k-1)\pi} \rangle$, $k\in \mathbb{Z}$;
$\boldsymbol{5.}$ rýdzo konkávna na $(-\infty ,0\rangle$ a na $\langle 0, \infty)$, ale nie na ich zjednotení;
$\boxed{358.}$ $\boldsymbol{1.}$ $p,q$ nepárne, $p>q>0$ (funkcia $x^{p/q}$ je pre $x<0$ definovaná len vtedy, keď $p$ je párne alebo $p$ aj $q$ sú nepárne; derivácia v bode $0$ je vlastná len pre $p\geq q$);
$\boldsymbol{2.}$ $9b^{2}-24ac>0$;
$\boldsymbol{3.}$ $\vert a \vert \leq 2$ (pre $\vert a \vert < 2$ je  $y''>0$ na $\mathbb{R}$, teda $y$ je rýdzo konvexná na $\mathbb{R}$; pre  $\vert a \vert =2$ má $y''$ nulovú hodnotu len v jednom bode, inde $y''>0$; dá sa dokázať, že aj v tomto prípade je  $y$ rýdzo konvexná (pozri pr. 462);
$\boxed{360.}$ $\boldsymbol{3.}$ stačí obidve strany vydeliť číslom $2$ a využiť konvexnosť funkcie $x \ln x$;
$\boldsymbol{4.}$ $y=\frac{2}{\pi}x$ je rovnica spojnice bodov $(0,0)$, $(\frac{\pi}{2},1)$ ležiacich na grafe funkcie $\sin $;
$\boldsymbol{6.}$ $y=x-1$ je rovnica spojnice bodov $(1,0)$, $(2,1)$ ležiacich na grafe konkávnej funkcie $y=\log _{2}x $;
$\boxed{361.}$ sporom; $f$ by musela byť rýdzo konkávna, a teda aj nekonštantná na $\mathbb{R}$; ak $a < b$, $f(a)\ne  f(b)$, tak sa spojnica bodov $(a,f(a))$, $(b,f(b))$ pretína s osou $Ox$, pritom na $\mathbb{R}-\langle a,b \rangle$ leží graf funcie $f$ pod touto spojnicou, to vedie k sporu s kladnosťou funkcie $f$ (každý z uvedených krokov treba samozrejme podrobne zdôvodniť);
$\boxed{362.}$ $\boldsymbol{1.}$ pretože z definície konvexnosti funkcie $f$  na $I$ vyplýva $\forall x,y,z\in I, \quad x<z<y: \, f(z) <f(x)+(f(y)-f(x)).\frac{z-x}{y-x}$ (v definícii stačí položiť  $z=px+qy$ a odtiaľ vyjadriť $q=\frac{z-x}{y-x}$, môžeme dokázať, že funkcia $F(x):=\frac{f(x)-f(a)}{x-a}$ je rastúca na $I\cap (a,+\infty)$ a na $I\cap (-\infty,a)$; pritom $f'(a)=\lim_{x \to a}F(x)$; odtiaľ vyplýva $\forall x\in I, \, x>a: \quad \frac{f(x)-f(a)}{x-a}>f'(a)$, $\forall x\in I, \, x<a: \quad \frac{f(x)-f(a)}{x-a}<f'(a)$; z toho po úprave  $\forall x\in I, \, x \ne a: \quad f(x)>f(a)+f'(a)(x-a)$ (úvahy pre krajné a pre vnútorné body intervalu $I$ treba robiť každú samostatne);
$\boxed{363.}$ $f'(0)=0$ (vypočíta sa priamo z definície); pre $x>0$ je  $f(x)>0$, pre  $x<0$  je $f(x)<0$; $f''(x)=12x+(6x-\frac{1}{x})\cos\frac{1}{x}+4\sin\frac{1}{x}$ je spojitá na $\mathbb{R}-\lbrace 0 \rbrace$ a  na každom z intervalov $(-\varepsilon ,0)$, $(0,\varepsilon)$, $\varepsilon>0$, nadobúda kladné aj záporné hodnoty (vyplýva to z faktu, že pre $x=\frac{1}{2k\pi}$ je $f''(x)=18x-\frac{1}{x}$, pre  $x=\frac{1}{(2k+1)\pi}$ je $f''(x)=6x+\frac{1}{x}$ a obidve tieto funkcie majú v bode $0$ jednostranné limity $+\infty$ a $-\infty$); preto každý z intervalov $(-\varepsilon ,0)$, $(0,\varepsilon)$, $\varepsilon>0$ obsahuje intervaly, na ktorých je $f$ rýdzo konvexná, aj intervaly na ktorých je  $f$ rýdzo konkávna;
$\boxed{364.}$ $\boldsymbol{1.}$ lok. minimum: $f(0)=0$, lok. maximum: $f(1)=f(-1)=1$;
$\boldsymbol{2.}$ lok. minimum: $f(-2)=-9$, $f(3)=-\frac{161}{4}$; lok. maximum: $f(0)=7$;
$\boldsymbol{3.}$ lok. minimum: $f(k\pi-\frac{\pi}{4})=-\frac{e^{(4k-1)\pi/4}}{\sqrt{2}}$, $k$ párne;  lok. maximum: $f(k\pi-\frac{\pi}{4})=\frac{e^{(4k-1)\pi/4}}{\sqrt{2}}$, $k$ nepárne $(\cos x+\sin x=\sqrt{2}\sin (x+\frac{\pi}{4}))$; 
$\boxed{365.}$ $\boldsymbol{1.}$ lok. minimum: $f(1)=f(3)=3$; lok. maximum: $f(2)=4$;
$\boldsymbol{2.}$ lok. minimum: $f(1)=0$, $f(-\frac{5+\sqrt{13}}{6})$; lok. maximum: $f(-\frac{\sqrt{13}-5}{6})$;
$\boldsymbol{3.}$ lok. minimum: $f(\frac{7}{5})=-\frac{1}{24}$; 
$\boldsymbol{4.}$ lok. minimum: $f(1)=0$; lok. maximum: $f(e^{2})=\frac{4}{e^{2}}$;
$\boldsymbol{5.}$ lok. minimum: $f(-1)=0$; lok. maximum: $f(9)=\frac{10^{10}}{9}$;
$\boxed{366.}$ $\boldsymbol{1.}$ nemá;
$\boldsymbol{2.}$ má v bode 0 lok. maximum $f(\frac{3}{2})=\frac{1}{4}$ (ak $f$ je spojitá, tak $\vert f \vert$ nadobúda lokálne extrémy v bodoch lokálnych extrémov funkcie $f$ a v nulových bodoch funkcie $f$);
$\boldsymbol{2.}$ lok. maximum: $f(-\frac{1}{\sqrt[4]{27}})=f(\frac{1}{\sqrt[4]{27}})=\frac{2}{3\sqrt{3}}$; lok. minimum: $f(0)=0$ (posledné tvrdenie nevyplýva z viet 16 a 17 - $f'(0)$ totiž neexistuje - a treba ho preto dokázať samostatne);
$\boldsymbol{3.}$ lok. maximum: $f(3)=3$ (bod 3 nie je vnútorný bod $D(f)$ uvedené tvrdenie teda nevyplýva z viet 15 - 17 a treba ho dokázať); 
$\boldsymbol{4.}$ lok. minimum: $f(\frac{3}{4})=-\frac{3}{4\sqrt[3]{4}}$ (pretože neexistuje vlastná $f'(1)$, preskúmame bod 1 zvlášť: buď priamo ukážeme, že $f$ nemá v bode $1$ lokálvy extrém, alebo dokážeme, že $f'(1)$ je nevlastná a použijeme vetu 15);
$\boxed{368.}$ $\boldsymbol{1.}$ $M=11$, $m=2$;
$\boldsymbol{2.}$ $M=100,01$; $m=2$;
$\boldsymbol{3.}$ $M=8$; $m$ neexistuje ($f$ je rastúca na $\langle 0,4 \rangle$; zrejme $0=\inf _{x\in (0,4\rangle}f(x)$);
$\boldsymbol{4.}$  $M=1$ ($f$ rastie na  $\langle 0,\frac{\pi}{4} \rangle$, klesá na $\langle \frac{\pi}{4} ,frac{\pi}{2})$), $m$ neexistuje ($\lim_{x \to \pi/2-}f(x)=-\infty$);
$\boldsymbol{5.}$ $M$ neexistuje,  $m=4$;
$\boldsymbol{6.}$ $M=\frac{\pi}{4}$; $m=0$;
$\boxed{369.}$ $\boldsymbol{3.}$ vyšetrite monotónnosť funkcie $\frac{x^{2}+1}{x^{2}+x+1}$ a nájdite jej limity v bodoch $+\infty $, $-\infty $; na základe toho si načrtnite jej graf; obrázok je návodom k vlastnému dôkazu;
$\boxed{370.}$ $\boldsymbol{1.}$ dokážte, že $P'$ musí pri prechode cez tento bod zmeniť znamienko; využite, že $P'$ ako polynóm nepárneho stupňa má v bodoch $+\infty$, $-\infty$ nevlastné limity opačných znamienok;
$\boxed{371.}$ nemôže; nepriamo: nech $a<b$ sú také stacionárne body; funkcia $f/\langle a,b \rangle$ nadobúda v niektorom bode $c$ kompaktu $\langle a,b \rangle$ svoje minimum; treba dokázať $c\in (a,b)$ a potom využiť vetu 15;
$\boxed{372.}$ $f'$ je spojitá na $\mathbb{R}-\lbrace 0 \rbrace$ a v každom ľavom, resp. pravom okolí bodu $O$ nadobúda kladné aj záporné hodnoty, ďalej pozri myšlienku návodu k pr. 363;
$\boxed{373.}$ $\boldsymbol{1.}$ výška  $v=\frac{2R}{\sqrt{3}}$, polomer podstavy $r=\frac{\sqrt{2} R}{\sqrt{3}}$ (existuje medzi valcami vpísanými do tejto gule valec s najmenším objemom?);
$\boldsymbol{2.}$  $v=\frac{2R\sqrt{5-\sqrt{5}}}{\sqrt{10}}$,  $r=\frac{R\sqrt{5+\sqrt{5}}}{\sqrt{10}}$;
$\boxed{374.}$ bod $(\frac{-1+\sqrt{33}}{8},\frac{1+3\sqrt{33}}{8})$;
$\boxed{375.}$ ak vzdialenosť bodov $(a,b)$ a  $(c,f(c))$ je maximálna (minimálna), tak  $f'(c)=\frac{c-a}{f(c)-b}$, to vyplýva z vety 15; pri dôkaze faktu, že dotyčnica v bode $(c,f(c))$	je kolmá na spojnicu bodov $(a,b)$ a $(c,f(c))$,  treba rozlíšíť prípady $a=c$,	$a\ne c$;
$\boxed{376.}$ vo vydialenosti $\sqrt{ab}$ ($cm$) ($\varphi =\arctan \frac{b}{x-\arctan \frac{a}{x}}$, $b>a$);
$\boxed{377.}$ platí $f(x)=f(x)-f(a)$, $g(x)=g(x)-g(a)$; potom zlomok rozšíriť výrazom $\frac{1}{x-a}$; pozor, vetu 18 nemožno použiť (nie je zaručené splnenie predpokladu 1);
$\boxed{378.}$ $\boldsymbol{2.}$ $2$;
$\boldsymbol{3.}$ $\frac{1}{3}$;
$\boldsymbol{4.}$ $0$; (l'Hospitalovo pravidlo použiť n-krát za sebou);
$\boldsymbol{5.}$ $0$ (najprv subst. $t=\frac{1}{x^{2}}$);
$\boldsymbol{6.}$ $-\frac{1}{2}$ (najprv zlomok rozšíriť členom x a využiť, že $\lim_{x \to 0}\frac{x}{\tan x}=1$; použitie l'Hospitalovho pravidla sa tak zjednoduší);
$\boldsymbol{7.}$ $4$ (napísať v tvare súčtu dvoch limít; využiť, že $\lim_{x \to 0}\frac{x^{3}\cos \pi/x}{x^{2}}=0$, l'Hospitalovo pravidlo použiť len na výpočet druhej limity);
$\boxed{379.}$ $\boldsymbol{1.}$ $0$;
$\boldsymbol{2.}$ $0$;
$\boldsymbol{3.}$ $2$;
$\boldsymbol{4.}$ $0$;
$\boxed{380.}$ $\boldsymbol{1.}$ $1$;
$\boldsymbol{2.}$ $e$;
$\boldsymbol{3.}$ $e^{1/6}$;
$\boldsymbol{4.}$ $e^{-1/2}$;
$\boldsymbol{5.}$ $1$ (po prvom použití l'Hospitalovo pravidla dostaneme výraz $\frac{x}{\ln \frac{1}{x}}$; ten - hoci v bode $0$ nie je typu $\frac{\infty}{\infty}$ ani  $\frac{0}{0}$ - vyhovuje všetkým predpokladom vety 18 (zvlášť si všimnite predpoklad 3), možno teda ešte raz použiť l'Hospitalovo pravidlo; $\lim_{x \to 0}(x/\ln (1/x))$ možno samozrejme nájsť aj bez toho, stačí použiť vetu 10 z kapitoly 2);
$\boxed{381.}$ $\boldsymbol{1.}$ $\frac{1}{2}$; (po úprave na spoločného menovateľa zlomok rozšírte členom $x$ a využite, že  $\lim_{x \to 0}\frac{e^{x}-1}{x}=1$);
$\boldsymbol{2.}$ $\frac{2}{3}$; (po úprave na spoločného menovateľa  využite, že  $\lim_{x \to 0}\frac{x^{2}}{\sin ^{2}x}=1$);
$\boldsymbol{3.}$ $\frac{1}{2}$ (najprv subst. $\frac{1}{x}=1$);
$\boxed{382.}$ $\boldsymbol{1,2.}$ nie, nie je splnená podmienka 4 z vety 18;
$\boldsymbol{3.}$ nie, nie je splnená podmienka 3 z vety 18 (výpočet jednotlivých limít: 1. zlomok rozšíriť členom $x$ a použiť pr. 131.1; 2. využiť, že $\lim_{x \to 0}x\sin \frac{1}{x}=0$; 3. dosadiť);
$\boxed{383.}$ správny je druhý postup; derivácia menovateľa je $\cos x.e^{\sin x}.(x+2\cos x+\cos x.\sin x$), tá má nulové hodnoty v každom okolí bodu  $+\infty$, nie je teda splnená podmienka 2 z vety 18, preto sme neboli oprávnení používať  l'Hospitalovo pravidlo;
$\boxed{384.}$ $\boldsymbol{1.}$ na výpočet $\lim_{x \to a}\frac{f(x)-f(a)}{x-a}$ použite l'Hospitalovo pravidlo;
$\boldsymbol{2a/}$  $0=\lim_{x \to 0}f'(x)=f'(0)$;
$\boldsymbol{2b/}$  $f'(1)=+\infty$, $f'(-1)=+\infty$ (použili sme unalógiu tvrdenia z pr. 384.1 pre jednostranné derivácie);
$\boldsymbol{2c/}$  $f'(0)=\frac{1}{\sqrt{2}}$ (nezabudnite najprv preveriť spojitosť funkcie $f$ v bode 0);
$\boxed{385.}$ $\boldsymbol{1.}$ $\frac{1}{\sqrt{2}}+\frac{1}{\sqrt{2}}(x+\frac{1}{2})-\frac{1}{2:\sqrt{2}}(x+\frac{1}{2})^{2}+\frac{1}{2:\sqrt{2}}(x+\frac{1}{2})^{3}-\frac{15}{4:\sqrt{2}}(x+\frac{1}{2})^{4}$;
$\boldsymbol{2.}$ $x+\frac{x^{3}}{3}$;
$\boldsymbol{3.}$ $\frac{1}{\sqrt{2}}+\sqrt{2}x-\sqrt{2}x^{2}-\frac{4\sqrt{2}}{3!}x^{3}+\frac{8\sqrt{2}}{4!}x^{4}+\frac{16\sqrt{2}}{5!}x^{5}$;
$\boldsymbol{4.}$ $\sum_{k=0}^N (-2)^{k}(x+1)^{k}$ ($f^{(k)}(-1)=(-1)^{k}k!2^{k}$;
$\boldsymbol{5.}$ Taylorovým polynómom stupňa i je funkcia $T(x)=\frac{1}{2}$; Taylorove polynómy stupňov $2m$ a $2m+1$ ($m \in \mathbb{R}$) sa zhodujú a majú tvar $\frac{1}{2}+\sum_{k=0}^N \frac{(-1)^{k+1}4^{2k-1}}{(2k)!}(x-\frac{\pi}{4})^{2k}$ (predpis funkcie $f$ možno upraviť pomocou vzorcov $\sin ^{2}\alpha =\frac{1-\cos 2\alpha}{2}$, $\cos ^{2} \alpha =\frac{1+\cos 2\alpha}{2}$);
$\boxed{386.}$ platí (*) $A_{0}+...+A_{n}(x-a)^{n}+o((x-a)^{n})=f(x)=f(a)+...+\frac{f^{(n)}}{n!}(x-a)^{n}+o((x-a)^{n})$ (druhá rovnosť vyplýva z vety 19, odtiaľ vyplýva rovnosť limity pravej a ľavej strany v bode $a$, preto $A_{0}=f(a)$, ak po odstránení $A_{0}$, resp. $f(a)$ z ľavej resp. pravej strany vydelíme obidve strany členom $(x-a)$ (a využijeme, že $\frac{o((x-a)^{n})}{x-a}=o((x-a)^{n-1})$), dostaneme $A_{1}+A_{2}(x-a)+...+A_{n}(x-a)^{n-1}+o((x-a)^{n-1})=f'(a)+...+\frac{f^{(n)}}{n!}(x-a)^{n-1}+o((x-a)^{n-1})$, odtiaľ opäť vyplýva rovnosť limít v bode a atď;
$\boxed{387.}$ $\boldsymbol{1.}$  $e^{\frac{x}{7}}=\sum_{k=0}^n \frac{x^{k}}{7^{k}.k!}+o(x^{n})$;
$\boldsymbol{3.}$ Maclavrinove polynómy stupňa 1, 2, ..., 5 sa zhodujú a majú tvar $T(x)=1$; Maclavrinove polynómy stupňa $n=6m, 6m+1,...,6m+5$ $(m \in \mathbb{N}$ sa zhodujú a majú tvar $T(x)=1+ \sum_{k=1}^m (-1)^{k}\frac{x^{6k}}{(2k!)}$, preto možno písať $\cos x^{3}=1+\sum_{k=1}^m (-1)^{k}.\frac{x^{6k}}{(2k!)}+o(x^{6m+5})$;
$\boldsymbol{5.}$ $\frac{1}{2+3x^{2}}=\frac{1}{2}. \frac{1}{1-(-\frac{3}{2}x^{2})}=\frac{1}{2}+\sum_{k=1}^m (-\frac{3}{2})^{k}x^{2k}+o(x^{2m+1})$,  Maclavrinov polynóm stupňa 1 má tvar $T(x)=\frac{1}{2}$ (využili sme vzorec V;  Maclavrinove polynómy stupňa $n=2m$ a $n=2m+1$ sa zhodujú);
$\boldsymbol{5.}$ $x^{3}\sin 3x= \sum_{k=1}^m (-1)^{k+1}\frac{3^{2k-1}x^{2k+2}}{(2k-1)!}+e(x^{2m+3})$ (teda Maclavrinov polynóm stupňa 1, 2, 3 majú rovnaký  tvar $T(x)=0$,  Maclavrinove polynómy stupňov  $n=2m+2, \, n=2m+3\quad (m\in \mathbb{N})$ sa zhodujú);
$\boxed{388.}$ $\boldsymbol{2.}$ $1+x-x^{2}+\frac{3}{2}x^{3} $ (najprv sme použili vzorec VI pre $\alpha=\frac{1}{3} $);
$\boldsymbol{3.}$ $-\frac{x^{3}}{8}$ (obidve strany rovnosti $\ln (1-\frac{x}{2})=-\frac{x}{2}+o(x)$ sme umocnili na tretiu);
$\boldsymbol{4.}$	$x+\frac{x^{2}}{2} -\frac{2x^{3}}{3}+\frac{x^{4}}{4} $ (položili sme  $ z=x^{2}+x $  a najprv použili vzorec IV; ak  $R(z)=o(z^{4})$, tak  $R(x^{2}+x)= o(z^{4}) $;
$\boldsymbol{5.}$ $\cos x-\ln (1+x)=(1-\frac{x^{2}}{2}+\frac{x^{4}}{24}+o(x^{4})). (x-\frac{x^{2}}{2}+\frac{x^{3}}{3}-\frac{x^{4}}{4}+\frac{x^{5}}{5}+ o(x^{5}))=x-\frac{x^{2}}{2}-\frac{x^{3}}{6}+\frac{3x^{5}}{40}+ o(x^{5}))$;
$\boxed{389.}$ $\boldsymbol{1.}$  $\sum_{k=0}^m\frac{e^{-2}}{k!}(x+1)^{2k} $ pre $n=2m$ aj pre $n=2m+1$;
$\boldsymbol{2.}$  $\sum_{k=0}^m\frac{(-1)^{k}e^{2}2^{k-2}}{k!}(k^{2}+3k+4)(x+1)^{k} $ $f(t+1)=e^{2}.(e^{-2t}-2te^{-2t}+t^{2}e^{-2t})=e^{2}[\sum_{k=0}^n\frac{(-1)^{k}2^{k}t^{k}}{k!} +o(t^{n})+\sum_{i=0}^{n-1}\frac{(-1)^{i+1}2^{i+1}t^{i+1}}{i!} +o(t^{n})+\sum_{j=0}^{n-2}\frac{(-1)^{j}2^{j}t^{j}}{j!} +o(t^{n})]$ $=e^{2}[\sum_{k=0}^n\frac{(-1)^{k}2^{k}t^{k}}{k!}+\sum_{k=1}^n\frac{(-1)^{k}2^{k}t^{k}}{(k-1)!}+\sum_{k=2}^n\frac{(-1)^{k}2^{k-2}t^{k}}{(k-2)!}+o(t^{n})]=$ (položili sme $i+1=k,j+2=k$)$=e^{2}[1-2t+\sum_{k=2}^n\frac{(-1)^{k}2^{k}t^{k}}{k!}-2t+\sum_{k=2}^n\frac{(-1)^{k}2^{k}t^{k}}{(k-1)!}+\sum_{k=2}^n\frac{(-1)^{k}2^{k}t^{k}}{(k-2)!}+o(t^{n})]$;
$\boldsymbol{3.}$   $\frac{x}{2}-\frac{x^{2}}{4}+\sum_{k=3}^n\frac{(-1)^{(k+1)}(k-1)x^{k}}{k(k-2)}$ pre  $n\geq 3$ Maclavrinov polynóm stupňa 1, resp. 2 je $\frac{x}{2}-\frac{x^{2}}{4}(\ln \sqrt{1+x}=\frac{1}{2}\ln (1+x))$;
$\boldsymbol{4.}$ $f(x)=1-x^{2}+\sum_{k=2}^m 2.(-1)^{k}x^{2k}+o(x^{2m+1})$;
$\boldsymbol{5.}$ $f(x)=-\frac{7}{6}+\sum_{k=1}^n (-\frac{5}{3})(\frac{10}{21})^{k}(x+\frac{1}{10})^{k}+o((x+\frac{1}{10})^{n})$  $(f(t-\frac{1}{10})=\frac{1}{2}-\frac{5}{3}.\frac{1}{1-\frac{10t}{21}})$;
$\boldsymbol{6.}$ $f(x)= \ln 6 - \sum_{k=1}^n \frac{1}{k}(\frac{1}{2^{k}}+\frac{1}{3^{k}})(x-1)^{k}+o(x-1)^{n})$ (pre $t<2$ platí $f(t+1)=\ln ((t-2)(t-3))=\ln (2-t)+\ln (3-t)=\ln 2 +\ln (1-\frac{t}{2}) +\ln 3+\ln (1-\frac{t}{3}))$;
$\boxed{390.}$ $\boldsymbol{1.}$ $\frac{1}{2} $  ($= \lim_{x \to 0}\frac{x^{3}/2+o(x^{3})}{x^{3}})$;
$\boldsymbol{3.}$ $\frac{1}{3} $;
$\boldsymbol{4.}$ $10 $ ($= \lim_{x \to 0}(-\frac{x^{3}/6+o(x^{3})}{x^{2}+o(x^{2})})$);
$\boldsymbol{5.}$ $\frac{4}{3} $;
$\boldsymbol{6.}$ $\frac{1}{2} $ ($\ln (\cos x +\frac{x^{2}}{2})= \ln (1+\cos x -1 +\frac{x^{2}}{2}))$, potom použiť vzorec IV; ak $R(z)=o(z^{2})$, tak $R( \cos x -1+\frac{x^{2}}{2})= o(x^{4}))$;
$\boxed{391.}$ $\boldsymbol{1.}$ $\frac{60!}{30!} $  ($f^{(60)}(0)$ možno vyjadriť pomocou koeficientu pri $x^{60}$ v Maclavrinovom polynóme stupňa $n\geq 60$ funkcie $f$);
$\boldsymbol{2.}$ $0$  ($\frac{1}{1+x+x^{2}}=\frac{1-x}{1-x^{3}}$ pre $x \ne 1$);
$\boxed{392.}$ $\boldsymbol{1.}$ $\bigtriangleup < \frac{e}{11!}$ (použili sme Lagrangeov tvar zvyšku);
$\boldsymbol{2.}$ $\bigtriangleup < \frac{\cos 0,5}{3840}$ (polynóm vpravo je Maclavrinovým polynómom stupňa 3 aj stupňa 4 funkcie $\sin$, preto sme odhadli zvyšky v obidvoch prípadoch a z odhadov vybrali ten menší);
$\boldsymbol{3.}$ $\bigtriangleup < 0,0015$;
$\boxed{393.}$ $\boldsymbol{2.}$ podľa vety 20 $e^{x}=1+x+\frac{e^{\upsilon (x)}}{2}x^{2}$, pritom $\upsilon (x)<x$;
$\boxed{394.}$ obr. 5, funkcia je nepárna;
$\boxed{395.}$ obr. 6;
$\boxed{396.}$ obr. 7;
$\boxed{397.}$ obr. 8;
$\boxed{Obr.5.}$  $y=3x-x^{3}$
$\boxed{Obr.6.}$  $y=x^{4}/(1+x)^{3}$
$\boxed{Obr.7.}$  $y=(x+1)^{3}/(x-1)^{2}$
$\boxed{Obr.8.}$  $y=(x^{2}+x-1)/(x^{2}-2x+1)$
$$*$$
$$*$$
$$*$$
$$*$$
(Pri popise obrázkov sme použili tieto označenia: n - bod, v ktorom sa funkčná hodnota rovná 0; m - bod, v ktorom funkcia nadobúda lokálne minimum; M - bod, v ktorom funkcia nadobúda lokálne maximum; i - inflexný bod; ak je os Ox (Oy) asymptotou, je k nej pripísaná jej rovnica y=0 (x=0), ostatné asymptoty sú vyznačené čiarkovane; s výnimkou obr. 17 a 23 sa jednotka dĺžky na Ox a na Oy zhodujú a jednotka dĺžky je označená len na Oy.)
$\boxed{398.}$ obr. 9; nepárna funkcia, $y_{-}'(-1)=y_{+}'(-1)=+\infty $; pozor: na intervale $\langle -\sqrt{\frac{3}{2}},0 \rangle$ $y$ nie je konvexná, hoci je konvexná na intervale $\langle -\sqrt{\frac{3}{2}},-1 \rangle$ a na intervale $\langle -1,0 \rangle$; podobná poznámka sa vzťahuje aj na interval $\langle 0,\sqrt{\frac{3}{2}} \rangle$;
$\boxed{399.}$ obr. 10; 
$\boxed{400.}$ obr. 11;	 nepárna fukcia;
$\boxed{401.}$ obr. 12;  párna funkcia, $y_{-}'(-1)=y_{-}'(1)=-\infty $, $y_{+}'(-1)=y_{+}'(1)=+\infty $; $y$ nie je konkávna na množine $\mathbb{R}$, hoci je konkávna na intervaloch $(-\infty ,-1\rangle , \langle -1,1 \rangle $ a $\langle 1, \infty)$;
$\boxed{402.}$ obr. 13;  $y_{-}'(0)=-\infty $, $y_{+}'(0)=+\infty $; $y$ nie je konkávna na intervale $\langle -2+\sqrt{3}, +\infty)$, hoci je konkávna na intervaloch $\langle -2+\sqrt{3},0 \rangle$ a $\langle 0,+\infty)$;
$\boxed{403.}$ obr. 14; funkcia s periódou $2\pi$;
$\boxed{404.}$ obr. 15; párna funkcia s periódou $2\pi$;
$\boxed{405.}$ obr. 16; $y_{-}'(0)=-\infty $, $y_{+}'(0)=+\infty $; $y$ nie je konkávna na intervale $\langle \frac{2-\sqrt{6}}{3},\frac{2+\sqrt{6}}{3} \rangle$;
$\boxed{Obr.9.}$  $y=x \sqrt{\vert}x^{2}-1 \vert)$
$\boxed{Obr.10.}$  $y=(x-2)/\sqrt{x^{2}+1}$
$\boxed{Obr.11.}$  $y=x/\sqrt[3]{x^{2}-1}$
$$*$$
$$*$$
$$*$$
$$*$$
$\boxed{406.}$ obr. 17;
$\boxed{407.}$ obr. 18;
$\boxed{408.}$ obr. 19; nepárna funkcia,  $y_{-}'(1)=y_{+}'(-1)=1 $,  $y_{+}'(1)=y_{-}'(-1)=-1 $; pozor $-1,1$ nie sú podľa nami používanej definície inflexnými bodmi (pri úprave výrazu pre $y'$ nezabudnite, že  $\sqrt{a^{2}}=\vert a \vert $, pri výpočte $y_{+}'(1),y_{+}'(-1)$ použite obdobu tvrdenia z pr. 384.1 pre jednostranné derivácie);
$\boxed{409.}$ obr. 20; $0 \notin D(y)$, existuje ale $\lim_{x \to 0-}y(x)=0$;
$\boxed{410.}$ obr. 21; 
$\boxed{411.}$ obr. 22; 0 je bod odstrániteľnej nespojitosti,  $\lim_{x \to 0}y(x)=1$;
$\boxed{413.}$ $(f.g)'(a)=f(a).g'(a)$;
$\boxed{414.}$ tvrdenie vyplýva zo vzťahu $\frac{f^{n}(x)-f^{n}(a)}{x-a}=\frac{f(x)-f(a)}{x-a}.(f^{n-1}(x)+f^{n-2}(x).f(a)+...+f^{n-1}(a))$;
$\boxed{417.}$ $\boldsymbol{1.}$ neplatí
$$*$$
$$*$$
$$*$$
$$*$$
$\boxed{Obr.12.}$  $y=(x+1)^{2/3}+(x-1)^{2/3}$
$\boxed{Obr.13.}$  $y=\sqrt[3]{x^{2}/(x+1)}$
$\boxed{Obr.14.}$  $y=\frac{1}{2}\sin 2x +\cos x$
($x+\sin x$ je neperiodická funkcia s periodickou deriváciou);
$\boldsymbol{2.}$ platí;
$\boxed{418.}$ $(f^{2})_{\pm}'(a)=2f(a).f_{\pm}'(a)$; nutná a postačujúca podmienka je $f(a)=0 $;
$\boxed{419.}$ $\boldsymbol{1.}$  $\sum_{j=1}^m$  ($\sum_{l=1}^n$ ($\varphi _{lj}'.\prod _{k=1,k\ne l}^n \varphi _{kj}))$;
$\boldsymbol{2.}$  $\sum_{j=1}^m$  ($\sum_{l=1}^n$ ($\varphi _{lj}'.\prod _{j=1}^m \sum_{j=1}^n \varphi _{kj}))$;
$\boxed{420.}$ $\boldsymbol{1.}$ $F'(x)=3x^{2}+3x+20$;
 $\boldsymbol{2.}$ $F'(x)=6x^{2}$;
$\boxed{421.}$ $\boldsymbol{1.}$ $f(x)=\vert x-a\vert$, $g(x)=x^{2}$;
$\boldsymbol{2.}$ $f(x)=( x-a)^{2}$, $g(x)=\vert x \vert$;	
$\boldsymbol{3.}$ $f(x)=2( x-a)+\vert x-a \vert $, $g(x)=2x - \vert x \vert$ (všimnite si, že jednostranné derivácie$ g\circ f$ možno v tomto prípade vypočítať pomocou $f_{+}'(a)$ a $g_{+}'(0)$, resp. $f_{-}'(a)$ a $g_{-}'(0))$;
$\boxed{422.}$ za daných predpokladov musí platiť
$$*$$
$$*$$
$$*$$
$$*$$
$\boxed{Obr.15.}$  $y=\cos x / \cos 2x$
$\boxed{Obr.16.}$  $y=x^{2/3}e^{-x}$
$f(a)=0$, $f'(a)\ne 0$, ďalej pozri pr. 412;
$\boxed{423.}$ 1. $f(x)=ax$ pre $x\ne \frac{1}{n}$, $x\ne n$ ($n \in \mathbb{N}$) (v tomto prípade využívame možnosť širšieho chápaania pojmu derivácia pre $f'(0))$, alebo 2. $f$ je nepárna funkcia definovaná pre $x\geq 0$ nasledovne: zoraďme spočitateľnú množinu ($\mathbb{Q}\cap \langle 0, \infty)) - \lbrace x \in \mathbb{R}:x=ap^{2}+ap, \, p\in \langle 0, \infty) \cap (\mathbb{Q}- \mathbb{N}) \rbrace $ (tá obsahuje aj všetky čísla tvaru $\frac{a}{n}$, kde $n \in \mathbb{N}$) do prostej postupnosti $\lbrace b_{n} \rbrace _{n=1} ^{\infty}$, potom položíme $f(x)=ax$ pre $x \in (0,\infty) - \mathbb{Q}$, $f(x)=ax^{2}+ax$ pre $x \in \langle 0,\infty )\cap (\mathbb{Q}- \mathbb{N})$, $f(n)=b_{n}$ $ n \in \mathbb{N}$ ( v tomto prípade nevyužívame možnosť širšieho chápania pojmu derivácia ani pre $f'(0)$ ani pre ($f^{-1})'(0))$;
$$*$$
$$*$$
$$*$$
$$*$$
$\boxed{Obr.17.}$  $y=\ln x / \sqrt{x}$
$\boxed{Obr.18.}$  $y=x \arctan x$
$\boxed{424.}$ $f''(0)=0$ (najprv treba odvodiť vzorec pre druhú deriváciu inverznej funkcie: ($f^{-1})''(x)=-\frac{f''(f^{-1}(x))}{(f')^{3}(f^{-1}(x))}$; stanovte predpoklady, za ktorých možno toto odvodenie vykonať);
$\boxed{425.}$ $\boldsymbol{1.}$ $y^{(n)}=\frac{3^{n}}{4}\cos (3x+\frac{n\pi}{2})+\frac{3}{4}\cos (x+ \frac{n\pi}{2}) $;
$\boldsymbol{2.}$ $y^{(n)}=2^{2n-3}\cos (4x+\frac{n\pi}{2}+2^{n-1}.\cos (2x+\frac{n\pi}{2})$;
$$*$$
$$*$$
$$*$$
$$*$$
$\boxed{Obr.19.}$  $y=\ln x / \arcsin (2x/(1+x^{2})$
$\boxed{Obr.20.}$  $y=(x+2)e^{1/x}$
$\boxed{Obr.21.}$  $y=\arccos((1-x)/(1-2x))$
$\boxed{Obr.22.}$  $y=x^{x}$
$\boldsymbol{3.}$ $y^{(n)}=72^{x}\ln ^{n-1} 72. ((2x-1)\ln 72+2n)$;
$\boldsymbol{4.}$ úpravu $\ln (x^{2}-3x+2)=\ln (x-2)+\ln(x-1)$ nemožno priamo použiť (porovnaj s pr. 26.2), treba si pomôcť nasledovne: pre všetky $x \in D((\ln f)'(x)=(\ln \vert f \vert )'(x)$; $y^{(n)}=(n-2)!(-1)^{n}.(\frac{x-2n}{(x-2)^{n}}+\frac{x-n}{(x-1)^{n}})$ pre $n\geq 2$; $y'=\ln (x^{2}-3x+2)+\frac{2x^{2}-3x}{x^{2}-3x+2}$;
$\boldsymbol{5.}$ $y^{(n)}=(\sqrt{2})^{n}e^{x} \sin (x+\frac{n\pi}{4})$ (treba použiť rovnosť $\sin x +\cos x=\sqrt{2} \sin (x+\frac{\pi}{4}) $ a indukciu);
$\boxed{426.}$ $f^{(n)}(a)=n\varphi ^{(n-1)}(a)$ (na výpočet $f^{(n)}(a)$ nemožno použiť vetu 7, existencia $\varphi ^{(n)}(a)$ nie je zaručená);
$\boxed{429.}$ pozri pr. 329;
$\boxed{430.}$ $a_{0}+...+a_{n}x^{n}=(a_{0}\frac{x}{1}+...+a_{n}\frac{x^{n+1}}{n+1})'$, polynóm v zátvorke vpravo má korene 0,  1;
$\boxed{431.}$ pozri pr. 332;
$\boxed{432.}$ pozri pr. 331, polynóm $(x^{2}-1)^{n}$ má n-násobné korene 1 a -1;
$\boxed{433.}$ sporom, použiť pr. 331;
$\boxed{435.}$ možno postupovať indukciou, pričom dôkaz pre dané n sa robí sporom; ak P je polynóm s n zmenami znamienok, tak $(x^{-k}P(x))'$ (kde $k\in \mathbb{N}$ je vhodne zvolené, pozri návod k pr. 434) má n-1 zmien znamienok, pritom $(x^{k+1}(x^{-k}P(x))'$ je polynóm, ktorý má rovnaké kladné korene ako $(x^{-k}P(x))'$ a tiež $n-1$ zmien znamienok;
$\boxed{437.}$ $f$ je polynóm stupňa $k$, kde $k\leq n-1$ (pozri návod k pr. 339);
$\boxed{438.}$ ak $f'(x_{0})=0$, tak $\frac{f(x)-f'(x_{0})}{f'(x_{0})} =\frac{f'(c(x))}{f'(x_{0})}(x-x_{0})$, pričom uvedený podiel derivácií je ohraničená funkcia;
$\boxed{441.}$ pozri pr. 343;
$\boxed{444.}$  možno použiť pr. 443 tak, ako sa používa Rolleho veta pri dôkaze Lagrangeovej vety o strednej hodnote;
$\boxed{445.}$ sporom	 keby to nebola pravda, tak by existoval interval $I\subset (a,b)$, na ktorom by funkcia $f$ mala nulovú deriváciu, čo je spor s jej injektívnosťou (treba si uvedomiť, že $f$ aj $f^{-1}$ sú monotónne spojité funkcie);
$\boxed{446.}$ treba použiť Cauchyho vetu pre funkcie $x^{2}.f(x),x^{2}$;
$\boxed{447.}$ (pozri tiež pr. 352)
$\boldsymbol{1.}$ ak $f(x):=\ln x$, $g(x):=\frac{x-1}{\sqrt{x}}$, tak $f(1)=g(1)$, $f'(x)< g'(x)$ pre $x>0, x\ne 1$, preto $f(x)< g(x)$ pre $x>1$,  $f(x)>g(x)$ pre $x\in (0,1)$;
$\boldsymbol{2.}$ najprv zlogaritmovať, potom položiť $x=\frac{1}{t}$;
$\boldsymbol{3.}$ môžeme predpokladať $x\leq y$; položme $t=\frac{x}{y}$, dostaneme ekvivalentnú nerovnosť $(1+t^{\alpha})^{1/\alpha}>(1+t^{\beta})^{1/\beta}$, $t\in (0,1 \rangle$, potom zlogaritmovať;
$\boxed{449.}$ $\boldsymbol{3.}$ funkcia $h:=f-g$ vyhovuje všetkým predpokladom tvrdenia z pr. 449.1;
$\boxed{450.}$ označme $O(c)$ to okolie bodu $c\in I$,v ktorom platí $\forall x,y \in O(c)$, $x<c<y:f(x)<f(c)<f(y)$; nech $a,b\in I$, $a<b$, potom $\langle a,b \rangle$ je kompatk, z otvoreného pokrytia $\lbrace O(c);c\in \langle a,b \rangle \rbrace $ vyberme konečné podpokrytie a z neho také konečné pokrytie $\lbrace O(c_{1}),...,O(c_{n})\rbrace$, že $O(c_{i})\cap O(c_{i+1})\ne \varnothing$ pre $i=1,...,n-1,$ $a<c_{1}<x_{1}<c_{2}<x_{2}<...<x_{n-1}<c_{n}<b$, pričom  $a\in O(c_{1})$, $b\in O(c_{n})$, $x_{i}\in O(c_{i})\cap O(c_{i+1})$ pre $i=1,...,n-1$; potom $f(a)<f(c_{1})<f(x_{1}<...<f(b))$; teda $f(a)<f(b))$, ak $a<b$;
$\boxed{451.}$ $\boldsymbol{1.}$ to, že $f$ je rastúca v bode 0, možno dokázať (aspoň) dvoma spôsobmi: a/ priamo z definície: pre $x\in (0,1)$ je $x+x^{2} \sin \frac{2}{x}\geq x-x^{2} = x(1-x)>0$, pre $x\in (-1,0)$ je $x+x^{2} \sin \frac{2}{x}< x-x^{2}<0$ využili sme nerovnosť $\vert \sin \frac{2}{x}\vert \leq 1$ pre $x\ne 0$); b/ na základe pr. 412, pretože $f'(0)=1$;
$\boldsymbol{2.}$ $f'$ je spojitá na $\mathbb{R}- \lbrace 0 \rbrace$ a v ľubovoľne malom okolí bodu 0 nadobúda kladné aj záporné hodnoty (stačí vypočítať $f'(\frac{2}{\pi/2+2k\pi})$ pre $k \in \mathbb{Z}$ a $f'(\frac{1}{k\pi})$ pre $k \in \mathbb{Z} - \lbrace 0 \rbrace $, preto $f$ nemôže byť rastúca na žiadnom okolí bodu 0 (porovnaj aj s pr.350);
$\boxed{452.}$ uvedieme dva návody: a/ ukázať, že $f$ je klesajúca v každom bode $x\in I$, ďalej porovnaj s pr.450; b/ ak $M$ je konečná, porovnaj s pr.354;  ak $M$ je nekonečná, tak má aspoň jeden hromadný bod, ten musí byť krajným bododm intervalu $I$ (označme krajné body intervalu $I$ ako a, b, $a<b$); ak je tým hromadným bodom len bod b, možno  $M$ zoradiť do rastúcej postupnosti  $\lbrace a_{n} \rbrace _{n=1}^{\infty}$, $\lim_{n \to \infty}a_n =b$, pričom $f$ je klesajúca na $\langle a_{n}, a_{n+1} \rangle$  $n\in \mathbb{N}$; úvahy pre prípad, že množina hromadných bodov množiny $M$ je $\lbrace a \rbrace$, resp. $\lbrace a,b \rbrace$ sú podobné;
$\boxed{453.}$ matematickou indukciou; $f(\lambda _{1}x_{1}+...+  \lambda _{n}x_{n}+  \lambda _{n+1}x_{n+1})=f((1-\lambda _{n+1})(\frac{\lambda  _{1}}{1-\lambda _{n+1}}x_{1}+...+\frac{\lambda _{n}}{1-\lambda _{n+1}}x_{n})+\lambda _{n+1}x_{n+1})$, pritom  $\lambda _{1}+...+\lambda _{n}=1-\lambda _{n+1}$;
$\boxed{454.}$ $\boldsymbol{1.}$ využiť Jensenovu nerovnosť pre funkciu $f(x)=x^{2}$ (zvoliť  $\lambda _{1}=...=\lambda _{n}=\frac{1}{n}$);
$\boldsymbol{2.}$ zlogaritmovať; funkcia - $\ln x$ je konvexná na $\mathbb{R}^{+}$;
$\boldsymbol{3.}$ vynásobiť obidve strany číslom $1/n^{r}$ a využiť Jensenovu nerovnosť pre funkciu $f(x)=x^{r}$, $r>1, x>0$;
$\boxed{456.}$ ak je polynóm P párnou funkciou, tak $P(x)=a_{0}+a_{2}x^{2}+...+a_{2m}x^{2m}$; súčet klaných násobkov konvexných funkcií je konvexná funkcia;
$\boxed{458.}$ stačí dokázať druhé z uvedených tvrdení; pre $x\in (c,b)$ je funkcia $F(x):=\frac{f(x)-f(c)}{x-c}$ neklesajúca (pozri návod k pr. 362.1) a zdola ohraničená (napr. číslom $=\frac{f(a)-f(c)}{a-c}$; pre $x\in (c,b)$ ležia body grafu funkcie $f$  nad alebo na spojnici bodov $(a,f(a)), (c,f(c))$, preto $\lim_{x \to c+}F(x)$  existuje a je konečná (z uvedeného tiež vyplýva, že $f_{-}'(c)\leq f_{+}'(c)$);
$\boxed{459.}$ stačí použiť pr. 362.2;
$\boxed{461.}$ inverzná k rastúcej konvexnej (konkávnej) je rastúca konvexná (konkávna), ku klesajúcej  konvexnej (konkávnej) je klesajúca konvexná (konkávna); pri  dôkaze použite definíciu konvexnosti prepísanú na tvar $\forall x,y,z\in I,\, x<z<y: f(z)\leq f(x) +\frac{f(y)-f(x)}{y-x}(z-x)$ (porovnaj s návodom k pr. 362.1) a inšpirujte sa obrázkom;
$\boxed{463.}$ $\boldsymbol{1.}$ lok. maximum $3\sqrt[3]{3}$ pre $x=-3$, lok. minimum $-\sqrt[3]{44}$ pre $x=2$ (stačí nájsť lokálne extrémy odmocnenca a využiť, že $f(z)=\sqrt[3]{z}$ je rastúca funkcia);
$\boldsymbol{2.}$  pre $n$ párne: lok. maximum $n^{n}e^{1-n}$ pre $x=n-1$, lok. minimum 0 pre  $x=-1$; pre $n$ nepárne: lok. maximum $n^{n}e^{1-n}$ pre $x=n-1$;
$\boldsymbol{3.}$ ak $n$ je  párne: lok. maximum  $1$ pre $x=0$; ak $n$ je  nepárne: nemá lok. extrémy;
$\boldsymbol{4.}$ lok. minimum 0 pre  $x=0$ (pri úprave výrazu pre $y'$ využite vzorce $\sin \alpha \pm \cos \alpha = \sqrt{2}\sin (\frac{\pi}{4}\pm \alpha )$; $y$ rastie pre $x\geq 0$, klesá pre $x\leq 0$);
$\boxed{464.}$ $\boldsymbol{2.}$  najprv vydeľte výrazom $(x+a)$, potom položte $z=\frac{x}{a}$  a vyšetrite priebeh funkcie $f(z)=\frac{\sqrt[n]{1+z^{n}}}{1+z}$ $z\geq 0$;
$\boxed{465.}$  $x=\frac{H}{2}$ (dostrek je $\sqrt{2gx}t$, pričom $t$ vypočítame z rovnice $\frac{gt^{2}}{2=H-x}$, tj. $t$ je čas, za ktorý teleso spadne voľným pádom z výšky $H-x$);
$\boxed{466.}$  $(a^{2/3}+b^{2/3})^{3/2}$;
$\boxed{467.}$  $\boldsymbol{1.}$ $-3$;
$\boldsymbol{2.}$ $\frac{1}{3}$;
$\boldsymbol{3.}$ $0$ (rozšírte členom x a využite, že $\lim_{x \to 0}\frac{\ln (1+x)}{x}=1)$;
$\boldsymbol{4.}$ $+\infty$ (rozšírte členom x, l'Hospitalovo pravidlo použite len na prvý člen v čitateli);
$\boldsymbol{5.}$ $0$ (v exponente vyňať $x^{n}$, l'Hospitalovo pravidlo použite len na druhý člen v zátvorke);
$\boldsymbol{6.}$ $0$ (upraviť na tvar $e^{f(x)}$, v $f(x)$ vyňať $\ln ^{2}x$ pred zátvorku, l'Hospitalovo pravidlo použite len na druhý člen v zátvorke);
$\boldsymbol{7.}$ $1$ pri úpravách využite, že $\lim_{x \to 0}x^{x}=1)$;
$\boldsymbol{8.}$ $-1$;
$\boldsymbol{9.}$ $+\infty$;
$\boldsymbol{10.}$ $-\frac{1}{6}$ (odpočítať a pripočítať $\sqrt{x^{2}+x+1}$; prvý rozdiel možno previesť na typ $\frac{0}{0}$ substitúciou $x=\frac{1}{t}$ a vyňatím $\frac{1}{t}$; druhý rozdiel možno úpraviť na tvar $\ln (1+\frac{x}{e^{x}}).\frac{\sqrt{x^{2}+x+1}}{x})$;
$\boldsymbol{11.}$ a (vyňať $(x+a)^{1/x}$ a upraviť na tvar a $(x+a)^{1/x}-x(\frac{x^{1/(x+a)}}{(x+a)^{1/x}}-1)$, prvý člen v zátvorke upraviť na tvar $e^{u}$ a využiť, že $\lim_{u \to 0}\frac{e^{u}-1}{u}=1)$;
$\boxed{468.}$ pri dôkaze neexistencie $\lim_{x \to +\infty}\frac{f(x)}{g(x)})$ využite, že v bodoch $k\pi - \frac{\pi}{4}$ má funkcia $\frac{f}{g}$ nevlastné jednostranné limity opačného znamienka, preto  $\frac{f}{g}$ nie je zhora ani zdola ohraničená v žiadnom okolí bodu  $+\infty$;
$\boxed{469.}$ $\boldsymbol{1.}$ $\sum_{k=0}^n (-1)^{k}(k+1)x^{3k+1}$ pre $n=3n+1, 3n+2, 3n+3$ ($n\in \mathbb{N} \cup \lbrace 0 \rbrace $);
$\boldsymbol{2.}$ $\sum_{k=1}^n c^{k-1}_{-1/2}.(-1)^{k-1}\frac{x^{2k}}{2^{2k-1}}+\sum_{k=0}^n c^{k}_{-1/2}.(-1)^{k}\frac{x^{2k+1}}{2^{2k}}$, kde $c_{\alpha}^{0}=1$, $c_{\alpha}^{k}=\frac{\alpha(\alpha -1).....(\alpha -k+1)}{k!}$, $k\in \mathbb{N}$ (pre $x \in (-2,2)$ platí $f(x)=x(2+x)(4-x^{2})^{-1/2}$);
$\boldsymbol{3.}$ $\sum_{k=0}^n \frac{e^{-27}3^{k}}{k!}(x+3)^{2k+1}$ pre $n=2n+1,2n+2$ ($n=0,1,...)$;
$\boldsymbol{4.}$ $ 1+ \sum_{k=1}^n \frac{(-1)^{k-1}3^{k}}{k.2^{k}\ln 2}(x-4)^{2k}$ pre $n=2n,2n+1$  ($n\in \mathbb{N}$), Taylorov polynóm stupňa 1 má tvar $T(x)=1$ (využite vzorec $\log _{a} b=\frac{\ln b}{\ln a}$);
$\boldsymbol{5.}$ $\sum_{k=1}^n c^{k-1}_{-1/3}.(-1)^{k-1}(x-1)^{2k}$, Taylorov polynóm stupňa 1 má tvar $T(x)=0$ (symbol $c^{i}_{\alpha}$ pozri v návode k pr.469.2);
$\boldsymbol{6.}$ $-\frac{\sin 2}{2}+\frac{\sin 2}{2}(x+1)+\sum_{k=0}^n \frac{(-1)^{k}2^{2k}}{(2k+1)!}(x+1)^{4k+3}+\sum_{k=0}^n \frac{(-1)^{k+1}2^{2k}}{(2k+1)!}(x+1)^{4k+2}$; ($\sin \alpha .\cos \beta =\frac{1}{2}(\sin (\alpha +\beta)+\sin (\alpha -\beta))$);
$\boxed{470.}$ hodnoty $f^{(k)}(0)$ potrebné na nájdenie Maclaurinovho polynómu funkcie $f$ možno vyjadriť z koeficientov Maclaurinovho polynómu funkcie $f'$;
$\boldsymbol{1.}$ $x+\sum_{k=1}^n \frac{(-1)^{k}.1.3.....(2k-1)}{(2k+1)2^{k}k!} x^{2k+1}$ pre $n=2n+1,2n+2$ ($n\in \mathbb{N}$), Maclaurinove polynómy stupňa 1 a 2 majú tvar $T(x)=x$;
$\boldsymbol{2.}$ $\sum_{k=0}^n \frac{(-1)^{k}x^{2k+1}}{2k+1}$ pre $n=2n+1,2n+2$ ($n=0,1,...)$;
$\boldsymbol{3.}$ $\frac{\pi}{2}-x- \sum_{k=1}^n \frac{1.3.....(2n-1)}{(2k+1)2^{k}k!}x^{2k+1}$ pre $n=2n+1,2n+2$ ($n\in \mathbb{N}$), Maclaurinove polynómy stupňa 1 resp. 2 má tvar $T(x)=\frac{\pi}{2}- x$;
$$*$$
$$*$$
$$*$$
$$*$$
$\boxed{Obr.23.}$  $y=e^{-2x}\sin ^{2} x$
$\boxed{Obr.24.}$  $y^{3}=6x^{2}-x^{3}$
$\boxed{Obr.25.}$  $y^{2}=x^{4}(x+1)$
$\boxed{471.}$  nech $e=\frac{p}{q}$ ($p\in \mathbb{Z}, q\in \mathbb{N}$); zvoľme $n\geq q$; platí $\frac{p}{q}=e=1+\frac{1}{1!}+...+\frac{1}{n!}+\frac{\vartheta}{n!}$, pričom $\vartheta \in(0,1)$; ak z tejto rovnosti vyjadríme $\vartheta$, zistíme, že $\vartheta \in \mathbb{Z}$, čo je spor;
$\boxed{472.}$ ak uvedenú rovnosť porovnáme s Taylorovým vzorcom $(n+1)$-vého stupňa so zvyškom v Peanovom tvare, dostaneme po úprave: $\vartheta (x).\frac{f^{(n)}(a+\vartheta (x)(x-a))-f^{(n)}(a)}{\vartheta (x)(x-a)}=\frac{1}{n+1}.f^{(n+1)}(a)+\frac{o(x-a)}{x-a}$;
$\boxed{473.}$ $\boldsymbol{1.}$ $\frac{19}{90}$;
$\boldsymbol{2.}$ $\frac{1}{2}$;
$\boldsymbol{3.}$ $e^{2/3}$ ($=\lim_{x \to 0}(1+\frac{2}{3}x^{3}+o(x^{3}))^{1/x^{3}}$);
$\boldsymbol{4.}$ $\frac{1}{2}$ (subst. $\frac{1}{x}=t$;
$\boldsymbol{5.}$ $\frac{1}{3}$ (pred zátvorku vyňať x, potom subst. $\frac{1}{x}=t$);
$\boldsymbol{6.}$ $-\frac{1}{4}$;
$\boxed{474.}$ hodnoty $f(0),f(1)$ zapíšte pomocou Taylorovho polynómu stupňa 1 so stredom v bode x a so zvyškom v Lagrangeovom tvare, potom jednu z týchto rovností odčítajte od druhej;
$\boxed{475.}$ $\boldsymbol{1.}$ pre $M_{2}=0$ dokážeme uvedené tvrdenie samostatne; postup pre $M_{2}\ne 0$: z rovnosti (*) $f(x+h)=f(x)+f'(x)h+f''(c)h^{2}/2$ vyjadríme $f'(x)$; odtiaľ dostaneme: pre každé $h>0$ platí $M_{1}\leq 2M_{0}/h+M_{2}h/2$; odtiaľ $M_{2}(h-M_{1}/M_{2})^{2}+4M_{0}-M_{1}^{2}M_{2}\geq 0$  pre $h>0$, potom položíme $h=M_{1}/M_{2}$;
$\boldsymbol{2.}$ postup je analogický ako v pr. 475.1, naviac využijeme ešte rovnosť $f(x-h)=f(x)-f'(x).h+f''(d)h^{2}/2$, ktorú odčítame od rovnosti (*);
$\boxed{476.}$ $\boldsymbol{1.}$ obr. 23; stačí zostrojiť graf na niektorom intervale dĺžky $\pi$ a potom využiť vzťah  $y(x+\pi)=e^{-2\pi}$ $y(x)\quad (x\in \mathbb{R}) $; 
$\boldsymbol{2.}$ obr. 24; $y_{-}'(0)=-\infty$, $y_{+}'(0)=+\infty$, $y'(6)=-\infty$; pozor: $y$ nie je konkávna na $(-\infty ,6)$; podľa nami používanej definície nie je bod 6 inflexným bodom;
$\boldsymbol{3.}$ obr. 25; dotyčnica v bode $(-1,0)$ má rovnicu $x=-1$;
$\boldsymbol{4.}$ obr. 26; (pri riešení nerovníc  $y'\ne 0, y''\ne 0$ 	nezabúdajte, že pred umocňovaním na druhú treba skontrolovať, či obidve strany nerovnice majú rovnaké znamienko, a že pred násobením číslom $\alpha $ treba zistiť znamienko tohto čísla);
$$*$$
$$*$$
$$*$$
$$*$$
$\boxed{Obr.26.}$  $y=1-x+\frac{x^{3}}{3+x}$
$\boldsymbol{5.}$ obr. 27; dotyčnica v bode $(1,0)$ má rovnicu $x=1$; ak $y=x\sqrt{1-x}/(1+x)$, tak $y''=(1+x)^{-3}(1-x)^{-3/2}(-x^{3}-6x^{2}+15x-12)/4$, pritom  $-x^{3}-6x^{2}+15x-12<0$ pre $x>-5$ (to yistíme, ak vyšetríme rast, klesanie a lokálne extrémy funkcie $-x^{3}-6x^{2}+15x-12$ na intervale $(-5,\infty$ );
$\boxed{477.}$ $\boldsymbol{1.}$ obr. 28;
$\boldsymbol{2.}$ obr. 29.$$*$$
$$*$$
$$*$$
$$*$$
$\boxed{Obr.27.}$  $y^{2}=x^{2}(1-x)(1+x)^{-2}$
$\boxed{Obr.28.}$
$\boxed{Obr.29.}$



\end{document}

% TODO
% bold a), b)
%
% teória čiara z vonkajšej strany
